\documentclass[makeidx]{article}
\usepackage{xspace}
\usepackage{epsfig}
\usepackage{xcolor}
% Russian color aliases to match translated identifiers
\providecolor{комментарийЦвет}{rgb}{0.5,0.5,0.5}
\providecolor{обоснованиеЦвет}{rgb}{0.5,0.5,0.5}
\providecolor{метаЦвет}{rgb}{0,0,1}
\providecolor{нормативныйЦвет}{rgb}{0,0,0}
\usepackage{syntax}
\usepackage[fleqn]{amsmath}
\usepackage{amssymb}
\usepackage{semantic}
\usepackage{dart}
% XeLaTeX Unicode fonts for Cyrillic
\usepackage{fontspec}
\setmainfont{Noto Serif}
\setsansfont{Noto Sans}
\setmonofont{DejaVu Sans Mono}
\usepackage[hidelinks]{hyperref}
\usepackage{makeidx}
\makeindex
\title{Dart Programming Language Specification\\
{6th edition draft}\\
{\large Version 2.13-dev}}
\author{}

% For information about Location Markers (and in particular the
% commands \LMHash and \LMLabel), see the long comment at the
% end of this file.

% CHANGES
% =======
%
% Significant changes to the specification. Note that the versions specified
% below indicate the current tool chain version when those changes were made.
% In practice, new features have always been integrated into the language
% specification (this document) a while after the change was accepted into
% the language and implemented. As of September 2022, the upcoming version of
% the language which is being specified is indicated by a version number in
% parentheses after the tool chain version.
%
% Note that the version numbers used below (up to 2.15) were associated with
% the currently released language and tools at the time of the spec change,
% they were not aligned with the actual version of the language which is being
% specified in that version of this document. This is highly misleading, so we
% changed the subtitles to be a month rather than a version number. Similarly,
% the 'Version' specified on the front page has been changed to indicate the
% version of the language which will actually be specified by the next stable
% release of this document.
%
% Jul 2025
% - Clarify that operator `[]` and `[]=` are included in the `void` allowlist
%   rule about formal parameters of type `void`, and so are setters.
%
% May 2025
% - Add `s?.length` as a constant expression in the case where `s` satisfies
%   some requirements.
%
% Apr 2025
% - Change the rules about constants to be more consistent: Every constant
%   expression is also a potentially constant expression. Also, eliminate
%   redundancies in the rules about constant collection literals and constant
%   object expressions (they are now defined just once, outside 'Constants').
%
% Sep 2024
% - Clarify the extension applicability rule to explicitly state that it
%   is concerned with an instance member with the same basename, not a
%   static member.
%
% Jun 2024
% - Add missing references to section 'Type dynamic' at the points where the
%   static analysis of Object member invocations is specified.
%
% Feb 2024
% - Correct the rule about deriving a future type from a type (which is used
%   by `flatten`).
%
% Dec 2023
% - Allow `~/` on operands of type `double` in constant expressions, aligning
%   the specification with already implemented behavior.
% - Broaden the grammar rule about `initializerExpression` to match the
%   implemented behavior. Specify that an initializer expression can not be
%   a function literal.
% - Specify in which situations it is an error to declare an initializing
%   formal parameter.
%
% Nov 2023
% - Specify that the dynamic error for calling a function in a deferred and
%   not yet loaded library will occur before actual argument evaluation, not
%   after.
%
% Oct 2023
% - Introduce the rule that an `extension` declaration cannot have the name
%   `type`. This is needed in order to disambiguate an `extension type`
%   named `on`.
%
% Aug 2023
% - Correct text about built-in identifier error and turn it into commentary
%   (the normative text is in grammar rules using `typeIdentifier`), and
%   correct the grammar for import prefixes to use `typeIdentifier`.
%
% Jul 2023
% - Land the null-safety updates for sections about variables.
% - Change terminology: A 'static variable' is now a variable whose declaration
%   includes `static` (which used to be called a 'class variable'); the old
%   meaning of 'static variable' is now never used (we spell out every time
%   that it is a 'library variable or a static variable', where the latter is
%   the new meaning of that phrase). This avoids the confusion that came up
%   just about every single time the phrase 'static variable' came up in a
%   discussion.
% - Change terminology: The notion of 'mutable' and 'immutable' variables has
%   been eliminated (it is spelled out each time which modifiers must be
%   present or not present). The concept is confusing now where a variable can
%   be late and final, and assignments to it are allowed by the static analysis.
% - Change the definition of the 'element type of a generator function', due to
%   soundness issue with the current definition.
%
% Mar 2023
% - Clarify how line breaks are handled in a multi-line string literal. Rename
%   the lexical token NEWLINE to LINE\_BREAK (clarifying that it is not `\n`).
% - Clean up grammar rules (avoid `T?` as a superclass/etc., which is an
%   error anyway; change extension names to `typeIdentifier`, avoiding
%   built-in identifiers).
%
% Feb 2023
% - Change the specification of constant expressions of the form `e1 == e2`
%   to use primitive equality.
%
% Dec 2022
% - Change the definition of the type function 'flatten' to resolve soundness
%   issue, cf. SDK issue #49396.
% - Introducing null-safety related changes to the sections about variables
%   and local variables.
% - Change 'primitive operator ==' to 'primitive equality', and include
%   constraints on `hashCode` corresponding to the ones we have on `==`.
%
% 2.15
% - Allow generic instantiation of expressions with a generic function type
%   (until now, it was an error unless the expression denoted a declaration).
% - Allow generic instantiation of the `call` method of a function object.
% - Clarify that `TypeLiteral.extensionMethod()` is an error, in line with the
%   error for `int.toString()`.
% - Add support for function closurization for callable objects, in the sense
%   that `o` is desugared as `o.call` when the context type is a function type.
% - Clarify the treatment of `covariant` parameters in the interface of a class
%   that inherits an implementation where those parameters are not covariant.
% - Adjust and clarify simple string interpolation (to allow `'$this'`, which
%   is already implemented and useful).
% - Add several lexical rules about identifiers, clarifying different kinds.
% - Clarify the conflicts between extension members and `Object` instance
%   members.
% - Correct <partDeclaration> to include metadata.
% - Clarify the section about assignable expressions.
%
% 2.14
% - Add constraint on type of parameter which is covariant-by-declaration in
%   the case where the method is inherited (that case was omitted by mistake).
% - Clarify symbol equality and identity.
%
% 2.12 - 2.13 (there was no 2.11)
% - Revert the CL where certain null safety features were removed (to enable
%   publishing a stable version of the specification).
% - Add rule that a top-level pair of declarations with the same basename
%   is a compile-time error except when it is a getter/setter pair.
% - Change grammar to enable non-function type aliases. Correct rule for
%   invoking `F.staticMethod()` where `F` is a type alias.
% - Add missing error for cyclic redirecting generative constructor.
%
% 2.8 - 2.10
% - Change several warnings to compile-time errors, matching the actual
%   behavior of tools.
% - Eliminate error for library name conflicts in imports and exports.
% - Clarify the specification of initializing formals. Fix error: Add
%   missing `?` at the end for function typed initializing formal.
% - Adjust specification of JS number semantics.
% - Revise and clarify the sections Imports and Exports.
% - Bug fix: Change <qualifiedName> to omit the single identifier, then add it
%   back when <qualifiedName> is used, as <typeIdentifier> resp. <identifier>.
% - Clarify that a function-type bounded receiver has a `.call` method.
% - Merge the `static' and `instance' scope of a class, yielding a single
%   `class' scope.
% - Add specification of `>>` in JavaScript compiled code, in appendix.
% - Integrate the specification of extension methods into this document.
% - Specify identifier references denoting extension members.
% - Remove a few null safety features, to enable publishing a stable
%   version of the specification which is purely about Dart 2.10.
% - Reorganize specification of type aliases to be in section `Type Aliases'
%   (this used to be both there, and in 'Generics').
% - Clarify the cyclicity error for type aliases ("F is not allowed to depend
%   on itself).
% - Add the error for a type alias $F$ used in an instance creation etc., when
%   $F$ expands to a type variable.
% - Correct lexical lookup rules to include implicit extension method
%   invocation.
%
% 2.7
% - Rename non-terminals `<...Definition>` to `<...Declaration>` (e.g., it is
%   'class declaration' everywhere, so `<classDefinition>` is inconsistent).
% - Clarify that <libraryDeclaration> and <partDeclaration> are the
%   start symbols of the grammar.
% - Clarify the notion of being `noSuchMethod forwarded': `m` is indeed
%   noSuchMethod forwarded if an implementation of `m` is inherited, but
%   it does not have the required signature.
% - Clarify static checks on `yield` and `yield*` to explicitly ensure that
%   assignability is enforced per element.
% - Update whitelist for expressions of type void to include `await e`,
%   consistent with decision in SDK issue #33415.
% - Re-adjust `yield*` rules: Per-element type checks are not supported,
%   the given Iterable/Stream must have a safe element type.
% - Clarify that an expression of type `X extends T` can be invoked when
%   `T` is a function type, plus other similar cases.
% - Specify actual type arguments passed to a generic function which is invoked
%   with no type arguments (so it must be a dynamic invocation).
%
% 2.6
% - Specify static analysis of a "callable object" invocation (where
%   the callee is an instance of a class that has a `call` method).
% - Specify that <libraryImport> string literals cannot use string
%   interpolation; specify that it is a dynamic error for `loadLibrary`
%   to load a different library than the one which was used for
%   compile-time checks.
% - Specify that a constant expression `e.length` can only invoke an instance
%   getter (in particular, it cannot execute/tear-off an extension member);
%   similar changes were made for several operators.
% - Specify the type arguments of the fields of an `Invocation` received by
%   `noSuchMethod`, when invoked in response to a failed instance member
%   invocation.
%
% 2.4
% - Clarify the section `Exports'.
% - Update grammar rules for <declaration>, to support `static late final`
%   variables with no initializer; for several top-level declarations,
%   to correct the existence and placement of <metadata>; for
%   <assignableSelectorPart>, to simplify the grammar (preserving the
%   derivable terms); for <topLevelDeclaration>, to allow top-level final
%   and const variables with no type, and to allow `late final` top-level
%   variables, to allow `late` on a top-level variable declaration; and
%   adding <namedParameterType> to allow `required` parameters.
% - Make lexical identifier lookups use the rules for 'Identifier Reference'
%   consistently; that is, always look up `id` as well as `id=`, and commit
%   to the kind of declaration found by that lookup.
% - Specify the signature of the `call` method of a function object.
% - Add the rule that it is an error for a type variable of a class to occur
%   in a non-covariant position in a superinterface.
% - Correct several grammar rules, including: added <functionExpressionBody> (to
%   avoid semicolon in <functionBody>), adjusted <importSpecification>.
% - Revise section on cascades. Now uses compositional grammar, and
%   specifies static type, compile-time errors, and includes `?..`.
% - Correct the grammar and lexical rules for string literals.
% - Change specification of `await` to be type-safe, avoiding cases like:
%       FutureOr<Future<S>> ffs = Future<S>.value(s);
%       Future<S> fs = await ffs;
%
% 2.3
% - Add requirement that the iterator of a for-in statement must have
%   type `Iterator`.
% - Clarify which constructors are covered by the section 'Constant
%   Constructors' and removed confusing redundancy in definiton of
%   potentially constant expressions.
% - Integrate the feature specification of collection literal elements
%   (aka UI-as-code).
%
% 2.2
% - Specify whether the values of literal expressions override Object.==.
% - Allow Type objects as case expressions and const map keys.
% - Introduce set literals.
% - Specify that a getter/setter and a method with the same basename is
%   an error, also in the case where a class obtains both from its
%   superinterfaces.
% - Specify the Dart 2.0 rule that you cannot implement, extend or mix-in
%   Function.
% - Generalize specification of type aliases such that they can denote any
%   type, not just function types.
% - Clarify that 'Constant Constructors' is concerned with non-redirecting
%   generative constructors only.
%
% 2.1
% - Remove 64-bit constraint on integer literals compiled to JavaScript numbers.
% - Allow integer literals in a double context to evaluate to a double value.
% - Specify dynamic error for a failing downcast in redirecting factory
%   constructor invocation.
% - Specify that type arguments passed in a redirecting factory constructor
%   declaration must be taken into account during static checks.
% - Disallow any expression statement starting with `{`, not just
%   those that are map literals.
% - Define a notion of lookup that is needed for superinvocations, adjust
%   specification of superinvocations accordingly.
% - Specify that it is a dynamic error to initialize a non-static variable
%   with an object that does not have the declared type (e.g., a failed cast).
% - Specify for constructor initializers that target variable must exist and
%   the initializing expression must have a type which is assignable to its
%   type.
% - Specify for superinitializers that the target constructor must exist and
%   the argument part must have a matching shape and pass type and value
%   arguments satisfying the relevant constraints.
% - Reword rules about abstract methods and inheritance to use 'class
%   interface'.
% - Specify that it is an error for a concrete class with no non-trivial
%   \code{noSuchMethod} to not have a concrete declaration for some member
%   in its interface, or to have one which is not a correct override.
% - Use \ref{bindingActualsToFormals} in 3 locations, eliminating 2 extra
%   copies of nearly the same text.
% - Add figure in \ref{bindingActualsToFormals} for improved readability.
% - Introduce a notion of lookup which is needed for superinvocations.
% - Use new lookup concept to simplify specification of getter, setter, method
%   lookup.
% - Introduce several `Case<SomeTopic>` markers in order to improve
%   readability.
% - Reorganize several sections to specify static analysis first and then
%   dynamic semantics; clarify many details along the way. The sections are:
%   \ref{variables}, \ref{new}, \ref{const}, \ref{bindingActualsToFormals},
%   \ref{unqualifiedInvocation}, \ref{functionExpressionInvocation},
%   \ref{superInvocations}, \ref{assignment}, \ref{compoundAssignment},
%   \ref{localVariableDeclaration}, and \ref{return}.
% - Corrected error involving multiple uses of the same part in the same
%   program such that it takes exports into account.
% - Eliminate all references to checked and production mode, Dart 2 does
%   not have modes.
% - Integrate feature specification on noSuchMethod forwarders.
% - Specify that bound satisfaction in generic type alias type parameters
%   must imply bound satisfaction everywhere in the body.
% - Specify that super-bounded generic type alias applications must trigger
%   a well-boundedness check on all types occurring in the denoted type.
% - Corrected corner case of rules for generation of noSuchMethod forwarders.
% - Integrate feature specification on parameters that are
%   covariant-by-declaration.
% - Integrate feature specification on parameters that are
%   covariant-by-class.
% - Correct section 'Type of a function', allowing for adjustments needed
%   for rules related to covariant parameters.
% - Specified the dynamic type of function objects in several contexts, such
%   that the special treatment of covariant parameters can be mentioned.
% - Specified what it means for an override relation to be correct, thus
%   adding the parts that are not captured by a function type subtype check.
% - Introduced the notion of member signatures, specified that they are the
%   kind of entity that a class interface contains.
% - Corrected super-boundedness check to take variance into account at the
%   top level.
%
% 2.0
% - Don't allow functions as assert test values.
% - Start running "async" functions synchronously.
% - It is a static warning and dynamic error to assign to a final local.
% - Specify what "is equivalent to" means.
% - Remove @proxy.
% - Don't specify the exact object used for empty positionalArguments and
%   namedArguments on Invocation.
% - Remove the, now unnecessary, handling of invalid overrides of noSuchMethod.
% - Add >>> as overridable operator.
% - If initializing formal has type annotation, require subtype of field type.
% - Constant `==` operations now also allowed if just one operand is null.
% - Make flatten not be recursive.
% - Disallow implementing two instantiations of the same generic interface.
% - Update "FutureOr" specification for Dart 2.0.
% - Require that a top-level "main" declaration is a valid script-entry
%   function declaration.
% - State that the return type of a setter or []= is void when not specified.
% - Clarify that "noSuchMethod" must be implemented, not just redeclared
%   abstractly, to eliminate certain diagnostic messages.
% - Add generic functions and methods to the language.
% - Don't cause warning if a non-system library import shadows a system library.
% - Update mixin application forwarding constructors to correctly handle
%   optional parameters and const constructors.
% - Specify `call` for Dart 2 (no function type given to enclosing class).
% - Clarify that an identifier reference denoting a top-level, static, or
%   local function evaluates to the closurization of that declaration.
% - Make `mixin` and `interface` built-in identifiers.
% - Make `async` *not* a reserved word inside async functions.
% - Add 'Class Member Conflicts', simplifying and adjusting rules about
%   member declaration conflicts beyond "`n` declared twice in one scope".
% - Specify that integer literals are limited to signed 64-bit values,
%   and that the `int` class is intended as signed 64-bit integer, but
%   that platforms may differ.
% - Specify variance and super-bounded types.
% - Introduce `subterm' and `immediate subterm'.
% - Introduce `top type'.
% - Specify configurable imports.
% - Specify the dynamic type of the Iterable/Future/Stream returned from
%   invocations of functions marked sync*/async/async*.
% - Add appendix listing the major differences between 64-bit integers
%   and JavaScript integers.
% - Remove appendix on naming conventions.
% - Make it explicit that "dynamic" is exported from dart:core.
% - Remove "boolean conversion". It's just an error to not be a bool.
% - Adjust cyclic subtype prevention rule for type variables.
% - Clarify that it is an error to use FutureOr<T> as a superinterface etc.
% - Eliminate the notion of static warnings, all program faults are now errors.
% - It is no longer an error for a getter to have return type `void`.
% - Specify that each redirection of a constructor is checked, statically and
%   dynamically.
% - Specify that it is an error for a superinitializer to occur anywhere else
%   than at the end of an initializer list.
% - Update the potentially/compile-time constant expression definitions
%   so that "potentially constant" depends only on the grammar, not the types
%   of sub-expressions.
% - Make `==` recognize `null` and make `&&` and `||` short-circuit in constant
%   expressions.
% - Add `as` and `is` expressions as constant expressions
% - Make `^`, `|` and `&` operations on `bool` constant operations.
% - Integrate subtyping.md. This introduces the Dart 2 rules for subtyping,
%   which in particular means that the notion of being a more specific type
%   is eliminated, and function types are made contravariant in their
%   parameter types.
% - Integrate instantiation to bound. This introduces the notions of raw
%   types, the raw-depends relation, and simple bounds; and it specifies
%   the algorithm which is used to expand a raw type (e.g., `C`) to a
%   parameterized type (e.g., `C<int>`).
% - Integrate invalid_returns.md. This replaces the rules about when it is
%   an error to have `return;` or `return e;` in a function.
% - Integrate generalized-void.md. Introduces syntactic support for using
%   `void` in many new locations, including variable type annotations and
%   actual type arguments; also adds errors for using values of type `void`.
% - Integrate implicit_creation.md, specifying how some constant expressions
%   can be written without `const`, and all occurrences of `new` can be
%   omitted.
%
% 1.15
% - Change how language specification describes control flow.
% - Object initialization now specifies initialization order correctly.
% - Specifies that leaving an await-for loop must wait for the subscription
%   to be canceled.
% - An await-for loop only pauses the subscription if it does something async.
% - Assert statements allows a "message" operand and a trailing comma.
% - The Null type is now considered a subtype of all types in most cases.
% - Specify what NEWLINE means in multiline strings.
% - Specified the FutureOf type.
% - Asserts can occur in initializer lists.
%
% 1.14
% - The call "C()" where "C" is a class name, is a now compile-time error.
% - Changed description of rewrites that depended on a function literal.
%   In many cases, the rewrite wasn't safe for asynchronous code.
% - Removed generalized tear-offs.
% - Allow "rethrow" to also end a switch case. Allow braces around switch cases.
% - Allow using `=` as default-value separator for named parameters.
% - Make it a compile-time error if a library includes the same part twice.
% - Now more specific about the return types of sync*/async/async* functions
%   in relation to return statements.
% - Allow Unicode surrogate values in String literals.
% - Make an initializing formal's value accessible in the initializer list.
% - Allow any expression in assert statements (was only conditionalExpression).
% - Allow trailing commas in argument and parameter lists.
%
% 1.11 - ECMA 408 - 4th Edition
% - Specify that potentially constant expressions must be valid expressions
%   if the parameters are non-constant.
% - Make "??" a compile-time constant operator.
% - Having multiple unnamed libraries no longer causes warnings.
% - Specify null-aware operators for static methods.
%
% 1.10
% - Allow mixins to have super-classes and super-calls.
% - Specify static type checking for the implicit for-in iterator variable.
% - Specify static types for a number of expressions where it was missing.
% - Make calls on the exact type "Function" not cause warnings.
% - Specify null-aware behavior of "e?.v++" and similar expressions.
% - Specify that `package:` URIs are treated in an implementation dependent way.
% - Require warning if for-in is used on object with no "iterator" member.
%
% 1.9 - ECMA-408 - 3rd Edition
%

\begin{document}





{заголовки}
 Спецификация языка программирования Dart


% начало Ecma boilerplate
{Сфера применения}


%
Стандарт Ecma определяет синтаксис и семантику
Язык программирования Dart.
В нем не указаны API библиотек Dart.
За исключением тех случаев, когда библиотечные элементы необходимы для
Правильное функционирование самого языка
(например, существование класса ) с помощью методов
, .


 Соответствие.


%
Соответствующая реализация языка программирования Dart
должен обеспечивать и поддерживать все API;
(библиотеки, типы, функции, геттеры, сеттеры,
будь то верхнего уровня, статический, экземпляр или локальный
предусмотренных в данной спецификации.

%
Допускается соответствующая реализация для предоставления дополнительных API,
но не дополнительный синтаксис,
За исключением экспериментальных особенностей.


 Нормативные ссылки


%
Следующие документы являются незаменимыми для
применения данного документа.
Для датированных ссылок применяется только цитируемое издание.
Для недатированных ссылок последнее издание ссылающегося документа
Применяются (включая любые поправки).



Стандарт Unicode версии 5.0 с изменениями, внесенными Unicode 5.1.0 или его преемником.

Ссылка на API Dart, https://api.dartlang.org/



 Термины и определения


%
Термины и определения, используемые в данной спецификации, приведены в
корпус собственно спецификации.
% Конец Ecma Boilerplate


{Примечание)


%
Различают нормативный и ненормативный тексты.
Нормативный текст определяет правила Дарта.
Он приведен в этом шрифте.
В настоящее время ненормативный текст включает:


[Обоснование]
Обсуждение мотивации языковых дизайнерских решений происходит курсивом.
{%
Отличие нормативного от ненормативного помогает уточнить
какая часть текста является обязательной и какая часть является просто разъяснительной.
?
 [Комментарий]
Комментарии такие как
{%
Внимательный читатель заметит
Имя Дарт имеет четыре символа %
? "
служит для иллюстрации или уточнения спецификации,
Они не соответствуют нормативному тексту.
<%
Разница между комментарием и обоснованием может быть тонкой.
?
{%
Комментарий более общий, чем обоснование.
могут включать иллюстративные примеры или пояснения.%
?


%
Зарезервированные слова и встроенные идентификаторы
(перенаправлено с «P000372»)
{ >> жирный шрифт}.

{%
Примерами могут служить \SWITCH{} или .%
?

%
Грамматические постановки приведены в общем варианте EBNF.
Левая сторона производства заканчивается на {::=}.
С правой стороны чередование представлено вертикальными полосами,
и секвенирование по интервалу.
Как и в ПЭГ, чередование отдает приоритет левым.
Факультативные элементы производства закрепляются вопросительным знаком
Например: .
Присоединение звезды к элементу производственного средства
Он может повторяться ноль или более раз.
Добавление знака плюс к произведению означает, что это происходит один или несколько раз.
Парентезы используются для группирования.
Отрицание представлено префиксом элемента производства с тильдой.
Отрицание подобно некомбинатору ПЭГ,
Но он потребляет вход, если он совпадает.
В контексте лексического производства он потребляет
один персонаж, если он есть;
В противном случае, один токен, если он есть.

{%
Примером может служить:%
?

\begin{grammar}\color{комментарийЦвет}
<aProduction> ::= <anAlternative>
  \alt <anotherAlternative>
  \alt <oneThing> <after> <another>
  \alt <zeroOrMoreThings>*
  \alt <oneOrMoreThings>+
  \alt <anOptionalThing>?
  \alt (<some> <grouped> <things>)
  \alt \gtilde<notAThing>
  \alt `aTerminal' 
  \alt <A\_LEXICAL\_THING>
\end{grammar}


%
Таким образом представлены как синтаксические, так и лексические произведения.
Лексические постановки отличаются своими названиями.
Названия лексических произведений состоят исключительно из
Верхние символы и подчеркивания.
Как всегда, в грамматических произведениях
Белое пространство и комментарии между элементами производства
неявно игнорируются, если не указано иное.
Токены препинания появляются в кавычках.

%
Производственные мощности, насколько это возможно,
При обсуждении представленных ими конструкций.

%
P000434 является синтаксической конструкцией.
Это может считаться частью текста, который можно вывести в грамматике.
Это может быть дерево, созданное таким образом.
 для данного термина  Это синтаксическая конструкция
которое соответствует непосредственному поддереву 
рассматривается как производное дерево.
 данного термина  является ,
или немедленный субтермин ,
или субтермин немедленного субтермина .

%
Список {x_1, , x_n} обозначает любой список
<<<p000010> Элементы формы .
Обратите внимание, что  Может быть нулевой, в этом случае список пуст.
Мы широко используем такие списки в этой спецификации.

%
%
Для ,
Пусть  будет атомным синтаксическим объектом (как идентификатор),
 составная синтаксическая сущность (например, выражение или тип);
и {E} снова составная синтаксическая сущность.
Запись
{}{[x1/y1, ..., xn/yn]E@$[x/y\ldots]E$]
Затем обозначает копию 
В котором каждое появление  Он был заменен на «P000020».

%
Эта операция также известна как P000437.
Это тот вариант, который позволяет избежать захвата.
То есть, когда  Содержит конструкцию, которая вводит  в
вложенный объем для некоторых ,
Замена не заменяет  в этой сфере.
И наоборот, если такая замена поставит идентификатор 
(перенаправлено с «P000025») где  объявляется,
Соответствующие декларации в  Они систематически переименовываются в новые имена.

{%
Короче говоря, свобода захвата гарантирует, что «смысл» каждого идентификатора
Сохраняется при замене.%
?

%
Иногда мы злоупотребляем списком или буквальным синтаксисом карты, написав  \List{o}{1}{n}
(соответственно :, , $k_n$:\ $o_n$\}}
где  и $k_i$ Это могут быть объекты, а не выражения.
Цель состоит в том, чтобы обозначить список (соответственно карту) объекта.
Комментарии к статье 
(соответственно, ключами которого являются ) и значения являются .

%
%
Спецификации операторов часто включают такие заявления, как:
  >>{op}  эквивалентно методу вызова
. \metavar{op}()}}
x.op(y)@\code{. \metavar{op}()}}.
Такие спецификации следует понимать как сокращение для:



   эквивалентно методу вызова
\code{.{op'}()},
Предполагая класс  фактически объявлен неоператорский метод, названный 
Определение функции оператора .


<%
Это обрезание необходимо потому, что{\rm\code{. {op}()}}, где op является оператором,
Это не юридический синтаксис.
Тем не менее, это болезненно многословно, и мы предпочитаем изложить это правило здесь,
Используйте краткую и четкую запись в спецификации.%
?

%
Когда спецификация относится к порядку, заданному в программе,
означает порядок исходного текста программы,
Сканирование слева направо и сверху вниз.

%
Когда спецификация относится к
{свежие переменные}{изменяемые}! свежий.
означает локальную переменную с именем, которое нигде не встречается.
в текущей программе.
Когда спецификация вводит новую переменную, связанную с объектом,
Свежая переменная неявно связана с окружающей областью.

%
Ссылки на иные неуказанные названия субъектов программы
(например, классы или функции)
Они интерпретируются как имена членов основной библиотеки Дарта.

{%
Примерами могут служить классы  и 
Представляя, соответственно, корень классовой иерархии и
Реификация типов Runtime.
%
Можно было бы заявить, например,
локальная переменная, называемая ,
Поэтому, как правило, неверно предполагать, что
Имя  На самом деле, мы решим, что это основной класс.
Однако мы обычно не упоминаем об этом, для краткости.%
?

% % TODO(Eernst): Мы должны избавиться от понятия «эквивалентно».
% см. Языковая проблема https://github.com/dart-lang/language/issues/227.
% % В этой КЛ в нескольких местах была введена фраза "обращались как",
%% и вышеупомянутый выпуск 227 приведет к полному пересмотру
% от этой части документа. В частности, следующий пункт будет
% исключить.

%
Когда в спецификации указано, что одна часть синтаксиса 
другой фрагмент синтаксиса, это означает, что он эквивалентен во всех отношениях;
и прежний синтаксис должен генерировать те же ошибки компиляции
и имеют такое же поведение во время выполнения, как и последние, если таковые имеются.
{%
Сообщения об ошибках, если таковые имеются, всегда должны относиться к исходному синтаксису.
?
Если считается, что выполнение или оценка конструкции
эквивалентно исполнению или оценке другой конструкции,
Тогда только время выполнения эквивалентно,
Ошибки времени компиляции применяются только к исходному синтаксису.

%
%
Когда в спецификации указано, что одна часть синтаксиса  это

Еще одна часть синтаксиса ,
Это означает, что статический анализ  Статический анализ 
{(в частности, происходят одни и те же ошибки компиляции)}.
Более того, если  не имеет ошибок времени компиляции
Поведение  Во время выполнения точно такое же поведение .

{%
Ошибка {сообщения}, если таковые имеются,
Всегда следует обращаться к исходному синтаксису. %
?

{%
Короче говоря, всякий раз, когда  Обозначается как ,
читатель должен немедленно перейти к разделу о 
Для получения дополнительной информации о
Статический анализ и динамическая семантика .%
?

<%
Понятие «обработка» аналогично понятию синтаксического сахара:
 Обозначается как  ''
Также можно было бы написать
 Он был преобразован в «P000067».
Конечно, тогда его действительно следует назвать «семантический сахар».
в силу применимости преобразования и построения 
Может опираться на информацию статического анализа.

Дело в том, что мы конкретизируем только статический анализ и динамическую семантику.
Основной язык, который является подмножеством Dart
(немного меньше, чем Дарт)
Преобразование любой программы Дарта в
Программа на этом базовом языке.
Это помогает поддерживать языковую спецификацию последовательной и понятной.Потому что он показывает прямо
что некоторые особенности языка вводят существенную семантику,
и другие лучше описывать как простые сокращения существующих конструкций.
?

%
Спецификация использует одну синтаксическую конструкцию,
{{) {let expression@\LET{} выражение
который не может быть выведен в грамматике
{(то есть ни один исходный код Dart не содержит такого выражения)}.
Это выражение полезно для определения определенных синтаксических форм.
которые рассматриваются как другие синтаксические формы,
поскольку она позволяет вводить и инициализировать одну или несколько свежих переменных,
и использовать их в выражениях.

{%
То есть,  Выражение вводится только как инструмент.
Определение оценочной семантики выражения
в терминах других выражений, содержащих  выражения.%
?

%
Синтаксис  Выражение выглядит следующим образом:


<letExpression> ::= \LET{} <staticFinalDeclarationList>{} <выражение>


%
% \BlindDefineSymbol removed due to translation artifacts
\marginpar{\textcolor{black}{\ensuremath{e_{{let}}, e_j, v_j, k}}}%
Пусть  будет \LET Выражение формы
{}{}{}{}{}.
Негласно предполагается, что  является свежей переменной, ,
Если только не указано обратное.

%
 Содержит  вложенные рамки,  S}{1}{k}}
Сфера охвата для  является текущим объемом для ,
и прилагаемая область для  является , .
Текущий объем  является текущим объемом ,
Текущий объем  составляет , ,
и текущий объем  является .
Для ,  вводит конечную локальную переменную в ,
Статический тип  как заявленный тип.

{%
Типовой вывод  и тип контекста, используемый для вывода 
не являются актуальными.
Обычно предполагается, что вывод типа уже произошел.
()%
?

%
Оценка  доходы от
оценка  >> к объекту  >> и связывание  с ,
где , в указанном порядке.
Наконец,  оценивается на объект  А потом
 Показатель .

%
Правая граница каждой страницы в этом документе используется для указания
Ссылки на объекты.

%
Документ содержит индекс в конце.
Каждая запись в индексе относится к номеру страницы .
На странице  На марже есть «$\diamond$»
при определении данной индексируемой фразы,
и сама фраза показана с использованием {этот шрифт}.
Настоящим мы представляем

себя.

%
Правый край также содержит символы.
Всякий раз, когда символ
% \BlindDefineSymbol removed due to translation artifacts
\marginpar{\textcolor{black}{\ensuremath{C}}}%
\marginpar{\textcolor{black}{\ensuremath{x_j}}}%
\commentary{(скажем,  или )}
вводится и используется более чем в нескольких строках текста;
Это показано на полях.

{%
Дело в том, что легко найти определение символа.
Сканирование текста до тех пор, пока этот символ не появится на полях.
Чтобы избежать бесполезной многословности, некоторые символы не упоминаются на полях.
Например, мы можем ввести {e}{1}{k},
Но показывают только  и  на полях.

Обратите внимание, что может потребоваться посмотреть несколько строк текста.
Надпись «P000115» или символ,
Потому что маржинальные маркеры могут быть сдвинуты вниз по одной линии.
когда имеется более одного маркера для одной строки.
?


 Обзор


%
Дарт — это классовое, однонаследное, чистоеОбъектно-ориентированный язык программирования.
Dart необязательно типизирован () и поддерживает модифицированные дженерики.
Тип времени выполнения каждого объекта представлен как
Пример класса  который может быть получен
Позвонив Геттеру  в классе ,
Корень иерархии классов Дарта.

%
Программы Дарта можно статически проверять.
Программы с ошибками времени компиляции не имеют заданной динамической семантики.
Данная спецификация не дает возможности ответить на дополнительные вопросы.
О библиотеке или программе в данный момент
где, как известно, имеется ошибка времени компиляции.

{%
Однако инструменты могут поддерживать выполнение некоторых программ с ошибками.
Например, компилятор может компилировать определенные конструкции с ошибками, такими, что
Динамическая ошибка возникает при попытке
осуществить такую конструкцию,
или интегрированное время выполнения IDE может поддерживать открытие
окно редактора при выполнении такой конструкции,
Это позволяет разработчикам исправить ошибку.
Предполагается, что такие признаки будут представлять собой естественное расширение
Динамическая семантика Дарта, как указано здесь, но, как уже упоминалось,
Эта спецификация не пытается точно указать, что это означает.
?

%
Как указано в этом документе,
динамические проверки гарантированно выполняются в определенных ситуациях;
и некоторые нарушения системы типа бросают исключения во время выполнения.

{%
Реализация может опускать такие проверки всякий раз, когда они
гарантированный успех, например, на основе результатов статического анализа.
?

{%
Сосуществование между факультативной типизацией и реификацией
основывается на следующем:



Информация Reified Type отражает типы объектов во время выполнения
и всегда могут быть запрошены динамическими конструкциями проверки типа
(аналоги exampleOf, casts, typecase и т. д.  на других языках).
Информация о типе включает в себя
Доступ к экземплярам класса  представлять типы,
тип времени выполнения (класс ака) объекта,
и фактические значения параметров типа
Конструкторы и общие вызовы функций.

Типовые аннотации декларируют типы
переменные и функции (включая методы и конструкторы);

% % TODO(Eernst): Изменения при интеграции instantiate-to-bounds.md.
Аннотации типа могут быть опущены, и в этом случае они обычно
Заполняется шрифтом \DYNAMIC{}
(P000375).
%
?

% % TODO(Eernst): Обновление при добавлении выводов.
{%
Реализованный Dart включает в себя обширную поддержку вывода опущенных типов.
Эта спецификация делает предположение, что вывод имел место,
Следовательно, предполагаемые типы уже присутствуют в программе.
Однако в некоторых случаях информация отсутствует.
чтобы сделать вывод о пропущенном типе аннотации,
И поэтому эта спецификация все еще должна указывать, как с этим бороться.
Будущая версия этой спецификации также будет указывать тип вывода.
?

%
Программы Dart организованы по модульному принципу.
единицы, называемые {библиотеки} (P000376).
Библиотеки являются единицами инкапсуляции и могут быть взаимно рекурсивными.

{%
Но это не первый класс.
Чтобы получить несколько копий библиотеки, работающей одновременно,
Для этого необходимо создать изолятору.%
?

%
Выполнение программы dart может происходить с включенными или отключенными утверждениями.
Метод, используемый для включения или отключения утверждений, является специфическим для реализации.


\subsection{Скопинг}


%
A {пространство имён времени компиляции}{namespace!compiletime}
Это частичная функция, которая отображает имена для значений пространства имен.
Пространства имён времени компиляции используются гораздо чаще, чем пространства имён времени выполнения
{(определено ниже в этом разделе)},
Когда слово  используется в одиночку,
Это означает пространство имён компиляции времени. является лексическим токеном, который является 
ан  Идентификатор, за которым следует {=}, или
ан {оператор},
;
 это
Декларация, пространство имен или специальное значение 
(P000377).

%
Если  Тогда мы говорим, что 
{maps}{namespace!maps a key to a value}
тот
{key}{namespace!key}

{value}{namespace!value}
,
и это 
{имеет связывание}{namespace!has a binding}
.
{%
Дело в том, что  Частичные функции означают, что
Каждое имя отображается как минимум до одного значения пространства имен.
То есть, если  имеет привязки
 и 
%
?

%
Пусть {} будет пространством имен.
Мы говорим, что имя   
Если  является ключом от .
Мы говорим декларацию  имеется в виду) 
Если ключ  Соответствует .

%
Область применения  имеет связанное пространство имен {0}.
Обязательства {0} указаны в этом документе тем, что
Данная декларация  Название: P000129
{introduces}{declaration!introtro предприятие в сферу применения.
конкретный субъект {V} в ,
что означает связывание  Добавлено в {0}.

{%
В некоторых случаях название декларации отличается от
идентификатор, который встречается в синтаксисе декларации, используемом для его объявления.
Сеттеры имеют имена, отличные от соответствующих геттеров.
потому что у них всегда есть  автоматически добавляется в конце,
и унарный оператор минус имеет специальное название .%
?

{%
Как правило,  Декларация  себя,
Но есть исключения.
Например,
переменная декларация вводит неявно индуцированную декларацию Геттера;
а в некоторых случаях и неявно индуцированной установочной декларацией в
Сфера применения.%
?

{%
Обратите внимание, что метки () не включены в пространство имен области.
Они решаются лексически, а затем просматриваются в пространстве имен.
?

%
Это ошибка времени компиляции, если существует более одного объекта с тем же именем.
заявленных в том же объеме.

{%
Поэтому невозможно, например, определить класс, который заявляет, что
Метод и геттер с тем же названием в Дарте.
Точно так же нельзя объявить функцию верхнего уровня
Название библиотеки переменная или класс
Объявлено в той же библиотеке.%
?

%
Мы вводим понятие а
.
Это частичная функция от имен до сущностей времени выполнения,
В частности, места хранения и функции.
Каждое пространство имен во время выполнения соответствует пространству имен с одинаковыми ключами.
но со значениями, соответствующими семантике значений пространства имен.

<%
Пространство имен обычно отображает имя в декларацию,
И его можно использовать статически, чтобы выяснить, что это имя означает.
Например,
переменная связана с фактическим местом хранения во время выполнения.
Мы вводим понятие пространства имен во время выполнения, основанного на пространстве имен,
так, чтобы динамическая семантика могла получить доступ к объектам времени выполнения
Как это место хранения.
Один и тот же код может выполняться несколько раз с одним и тем же пространством имен во время выполнения.
или с различными пространствами имен времени выполнения для каждого исполнения.
Например, локальные переменные, объявленные внутри функции.
специфичны для каждого вызова функции,
и переменные экземпляра специфичны для объекта.%
?

%
Дарт лексически ограничен.
Сферы могут гнездиться.
Имя или декларация   Если  находится в пространстве имен, индуцируемом  или
Если  доступен в лексически закрытом диапазоне .
Мы говорим, что имя или декларация  является 
Если  доступен в текущем объеме.

%
Если заявление  Оригинальное название: P000143 находится в пространстве имен, индуцируемом областью ,
  Любое заявление под названием  Это доступно
в лексически замкнутой области .

{%
Следствием этих правил является то, что можно скрыть тип
Метод или переменная.
Соглашения об именах обычно предотвращают такие злоупотребления.
Тем не менее, следующие программы являются законными:
?



%
Наименования могут вводиться в сферу действия посредством заявлений в пределах сферы действия.
или другими механизмами, такими как импорт или наследование.

{%
Взаимодействие лексики и наследования является тонким.
В конечном счете, вопрос заключается в том, является ли лексический анализ
преобладает над наследством или наоборот.
Дарт выбирает первое.

Позволить наследственным именам иметь приоритет над местными именами
Они могут создавать неожиданные ситуации по мере развития кода.
В частности, поведение кода в подклассе может незаметно измениться.
Если в суперклассе вводится новое имя.
Рассмотрим:%
?



{%
Теперь предположим метод  Добавлено в .%
?



{%
Если наследование превалирует над лексическим размахом,
Поведение  Изменится неожиданным образом.
Ни один из авторов  Автор статьи 
Они обязательно осознают это.
В Дарте, если существует лексически видимый метод ,
Это всегда будет называться.

Теперь рассмотрим противоположный сценарий.
Начнем с версии  который содержит ,
Но не объявляйте  В библиотеке .
Опять же, есть изменения в поведении, но автор 
Это тот, кто ввел несоответствие, которое влияет на их код.
Новый код является лексически видимым.
Оба эти фактора повышают вероятность обнаружения проблемы.

Эти соображения становятся еще более важными.
Если ввести такие конструкции, как вложенные классы,
которые могут быть рассмотрены в будущих версиях языка.

Хорошие инструменты должны, конечно, стремиться информировать программистов.
В таких ситуациях (с осторожностью).
Например, идентификатор, который является как наследственным, так и лексически видимым.
Можно выделить (через подчеркивание или окраску).
Более того, тесная интеграция контроля источника с инструментами языкового знания
Выявляет такие изменения, когда они происходят.%
?


 Конфиденциальность


%
Dart поддерживает два уровня P000449: публичной и частной.
Декларация {private}{private!declaration}
Если его имя является частным,
В противном случае это {public}{public!declaration}.
Имя  \IndexCustom{private}{private!name}
f любой из идентификаторов, содержащих  является частным,
В противном случае это {public}{public!name}.
Идентификатор: {private}{private!identifier}
если оно начинается с подчеркивания () характер
В противном случае это {public}{public!identifier}.

%
Декларация  является \Index{accessible to a library} 
Если  Объявляется в  Если  является публичным.

<%
Это означает, что частные декларации могут быть доступны только в пределах
Библиотека, в которой они объявлены.%
?

{%
Конфиденциальность распространяется только на декларации в библиотеке.
Не в библиотечную декларацию в целом.
Это потому, что библиотеки не ссылаются друг на друга по имени.
Идея частной библиотеки бессмысленна.
(P000379).Если название библиотеки начинается с символа,
не влияет на доступность библиотеки или ее членов.
?

{%
Конфиденциальность на данный момент является статическим понятием, связанным с
определенный фрагмент кода (библиотека).
Он предназначен для поддержки проблем разработки программного обеспечения
Вместо проблем безопасности.
Ненадежный код всегда должен работать в другом Изоляторе.

Конфиденциальность указывается именем декларации - отсюда
Приватность и имена не являются ортогональными.
Это имеет то преимущество, что и люди, и машины
может признавать доступ к частным декларациям в точке использования;
без знания контекста, из которого выводится декларация.
?


 Конкуренция


%
Код Дарта всегда однопоточный.
В Дарте нет параллели между общими государствами.
Конкурентность поддерживается через актероподобные сущности, называемые .

%
Изолировать Термин — это единица параллелизма.
У него есть собственная память и собственная нить управления.
изолят осуществляет связь посредством передачи сообщений ().
Ни одно государство никогда не делится между собой.
Изолятор создается путем нереста ().


 Ошибки и предупреждения


%
Эта спецификация различает несколько видов ошибок.

%
 Ошибки времени компиляции
Ошибки, которые исключают исполнение.
Ошибка времени компиляции должна сообщаться компилятором Dart.
До того, как будет выполнен ошибочный код.

<<%
Реализация Дарта имеет значительную свободу.
Когда происходит компиляция.
Современные реализации языков программирования
часто переплетаются компиляция и исполнение,
так, чтобы компиляция способа могла быть отложена, например,
до тех пор, пока он не будет впервые использован.
Следовательно, ошибки времени компиляции в методе  может быть сообщено
Это произошло во время первого вызова P000156.

Дарт часто загружается непосредственно из источника.
без промежуточного бинарного представления.
В интересах быстрой загрузки реализация Dart
Например, можно избежать полного разбора методических тел.
Это можно сделать путем токенизации ввода и
проверка сбалансированных кудрявых брекетов при входе в тело.
В такой реализации будут обнаружены даже синтаксические ошибки
только в том случае, если метод необходимо выполнить,
В какое время он будет собран (смеется).

В среде разработки компилятор должен
сообщать об ошибках компиляции, чтобы наилучшим образом служить программисту.

Среда развития Дарта может выбрать поддержку
устранение ошибок при преобразовании программ, например,
Замена ошибочного выражения вызовом отладчика.
выходит за рамки настоящего документа
Чтобы определить, как работают такие преобразования,
где они могут быть использованы.%
?

%
Если происходит непойманная ошибка времени компиляции
в коде запущенного изолированного термина ,  немедленно приостанавливается.
Единственным обстоятельством, при котором может быть обнаружена ошибка компиляции времени, является
С помощью кода работает отражательно, где зеркальная система может поймать его.

<<%
Как правило, после ошибки компиляции и  приостановлено,
 будет прекращен.
Однако это зависит от общей окружающей среды.
Двигатель Dart работает в контексте ,
Программа, которая взаимодействует между двигателем и
окружающей вычислительной среды.
Время выполнения может быть, например, программой C++ на сервере.
Когда изолятор терпит неудачу с ошибкой времени компиляции, как описано выше,
контроль возвращается во время выполнения,
За исключением случая, описывающего проблему.
Это необходимо для того, чтобы время выполнения могло очищать ресурсы и т.д.
Именно тогда время выполнения решает, прекращать ли изоляторе или нет.%
?

%
 Статические предупреждения
ситуации, которые не препятствуют исполнению,
которые, однако, вряд ли будутОни могут вызвать проблемы или неудобства.
А. Компилятор Dart может сообщать обо всех,
или ни одно из указанных статических предупреждений до выполнения соответствующего кода.
А. Компилятор Dart также может сообщать о дополнительных предупреждениях.
не определены данной спецификацией.

%
Когда эта спецификация говорит, что  происходит,
Это означает, что соответствующий объект ошибки выбрасывается.
Когда говорится, что  происходит,
представляет собой неудавшуюся проверку типа во время выполнения,
и предмет, который брошен орудиями .

%
Всякий раз, когда мы говорим, что исключение  это
\IndexCustom{thrown}{throwing an exception},
Он действует так, как будто его бросило выражение ().
 как объект исключения и со следом стека
соответствующее текущему состоянию системы.
Когда мы говорим, что  {выбрасывается}{бросает класс},
где  Это класс, мы имеем в виду, что бросается экземпляр класса .

%
Если непойманное исключение выкинуто бегущим изолят ,
 немедленно приостанавливается.


 Переменные.


%
Переменные — это места хранения в памяти.


<finalConstVarOrType>:=  {} <тип>
 {} <тип>
  <varOrType>

<varOrType> ::= <
 <type>

<initializedVariableDeclaration> ::= <
<declaredIdentifier> ('=' <expression>)? (',' <initialized) Идентификатор> *

<initializedIdentifier> ::= <identifier> ('=' <expression>)?

<initializedIdentifierList> ::=
<initializedIdentifier> (',' <initialized) Идентификатор> *


%
{инициализированная переменная декларация}
Объясняет две или более переменных
эквивалентны множественным переменным декларациям, объявляющим
одинаковый набор переменных имен, в том же порядке,
с той же инициализацией, типом и модификатором.

{%
Например,
\code{{} x, y;}
является эквивалентным
\code{{} x; {} y;}
и
\code{\STATIC\,\,\LATE\,\,\FINAL String\ s1,\ s2\ =\ "foo";}
эквивалентно тому, чтобы иметь оба
\code{\STATIC\,\,\LATE\,\,\FINAL String\ s1;}
и
\code{\STATIC\,\,\LATE\,\,\FINAL String\ s2\ =\ "foo";}%
?

%
Можно включить в переменную декларацию модификатор .
Эффект от выполнения этого с переменной экземпляра описан в другом месте.
(P000383).
Это {ошибка времени компиляции} для объявления переменной.
которая не является переменной экземпляра для включения модификатора .

{%
Формальное объявление параметров индуцирует локальную переменную в область,
Но формальные декларации о параметрах не являются вариативными.
Не допускайте вышеописанных ошибок.
Влияние наличия модификатора {} на формальный параметр
описывается в другом месте
()%
?

%
В переменной декларации одной из форм

или
\code{$N$\,\,\id{}\ =\ $e$;}
где $N$ происходит от
\syntax{<metadata> <finalConstVarOrType>} и \id{} является идентификатором,
Мы говорим, что  является  идентификатора.
Для каждого идентификатора, который не является декларирующим событием,
Мы говорим, что это .
Мы также сокращаем это, чтобы сказать, что идентификатор
а  соответственно .

{%
В выражении формы  Возможно, что
 имеет статический тип \DYNAMIC{{} и  не могут быть связаны с
любое конкретное заявление под названием  во время составления,Но в этой ситуации  Идентификатор ссылки.%
?

%
Для краткости мы будем ссылаться на переменную, используя ее название.
Даже если имя переменной и сама переменная
Это очень разные понятия.

{%
Итак, мы поговорим о переменной .
Вместо введения переменной  Оригинальное название: P000570 ,
Для того, чтобы потом можно было сказать "переменная ".
Это не должно создавать двусмысленности.
{поскольку} понятие имени и понятие переменной
Они очень разные.%
?

%

является переменной декларацией, идентификатор декларирования которой
Далее следуют  и .

%
Переменная, объявленная на верхнем уровне библиотеки, называется
{library variable}{variable!library}
{Вариабельный!Верхний уровень}.
Это {ошибка времени компиляции} при объявлении переменной библиотеки.
Имеет модификатор .

>>>>>>>>>>>>>>>>>>>>>>>>>>>>>>>>>>>>>>>>>>>>>>>>>>>>>>>>>>>>>>>>>>>>>>>>>>>>>>>>>>>>>>>>>>>>>>>>>>>>>>>>>>>>>>>>>>>>>>>>>>>>>>>>>>>>>>>>>>>>>>>>>>>>>>>>>>>>>>>>>>>>>>>>>>>>>>>>>>>>>>>>>>>>>
A {статическая переменная}{вариабельная!статическая}
переменной, заявление которой является
немедленно вложенный внутрь , , , или  заявление
Включает в себя модификатор .

%
{ошибка времени компиляции} возникает, если статическая или библиотечная переменная имеет
отсутствие инициализации выражения и тип, который не является нулевым
(),
Если в переменной декларации нет
модификатор {} или модификатор .

%
A {нелокальная переменная}{вариабельная!нелокальная}
Это библиотечная переменная, статическая переменная или переменная экземпляра.
{%
То есть любой вид переменной, который не является локальной переменной.
?

%
A {постоянная переменная}{вариабельная!постоянная}
переменная, в которую включается модификатор .
Постоянная переменная должна быть инициализирована в постоянное выражение.
(),
{ошибка времени компиляции}.

{%
Начальное выражение  с постоянной переменной декларацией
происходит в постоянном контексте
(P000387).
Это означает, что  Модификаторы в  Не обязательно уточнять.%
?

{%
Это грамматически невозможно для постоянной переменной декларации.
Модификатор .
Даже если бы это было грамматически возможно,
a  Постоянная переменная все равно должна быть ошибкой времени компиляции.
Потому что быть постоянной времени компиляции по своей сути несовместимо.
быть просчитанным поздно.

Аналогично, переменная экземпляра не может быть постоянной.
()%
?


 Неявно индуцированные геттеры и сеттеры


% % TODO(Eernst): Когда будет сделан вывод, мы сможем сделать вывод
% %, что случаи без заявленного типа не существуют после вывода типа
% (например, «var x» или «var x = e»); затем мы можем заменить все правила;
% % о таких случаях, комментируя, что они могут существовать во входе;
%, но они ушли после вывода типа.
% %
% % В настоящее время мы полагаемся на предположение, что вывод типа уже
Произошел %%, что означает, что мы можем ссылаться на объявленный тип переменной.
% без упоминания типа вывода.

%
Следующие правила о неявно индуцированных геттерах и сеттерах
распространяется на все неместные переменные декларации.
{%
Местные переменные декларации
(перенаправлено с «P000389»)
не вызывают геттеров или сеттеров.%
?

%
 Геттер: Переменная с объявленным типом
Рассмотрим переменную декларацию одной из форм


 
  \,\,  \,\,$T$\,\,\id{} = $e$;  \,\,\CONST\,\,$T$\,\,\id{} = $e$;



где  Тип:  является идентификатором,
и {?} указывает, что данный модификатор может присутствовать или отсутствовать.
Каждая из этих деклараций неявно вызывает геттер.
()
с заголовком ,
чье обращение оценивается как описано ниже
(P000391).
В этих случаях заявленный тип  является .


%
 Геттер: переменный без декларируемого типа
Переменная декларация одной из форм


 
   \,\,\VAR\,\,\id{}=$e$;
 
   \,\,\FINAL\,\,\id{}=$e$;
  \,\,\CONST\,\,\id{}=$e$;


<<<p001010>
неявно вызывает геттер с заголовком, который
Содержит \STATIC{} если {}f декларация содержит \STATIC{} После чего следует
,
где  Получено из вывода типа
в случае, когда  существует,
и  является {} в противном случае.
Призыв этого геттера оценивает, как описано ниже.
(P000392).
В этих случаях заявленный тип  является .


%
 Сеттер: Переменная переменная с объявленным типом
Переменная декларация одной из форм


 
   \,\,$T$\,\,\id{}=$e$;



неявно вызывает сеттер () с заголовком
,
чье исполнение устанавливает значение  для входящего аргумента .

%
 Сеттер: Переменная переменная без объявленного типа с инициализацией
Переменная декларация формы
  \,\,\VAR\,\,\id{}=$e$;
неявно вызывает сеттер с заголовком
,
чье исполнение устанавливает значение  для входящего аргумента .

{%
Вывод типа мог бы обеспечить тип, отличный от 
()%
?


%
 Сеттер: изменчивая переменная без объявленного типа, без инициализации
Переменная декларация формы

неявно вызывает сеттер с заголовком
,
чье исполнение устанавливает значение  для входящего аргумента .

{%
Тип вывода еще не определен в этом документе
().
Обратите внимание, что тип вывода может меняться, например,
,
Что приведет нас к более раннему случаю.%
?


%
 Сеттер: Поздно-финальная переменная с объявленным типом
Переменная декларация формы

неявно вызывает сеттер () с заголовком
.
Если этот сеттер выполнен
в ситуации, когда переменная  не был связан,
Он связывает  с объектом, который  Это неизбежно.
Если этот сеттер выполнен
в ситуации, когда переменная  был привязан к объекту,
Возникает динамическая ошибка.


% Сеттер: переменная с поздним окончанием без объявленного типа, без инициализации
Переменная декларация формы

неявно вызывает сеттер с заголовком
.
Казнь указанного сеттера
в ситуации, когда переменная  не был связан
связываться  к объекту, что аргумент  Это неизбежно.
Исполнитель The Setter
в ситуации, когда переменная  был привязан к объекту
Это приведет к динамической ошибке.


%
Сфера, в которую вводятся неявные геттеры и сеттеры
Это зависит от вида варьируемой декларации.

%
Переменная библиотеки вводит геттер в
библиотечный объем ограждающей библиотеки.
Статическая переменная вводит статический геттер в
сфера действия корпуса немедленно закрывающего
класс, микширование, enum или расширение декларации.
Переменная экземпляра вводит в
сфера действия корпуса немедленно закрывающего
класс, микширование или декларирование enum
{(расширение не может декларировать переменные экземпляра)}.

%
Нелокальная переменная вводит сеттер, если {}f
не имеет модификатора \FINAL{} или модификатора \CONST{},
или {} и , но не имеет инициализирующего выражения.

%
Переменная библиотеки, которая вводит сеттер, введет
библиотечный сеттер в объем библиотеки.
Статическая переменная, которая вводит сеттер, введет
статический сеттер в область тела немедленно закрывающего
класс, микширование, enum или расширение декларации.
Переменная экземпляра, которая вводит установщик, введет
инстанция, установленная в объем тела немедленно закрывающего
класс, микширование или декларирование enum
{(расширение не может декларировать переменные экземпляра)}.

%
 быть переменной, объявленной переменной декларацией
имеет начальное выражение .
Это {ошибка времени компиляции}, если статический тип 
не присваивается заявленному типу .
Это {ошибка времени компиляции}, если переменная конечного экземпляра
заявление которого имеет инициализаторное выражение
Он также инициализируется конструктором.
либо путем инициализации формального или инициализатора списка.

{%
Ошибка времени компиляции, если переменная последнего экземпляра
который был инициализирован посредством
форма инициализации конструктора 
Также инициализируется в списке инициализаторов 
(P000397).

Декларация непоздней статической переменной  Оригинальное название: P000591
имеет модификатор \FINAL{} или модификатор \CONST{}
не вызывает сеттера.
Однако назначение на  в месте, где  находится в сфере
Это не обязательно ошибка времени компиляции.
Например, сеттер по имени  Можно найти по лексическому поиску
(P000398).

Аналогично, переменная непоздней конечной инстанции  не вызывает сеттера,
Но назначение может быть призывом к
Сеттер, который предоставляется другим способом.
Например, может быть, что лексический поиск ничего не дает.
и место назначения имеет доступ к ,
Интерфейс ограждающего класса имеет сеттер под названием 
(в этом случае и геттер, и сеттер наследуются).%
?

%
Рассмотрим переменную 
Декларация которого не имеет модификатора .
Предположим, что  не имеет начального выражения,
И это не инициализируется инициализацией формального
(),
ни элементом в списке инициализаторов конструктора
(перенаправлено с «P000400»).
Начальное значение  Тогда нулевой объект
().

{%
Обратите внимание, что существует множество ситуаций, когда такое заявление переменной
Ошибка времени компиляции,В этом случае начальное значение, конечно, не имеет значения.%
?

%
В противном случае переменная инициализация осуществляется следующим образом:

%
Объявление статической или библиотечной переменной
с инициализирующим выражением инициализируется лениво
().

{%
Ленивая семантика дается потому, что мы не хотим языка.
где имеет тенденцию определять дорогостоящие вычисления инициализации,
длительное время запуска приложений.
Это особенно важно для Дарта.
который должен поддерживать кодирование клиентских приложений.%
?

{%
Инициализация переменной экземпляра без инициализации выражения
Происходит во время выполнения конструктором
() %
?

%
%
Инициализация переменной экземпляра 
с инициализирующим выражением 
идет следующим образом:
 оценивается на предмет 
и переменной  Связано с .

{%
Он указывается в другом месте, когда происходит эта инициализация.
и в какой среде
(p.\,\pageref{исполнение Генеративных Конструкторов}),
,
%
?

{%
Если инициализирующее выражение бросает тогда
доступ к неинициализированной переменной предотвращен;
Потому что создание экземпляра
Это привело к тому, что инициализация
бросок.%
?

%
% Эта ошибка может произойти из-за неявного понижения от .
Ошибка динамического типа, если динамический тип  не является
подтип фактического типа переменной 
().


 Оценка неявных переменных Геттеров


%
Мы вводим два конкретных типа геттеров, чтобы указать
Поведение различных типов переменных
без дублирования спецификации их общего поведения.

%
А.
<{позднеинициализированный} Getter! (с поздней инициализацией)
является геттером  который неявно индуцируется нелокальной переменной 
имеет начальное выражение .
{%
Ниже описано, какие объявления вызывают поздний инициализ Геттера.%
?
Призыв  идет следующим образом:

%
Если переменная  Не был привязан к объекту тогда
 оценивается на предмет .
Если  Теперь он привязан к объекту и  является окончательным,
Возникает динамическая ошибка. В % реализации бросают «ошибку поздней инициализации».
В противном случае  Обязателен для ,
и оценки  Полное возвращение .
Если оценка  бросок тогда
Обращение к  завершает бросание одного и того же предмета и следа стека,
и не изменяет связывания .

%
Призыв  В ситуации, когда  был привязан к объекту 
Завершается немедленно, возвращаясь .

{%
Рассмотрим нелокальную переменную декларацию формы
,
чей неявно индуцированный геттер запоздало инициализируется.
Возможно, удивительно,
Если переменная  был привязан к объекту
Когда впервые вспыхивает его геттер,
 Они никогда не будут казнены.
Другими словами, выражение инициализации может быть предупредительным.
по заданию.%
?

{%
Также обратите внимание, что инициализирующее выражение может иметь побочные эффекты.
Это важно при инициализации.
Например: %
?



{%
В этом примере  призывать
неявно индуцированный геттер по имени ,
и переменной  На данный момент он не связан.
Таким образом, оценка инициализации выражения продолжается.
Это приводит к переключению  на ,
Это приводит к повторной оценке .Это вызывает геттер  Чтобы вновь призвать,
И все же верно, что переменная не была связана,
Таким образом, выражение инициализации оценивается повторно.
Это переключает  на ,
что вызывает оценку ,
который вызывает неявно индуцированный установщик, названный ,
и последний вызов геттера 
Возвращение 10.
Это означает, что  Оценить 11,
Затем переменная привязывается к 11.
Наконец, вызов геттера \code{i} в \code{main}
Возвращение 11.%
?

{%
Это изменение от семантики старых версий Дарта:
Выбрасывание исключения во время оценки инициализатора больше не устанавливает
переменная на ,
и считывание переменной во время оценки инициализатора
больше не вызывает динамических ошибок.
?

%
А.
{позднеинициализированный} {getter!late-uninitialized}
Это геттер (P000249). который неявно индуцируется нелокальной переменной 
Это не имеет начального выражения.
{%
Опять же, только {некоторые} нелокальные переменные без инициализации выражения
индуцировать поздне-неинициализированный геттер, как указано ниже.%
?
Призыв  идет следующим образом:
Если переменная  Он не привязан к объекту, поэтому возникает динамическая ошибка.
Если  Привязан к объекту  затем
Ссылка на  Завершается сразу, возвращаясь .

%
%
 Объявление нелокальной переменной, называемой .
Для краткости, мы будем ссылаться на него ниже как
пункт переменная  или , или просто как  или .
Исполнение неявно индуцированного геттера  идет следующим образом:

%
 Переменная непоздней инстанции}
Если  объявляет переменную экземпляра, не имеющую модификатора ,
затем вызов неявного геттера  оценивать на
Объект, который  Это неизбежно.

{%
Невозможно вызвать на объекте геттер экземпляра
до того, как объект был инициализирован.
Следовательно, переменная непоздней инстанции всегда связана
когда используется кеттер.%
?


%
{Поздняя, инициализированная переменная экземпляра}
Если  объявляет переменную экземпляра  который имеет модификатор 
и начальное выражение,
Неявно индуцированный геттер  является
позднеинициализированный геттер.
Это определяет семантику призыва.


%
{Поздняя, неинициализированная переменная экземпляра}
Если  объявляет переменную экземпляра  который имеет модификатор 
не имеет начального выражения,
Неявно индуцированный геттер  является
поздне-неинициализированный геттер.
Это определяет семантику призыва.

{%
В этом случае это возможно для  быть несвязанным,
Но есть также несколько способов связать  Для объекта:
Конструктор может иметь инициализацию
()
или запись в списке инициализаторов
()
Это позволит инициализировать ,
и неявно индуцированный сеттер по имени  мог бы
Завершено и выполнено нормально.%
?


%
 Статическая или библиотечная переменная
Если  объявляет статическую или библиотечную переменную,
Неявно индуцированный геттер  выполняется следующим образом:


  Непостоянная переменная с инициализатором.
В случае, когда  имеет инициализирующее выражение и не является постоянным;
Неявно индуцированный геттер  Это позднеинициализированный геттер.
Это определяет семантику призыва.{%
Обратите внимание, что эти статические или библиотечные переменные могут быть 
поздней инициализации в том смысле, что они не
Модификатор .%
?
  Постоянная переменная.
Если  объявляет постоянную переменную с инициализацией выражения ,
Результатом выполнения неявно индуцированного геттера является
значение постоянного выражения .
{%
Обратите внимание, что постоянное выражение не может зависеть от самого себя.
Так что циклических ссылок быть не может.%
?
  Переменная без инициализатора.
Если  Объявляет переменную  без инициализации выражения
и не имеет модификатора ,
a вызов неявно индуцированного геттера  оценивать на
Объект, который  Это неизбежно.

{%
В этом случае переменная всегда связана с объектом.
Это может быть нулевой объект.
что является начальным значением некоторых переменных объявлений
В данном случае.%
?

Если  объявляет переменную  без инициализации выражения
и имеет модификатор ,
неявно индуцированный геттер является поздне-неинициализированным геттером.
Это определяет семантику призыва.


Уменьшите белое пространство после подробного списка: Это всего лишь конечный символ.
\vspace{\baselineskip}\EndCase


 Функции.


%
Функции абстрагируются над выполняемыми действиями.


<функциональная подпись> ::= \gnewline{}
<type>? <identifier> <formalParameterPart>

<formalParameterPart> ::= <typeParameters>? <formalParameterList>

<functionBody> ::= ? <=> <выражение>> "
 (\ASYNC{} '*'? | \SYNC{} '*')) <блокировать>

<block> ::= '{' <statements> '' "


%
Функции могут быть введены посредством деклараций функций.
(),
Методические декларации (, ,
Геттерские декларации (),
устанавливающие декларации (),
и конструкторские декларации ();
Они могут быть введены буквальными функциями ().

%
Функция {asynchronous}{function!asynchronous}
если его тело обозначено {} или  Модификатор.
В противном случае функция {синхронная}{функция!синхронная}.
Функция - это {генератор}{функция!генератор}
если его тело обозначено  или  Модификатор.
Более подробная информация об этих концепциях приводится ниже.

{%
Является ли функция синхронной или асинхронной, является ортогональной к
Является он генератором или нет.
Генераторные функции — это сахар для функций
которые производятся систематически,
при помощи ленивой функции, которая {генерирует}
отдельные элементы коллекции.
Дарт обеспечивает такой сахар в обоих синхронных случаях,
Когда человек возвращается,
и в асинхронном случае, когда один возвращает поток.
Dart также позволяет выполнять синхронные и асинхронные функции.
Производит однозначную величину.%
?

%
Каждое заявление, которое вводит функцию, имеет подпись, которая указывает
его тип возврата, название и формальная параметрическая часть,
При этом возврат может быть невозможен,
Геттеры никогда не имеют формального параметра.
Функциональные буквы имеют формальную часть параметра, но не имеют типа возврата и имени.
Формальный параметр необязательно указывает
формальный список параметров типа функции,
Он всегда указывает свой официальный список параметров.
Функциональным органом является либо:



Заявление блока () содержащий
Заявления () выполняется функцией,
необязательно помеченный одним из модификаторов:
,  или .
%Если не известно статически, что тело функции
не может нормально
{%
(то есть он не может дойти до конца и «пройти сквозь»);
cf.~\ref{statementCompletion}%
?
Ошибка времени компиляции, если
Добавление  В конце тела
Это будет ошибка времени компиляции.
<{%
Например, это ошибка, если
Тип возврата синхронной функции ,
Тело может нормально функционировать.
Точные правила приведены в разделе ~.%
?

{%
Дарт поддерживает динамические функции,
Мы не можем гарантировать, что функция не возвращает объект.
не будет использоваться в контексте выражения.
Поэтому каждая функция должна либо бросать, либо возвращать объект.
Функциональное тело, которое заканчивается без броска или возврата
заставит функцию вернуть нулевой объект (),
Как и в случае с {} без выражения.
Для функций генератора ситуация более тонкая.
Подробнее см. в разделе ~\ref{return}.%
?

или

Форма  или форму \code{\ASYNC{} => $e$},
которые возвращают значение выражения  как если бы а
.
{%
Другие модификаторы здесь не применяются.
Поскольку они применяются только к генераторам, описанным ниже.
Генераторам не разрешается явно возвращать что-либо.
объекты добавляются в генерируемый поток или итерируемые с использованием
\YIELD{{} или .%
?
 быть заявленным типом возврата функции, которая имеет это тело.
Это ошибка времени компиляции, если выполняется одно из следующих условий:


 Функция синхронная,  не является ,
Это была бы ошибка компиляции времени.
Объявить функцию с телом
 \RETURN{} ; 
Вместо «P000541».
{%
В частности,  может иметь тип {any}
когда тип возврата .%
?
{%
Это позволяет сжатые декларации функций {}.
Легко понять такую функцию,
потому что тип возврата текстуально близок к возвращенному выражению .
Напротив, \code{\RETURN{} ;} в корпусе блока разрешается только
 с одним из нескольких конкретных статических типов,
менее вероятно, что разработчик понимает
что возвращенный объект не будет использоваться
()%
?
 Функция асинхронна, \flatten{T} не является \VOID,
Это была бы ошибка компиляции времени.
Объявить функцию с телом
\code{{)  < ; 
Вместо {\ASYNC{} .
{%
В частности,  может иметь тип 
когда сплющенный тип возврата ,%
?
{%
и обоснование аналогично синхронному случаю.%
?



%
Это ошибка времени компиляции
Если ,  или  Модификатор прилагается к
Тело сеттера или конструктора.

{%
Асинхронный сеттер будет мало полезен,
Поскольку сеттеры могут использоваться только в контексте назначения
(),
и выражение назначения всегда оценивает
значение правой стороны задания.
Если сеттер действительно выполнял свою работу асинхронно,
Можно было бы представить, что кто-то вернет будущее, которое решило вернуться к жизни.
Правая сторона задания после того, как сеттер выполнил свою работу.

Асинхронный конструктор по определению никогда не вернется.Пример класса, который он намеревается построить,
Вместо этого они возвращают будущее.
Назвать такого зверя через \NEW{} Это было бы очень запутанно.
Если вам нужно произвести объект асинхронно, используйте метод.

Можно было бы разрешить модификаторы для заводов.
Производитель:  может быть изменено на ,
Производитель:  может быть изменено на ,
и фабрика для  Он может быть изменен на .
Никакой другой сценарий не имеет смысла, потому что
Объект, возвращенный заводом, будет неправильного типа.
Эта ситуация очень необычна, поэтому не стоит делать исключение.
к общему правилу для конструкторов, чтобы позволить это.%
?

%
Ошибка времени компиляции, если заявленный тип возврата
Функция, отмеченная {} не является
Супертип  для некоторых типов .
Ошибка времени компиляции, если заявленный тип возврата
Функция с пометкой  не является
Супертип  для некоторых типов .
Ошибка времени компиляции, если заявленный тип возврата
Функция с пометкой  не является
Супертип  для некоторых типов .
Ошибка времени компиляции, если заявленный тип возврата
Функция, отмеченная  или , является .

%
Мы определяем

тип  следующим образом:
Если  имеет форму  или форму \code{FutureOr<>}
Тогда безсоюзный тип, полученный из  это
Безсоюзный тип, полученный из .
В противном случае, безсоюзный тип, полученный из , является .
{%
Например, свободный от союза тип, полученный из
 .%
?

%
Мы определяем
{тип элемента функции генератора}{%
Функция! Генератор! Тип элемента
 следующим образом:
%
 быть свободным от союза типом, полученным из объявленного типа возврата .
%
Если  является синхронным генератором и
 реализация  
()
Тип элемента  является .
%
Если  является асинхронным генератором и
 реализация  
Тогда тип элемента  является .
%
В противном случае, если  является генератором (синхронным или асинхронным)
 Это супертип 
{, которое включает 
Тип элемента  является .
 Дальнейшие случаи невозможны.


 Декларация функций Срок


%
 Это функция, которая
Он не является ни членом класса, ни функцией в буквальном смысле.
Декларация функций включает в себя следующее:
{library functions}{function!library}
которые являются функциональными декларациями
% (включая геттеров и сеттеров)
на верхнем уровне библиотеки и
{местные функции}{функция!местные}
которые являются функциональными декларациями, объявленными внутри других функций.
Библиотечные функции часто называют просто функциями верхнего уровня.

%
Декларация функции состоит из идентификатора, указывающего имя функции,
Возможно, предваряется возвратным типом.
Название функции сопровождается подписью и телом.
Для геттеров подпись пуста.
Тело пусто для внешних функций.

%
Сфера библиотечной функции — это сфера охвата библиотеки.
Описан объем локальной функции
В разделе .
В обоих случаях название функции находится в сфере действия.
в своей формальной области действия
(P000426).

%
Ошибка времени компиляции для предисловия декларации функции
со встроенным идентификатором .

%
Когда мы говорим, что функция   в другую функцию ,
Мы имеем в виду, что вызов  Причин, по которым выполняется с теми же аргументами и/или приемником, что и ,
и возвращает результат выполнения  абоненту ,
Если только  Исключение бросает,
В этом случае  И делает то же исключение.
Мы используем только термин для
Синтетические функции, введенные спецификацией.


 Формальные параметры


%
Каждая декларация функции non-getter включает
,
Состоит из списка необходимых позиционных параметров.
(),
за любыми дополнительными параметрами (P000428).
Дополнительные параметры могут быть определены либо как
набор названных параметров или как список позиционных параметров;
Но не оба.

%
Некоторые функции включают в себя
 (),
В этом случае мы говорим, что это
{общая функция}{функция!общая}.
A {негенерическая функция}{функция!негенерическая}
Это функция, которая не является общей.

%
 в функциональной декларации
Формальный список параметров типа, если таковой имеется, и формальный список параметров.

{%
Следующие виды функций не могут быть общими:
Геттеры, сеттеры, операторы и конструкторы.%
?

%
Формальный список параметров типа декларации функций вводит
Новая область, известная как функция
{Сфера действия параметра типа}{Scope!type parameter}.
Сфера действия параметра типа общей функции  Он заключен в
Сфера, в которой объявлен .
Каждый формальный параметр типа вводит тип в область параметров типа.

%
Если он существует, то размер параметра типа функции  это
текущий объем подписи ,
и для самого формального перечня параметров типа;
В противном случае область, в которой  Объявлено, что
существующий объем подписи .

{%
Это означает, что формальные параметры типа находятся в области применения.
в границах деклараций параметров,
позволяет использовать так называемые F-параметры типа


\code{класс C<X\EXTENDS{} Сравнительный <X>{}  \ldots\ \},


и формальные параметры типа находятся в области действия друг для друга,
Разрешение зависимостей типа
\code{класс D<X\EXTENDS{} Y >  \ldots\ \}.%
?

%
Формальный список параметров декларации функций
Новая область, известная как функция
{формальный параметр}{схема!формальный параметр}.
Формальный диапазон параметров негенерической функции  Он заключен в
Сфера, в которой объявлен .
Формальный параметр общей функции  Он заключен в
Сфера действия параметра типа .
Каждый формальный параметр вводит локальную переменную в
Формальный диапазон параметров.
Текущий объем подписи функции является
Сфера охвата, которая охватывает формальный диапазон параметров.

{%
Это означает, что в общей функциональной декларации
Тип возврата и аннотации типа параметра могут использовать
формальные параметры типа,
но формальные параметры не входят в объем подписи.%
?

%
Тело функциональной декларации вводит
Новая область, известная как функция
{область тела}{сфера!функциональное тело}.
Объем тела функции  включено в объем, введенный
Формальный диапазон параметров .

%
Ошибка времени компиляции, если формальный параметр
Объявляется постоянной переменной ().


<formalParameterList> ::= ''' "
 <<нормальные Формальные Параметры>> "
 <(<normalFormalParameters>>, <optionalOrNamedFormalParameters>> "
 <("Необязательные Формальные Параметры") "

<нормальные формальные параметры> ::= <<normalFormalParameter> (>normalFormalParameter>)

<optionalOrNamedFormalParameters> ::= <optionalPositionalFormalParameters>
 <namedFormalParameters>

<optionalPotionalFormalParameters> ::= <
<defaultFormalParameter> ('', <defaultFormalParameter>)* ''?'' "

<namedFormalParameters> ::= <
'{' <defaultNamedParameter> (',' <defaultNamedParameter>)* ','?'' "


%
Формальные списки параметров позволяют необязательно запятую
после последнего параметра (\syntax{','?}).
Список параметров с такой задней запятой
Во всех отношениях эквивалентен одному и тому же списку параметров без запятой.
Все списки параметров в этой спецификации показаны без задней запятой,
Правила и семантика в равной степени применимы к
соответствующий список параметров с задней запятой.


 Необходимые формы


%
 может быть определена одним из трех способов:



С помощью подписи функции, которая именует параметр и
Описывает его тип как тип функции ().
Это ошибка времени компиляции, если указаны значения по умолчанию
в подписи такого функционального типа.

В качестве инициализации формальной, которая действительна только как параметр
Генеративный конструктор ().

Обычная переменная декларация ().



<normalFormalParameter> ::= <
<metadata> <normalFormalParameterNoMetadata>

<normalFormalParameterNoMetadata> ::= <functionFormalParameter>
 <fieldFormalParameter>
 <Простой Формальный Параметр>

<functionFormalParameter> ::= {}
? <type>? <identifier> <formalParameterPart> '?'?

<simpleFormalParameter>:= <декларированный идентификатор>
  <идентификатор>

<declaredIdentifier> ::= ? <finalConstVarOrType> <identifier>

<fieldFormalParameter> ::= \gnewline{}
<finalConstVarOrType> \THIS{}<identifier> (<formalParameterPart> '?'?)?


%
Это ошибка времени компиляции, если формальный параметр имеет модификатор 
Модификатор .
Ошибка времени компиляции  происходит как
Первый токен {fieldFormalParameter}.

%
Ошибка времени компиляции, если параметр
<<Формальное поле Параметр возникает как параметр функции
который не является неперенаправляющим генеративным конструктором.

{%
A {fieldFormalParameter} объявляет инициализацию формальной,
который описан в другом месте
()%
?

%
Можно включить модификатор \COVARIANT{}
в некоторых формах декларирования параметров.
Эффект от этого описан в отдельном разделе.
(P000458).

<%
Обратите внимание, что также используется не терминальный {нормальный формальный параметр}
в правилах грамматики для факультативных параметров,
Это означает, что такие параметры также могут быть ковариантными.
?

%
Это ошибка времени компиляции, если модификатор  происходит
параметр функции, который не является
Способ экземпляра, установщик экземпляра или оператор экземпляра.


 Факультативные формы


%
Дополнительные параметры могут быть указаны и снабжены значениями по умолчанию.


<defaultFormalParameter> ::= <normalFormalParameter> ('=' <expression>)

<defaultNamedParameter> ::= 
<metadata> ? <normalFormalParameterNoMetadata>
\gnewline{}('=' | ':') <выражение>)?


Форма {<нормальный Формальный Параметр>:> <выражение>}
Это эквивалентно форме
{<normalFormalParameter> '=' <expression>}.
Колонно-синтакс включен только для обратной совместимости.
Он обесценивается и удаляется в
Более поздняя версия спецификации языка.

%
Это ошибка времени компиляции, если значение по умолчанию необязательного параметра
Не является постоянным выражением (P000459).
Если по умолчанию явно не указан необязательный параметр
Предусмотрен неявный дефолт {}.

%
Ошибка времени компиляции, если имя названного опционального параметра
Начинается с символа «_».

<%
Необходимость такого ограничения является прямым следствием
Дело в том, что именование и приватность не являются ортогональными.
Если мы позволим названным параметрам начаться с подчеркивания,
Они считаются частными и недоступными для
Абоненты из-за пределов библиотеки, где она была определена.
Если метод вне библиотеки переопределяется
метод с частным факультативным названием,
Это не будет подтипом оригинального метода.
Статическая шашка, конечно, будет отмечать такие ситуации.
Но последствием будет то, что добавление частного, названного формальным, сломается.
клиенты за пределами библиотеки таким образом, что они не могли легко исправить.
?


{Ковариантные параметры}


%
Дарт допускает формальные параметры примеров,
В том числе сеттеры и операторы,
Заявлено .

{%
Синтаксис для этого указан в предыдущем разделе.
()%
?

%
Ошибка времени компиляции, если модификатор  происходит
в декларации формального параметра функции
который не является методом экземпляра, множителем экземпляров или оператором.

{%
Как указано ниже, параметр также может быть ковариантным по другим причинам.
Общий эффект наличия ковариантного параметра 
в подписи данного метода 
- позволяет ковариантно перекрывать тип ,
Это означает, что тип, требуемый во время выполнения для данного фактического аргумента.
может быть подходящим подтипом типа, который известен во время компиляции.
на сайте call.%
?

{%
Этот механизм позволяет разработчикам явно запрашивать
Гарантия времени компиляции, которая в противном случае поддерживается(а именно: что фактический аргумент, статический тип которого удовлетворяет требованию)
Это также будет сделано во время выполнения
Заменяются динамическими проверками.
В обмен на принятие этих динамических проверок типа,
Разработчики могут использовать ковариантные параметры для выражения дизайна программного обеспечения.
где динамические проверки типа известны (или, по крайней мере, доверены) для успеха;
на основании рассуждений о том, что анализ статического типа не захватывает.%
?

%
%
 Сигнатура метода с формальными параметрами типа
 X}{1}{}
Формальные позиционные параметры {p}{1}{n}
%
именованные формальные параметры {q}{1}{k}.
%
 быть подписью метода с формальными параметрами типа
 XX \!}{1}{}
Формальные позиционные параметры <p> '\!}{1}{n},
%
именованные формальные параметры < '\!}{1}{k'}.
%
Предположим, что  и ;
Мы говорим:  является параметром в  это
{соответствует параметру}
Формальный параметр  в .
Предположим, что  и ;
Мы говорим, что  является параметром в  это
{соответствует} формальному параметру
 в , если .
%
Точно так же мы говорим, что формальный параметр типа
 
{соответствует} формальному параметру типа
 от , для всех .

{%
Сюда относится случай, когда  соответственно  иметь
необязательные позиционные параметры,
в каком случае  соответственно  Должен держать,
Но мы можем иметь .
Случай, когда числа формальных параметров типа отличаются, не имеет значения.
?

% Быть ковариантным — свойство параметра интерфейса класса;
% это означает, что мы говорим только о исходном ключевом слове 
% и класс, содержащий соответствующую декларацию, когда мы обнаруживаем
% первый раз, когда данный параметр является ковариантным. С этого момента и на
% оно «переносится» по ссылкам подтипа, связанным с интерфейсами классов;
% так, что мы можем получить его индуктивно из непрямого суперинтерфейса
% путем проверки того, имеют ли прямые суперинтерфейсы сигнатуру метода
% с соответствующим именем и соответствующим параметром, а затем проверка
% от этого параметра. Такой же подход применим и к ковариантам по классам.

%
%
 Это класс, который объявляет
Метод  который имеет
параметр  заявление которого имеет модификатор ;
В этом случае мы говорим, что параметр  это
{covariant-by-declaration}{parameter!covariant-by-declaration}.
%
В этом случае интерфейс  имеет сигнатуру метода ,
и эта подпись имеет параметр ;
Мы также говорим, что параметр  В этом методе подпись
{covariant-by-declaration}
%
Наконец, параметр  сигнатуры метода 
интерфейс класса  это
{ковариант по декларации}
если прямой суперинтерфейс 
имеет доступную сигнатуру метода  с таким же названием, как ,
который имеет параметр  что соответствует ,
 является ковариантным по декларированию.

%
Предположим, что  Является общим классом
с официальными декларациями параметров типа
,
%
 быть декларацией метода экземпляра в 
(который может быть методом, установщиком или оператором),
 быть параметром, объявленным , и
Пусть  будет объявленным типом .
%
Параметр  это{ковариант по классу}{параметр!ковариант по классу}
Если для любого ,
 происходит в ковариантном или инвариантном положении в .
%
В этом случае интерфейс  также имеет сигнатуру метода ,
и эта подпись имеет параметр ;
Мы также говорим, что параметр  В этом методе подпись
{ковариант по классу}.
Наконец, параметр  сигнатура метода 
интерфейс класса  это
{ковариант по классу}
если прямой суперинтерфейс 
имеет доступную сигнатуру метода  с таким же названием, как ,
который имеет параметр  что соответствует ,
 является ковариантным по классу.

%
Формальный параметр  это
{ковариант}{параметр!ковариант}
Если  является ковариантным по декларированию, или  является ковариантным по классу.

{%
Параметр может быть одновременно
Ковариант по классу и ковариант по классу.
Обратите внимание, что параметр может быть
ковариант по классу или ковариант по классу
на основании заявления в любом прямом или косвенном суперинтерфейсе,
В том числе любой суперкласс:
Приведенные выше определения распространяют эти свойства.
Интерфейс от каждого из его прямых суперинтерфейсов
Но они в свою очередь получат
свойство от их прямых суперинтерфейсов,
и так далее.%
?


 Тип функции.


%
Этот раздел определяет статический тип, который приписывается
функция, обозначенная декларацией функции,
и динамический тип соответствующего функционального объекта.

%
В данной спецификации,
обозначение, используемое для обозначения типа функции,
То есть, ,
следует синтаксису языка,
За исключением того, что 
.
Это означает, что каждый тип функции имеет одну из форм.
 T_0
 T_0


где  является возвратным типом,
 Формальные параметры типа с границами , ,
 Формальные типы параметров для ,
 Названия параметров для .
Типы негенерических функций охватываются случаем ,
где перечень типовых параметров
\code{\ldots{}>>
В целом опущен.
%
Аналогичным образом факультативные скобки  и \{\} опущены
когда нет дополнительных параметров.

% Мы обещаем, что обе формы всегда получают одинаковую обработку для k=0.
{%
Обе формы с опционами охватывают типы функций без опционов.
когда 
Каждое правило данной спецификации таково, что
Любая из двух форм может использоваться без двусмысленности.
Определить лечение типов функций без опций.%
?

%
Если декларация функции не объявляет тип возврата явно,
Тип возврата: \DYNAMIC{} (),
Если это не конструктор,
В этом случае он не считается возвратным,
или является сеттером или оператором ,
В этом случае тип возврата составляет .

%
Декларация функции может объявлять формальные параметры типа.
Тип функции включает имена параметров типа.
и для каждого параметра типа верхней границы,
который считается встроенным классом 
если не указана граница.
При последовательном переименовании параметров типа
может сделать два типа функций идентичными;
Они считаются одного и того же типа.

{%
Удобно включать формальные имена параметров типа в типы функций.
Они необходимы для того, чтобы выразить такие вещи, как отношения между
различные параметры типа, F-связи и типы формальных параметров.Однако мы не хотим различать два типа функций.
если они имеют одинаковую структуру и отличаются только выбором названий.
Эта обработка имен также известна как альфа-эквивалентность.
?

%%
В следующих трех пунктах,
Если число {s} формальных параметров типа равно нулю, то
Перечень параметров типа в типе функции опущен.

%
 Быть функцией с
Типовые параметры ,
Требуемые формальные типы параметров  T}{1}{n}
Тип возврата ,
без дополнительных параметров.
Статический тип  это
{T_0}.

%
 Быть функцией с
Типовые параметры ,
Требуемые формальные типы параметров  T}{1}{n}
Тип возврата 
и позиционные факультативные типы параметров {T}{n+1}{n+k}.
Тогда статический тип  это
{T_0}.

%
 Быть функцией с
Типовые параметры ,
Требуемые формальные типы параметров  T}{1}{n}
Тип возврата ,
именованные параметры {T}{x}{n+1}{n+k}
где   может иметь или не иметь значение по умолчанию.
Статический тип  это
{T_0}.

%
 Статический тип функциональной декларации .
 быть типом времени выполнения объекта функции  полученный
функция клозурита
()
или пример метода клозурирования
()
Применительно к ,
 быть фактическим типом, соответствующим 
В случае, когда  был создан
().
<<<<<<<<<<<<<<<<<<<<<<<<<<<<<<<<<<<<<<<<<<<<<<<<<<<<<<<<<<<<<<<<<<<<<<<<<<<<<<<<<<<<<<<<<<<<<<<<<<<<<<<<<<<<<<<<<<<<<<<<<<<<<<<<<<<<<<<<<<<<<<<<<<<<<<<<<<<<<<<<<<<<<<<<<<<<<<<<<<<<<<<<<<<<<<<<<<<<<<<<<<<<<<<<<<<<<<<<<<<<<<<<<<<<<<<<<<<<<<<<<<<<<<<<<<<<<<<<
 Может содержать переменные свободного типа, но  Содержит фактические значения.%
?
После этого следует придерживаться следующего:
 - класс, реализующий встроенный класс ;
 является подтипом ;
 не является подтипом какого-либо функционального типа, который является соответствующим подтипом .
{%
Если бы мы пропустили последнее требование,
 < int\,\FUNCTION[int]]
можно оценить по  с заявлением
\code{{} f()\{\}},
Что, очевидно, не является целью.%
?

{%
Реализация зависит от выбора
соответствующее представление для функциональных объектов.
Например, считать, что
Функционный объект, полученный посредством извлечения свойств
рассматривает равенство иначе, чем другие функциональные объекты;
Скорее всего, это другой класс.
Реализации могут также использовать различные классы для функциональных объектов.
В зависимости от типа и типа.
Арити может косвенно зависеть от того, является ли функция
метод экземпляра (с неявным параметром приемника) или нет.
Вариации многообразны и, например,
Нельзя предположить, что два различных объекта функции
обязательно будет иметь тот же тип времени выполнения.%
?


 Внешние функции.


%
{внешняя функция}{функция!внешняя функция}
является функцией, орган которой предоставляется отдельно от его декларации.
Внешняя функция может быть
Функция верхнего уровня (),
метод (, ,
геттер (),
Сеттер (),
Неперенаправляемый конструктор
(, .
Внешние функции вводятся через встроенный идентификатор. <
()
Затем следует подпись функции.

<<<<<<<<<<<<<<<<<<<<<<<<<<<<<<<<<<<<<<<<<<<<<<<<<<<<<<<<<<<<<<<<<<<<<<<<<<<<<<<<<<<<<<<<<<<<<<<<<<<<<<<<<<<<<<<<<<<<<<<<<<<<<<<<<<<<<<<<<<<<<<<<<<<<<<<<<<<<<<<<<<<<<<<<<<<<<<<<<<<<<<<<<<<<<<<<<<<<<<<<<<<<<<<<<<<<<<<<<<<<<<<<<<<<<<<<<<<<<<<<<<<<<<<<
Внешние функции позволяют вводить информацию о типе кода
что не известно компилятору Dart.%
?

{%
Примерами внешних функций могут быть внешние функции.
(определяется в C, Javascript и т.д.),
примитивы реализации (как определено системой времени выполнения Dart),
код, который генерируется динамическиИнтерфейс которого статически известен.
Однако абстрактный метод отличается от внешней функции.
как имеет {no} тело.%
?

%
Внешняя функция связана с телом посредством
Конкретный механизм реализации.
Попытка вызвать внешнюю функцию
которая не связана со своим телом.
Выбрасывает  или какой-либо его подкласс.

%
Может быть повышена конкретная ошибка компиляции
в объявлении функции {}.

<<%
Такие ошибки указывают на то, что
Каждый вызов этой функции будет бросать, например,
потому что известно, что он не будет связан с телом.
?

%
Настоящий синтаксис приведен в
разделы  и  ниже.


 Классы


%
 определяет форму и поведение набора объектов, которые являются его
\IndexCustom{instances}{instance}.
Классы могут быть определены классовыми декларациями, как описано ниже.
или через приложения для смешивания ().


<classDeclaration> ::=
  >> <typeIdentifier> <typeParameters>
{} <суперкласс>? <Интерфейсы>.
\gnewline{} '{' (<metadata> <classMemberDeclaration>)* '} "
  \CLASS{{}<mixinApplication Класс

<typeNotVoidList> ::= <typeNotVoid> (',' <typeNotVoid>)*

<classMemberDeclaration> ::= <declaration>; "
 <методическая подпись> <Функциональное тело>

<methodSignature> ::= <constructorSignature> <initializers>
 <factoryConstructorSignature>
  <Функциональная подпись>
  <getterSignature>
  <setter подпись>
 <подпись оператора>

<декларация> ::= \EXTERNAL{{}<factoryConstructorSignature>
 \EXTERNAL{{}<constantConstructorSignature>
 {}<ConstructorSignature>
 ({) ? <getterSignature>
 (\EXTERNAL) ? <setter подпись>
 (\EXTERNAL) ? <Функциональная подпись>
 ? <operatorSignature>
 < <type>? <staticFinalDeclarationList>
 \STATIC{} <type>? <staticFinalDeclarationList>
 < < {} <type>? <инициализированный ИдентификаторСписок
   <varOrType> <инициализированный ИдентификаторСписок
 < < <type>? <identifierList>
 <  <varOrType> <инициализированный ИдентификаторСписок
   >>{} <type>? <инициализированный ИдентификаторСписок
  <varOrType> <инициализированный ИдентификаторСписок
 <redirectoryConstructorSignature>
 <constantConstructorSignature> (<redirection> | <initializers>)?
 <constructorSignature> (<redirection> | <initializers>)?

<staticFinalDeclarationList>:= <
<staticFinalDeclaration> (<staticFinalDeclaration>) *

<staticFinalDeclaration> ::= <identifier> <expression>


%
Можно включить модификатор \COVARIANT{}
в некоторых формах декларации.
Эффект от этого описан в другом месте.
(P000476).

%
Класс имеет конструкторов, экземпляров и статических членов.
{члены организации}{члены организации!) класса
являются его методы экземпляра, геттеры, множители и переменные экземпляра.
{статические члены}{члены!статические} класса
Это его статические методы, геттеры, сеттеры и переменные.
{члены}{члены} класса
являются его статическими и экземплярными членами.

%
Классовая декларация вводит две сферы:


A {type-parameter scope}{scope!type parameter},
который является пустым, если класс не является общим ().
Включающая сфера охвата типового параметра декларации класса является
библиотечный объем текущей библиотеки.

A {область тела}{сфера действия} Классное тело.
Включающая сфера охвата органа сфера охвата классовой декларации является
Сфера охвата параметра типа декларации класса.


%
текущая сфера действия декларации члена инстанции,
Статическая декларация,
или декларация застройщика является
Область применения класса, в котором он объявлен.

%
Нынешний случай
{(и, следовательно, его члены)}
Доступ можно получить только в определенных местах в классе:
Мы говорим, что местоположение 
{имеет доступ к }{имеет доступ к этому@имеет доступ к }
f  находится внутри тела декларации
элемент экземпляра или генеративный конструктор,
или в начальном выражении  Пример переменной декларации.

{%
Обратите внимание, что начальное выражение для не-{} переменная переменная
не имеет доступа к ,
и ни одна из частей декларации с пометкой .%
?

%
Каждый класс имеет один суперкласс.
кроме классов  Который не имеет суперкласса.
Класс может реализовывать несколько интерфейсов.
Объявив их в своем имплементационном пункте ().

%
{абстрактная декларация класса}{классовая декларация}! абстрактный
Классовая декларация, которая явно декларируется
{} Модификатор.
A {конкретная декларация класса}{классовая декларация}! конкретный
Классовая декларация, которая не является абстрактной.
{абстрактный класс}{класс!абстракт} - это класс
заявление которого является абстрактным, и
а {конкретный класс}{класс!конкретный} - это класс
Чья декларация конкретна.

{%
Мы хотим различного поведения для конкретных классов и абстрактных классов.
Если  Он должен быть абстрактным,
Мы хотим, чтобы статическая шашка предупреждала о любой попытке инстанцировать ,
Мы не хотим, чтобы чекеры жаловались на нереализованные методы в P000121.
В отличие от этого, если  Он должен быть конкретным,
Проверяющий должен предупредить обо всех нереализованных методах,
Позволяет клиентам делать это свободно.%
?

{%
Интерфейс класса  это
неявный интерфейс, который объявляет подписи членов экземпляра
которые соответствуют членам экземпляра, заявленным ,
и чьи прямые суперинтерфейсы
Прямые суперинтерфейсы 
(, %
?

%
Когда имя класса появляется как тип,
Это имя обозначает интерфейс класса.

%
Это ошибка времени компиляции, если класс называется  объявлять
член с базовым именем () .
Если общий класс называется  объявляет переменную типа ,
Это ошибка времени компиляции
Если  равно ,
или если  имеет член, базовое имя которого ,
или если  Имеет конструктор под названием .

{%
Вот простые примеры, иллюстрирующие разницу между
"имеет члена" и "объявляет его членом".
Например,  \IndexCustom{declares}{declares member}
один член, названный ,
Но это {has}{has Member} два таких члена.
Правила наследования определяют, какие члены класса имеют.
?




 Полностью реализовать интерфейс


% Обратите внимание, что правила здесь и в  Они пересекаются, но они
% необходимое: Этот раздел посвящен конкретным методам, включая
% унаследованных и  Заинтересован в случае% членов, заявленных в , включая как конкретные, так и абстрактные.

%
% Использование "конкретного члена" ниже может показаться избыточным, поскольку класс
% не наследует абстрактных членов от своего суперкласса.
% подчеркивает тот факт, что даже при абстрактной декларации  это
%, заявленных в ,  Не имеет значения  которых здесь будет достаточно.
Конкретный класс должен полностью реализовать свой интерфейс.
%
 быть конкретным классом, объявленным в библиотеке , с интерфейсом .
Предположим, что  имеет подпись участника  Доступен по адресу .
Ошибка времени компиляции  не иметь
конкретный член с тем же названием, что и  Доступен по адресу ,
Если только  имеет нетривиальный 
(P000484).

%
Каждый конкретный член должен иметь соответствующую подпись:
Предположим, что  имеет конкретный член с
Оригинальное название:  Доступен по адресу ,
Пусть {m''} будет его подписью.
{%
Конкретный член может быть объявлен в  или унаследованы от суперкласса.
?
Пусть {m'} является подписью участника, полученной от 
Добавляя, если не присутствует, модификатор 
()
для каждого параметра   где
Соответствующий параметр в  Имеет модификатор .
Ошибка времени компиляции  не является правильным переопределением 
(),
Если этот конкретный элемент не является  экспедитор
(P000487).

{%
Рассмотрим конкретный класс ,
Предположим, что  Объявлять или наследовать
Реализация с таким же названием
для каждой подписи в интерфейсе.
Это все еще ошибка, если одна или несколько из этих реализаций
имеет такие параметры или типы, которые не удовлетворяют
подпись соответствующего члена в интерфейсе.
Для этой проверки не хватает  >> Модификаторы добавляются неявно
к подписи наследуемого члена (так мы получаем  от ).
При добавлении модификатора {} к одному или нескольким параметрам
(что произойдет только тогда, когда конкретный член наследуется),
реализация может выбрать неявно индуцировать метод пересылки
с такой же подписью, как ,
для выполнения требуемой динамической проверки типа,
Используйте унаследованный метод.%
?

%
Это конкретный выбор реализации, независимо от того, является ли он или нет.
используется неявно индуцированный метод пересылки
когда модификатор {} добавляется
один или несколько параметров в .

{%
Это верно, несмотря на то, что
Такие методы пересылки можно наблюдать.
Например, мы можем сравнить
время выполнения типа отрыва метода от
Получатель типа 
к типу времени выполнения отрыва суперметода от
местоположение в теле .%
?

%
с использованием или без использования метода,
Подпись участника в интерфейсе  является .

{%
Способ пересылки не изменяет интерфейс ,
Это деталь реализации.
В частности, это касается даже тех случаев, когда
четкое заявление о методе пересылки имело бы
Изменился интерфейс , потому что $m'$ Это подтип .

Когда класс имеет нетривиальный ,
класс может оставить некоторых членов нереализованными,
и классу разрешается иметь  экспедитор
который не удовлетворяет интерфейсу класса
(в этом случае он будет отменен)
Еще один экспедитор P000573.

Вот пример: %
?



%
Параметры, которые являются ковариантными по декларированиюдолжны также удовлетворять следующим ограничениям:
Предположим, что параметр   Имеет модификатор .
Допустим, что прямой или косвенный суперинтерфейс  иметь
Подпись метода  с таким же названием, как  Доступен по адресу ,
 имеет параметр  Это соответствует P000178.
В этой ситуации происходит ошибка компиляции времени.
Если тип  Это не подтип и не супертип типа .

{%
Это гарантирует, что унаследованный метод удовлетворяет тем же ограничениям.
для каждого формального параметра, который является ковариантным по декларированию
как ограничение, указанное для декларации в 
() %
?


 Методы инстанций


%
\IndexCustom{Методические методы}{метод!инстанция}
Функции ()
декларация которых непосредственно содержится в классовой декларации
И это не объявлено .
  есть
Методы экземпляра, объявленные 
и методы экземпляра, унаследованные  Из своего суперкласса
().

%
%
Рассмотрим класс 
и заявление члена экземпляра  в $C$, с подписью члена $m$
().
Ошибка времени компиляции  отменяет декларацию
% Обратите внимание, что  является доступной в силу определения "переопределений".
с подписью члена 
Прямой суперинтерфейс 
(),
Если  не является правильным членом, переопределяющим 
(P000493).

{%
Это не единственный вид конфликта, который может существовать:
Заявление члена суда  может противоречить другому заявлению ,
Даже в тех случаях, когда они не имеют одинаковых названий.
Или это не одно и то же заявление.
Например,  Это может быть пример Геттера и  статический сеттер
() %
?

%
Для каждого параметра $p$ $m$, где присутствует \COVARIANT{},
Ошибка времени компиляции при наличии
прямой или косвенный суперинтерфейс  который имеет
Подпись доступного метода  с таким же названием, как ,
 имеет параметр  что соответствует 
(),
Если только тип  Это подтип или супертип типа .

{%
Это означает, что
параметр, который является ковариантным по декларированию, может иметь тип
который является супертипом или подтипом типа
соответствующий параметр в суперинтерфейсе,
Но эти два типа не могут быть несвязанными.
Обратите внимание, что это требование должно быть выполнено
для каждого прямого или косвенного суперинтерфейса отдельно,
Потому что эти отношения не транзитивны.
?

{%
Суперинтерфейс может быть статически известным типом приемника.
Это означает, что мы расслабляем потенциальные отношения типизации.
между статически известным типом параметра и
Тип, который действительно требуется во время выполнения
отношения подтипа или супертипа,
Вместо строгих отношений супертипа
Это относится к параметру, который не является ковариантным.
Следует отметить, что статически он не известен.
на месте вызова, является ли любой данный параметр ковариантным,
Поскольку ковариация может быть введена в
надлежащий подтип статически известного типа приемника.
Мы решили отдать приоритет гибкости, а не безопасности здесь.
Потому что весь смысл ковариантных параметров в том, что разработчики
Вы можете сделать выбор, чтобы увеличить гибкость.
в компромиссе, где некоторая безопасность статического типа теряется.
?


 Операторы


%
{Операторы}{операторы} - это методы экземпляров со специальными именами,За исключением оператора \lit{[]] который является экземпляром Геттера
и оператор <{[]=}, который является набором экземпляров.


<оператор подпись> ::= \gnewline{}
<type>? \OPERATOR{} <operator> <formalParameterList>

<оператор>:=~ "
 <бинарный оператор>
 "
<[] "

<binaryOperator> ::= <multiplicativeOperator>
<<p001031> <additive-оператор>
 <shiftOperator>
 <relationOperator>
== "
 <bitwiseOperator>


%
Декларация оператора идентифицируется с помощью встроенного идентификатора.
()
.

%
Для операторов, определяемых пользователем, допускаются следующие имена:
<<,
\lit{>>,
<<==
{>==
 >>{=======================================================================================================================================================================================================================================================
{-}
{+}
{/}
/>
\lit{*},





,

\lit{[=],
<{[]],
.

%
Ошибка времени компиляции, если оператор объявлен пользователем
{[]=} не является 2.
Это ошибка времени компиляции, если активность оператора, объявленного пользователем.
с одним из имен:
<<,
\lit{>>,
<<==
{>==>
<\lit{======================================================================================================================================================================================================================================================
{-}
\lit{+}
\lit{/}
{/}
\lit{*},

{ |}}




\lit{\gtgt,
\lit{[]]
Это не 1.
Ошибка времени компиляции, если оператор объявлен пользователем
{-}
Это не 0 или 1.

{%
{-} Оператор уникален
Допускаются две перегруженные версии.
Если у оператора нет аргументов, он обозначает унарный минус.
Если у него есть аргумент, он обозначает двоичное вычитание.%
?

%
Имя унарного оператора \lit{-} — .

{%
Это устройство позволяет различать два метода.
для целей метода поиска, переопределения и отражения.%
?

%
Ошибка времени компиляции, если оператор объявлен пользователем

Это не 0.

%
Ошибка времени компиляции для объявления необязательного параметра в операторе.

%
Ошибка времени компиляции, если оператор объявлен пользователем \lit{[=]
Объявляет тип возврата, отличный от .

{%
Если для оператора, объявленного пользователем, не указан тип возврата
{[=],
Тип возврата - \VOID{} ()%
?

{%
Тип возврата {} Потому что
заявление о возврате в реализации оператора
\lit{[=]
Не возвращает объект.
Рассмотрим оценку выражения  в форме
,
и предположить, что оценка  Дает объекту .
 После этого будет оцениваться ,
и даже если исполняемый орган оператора
\lit{[=]
завершается с помощью объекта ,
То есть, если  Возвращается, его просто игнорируют.
Обоснование такого поведения заключается в том, что
Задания должны быть гарантированы для оценки назначенному объекту.%
?


 Метод 


%
Метод  Имплицитно упоминается во время исполнения
В ситуациях, когда один или несколько членов поиска терпят неудачу
(,
,
.

{%
Мы можем думать о P000578. как резервная копия
который вступает в силу при обращении члена  Попытка,но нет члена, названного ,
или существует,
Но данный призыв имеет форму списка аргументов.
не соответствует декларации 
(проведение меньшего количества позиционных аргументов, чем требуется, или больше, чем поддерживается)
или передавать именные аргументы с именами, не объявленными .
% Следующее предложение охватывает как функциональные объекты, так и экземпляры
% класс с методом, названным , потому что
% ошибка времени компиляции  с неправильно оформленным аргументом
Список %, если тип не является \DYNAMIC{} или \FUNCTION.
Это может произойти только для обычного вызова метода.
когда приемник имеет статический тип ,
или для выполнения функции, когда
Вызванная функция имеет статический тип {} или .
%
Метод  могут быть использованы и другими способами, например,
Его можно назвать как любой другой метод.
На него можно ссылаться из  экспедитор,
Как описано ниже.%
?

%
Мы говорим, что класс  \Index{имеет нетривиальный \code{noSuchMethod}}
Если  имеет конкретный член, названный 
который отличается от заявленного во встроенном классе .

{%
Обратите внимание, что это должен быть метод, который принимает один позиционный аргумент.
Чтобы правильно переопределить  в .
Например, он может иметь подпись.
 или
,
но не
.
Это означает, что ситуация, при которой  на него ссылаются
(явно или неявно)
С одним реальным аргументом нельзя не согласиться по той причине, что
«Такого метода нет»
Чтобы мы вошли в бесконечный цикл
Пытаясь вызвать .
Однако можно столкнуться с динамической ошибкой 
Во время ссылки на 
Поскольку фактический аргумент не удовлетворяет проверке типа,
Но эта ситуация приведет к динамическим ошибкам типа.
Вместо того, чтобы повторять попытку 
().
Вот пример, когда ошибка динамического типа возникает из-за
Предпринята попытка пройти 
где принимается только нулевой объект: %
?



%
%
 Будьте конкретным классом,
 быть библиотекой, содержащей декларацию ,
 Будь именем.
Тогда $m$ \Index{noSuchMethod forwarded} в $C$ если {f}
Одно из следующего верно:


  Запрошено в программе: }
 имеет нетривиальный ,
Интерфейс  содержит подпись члена  ,
 Нет конкретного члена, названного  Доступно для 
который правильно перекрывает 
{%
(то есть, ни один член не назван ) Объявляется или наследуется по адресу ,
или наследуется, но не имеет требуемой подписи.
}
В этом случае мы также говорим, что  пересылается noSuchMethod.

 Принуждение к конфиденциальности:
Существует прямой или косвенный суперинтерфейс
  Объявляется в библиотеке  Отличие от ,
Интерфейс  содержит подпись члена  
 Это частное имя,
Без суперкласса  иметь
Конкретный член, названный  Доступно для 
который правильно перекрывает .
В этом случае мы также говорим, что  пересылается noSuchMethod.


%
Для конкретного класса , a

неявно индуцируется для каждой подписи
Который не является ниспосланным.

%
Ошибка времени компиляции, если имя  NoSuchMethod скачать
в конкретном классе ,и суперкласса  имеет доступную конкретную декларацию 
который не является форвардом.

%
% % TODO(Eernst): Мы привыкли говорить «интерфейс  или ; но мы
% необходимо изменить «доступный» на схему, аналогичную слиянию имен
%, чтобы можно было точно сказать.
Экспедитор noSuchMethod является конкретным членом 
с подписью, взятой из интерфейса ,
и с одинаковым значением по умолчанию для каждого дополнительного параметра.
Это может быть вызвано в обычном призыве и в сверхвызове.
и когда  Это метод, который можно клозировать.
()
Использование извлечения имущества
().

{%
Ошибка, связанная с неявно вызванным экспедитором
Это отменяет написанную человеком декларацию.
Это может произойти только в том случае, если конкретное заявление не
правильно отменять . Рассмотрим следующий пример: %
?



<{%
В этом примере,
Реализация с подписью  Унаследовано от ,
и суперинтерфейс  объявлять
Подпись .
Это ошибка времени компиляции, потому что  не иметь
Реализация метода с подписью .
Мы не хотим косвенно побуждать
 экспедитор с подписью 
Потому что он будет перевешивать ,
Это, вероятно, будет очень запутанным для разработчиков.
%
В частности, это вызов, как 
призывать ,
Хотя это призыв, который является правильным для
Объявление  в .
%
Поэтому мы требуем от разработчиков четкого разрешения конфликта.
когда неявно индуцируется  экспедитор
будет отменять явно объявленное унаследованное осуществление.
%
Однако это не проблема,
Пусть будет  отклонить
 экспедитор,
В такой ситуации ошибки нет.%
?

{%
Это означает, что  У экспедитора такая же
свойства как явно заявленный конкретный член,
За исключением, конечно, того, что  экспедитор
не препятствует индуцированию себя или другого экспедитора noSuchMethod.
Мы не указываем тело  экспедитор,
Но он будет ссылаться на ,
и динамическая семантика ее выполнения указана ниже.%
?

{%
В начале этого раздела мы упоминали о неявных призывах.
 может произойти только
с приемником статического типа \DYNAMIC{}
или функция статического типа \DYNAMIC{} или .
 экспедитор,
 Также может быть использован
на приемнике, статический тип которого не является .
Аналогичной ситуации для функций не существует.
потому что невозможно вызвать  экспедитор
в класс функционального объекта.%
?

{%
Для конкретного класса ,
Мы можем думать о нетривиальном [P000619].
(объявлено или унаследовано )
в качестве просьбы о "автоматическом осуществлении" всех невыполненных членов
в интерфейсе  в качестве экспедиторов .
Кроме того, имеется неявная просьба о
Автоматическая реализация всех нереализованных
Недоступные представители любого конкретного класса,
Существует ли нетривиальный P000621?
Обратите внимание, что последний не может быть написан явно в Дарте.
потому что их имена недоступны;
Но язык все еще может указать, что они индуцированы неявно.
потому что компиляторы контролируют обработку личных имен.
?

%
%
Для динамической семантики,
Предположим, что класс  имеет неявно индуцированный
 экспедитор по имени ,
Формальные параметры типа
,
Формальные позиционные параметры
%
\commentary{(некоторые из которых могут быть необязательными, когда )},
именованные формальные параметры с именами

{(с приведенными выше значениями по умолчанию)}.

{%
Для этого не нужно различать
подпись, которая имеет необязательные позиционные параметры и
подпись, в которой указаны параметры,
потому что первое покрывается .%
?

%
Исполнение тела  создавать
экземпляр {im} предопределенного класса 
такой, что:


  Показатель  f  Это метод.
  Показатель  Если f  является геттером.
  Показатель {} f  Это сеттер.
  Оценивает символ .
  Оценка в неизменяемый список
чей динамический тип ,
содержат те же объекты, что и список, полученный в результате оценки
.
  Оценка на неизменяемую карту
чей динамический тип ,
с теми же ключами и значениями, что и карта, полученная в результате оценки

\code{<Symbol, Object>$\#x_1$: : .
  Оценить неизменяемый список
чей динамический тип ,
содержит те же объекты, что и список, полученный в результате оценки
.


%
Далее  Называется  Как реальный аргумент,
и полученный оттуда результат возвращается исполнением .

{%
Это обычный метод вызова 
().
То есть  экспедитор в классе  может вызвать
Реализация  Об этом говорится в
Подкласс .

Проверка динамического типа (DVD) - Проверка динамического типа (DVD) - Проверка динамического типа (DVD) - Проверка динамического типа (DVD) - Проверка динамического типа (DVD) - Проверка динамического типа (DVD) - Проверка динамического типа (DVD) - Проверка динамического типа (DVD) - Проверка динамического типа (DVD) - Проверка динамического типа (DVD) - Проверка динамического типа (DVD)
выполняется таким же образом, как и для вызова
Явно заявленный метод.
В частности, фактический аргумент перешел к ковариантному параметру.
Они будут проверяться динамически.

Кроме того, как и другие обычные методы,
Ошибка динамического типа, если результат возвращается
 Форвард имеет тип, который не является подтипом.
Тип возврата экспедитора.

Один особый случай, о котором следует знать, это когда экспедитор оторван.
А затем ссылается на фактический список аргументов, который не соответствует
Официальный список параметров.
В этой ситуации мы получим призыв

Вместо «P000646» в оригинальном приемнике,
потому что это вызов объекта функции
(и они не перекрывают ):%
?




 Оператор «==» и Примитивное равенство


%
Оператор {==} используется неявно в определенных ситуациях,
В частности, постоянные выражения
()
Это создает ограничения для данного оператора.
Аналогичная ситуация с геттером .
Для того, чтобы определить эти ограничения, как только мы вводим понятие
.

<<<p001145>{%
Известно, что некоторые постоянные выражения имеют значение
Равенство которых примитивно.
Это полезно знать, так как позволяет
Значение выражения равенства
и значение вызовов 
вычисляется во время компиляции.
В частности, его можно использовать для строительства
постоянные сборы во время компиляции,
Его можно использовать для проверки
Все элементы в постоянном наборе различны,
и все ключи в постоянной карте различны.%
?


 Нулевой объект имеет примитивное равенство.
().
 Каждый экземпляр типа ,  и 
примитивное равенство. Каждый экземпляр типа <
который был первоначально получен путем оценки буквального символа или
постоянное обращение к конструктору  класс
примитивное равенство.

%% TODO (Eernst): С Dart 2.15 мы должны упомянуть «Список» здесь,
% %, если «постоянный тип» не включает этот случай.
Каждый экземпляр типа 
которая была первоначально получена путем оценки постоянного типа
()
примитивное равенство.

 быть объектом, полученным путем оценки постоянного списка
(),
Постоянная карта буквально
(), или
постоянный набор буквальных
(),
 примитивное равенство.

% % TODO (Eernst): С Dart 2.15 рассмотрите «f<int>».
Объект функции, полученный путем клосуризации функции
Статический метод или функция верхнего уровня
()
как значение постоянного выражения
примитивное равенство.

Пример  примитивное равенство
Если динамический тип  класс ,
 примитивное равенство.

Класс  примитивное равенство.

Класс  примитивное равенство
если он не имеет реализации оператора 
который перевешивает унаследованный от ,
и не имеет реализации геттера 
Это отменяет унаследованное от P000659.


%
Когда мы говорим, что данный случай или класс
{не имеет примитивного равенства}{%
Примитивное равенство! не имеет,
Это означает, что это неправда, что упомянутый экземпляр или класс
примитивное равенство.


 Геттерс.


%
Геттеры — это функции (). которые используются
для получения значений свойств объекта.


<getterSignature> ::= <type>? \GET{} <идентификатор>


%
Если тип возврата не указан, тип возврата геттера .

%
Определение Геттера, прикрепленное к {} Модификатор определяет
Статический геттер.
В противном случае, он определяет пример геттера.
Имя геттера дается идентификатором в определении.

%
  есть
те экземпляры, которые объявлены ,
явно или косвенно,
и экземпляров, унаследованных  из своего суперкласса.
  есть
Эти статические геттеры объявлены .

{%
Декларация Геттера может противоречить другим декларациям.
().
В частности, геттер никогда не может переопределить метод.
и метод никогда не может переопределить геттер или переменную экземпляра.
Правила, когда геттер правильно перевешивает другого члена
Дано в другом месте
() %
?


 Сеттеры.


%
Сеттеры — это функции (). которые используются для установки
Значения свойств объекта.


<setterSignature> ::= <type> \SET{} <identifier> <formalParameterList>


{%
Если тип возврата не указан, тип возврата задаваемого устройства {}
()%
?

%
Определение сеттера, приставленное к {} Модификатор определяет
Статический сеттер.
В противном случае он определяет набор экземпляров.
Имя сеттера получается путем добавления строки «=» к
идентификатор, указанный в его подписи.

{%
Следовательно, имя сеттера никогда не может конфликтовать с, переопределять или переопределяться.
Геттер или метод.%
?

%
  есть
Настройщики экземпляров, объявленные 
явно или косвенно,
и наборы экземпляров, унаследованные  из своего суперкласса.
 естьЭти статические установки объявлены ,
либо неявно, либо явно.

%
Это ошибка времени компиляции, если формальный список параметров множителя
не состоит из одного необходимого формального параметра. .
{%
Это можно сделать с помощью грамматики,
Но в этом случае мы должны указать правила оценки.%
?

%
Это ошибка времени компиляции, если установщик объявляет тип возврата, отличный от .
Ошибка времени компиляции, если класс имеет
Сеттер по имени  Тип аргумента  и
Геттер по имени  с возвратным типом ,
 Не может быть назначено .

{%
Правила, когда сеттер правильно перевешивает другого участника
Дано в другом месте
(P000517).
Декларация учредителя может противоречить и другим декларациям.
() %
?


 Абстрактные участники


%
{абстрактный метод}{метод!абстракт}
(соответственно,
{abstract getter}{getter!abstract} или
{abstract setter}{setter!abstract}
является методом экземпляра, геттера или сеттера, который не объявлен 
и не обеспечивает реализацию.
A {конкретно} Метод {метод!конкретный}
(соответственно,
{конкретно} Геттер (getter!conrete)
\IndexCustom{concrete setter}{setter!concrete}
Это метод экземпляра, геттера или сеттера, который не является абстрактным.

{%
Абстрактные экземпляры полезны из-за их взаимодействия с классами.
Каждый класс Dart создает неявный интерфейс.
Dart не поддерживает явное определение интерфейсов.
Использование абстрактного класса вместо традиционного интерфейса
имеет важные преимущества.
Абстрактный класс может обеспечить реализацию по умолчанию.
Также можно использовать статические методы.
Устранение потребности в классах обслуживания
 или ,
цель которой состоит в том, чтобы группировать утилиты, относящиеся к данному типу.
?

{%
Призыв абстрактного метода, геттера или сеттера не может произойти.
Поиск () Никогда не уступать
Абстрактный член в результате.
Один из способов думать об этом — это
Абстрактное заявление члена в подклассе
не отменяет и не отменяет унаследованную реализацию.
Он служит только для указания подписи данного члена, который
каждый конкретный подтип должен иметь реализацию;
То есть он вносит свой вклад в интерфейс класса,
Не для самого класса.%
?

{%
Целью абстрактного метода является предоставление декларации.
Для таких целей, как проверка типа и отражение.
В рамках таких формулировок часто используются такие формулировки, которые
ожидаемое смешивание будет обеспечено суперклассом, к которому применяется смесь.%
?

{%
Мы хотим определить, объявляется ли конкретный класс абстрактными членами.
Тем не менее, код должен работать следующим образом: %
?



{%
Во время выполнения конкретный метод  Заявлено в 
будут казнены,
И никаких проблем не должно возникнуть.
Поэтому никакой ошибки не должно быть
если соответствующий конкретный член существует в иерархии.
?


 Переменные инстанций


%
{Вариабельные значения}{вариабельные значения!
переменные, декларации которых
немедленно содержатся в классовой декларации
И это не объявлено .
  есть
Переменные экземпляра, объявленные 
и переменные экземпляра, унаследованные  из своего суперкласса.

%
Это ошибка времени компиляции, если переменная экземпляра объявлена постоянной.

{%
Понятие постоянной переменной экземпляра
Это сложно и запутанно для программистов.
Переменная экземпляра предназначена для изменения в каждом экземпляре.Постоянная переменная экземпляра будет иметь одинаковое значение для всех экземпляров.
И как таковой уже является сомнительной идеей.

Язык может интерпретировать заявления с переменной константой как
Геттеры, возвращающие константу.
Однако постоянная переменная экземпляра не может рассматриваться как
Истинная постоянная времени компиляции,
И как бы то ни было, его геттер будет подвержен переопределению.

Учитывая, что стоимость не зависит от конкретного случая,
Лучше использовать статическую переменную.
Геттер экземпляра для него всегда может быть определен вручную, если это необходимо.
?

%
Это возможно для объявления переменной экземпляра
Включает модификатор {}
().
Результатом является то, что формальный параметр
соответствующий неявно индуцированный сеттер
считается ковариантным по декларации
(P000521).

{%
Модификатор {} на переменной экземпляра не имеет других эффектов.
В частности, тип возврата неявно индуцированного геттера
может быть преодолена ковариантно без ,
Он никогда не может быть переопределен на супертип или несвязанный тип.
независимо от того, является ли модификатор  Присутствует.%
?


 Конструкторы.


%
 Специальная функция, которая используется
в выражениях создания () для получения объектов,
Обычно путем их создания или инициализации.
Конструкторы могут быть генеративными ().
Или это могут быть заводы (P000524).

%
 Всегда начинается с
название его непосредственного замкнутого класса,
и необязательно может сопровождаться точкой и идентификатором .
Ошибка времени компиляции, если имя конструктора
Это не имя конструктора.

%
The
{функциональный тип конструктора}{функциональный тип! из конструктора.
 является типом функции
класс, который содержит декларацию ,
и чьи формальные типы параметров, факультативность и имена названных параметров
соответствует декларации .

{%
Обратите внимание, что тип функции  Конструктор  может содержать
переменные типа, объявленные прилагаемым классом .
В этом случае мы можем применить замену на , как в
,
где  Формальные параметры типа 
 определяются в данном контексте.
Мы также можем опустить такую замену, когда данный контекст
Объем корпуса , где  в объеме.%
?

{%
Декларация конструктора может противоречить статическим декларациям членов
()%
?

{%
Декларация конструктора не вводит имя в объем.
Если функция выражения вызова
()
или создание экземпляра
()
обозначает строителя как
, , , или \code{\metavar{prefix}.. 
разрешение зависит от объема библиотеки для определения класса
(возможно, через префикс импорта).
Классовая декларация проверяется непосредственно
Имеет ли он конструктор под названием  соответственно .
Невозможно, чтобы идентификатор непосредственно ссылался на конструктор,
Поскольку конструктор не находится ни в одной области, используемой для разрешения идентификаторов.
?

%
Если для класса  не указан конструктор,
имеет конструктор по умолчанию \code{C():\ \SUPER()\ \{\}},
Если только  Это встроенный класс .
% % TODO(Eernst): С нулевой безопасностью добавьте «или .


 Генеративные конструкторы.


%
A {генеративный конструктор}{конструктор!генератор}
Декларация состоит из имени конструктора, формального списка параметров.
(),и либо пункт перенаправления, либо список инициализаторов и факультативный орган.


<constructorSignature> ::= <constructorName> <formalParameterList>

<constructorName> ::= <typeIdentifier> ('.' <identifier>)?


<{%
Видишь? \synt{declaration} и {metodSignature}
введение списка перенаправления или инициализатора и тела.%
?

% % TODO(Eernst): Добавить {...}> ниже, когда эта команда добавлена.
%
Ошибка времени компиляции возникает, если декларацией генеративного конструктора
Имеет корпус формы .

<<%
В других декларациях функций,
Этот тип тела подразумевает, что значение  возвращается,
Но генеративные конструкторы ничего не возвращают.
?

%
Если официальное заявление о параметрах  происходит от
{поле Формальный Параметр}
Объявлено: .
Термин формы  Содержится в ,
 {имя} .
Ошибка времени компиляции  Не является также именем
переменная экземпляра немедленно закрывающего класса или числа.

%
Ошибка времени компиляции для инициализации формального параметра
Это происходит в любой функции, которая не является генеративным конструктором.
Кроме того, это ошибка компиляции для инициализации формального параметра.
Происходит в перенаправлении или внешнем конструкторе.
<<<p001213>> В Particuar всегда есть закрытый класс или enum. ?

%
Предположим, что  Объявление параметра инициализации, названного .
 быть типом переменной экземпляра, называемой  в
Немедленно включающий класс или число.
Если  имеет аннотацию типа  Заявленный тип  является .
В противном случае заявленный тип  является .
Ошибка времени компиляции, если заявленный тип 
Не является подтипом .

%
Инициализация формальностей является исключением из правила, согласно которому
Каждый формальный параметр вводит локальную переменную в
Формальный диапазон параметров ().
Когда формальный список параметров неперенаправляющего генеративного конструктора
содержит любые инициализирующие формы, вводится новый объем,
{формальный параметр начального диапазона}{%}
Формальный параметр инициализатор.
который является текущим объемом списка инициализаторов конструктора,
и который заключен в область, в которой объявлен конструктор.
Каждая инициализация в официальном списке параметров
вводит конечную локальную переменную в формальную область параметрического инициализатора;
но не в формальный параметр;
любой другой формальный параметр вводит локальную переменную в
Формальный параметр и формальный параметр инициализатор.

<{%
Это означает, что формальные параметры, включая инициализацию формальностей,
должны иметь различные названия, и что инициализация формальностей
находятся в пределах списка инициализаторов,
Но они не подходят для корпуса конструктора.
Когда формальный параметр вводит локальную переменную в два объема,
Это все еще одна переменная и, следовательно, одно место хранения.
Тип конструктора определяется по его формальным параметрам.
Включая инициализацию форм.%
?

%
Инициирование формальностей выполняется во время
Строительные конструкции, описанные ниже.
Выполнение инициализации формально 
вызывает переменную экземпляра  непосредственно окружающего класса
для присвоения значения соответствующего фактического параметра,
% Это может произойти из-за неудачного неявного слепка.
Если объект не имеет динамического типа
который не является подтипом заявленного типа переменной экземпляра ,
В этом случае возникает динамическая ошибка.

{%Вышеприведенное правило позволяет использовать инициализацию форм в качестве дополнительных параметров: %
?



<<%
является законным и имеет тот же эффект, что и %
?



%
 Это случай, личность которого отличается от
любой ранее выделенный экземпляр своего класса.
Генеративный конструктор всегда работает на свежем примере
сразу же закрывающий класс.

{%
Вышеприведенное удерживает, если конструктор фактически запущен, как это происходит на .
Если конструктор  На него ссылается ,  не может быть запущен;
Вместо этого можно посмотреть на канонический объект.
См. раздел о создании экземпляра ().%
?

%
Если генераторный конструктор  Не является перенаправляющим конструктором
и тело не предоставлено, тогда  Имеет пустое тело \code{\{\}}.


 Перенаправление генерирующих конструкторов


%
Генератор может быть
{redirecting}{constructor!redirecting}
В этом случае его единственным действием является вызов другого генеративного конструктора.
У перенаправляющего конструктора нет тела.
Вместо этого он имеет пункт перенаправления, который указывает
К кому обращен призыв,
с какими аргументами.


<redirection> ::= ':' \THIS{} ('.' <identifier>)? <аргументы>


\def\ConstMetavar{{\CONST?

%
%
Предположим, что

- наименование и формальные параметры типа ограждающего класса,
%
 обозначает либо  или ничего,
 -  или  для некоторого идентификатора ,
%
 - идентификатор.
Рассмотрим декларацию перенаправления генеративного конструктора
одной из форм





 $N$($T_1\ x_1 \ldots,\ T_n\ x_n,\ $\{$T_{n+1}\ x_{n+1} = d_1 \ldots,\ T_{n+k}\ x_{n+k} = d_k$\}):\ $R$;


где  является одной из форм





.

%
The
{redirectee constructor}{constructor!redirectee}}
Для этого конструктор обозначается
 соответственно .
Ошибка времени компиляции, если тип списка статических аргументов
()

не является подходящим совпадением для официального списка параметров перенаправления.

{%
Обратите внимание, что в случае, когда не пропущены именованные параметры
Покрывается разрешением  быть нулевым,
В случае, когда  является негенерическим классом
Покрывается разрешением  быть нулевым,
В этом случае формальный список параметров типа и фактический список аргументов типа
опущены
()%
?

{%
Нам требуется назначаемый матч, а не более строгий матч подтипа.
потому что конструктор генеративного перенаправления  ссылается на свой редирект 
в манере, которая напоминает вызов функции в целом.
Например,  Можно ли принять аргумент 
и передать выражение , используя , такое как , ,
И это было бы удивительно
Если  были подвержены более строгим ограничениям, чем те, которые применялись к
фактические аргументы для функционирования призывов в целом.%
?

%
Перенаправленный генеративный конструктор  
от перенаправления генеративного конструктора  если {f}
 является конструктором перенаправления ,
 Конструктор перенаправления 
 Перенаправление доступно из .
Ошибка времени компиляции при перенаправлении генеративного конструктора
Это перенаправление-доступно от самого себя.

%
Когда  ,Это ошибка времени компиляции, если перенаправление не является постоянным конструктором.
Более того, когда  , каждый
,
Должно быть потенциально постоянное выражение ().

%
% Эта ошибка может произойти из-за неудачного неявного слепка.
Это динамическая ошибка типа, если фактический аргумент принят.
перенаправление генеративного конструктора 
не является подтипом фактического типа ()
соответствующего формального параметра в декларации .
% Эта ошибка может произойти из-за неудачного неявного слепка.
Ошибка динамического типа, если фактический аргумент принят
перенаправление  перенаправления генеративного конструктора
не является подтипом фактического типа
()
соответствующего формального параметра в декларации перенаправления.


 Список инициализаторов.


%
Список инициализаторов начинается с толстой кишки,
и состоит из запятой отдельного списка отдельных .

{%
Существует три типа инициализаторов.


[] A {суперинициализатор} идентифицирует
--, то есть
Конкретный конструктор суперкласса.
Исполнение суперинициализаторных причин
Список инициализаторов суперконструктора, который должен быть выполнен.
[] {фактор инициализации переменной напряжения}
присваивает объект индивидуальной переменной экземпляра.
[] Утверждение.
%
?


<initializers> ::= ':' <initializer ListEntry> (',' <инициализатор) ListEntry> *

<инициализатор ListEntry> ::= \SUPER{} <аргументы>
 {} <идентификатор> <аргументы>
 <fieldInitializer>
 <утверждение>

<fieldInitializer> ::= <
({} '.')? <идентификатор> '=' <инициализатор Выражение

<initializerExpression> ::= \gnewline{}
<assignableExpression> <assignmentOperator> <expression>
 <условное выражение>
 <cascade>
 <rowExpression>


% % TODO(Eernst): Когда #54262 будет решен, удалите следующий абзац.
% % Если <выражение> изменяется таким образом, что получается <выражение функции>,
% % следующая ошибка будет просто свойством правила грамматики
%, потому что <initializerExpression> _not_ будет выведено <выражение функции>.
%
В качестве особого правила,
<инициализатор Выражение не может вывести  Выражение.

{%
Это решает почти двусмысленность:
В \code{A():\,\,x\,\,=\,\,\,\{\,\ldots\,\}},
 могут быть инициализированы до пустой записи,
Блок может быть телом конструктора.
Альтернативно,  можно инициализировать в объект функции,
Тогда у строителя не было бы тела.
Было бы известно только, какой случай мы имеем, когда встречаем
(или не встречать)
Семинария в самом конце.
Это считалось нечитаемым.
Следовательно, парсеры могут взять на себя обязательство не анализировать выражение функции.
в этой ситуации.
Обратите внимание, что это все еще возможно для  Выражение} для получения
термин, который содержит выражение функции в качестве подтермина, например,
\code{A():\,\,x\,\,=\,(()\,\{\,\ldots\,\});}.%
?

%
Начальник формы  является эквивалентным
инициализатор формы ,
Обе формы называются P000553.
Ошибка времени компиляции, если класс
не объявляет переменную экземпляра, названную .Это ошибка времени компиляции, за исключением статического типа 
присваивается заявленному типу .

%
Рассмотрим  <<<p000010> в форме



соответственно


.

\noindent{%}
 быть суперклассом замкнутого класса .
Ошибка времени компиляции в классе  не объявлять
Генеративный конструктор по имени  (перенаправлено с «P000935»).
В противном случае, статический анализ  выполняется
как указано в разделе ~,
Как будто  соответственно 
имел функциональный тип обозначенного конструктора,
% % TODO(Eernst): Следующее очень неточно, оно просто служит для запоминания.
% %, что мы должны указать, как обращаться с переменными типа в этом параметре
% % часть. Одна вещь, которую мы лелеем, заключается в том, что суперкласс может быть смесью.
%-ное применение.
Замена формальных переменных типа суперкласса
для соответствующего фактического типа аргументов, переданных суперклассу
в заголовке текущего класса.

%
Давай. {k} — генеративный конструктор.
 может включать в свой список инициализаторов не более одного суперинициализатора
или возникает ошибка времени компиляции.
Если нет суперинициализатора,
добавлен неявный суперинициализатор формы \SUPER{}()
в конце списка инициализаторов ,
Если только класс не является классом .
Это ошибка времени компиляции, если появляется суперинициализатор.
в Список инициализаторов P000023 в любой другой позиции, чем в конце.
Ошибка времени компиляции, если более одного инициализатора соответствует
Для данного экземпляра переменная появляется в списке инициализаторов .
Это ошибка времени компиляции, если список инициализаторов  содержит
инициализатор переменной, инициализированный посредством
форма инициализации .
Это ошибка времени компиляции, если список инициализаторов  содержит
инициализатор для конечной переменной  чье заявление включает
Выражение инициализации.
Ошибка времени компиляции  включает в себя инициализацию формального
для конечной переменной  чье заявление включает
Выражение инициализации.

%
Пусть {f} будет переменной конечного экземпляра, объявленной в
Немедленно включающий класс или число.
Ошибка времени компиляции происходит, если  инициализируется
одним из следующих средств:


  объявляется инициализацией переменной декларации.
  инициализируется посредством инициализации формы .
  имеет инициализатор в списке инициализаторов .


%
Это ошибка времени компиляции, если список инициализаторов  содержит
инициализатор для переменной, которая не
переменная экземпляра, объявленная в непосредственно окружающем классе.

{%
Список инициализаторов может содержать инициализатор для
любой переменной экземпляра, объявленной непосредственно окружающим классом,
Даже если это не окончательно.%
?

%
Это ошибка времени компиляции, если генеративный конструктор класса 
Включает в себя суперинициализатор.


 Выполнение генеративных конструкторов.


%
%
Исполнение генеративного конструктора  Тип 
для инициализации нового экземпляра 
всегда выполняется в отношении набора привязок для его формальных параметров.
и типовые параметры непосредственного замкнутого класса или перечня, связанного с
набор фактических аргументов типа , \DefineSymbol{\List{t}{1}{m}}.

{%
Эти связывания обычно определяются выражением создания экземпляра.
Строитель (прямо или косвенно).
Тем не менее, они также могут быть определены отражающим вызовом.
?

%
Если  перенаправляет, тогда его пункт перенаправления имеет форму



где \DefineSymbol{g} идентифицирует другой генеративный конструктор
непосредственно окружающего класса.
Затем выполняется  Для инициализации  доходы от
Оценка списка аргументов

Список аргументов  в форме

В среде, где
Типовые параметры ограждающего класса связаны с
.

%
Далее, тело  Выполняется для инициализации 
В отношении переплетов, которые составляют карту
Формальные параметры  для соответствующих объектов
в реальном списке аргументов ,
<< Ссылаясь на ,
и типовые параметры непосредственного замкнутого класса или перечня, связанного с
\List{t}{1}{m}.

%
В противном случае  не перенаправляется.
После этого казнь осуществляется следующим образом:

%
Заявления о переменных экземплярах непосредственно прилагаемого класса или перечня
Посещаются в том порядке, в котором они появляются в тексте программы.
Для каждого такого заявления , если  имеет форму
\code{\synt{финалConstVarOrType)  = ;
 оценивается на предмет 
и переменная экземпляра $v$  Ссылка: .

%
Любые инициализирующие формы, объявленные в списке параметров 
выполняются в том порядке, в котором они появляются в тексте программы.
% На самом деле этот порядок ненаблюдаем, это можно сделать в любое время.
% до начала работы организма, так как эти .
Затем инициализаторы списка инициализаторов  выполняются для инициализации 
в порядке их появления в программе, как описано ниже
(p.\,\pageref{executionOfInitializerLists}).

<{%
Мы могли бы наблюдать порядок, выполняя побочные действия, называемые внешними процедурами.
Поэтому необходимо указать порядок.%
?

%
Если есть переменная  объявленный
непосредственно включающий класс или
Он не привязан к объекту,
Все такие переменные инициализируются нулевым объектом ().

%
Затем, если прилагаемый класс не является , явно указанный или
Неявно добавленный суперинициализатор () выполняется для
дальнейшей инициализации .

%
После завершения суперинициализатора тело  исполняется
в сфере, где {} связано с .

<{%
Этот процесс гарантирует, что ни одна неинициализированная переменная конечного экземпляра не
Его всегда можно увидеть по коду.
Обратите внимание, что \THIS{} не находится в области действия на правой стороне инициализатора.
(см. )
Таким образом, ни один метод экземпляра не может выполняться во время инициализации:
метод не может быть использован непосредственно,
не может быть \THIS{} передаваться в любой другой код, на который ссылаются
в инициализаторе.%
?


<<<p001212>> Исполнение списков инициализаторов.


%
Во время выполнения генеративного конструктора инициализировать экземпляр
\DefineSymbol{i)
выполнение инициализатора формы 
идет следующим образом:

%
Во-первых, выражение  оценивается на предмет .
Затем переменная экземпляра   Обозначается как .
% Эта ошибка может произойти из-за неявного слепка.
Ошибка динамического типа, если динамический тип  не является
Подтип фактического типа
()
переменной .

%
Исполнение инициализатора, который является утверждением, происходит путем
Выполнение утверждения ().

%
Рассмотрим суперинициализатор {s} формы



соответственно


.

%
%
 класс, в котором  Появляется и позволяет  Это суперкласс P000096.Если  является общим (),
давать  - фактический тип аргументов, переданных ,
Получено путем замены действительных привязок \List{t}{1}{m)
Формальные параметры типа 
в суперклассе, как указано в заголовке .
Давай. \DefineSymbol{k} быть конструктором заявлено в $S$ названный
 соответственно .

%
Выполнение  идет следующим образом:
Список аргументов

оценивается по фактическому списку аргументов  в форме
.
Затем корпус суперконструктора  исполняется
в среде, где формальные параметры  связаны с
соответствующие фактические аргументы из ,
Формальные параметры типа  Обозначается как .


 Фабрики


%
A {завод}{конструктор!завод}
Конструктор, предваряемый встроенным идентификатором
()
.


<factoryConstructorSignature> ::= \gnewline{}
 \FACTORY{{}<constructorName> <formalParameterList>


%
Тип возврата фабрики, подпись которой
форма \FACTORY{}  или
форма \FACTORY{} 
, если  не является общим типом;
В противном случае тип возврата

где  являются типовыми параметрами ограждающего класса.

%
Ошибка времени компиляции  Это не имя
% % TODO (Eernst): Come Dart 3.0, add 'mixin'.
Немедленно включающий класс или число.

%
% Эта ошибка может произойти из-за неявного слепка.
Это динамическая ошибка типа, если завод возвращает ненулевой объект.
тип которого не является подтипом его фактического
()
Тип возврата.

{%
Бесполезно позволять фабрике возвращать нулевой объект (P000815).
Но это более единообразно, чтобы позволить это, как это делают правила в настоящее время.
?

{%
Заводы устраняют классические недостатки, связанные с конструкторами
на других языках.
Заводы могут выпускать экземпляры, которые не были недавно выделены:
Они могут прийти из кэша.
Кроме того, фабрики могут возвращать экземпляры разных классов.
?


 Перенаправление заводских строителей.


%
A {перенаправление заводского конструктора}{конструктор!перенаправление завода}
указывает вызов конструктору другого класса, который должен быть использован
Всякий раз, когда вызывается перенаправляющий конструктор.


<redirectoryConstructorSignature> ::= \gnewline{}
 \FACTORY{} <constructorName> <formalParameterList> \gnewline{}
<конструкторское проектирование>

<constructorDesignation> ::= <typeIdentifier>
 <qualifiedName>
 <typeName> <typeArguments> ('.' <identifier>)?


Предположим, что
%

- наименование и формальные параметры типа ограждающего класса,
%
  или пустой,
 -  или  для некоторого идентификатора ,
и \DefineSymbol{ является идентификатором,
Затем рассмотрите декларацию перенаправляющего заводского конструктора.
{k} одной из форм




%
где  является одной из форм
 или
.

%
Ошибка времени компиляции  не обозначает
класс, доступный в текущем объеме.
Если  обозначает такой класс ,
Ошибка времени компиляции  не обозначает конструктора.
% Этой индукции достаточно, чтобы проверить абстрактность на один уровень ниже.
%, потому что это ошибка перенаправления, если это происходит после нескольких
% перенаправления:
В противном случае это ошибка времени компиляции.Если  обозначает генеративный конструктор и  является абстрактным.
В противном случае,
{redirectee constructor}{constructor!redirectee}}
Для этого декларация является конструктором {k'} обозначается .

%
Перенаправляющий заводской конструктор  
от перенаправляющего заводского конструктора  если {f}
 является конструктором перенаправления ,
 Конструктор перенаправления 
 Перенаправление доступно из .
Ошибка времени компиляции при перенаправлении конструктора завода
Это перенаправление-доступно от самого себя.

%
Пусть  Статический тип списка аргументов
()

когда  не приводит никаких аргументов, и

когда  Приведем несколько названных аргументов.
Ошибка времени компиляции 
не является подтипом для формального списка параметров перенаправления.

{%
Требуется подтип матча
(в отличие от более прощального матча)
который используется с конструктором генеративного перенаправления,
потому что завод перенаправляет конструктора  всегда призывает
его перенаправление  с
Те же самые аргументы, что и  получено.
Это означает, что понижение фактического аргумента
 и 
не используется, так как фактический аргумент
тип, требуемый ,
или это будет означать динамическую ошибку, которая просто задерживается на один шаг.
?

{%
Обратите внимание, что необычный случай охватывается
 или  или оба равны нулю,
в этом случае формальный список параметров типа класса 
и/или фактический список аргументов типа
Конструктор перенаправления опущен
()%
?

%
Ошибка времени компиляции  явно указывает
значение по умолчанию для необязательного параметра.

{%
Значения по умолчанию, указанные в  будет игнорироваться,
Так как это параметры .
Следовательно, значения по умолчанию не допускаются.%
?

%
Ошибка времени компиляции, если формальный параметр  имеет значение по умолчанию
тип которого не является подтипом аннотации типа
по соответствующему формальному параметру в .

{%
Обратите внимание, что невозможно изменить аргументы, передаваемые .%
?

{%
На первый взгляд, можно подумать, что
Обычные заводские конструкторы могли бы просто создать
других классов и вернуть их,
Перенаправлять фабрики не нужно.
Однако перенаправление заводов имеет ряд преимуществ:



Абстрактный класс может быть постоянным конструктором.
Использует константный конструктор другого класса.

Перенаправляющий завод-конструктор избегает необходимости в экспедиторах
повторять формальные параметры и их значения по умолчанию.
%
?

%
Ошибка времени компиляции  Приставка с приставкой {} модификатор
 Не является конструктором (P000818).

%
Пусть \DefineSymbol{\List T}{1}{m}} — фактический тип аргументов, переданный 
в декларации .
Пусть \DefineSymbol{{X}{1}{m}} являются формальными параметрами типа, объявленными
Класс, содержащий декларацию .
 F' - тип функции  (P000819).
Ошибка времени компиляции 
не является подтипом функции типа .

{%
В случае, когда два класса не являются общими
Это просто проверка подтипа на типы функций двух конструкторов.
В общем, это означает, что полученный объект соответствует% % TODO: Come Dart 3.0, добавьте «миксин».
Интерфейс класса или перечня, который непосредственно включает .%
?

%
Для динамической семантики,
предположить, что {k} - перенаправляющий заводской конструктор
и \DefineSymbol{k'} является перенаправлением .

%
% Эта ошибка может произойти из-за неявного слепка.
Это ошибка динамического типа, если фактический аргумент принят в вызове 
не является подтипом фактического типа ()
соответствующего формального параметра в декларации .

%
При перенаправлении  является фабрикантом,
Выполнение  Выполнение 
Фактические аргументы перешли к .
Результат исполнения  является результатом .

%
При перенаправлении  является генеративным конструктором,
 Будьте свежим примером ()
Класс, содержащий .
Выполнение  Тогда это равносильно исполнению  инициализировать ,
регулируется теми же правилами, что и выражение создания экземпляра
().
Если  Затем выполняется обычное выполнение 
обычно возвращается ,
иначе  Завершает, бросая исключение и след стека
Выброшено .


 Постоянные строители.


%
A {постоянно} Конструктор (Constructor!Constant)
Может использоваться для создания объектов постоянной времени компиляции ().
Константный конструктор закрепляется зарезервированным словом .


<constantConstructorSignature> ::= \gnewline{}
\CONST{{}<constructorName> <formalParameterList>


{%
Постоянные конструкторы имеют более сильные ограничения, чем другие.
Например, вся работа
Конструктор генеративной постоянной без перенаправления
должно быть выполнено в его инициализаторах
При инициализации выражений
Примерные переменные закрывающего класса
(Возможно, это произошло раньше).
Потому что инициализация должна быть постоянной.%
?

%
Постоянные перенаправления генеративных и заводских конструкторов указаны в другом месте.
(p.\,\pageref{redirectingGenerativeConstructors}}
p.\,\pageref{redirectingFactoryConstructors}.
Отныне этот раздел посвящен
Неперенаправляющие константные генераторы.

%
Это ошибка времени компиляции, если конструатор генеративной константы не перенаправляет
Объявляется классом, который имеет переменную экземпляра, которая не является окончательной.

{%
Вышеуказанное относится как к локально объявленным, так и к унаследованным переменным экземпляров.
?

%
Если неперенаправленный константный конструктор \DefineSymbol{k)
Объявляется классом ,
Это ошибка времени компиляции
переменная, объявленная в 
иметь инициализирующее выражение, которое не является постоянным выражением.

{%
Суперкласс  не может иметь такого инициализирующего выражения  тоже.
Если у него есть неперенаправляющий константный конструктор
 Это ошибка,
Если у вас нет такого конструктора
затем (явный или неявный) суперинициализатор в  Ошибка.%
?

%
Суперинициализатор, который появляется, явно или неявно,
в списке инициализаторов постоянного конструктора
Необходимо указать конструктор генеративной постоянной
Сверхкласс немедленно закрывающегося класса,
или возникает ошибка времени компиляции.

%
Любое выражение, которое появляется внутри
Список инициализаторов постоянного конструктора
Должно быть потенциально постоянное выражение
(),
или возникает ошибка времени компиляции.

%
Когда константный конструктор \DefineSymbol{k) На него ссылаются
Постоянное выражение объекта,Ошибка времени компиляции, если
Обращение к  Во время пробежки было бы исключение,
Ошибка времени компиляции, если
Подмена фактических аргументов формальными параметрами
дает инициализирующее выражение  в списке инициализаторов 
Это не постоянное выражение.

{%
Например, если  
где  является формальным аргументом  с типом ,
 потенциально постоянен и может быть использован в списке инициализаторов .
Ошибка вызова  С аргументом типа 
Если  Отличается от класса ,
Даже если   Геттер,
и то же выражение будет оцениваться без ошибок во время выполнения.
?


 Статические методы


%
{Статические методы}{метод!статика}
функции, кроме геттеров или сеттеров,
декларация которых непосредственно содержится в классовой декларации
Об этом говорится в статье .
Статические методы класса  Статические методы, объявленные .

{%
Наследование статических методов мало полезно в Дарте.
Статические методы не могут быть отменены.
Любой необходимый статический метод можно получить из его декларирующей библиотеки.
И нет необходимости вносить его в сферу действия через наследование.
Опыт показывает, что разработчики запутались в
Идея наследственных методов, которые не являются методами экземпляра.

Конечно, вся концепция статических методов является спорной.
Но он сохранился здесь, потому что многие программисты знакомы с ним.
Статические методы Дарта можно рассматривать как функции библиотеки.
?

{%
Заявления статического метода могут противоречить другим заявлениям
()%
?


<<<p001280>> Суперклассы.


% % TODO(Eernst): Надо сказать, что сверхкласс, который получается
% % по применению смешивания является общим, когда  является общим или, по крайней мере,
%, когда классы используют одну или несколько переменных типа 
% % в \EXTENDS{} или \WITH{} Статья . Ниже говорится, что
% % этих оговорок находятся в области параметров типа , но это
%% не позволяют нам говорить о суперклассе как о реальном, автономном
Класс %% (если только мы не начнем определять вложенные классы так, что
Процентный суперкласс может быть объявлен в этой области.

%
Суперкласс  класса  чье заявление имеет пункт
\code{{} 
и пункт о продлении
\code{\EXTENDS{} 
абстрактный класс, полученный путем применения
Композиция микширования () .
Имя  является новым идентификатором.
Если нет  Затем указывается пункт {} пункт
Класс  Определяет свой суперкласс.
Если нет  пункт уточняется, затем либо:


  – это , не имеющий суперкласса. или
 Класс  Считается, что имеет  пункт о форме
\code{{) Объект, и вышеприведенные правила применяются.


%
Ошибка времени компиляции для указания {} пункт
для класса .


<суперкласс> ::= < <typeNotVoid> <mixins>
 <mixins>

<mixins> ::= \WITH{{<typeNotVoidList>>


%
Область применения \EXTENDS{} и \WITH{} класс  это
Типовой параметр .

%
Ошибка времени компиляции, если тип
в {} класс  это
переменная типа (),
псевдоним типа, не обозначающий класс (),
Перечисленный тип (),Отложенный тип (), тип \DYNAMIC{} (),
или типа   (P000832).

{%
Обратите внимание, что {} является зарезервированным словом,
что подразумевает, что те же ограничения применяются для типа ,
Аналогичные ограничения установлены для других типов, таких как:
 () и
 ()%
?

{%
Параметры типового класса доступны в
лексический охват статьи о суперклассе,
Потенциально затеняющие классы в окружающем пространстве.
Поэтому следующий кодекс является незаконным.
и должно вызывать ошибку времени компиляции: %
?



%
% % TODO(Eernst): Подумайте о замене всех случаев «суперкласса».
«прямой или косвенный суперкласс», потому что это слишком запутанно.
Класс   класса  Если {f} либо:


  является суперклассом , или
  является суперклассом класса ,
 Это суперкласс P000245.


%
Ошибка времени компиляции в классе  Это суперкласс сам по себе.


 Наследование и преобладание


%
 Будьте классом, пусть  быть суперклассом , и
 быть суперклассами  Это также подклассы .
  все конкретные, доступные экземпляры 
которые не были отменены конкретной декларацией в 
или по меньшей мере в одном из .

<%
Было бы более привлекательно дать чисто местное определение наследственности.
Это зависело только от членов прямого суперкласса .
Класс  может наследовать член  это
Не является членом своего суперкласса. .
Это может произойти, когда член  является частным для библиотеки  ,
тогда как  Из другой библиотеки ,
Сеть суперклассов  включает класс, заявленный в .%
?

%
Класс может переопределить членов экземпляра
В противном случае они были бы унаследованы от суперкласса.

%
 быть классом, объявленным в библиотеке , и
 быть набором всех суперклассов ,
где  Это суперкласс  для .
 Это встроенный класс .
 декларировать конкретный член , и
 быть конкретным членом , который имеет
Имя: ,
 Доступен по адресу .
Далее  переопределение 
Если  не перевешивается конкретным членом
По крайней мере, один из 
ни того, ни другого  или  являются переменными экземпляров.

{%
Инстанционные переменные никогда не превосходят друг друга.
Геттеры и сеттеры, индуцированные переменными экземпляров do.%
?

{%
Опять же, было бы предпочтительнее местное определение переопределения.
Но не учитывает конфиденциальность библиотеки.%
?

{%
Является ли оверрайд законным или нет, определяется относительно
все прямые суперинтерфейсы, а не только интерфейс суперкласса;
Это описано в другом месте
(P000835).
Статические члены никогда ничего не отменяют,
Они могут участвовать в некоторых конфликтах.
Включение деклараций в суперинтерфейсы
()%
?

{%
Для удобства, вот краткое изложение соответствующих правил,
Использование слова «ошибка» для обозначения ошибок времени компиляции.
Помните, что это не нормативно.
Язык управления находится в соответствующих разделах спецификации.

 Существует только одно пространство имен
для геттеров, сеттеров, методов и конструкторов ().
Нелокальная переменная  Вводит геттер ,
и нелокальной переменной 
Также вводится сеттер
Если оно не является постоянным и окончательным,
или оно является поздним и окончательным и не имеет начального выражения
 (, .
Когда мы говорим об участниках, мы имеем в виду
доступные экземпляры, статические или библиотечные переменные,
Геттеры, сеттеры и методы
(P000840).
 Вы не можете иметь двух членов с одинаковым именем в одном классе.
Они объявили или унаследовали (). .
 Статические члены никогда не наследуются.
 Это ошибка, если у вас есть статический элемент под названием  В вашем классе
и экземпляр с тем же именем
().
 Это ошибка, если у вас есть статический установщик ,
и экземпляр член  (P000844).
 Это ошибка, если у вас есть статический геттер 
и набором экземпляров  (P000845).
 Если вы определили экземпляр члена по имени ,
и ваш суперкласс имеет одноимённый экземпляр,
Они преобладают друг над другом.
Это может быть или не быть законным.
 
Если два члена перевешивают друг друга,
Это ошибка, если только она не является правильной.
(P000846).
 Сеттеры, геттеры и операторы никогда не имели
факультативные параметры любого вида;
Это ошибка (, , .

Это ошибка, если член имеет то же имя, что и его класс.
(P000850).
 Класс имеет неявный интерфейс ().
 Члены суперинтерфейса не наследуются классом.
Они наследуются по неявному интерфейсу.
Интерфейсы имеют собственные правила наследования
(P000852).
 Член является абстрактным, если
не имеет тела и не помечена 
(, .
 Класс является абстрактным, если он явно обозначен .
 Это ошибка, если конкретный класс не реализует какой-либо член.
его интерфейса, и нет нетривиального 
(P000855).
 Ошибочно называть нефабрику конструктором абстрактного класса.
использование выражения создания экземпляра (),
Такой конструктор может быть вызван только от другого конструктора.
Использование суперинвокации ().
 Если класс определяет экземпляр члена по имени ,
и любой из его суперинтерфейсов имеет подпись участника под названием ,
Интерфейс класса содержит  от самого класса.
 Интерфейс наследует всех членов своего суперинтерфейса.
Они не перегружены и не являются членами нескольких суперинтерфейсов.
 Если несколько суперинтерфейсов интерфейса
Определение члена с таким же названием, как ,
В большинстве случаев один из членов наследуется.
Этот член (если он существует) является тем, чей тип является подтипом.
всех остальных.
Если такого члена нет, возникает ошибка.
(P000858).
 Правило  относится к интерфейсам, а также классам
(P000860).
 Это ошибка, если конкретный класс не имеет реализации.
Для метода в его интерфейсе
если он не имеет нетривиального 
(P000861).
 Идентификатор поименованного конструктора не может быть таким же, как
Базовое имя статического члена, объявленного в том же классе
().
%
?


 Суперинтерфейсы


%
Класс имеет набор .Этот набор содержит интерфейс своего суперкласса.
и интерфейсы классов, указанных в
 Положение о классе.


<интерфейсы> ::= \IMPLEMENTS{{<typeNotVoidList>>


%
Область применения  класс  это
Типовой параметр .

%
Ошибка времени компиляции, если элемент
в списке типов  класс  это
переменная типа (),
псевдоним типа, не обозначающий класс (),
Перечисленный тип (),
Отложенный тип (), тип < (),
или типа   (P000868).
Ошибка времени компиляции, если два элемента в списке типов
% % TODO(Eernst): См. понятие nnbd «один и тот же тип».
 класс  обозначает один и тот же тип .
Это ошибка времени компиляции, если суперкласс класса  это
Один из элементов списка типов  Положение .

{%
Можно утверждать, что это безвредно.
повторение типа в списке суперинтерфейсов;
Так зачем делать это ошибкой?
Дело не столько в том, что ситуация ошибочна,
Но это бессмысленно.
Таким образом, это признак того, что программист вполне мог иметь в виду.
Сказать что-то еще - и это ошибка, которая должна быть
Привлекает к себе внимание.%
?

%
Ошибка времени компиляции в классе  Имеет два суперинтерфейса
Это разные экземпляры одного и того же общего класса.
{%
Например, класс не может иметь
оба  и  в качестве суперинтерфейсов
прямо или косвенно.%
?

%
При общем классе  объявляет параметр типа ,
Ошибка времени компиляции  происходит в нековариантном положении
% Можно сказать «прямой суперинтерфейс», но легко понять, что это
% достаточно для проверки прямых суперинтерфейсов, и это верно
% для косвенных.
в типе, который определяет суперинтерфейс .
{%
Например, класс не может иметь
\code{Список\VOID{} ()
в его \EXTENDS{} или {} клаузулы.%
?

%
Это ошибка времени компиляции, если интерфейс класса  это
Суперинтерфейс самого себя.

{%
Класс не наследует членов от своих суперинтерфейсов.
Однако его неявный интерфейс делает.%
?


 Конфликты членов класса.


%
Некоторые пары деклараций членов класса, смеси, перечня и расширения
не может сосуществовать,
Хотя они не вводят одно и то же имя в один и тот же объем.
В этом разделе указаны эти ошибки.

%
 из геттера или метода под названием  означает ;
Базовое имя сеттера  является .
Базовое имя оператора  - ,
За исключением оператора  Имя пользователя: .

%
 Будьте классом.
Ошибка времени компиляции 
Объявлен конструктор по имени  и
Статический элемент с базовым именем .
Ошибка времени компиляции 
Объявляет статический элемент с именем базы  и
Интерфейс  имеет элемент экземпляра с именем базы .
Это ошибка времени компиляции, если интерфейс 
имеет метод экземпляра под названием  Настройщик экземпляров с именем базы .
Ошибка времени компиляции  Статический метод, названный 
и статический сеттер с именем базы .

%
Когда {C} представляет собой смесь или расширение,
Ошибки времени компиляции происходят по тем же правилам.
{%
В некоторых случаях это избыточно.Например, это уже ошибка для смесителя объявить конструктор.
Но существуют и полезные случаи, например, конфликт между статичным членом.
и примерный член.%
?

%
Эти ошибки возникают, когда геттеры или сеттеры определены явно.
а также когда они вызваны переменными декларациями.

{%
Обратите внимание, что другие ошибки, которые похожи по своей природе, рассматриваются в другом месте.
Например, если  Класс, который имеет два суперинтерфейса  и ,
где  Метод получил название 
 имеет геттер с именем ,
Тогда это ошибка, потому что вычисление интерфейса 
включает расчет объединенной подписи члена
()
из этого геттера и этого метода,
Это ошибка для объединенной подписи члена.
Включая геттер и негетер.%
?


 Интерфейсы


%
Этот раздел вводит понятие интерфейсов.
Прежде всего, мы определяем понятие подписей членов.
Эта концепция необходима в определении интерфейсов.

% Нам нужна отдельная концепция подписей членов,
%, так что мы можем получить чистую обработку того, как вычислить
% интерфейс класса и общий интерфейс набора
% суперинтерфейсов, например, суперинтерфейсов класса или
% общий интерфейс, представленный оговоркой «on» в
% «миксин». Например, мы не хотим указывать каждый раз.
% мы проверяем наличие конфликтов, которые не имеют значения для организма;
% что метаданные не имеют значения, интерфейс класса
% никогда не должно содержать тела какого-либо члена в
% первое место. Также четкое разделение синтаксического
Декларация % и подпись участника обеспечивают четкое определение
% возможность введения преобразований, например, для замены
% некоторые типы параметров другими, которые необходимы для
% спецификация правильного отношения переопределения.

%
 
может быть получено из декларации члена экземпляра класса .
Он содержит ту же информацию, что и ,
За исключением того, что  опустить тело, если оно есть;
содержит тип возврата и типы параметров
даже если они подразумеваются в ;
опускает названия позиционных параметров;
не содержит модификатора  по каждому параметру, если таковой имеется;
Опускает метаданные
();
и опускает информацию о том, является ли член
, , , или .
Не имеет значения, будет ли  дается как явный синтаксис
или индуцируется имплицитно, например, посредством переменной декларации.
Наконец, если  имеет формальные параметры,
Каждый из них имеет модификатор 
()
если и только если этот параметр является ковариантным
(P000872).

%
Мы используем синтаксис, аналогичный синтаксису абстрактной декларации участника.
Указать подписи членов.
Разница в том, что названия позиционных параметров опущены.
Этот синтаксис используется только для целей спецификации.

{%
Подписи участников являются синтетическими объектами, то есть
Они не поддерживаются как конкретный синтаксис в программе Dart.
Это вычисленные объекты, используемые во время статического анализа.
Однако полезно уметь указать на
Свойства подписи члена в данной спецификации
через синтаксическое представление.
Подпись делает его явным
является ли параметр ковариантным по декларированию,
Но остается неявным, является ли он ковариантным по классу.
(P000873).
Причина этого заключается в том, что правило для определения того,
Данное отношение override правильно
()
зависит от первого, а не от второго.
?

%
Пусть {m} будет сигнатурой метода формы





  T}{1}{n},}

\qquad   T {\displaystyle T} d}{n+1}{n+k})}.


{функциональный тип}{методическая подпись! тип функции
 тогда


{T_0}.

%
Пусть {m} будет сигнатурой метода формы





  T}{1}{n},}


<\code\qquad\qquad\{\TripleList\COVARIANT T}{x}{= d}{n+1}{n+k}<<< P001613>>>).


 тогда


{T_0}.

%
Пусть \DefineSymbol{m} будет сеттерской подписью формы
.
{функциональный тип}  тогда
\FunctionTypeSimple{\VOID}{$T$}.

%
Тип функции подписи члена остается неизменным, если
Некоторые или все значения по умолчанию опущены.

<<%
Мы не указываем тип функции подписи Геттера.
Для таких подписей мы вместо этого будем напрямую ссылаться на тип возврата.
?

%
 Синтетическая сущность, которая определяет
Как можно взаимодействовать с объектом.
Интерфейс имеет сигнатуры метода, Getter и Setter.
и набор суперинтерфейсов,
Это опять же интерфейсы.
Каждый интерфейс является неявным интерфейсом класса.
В этом случае мы называем это
{class interface}{interface!class}
или сочетание нескольких других интерфейсов,
В этом случае мы называем это
{комбинированный интерфейс}{интерфейс!комбинированный}.

%
 Будьте классом.
   Интерфейс, который заявляет
подпись, полученная от
каждый экземпляр, заявленный .
  Это прямые суперинтерфейсы 
(P000875).

{%
Мы говорим, что интерфейс класса "объявляет" эти подписи членов,
Таким образом, мы можем сказать, что интерфейс «объявляет» или «имеет» участник,
Так же, как и для классов.
Обратите внимание, что подпись участника  интерфейса класса 
может иметь параметр  с модификатором ,
Хотя  была получена из декларации  в 
и параметр, соответствующий  в  не
Есть такой модификатор.
Это связано с тем, что  может иметь «наследственное» " "
свойство быть ковариантным по декларации
Один из его суперинтерфейсов
()%
?

%
С целью выполнения статических проверок по обычным методам вызовов
()
и извлечение имущества
(),
любой тип  Что такое P000378 ограниченный
(),
где  класс с интерфейсом ,
Также считается, что интерфейс .
Точно так же, когда   где  является функциональным типом,
 Считается, что метод называется \CALL{} с подписью ,
Тип функции  является .

%
{I, \List{I}{1}{k}}%
  Список интерфейсов  >> I}{1}{k)
является интерфейсом, который объявляет набор подписей членов ,
где  определяется как указано ниже.
  Набор \List I}{1}{k}.

%
 быть совокупностью всех подписей членов, объявленных  I}{1}{k}.
{M} - наименьший набор, удовлетворяющий следующему:


 Для каждого имени  Библиотека   содержит
подпись участника под названием  который доступен по адресу , быть объединенной подписью участника под названием 
из \List{I}{1}{k} в отношении .
Это ошибка времени компиляции
если вычисление этой объединенной подписи не удалось.
В противном случае  Содержит .


{%
Интерфейсы должны содержать недоступные подписи участников.
Они могут быть доступны из интерфейсов, связанных с
Декларации подтипов.%
?

{%
Например, класс  В библиотеке  может объявлять частного члена по имени
,
Класс  В другой библиотеке  может расширяться ,
Класс  В библиотеке  может расширяться ;
 может объявить члена, который переопределяет  из ,
и это отношение переопределения должно проверяться на основе интерфейса .
Таким образом, мы не можем разрешить интерфейс 
«Забыть» недоступных членов .

Для конфликтов ситуация еще более требовательна:
Классы  и  В библиотеке  могут объявлять частных
 и ,
Подтип  В другой библиотеке  может иметь
a < Перечисление обоих  и .
В этом случае мы должны сообщать о конфликте, даже если конфликтующие стороны
Заявления не доступны для ,
потому что эти подписи членов затем не пересылаются
(),
и призыв к  на примере $D$ в $L$
должен вернуть "int" в соответствии с подписью первого члена;
и он должен возвращать объект функции согласно второму,
и призыв к 
Должен вернуть  с первой подписью, и она должна быть
(время компиляции или, для динамического вызова, время выполнения) со вторым.%
?

{%
Возможно, невозможно одновременно удовлетворить эти ограничения.
Это неизбежно будет сложная семантика.
Поэтому мы решили сделать это ошибкой.
К сожалению, добавление частной декларации
В одной библиотеке может быть нарушен существующий код в другой библиотеке.
Но следует отметить, что конфликты могут быть выявлены локально.
в библиотеке, где существуют частные декларации,
Они возникают только для частных лиц, имеющих
одно и то же имя и несовместимые подписи.
Переименовать этого частного члена во что-нибудь, что не использовалось в этой библиотеке
Устранит конфликт и не сломает ни одного клиента.
?


 Комбинированные подписи членов.


%
Этот раздел определяет, как вычислить подпись участника, которая
должным образом обозначать приоритетный набор из нескольких подписей членов;
взяты из заданного списка интерфейсов.

{%
Комбинированная подпись имеет тип, который
Подтип всех типов, указанных для этого члена.
Это необходимо для того, чтобы гарантировать, что тип
Член  класса  является четко определенным,
Даже в том случае, когда  наследник
несколько различных деклараций 
и не переопределяет .
В случае неудачи он служит для конкретизации ситуаций.
разработчик должен добавить декларацию, чтобы разрешить двусмысленность.
Подписи участников являются приоритетными в том смысле, что мы выберем
подпись участника интерфейса с наименьшим возможным индексом;
в случае, когда несколько подписей членов одинаково подходят
быть выбрана в качестве объединенной подписи.
То есть «выигрывает первый интерфейс».%
?

%
Для целей вычисления объединенной подписи члена,
Нам нужно особое понятие
\IndexCustom{равенство}{равенство подписей членов}
подписей членов.
Две подписи членов  и  равнымиЕсли они имеют одно и то же имя,
доступ к одному и тому же набору библиотек;
имеют одинаковый тип возврата (для геттеров),
или один и тот же тип функции и одни и те же случаи 
(для методов и сеттеров).

{%
В частности, частные методы из разных библиотек никогда не бывают равными.
Лучшие типы также различаются.
Например,
\FunctionTypeSimple{}{} и {}{}
Они не равны, хотя и являются подтипами друг друга.
Мы нуждаемся в этом различии, потому что управление несоответствиями высшего типа является
одна из целей вычисления комбинированного интерфейса.%
?

%
%
Теперь мы определили состав подписей.
Пусть  быть идентификатором  библиотека,
{I}{1}{k} список интерфейсов,
и {M_0} набор
Все подписи участников от \List{I}{1}{k}
Доступен по адресу .
The
{комбинированная подпись участника
< \List{I}{1}{k} в отношении }{%
объединенная подпись члена
является подписью участника, которая получается следующим образом:

%
Если  пуст, вычисление объединенной подписи члена не удалось.

%
Если  содержит ровно одну подпись ,
Комбинированная подпись участника .

%
В противном случае  содержит более одной подписи
\DefineSymbol{\List{m}{1}{q}}

%
 Неудачные смеси
Если  содержит по меньшей мере одну подпись
и, по меньшей мере, одной подписью,
Вычисление объединенной подписи не увенчалось успехом.


%
 Геттерс.
Если  содержит только геттерские подписи,
расчет объединенной подписи члена продолжается, как описано ниже
Методы и методы,
За исключением того, что он использует тип возврата подписи геттера.
где методы и установщики используют тип функции подписи участника.


%
 Методы и установки
В данном случае  состоит только из подписей,
или только подписи,
Потому что имя  в первом случае всегда заканчивается на 
Что никогда не бывает правдой в последнем случае.

%
%
Определить, существует ли непустой набор  такой, что
для каждого ,
Тип функции  является подтипом
Тип функции  Для каждого .
%
Если такого набора не существует, вычисление объединенной подписи члена не удается.
{%
Полезная интуиция в этой ситуации заключается в том, что
Подписи участников не согласованы
Какой тип подходит для члена по имени .
В противном случае у нас есть набор подписей членов, которые являются наиболее конкретными. " "
в том смысле, что их типы функций являются подтипами их всех.
?

{% Сфера действия для N_min, M_min, N_firstMin.

\def\Nall{\ensuremath{N_{\mbox{\scriptsize{}all}}}}
\def\Mall{{M_{\mbox{\scriptsize{}all}}}}
\def\MallFirst{{M_{{\scriptsize{}first}}}}

%
В противном случае, когда набор  Как указано выше, существует
Пусть {} будет наибольшим набором, удовлетворяющим требованию ,
%
И пусть .
{%
То есть, \Mall{} содержит все подписи участников, названные 
с наиболее специфическим типом.
Dart subtyping - это частичный предварительный заказ.
который обеспечивает существование такого наибольшего набора наименьших элементов,
При наличии непустого набора элементов.
Мы можем иметь несколько таких подписей, потому что подписи членов
Они могут быть такими, что не равны.
И все же их типы функций являются подтипами друг друга.
Мы должны вычислить одну подпись участника из ,и мы делаем это, используя порядок заданных интерфейсов.%
?

%
%
 наименьшее число, такое, что
 Он не пустой.
 Один элемент, который  содержит.
{%
Этот набор содержит ровно один элемент, потому что он не пуст.
и ни один интерфейс не содержит более одной подписи участника под названием .
Другими словами, мы выбираем  В качестве подписи от
Первый возможный интерфейс
Среди наиболее специфических подписей %
?
?

%
Комбинированная подпись участника ,
который получен из  Добавление модификатора {}
для каждого параметра  (если его еще нет)
когда существует 
параметр, соответствующий 
()
Имеет модификатор .
{%
Другими словами,
каждый параметр в объединенной подписи члена помечается ковариантным
если любой из соответствующих параметров отмечен ковариантным,
И не только среди самых известных,
Но среди подписей \emph{all} с именем \id{} (перенаправлено с «P000455»)
в приведенном списке интерфейсов.%
?



 Суперинтерфейсы


%
Интерфейс имеет набор прямых суперинтерфейсов
(P000882).
Интерфейс  представляет собой  интерфейс 
Если {f} либо  Это прямой суперинтерфейс 
 является суперинтерфейсом прямого суперинтерфейса .

%
Когда мы говорим, что тип 
 {implements}{type!implements a type}
другой тип ,
Это означает, что  является суперинтерфейсом ,
  ограниченный каким-либо типом 
(),
 Это суперинтерфейс P000470.
Предположим, что  Это сырой тип
()
чье заявление объявляет  Типовые параметры.
Когда мы говорим, что тип 
 {implements}{type!implements a raw type}
,
Это означает, что существуют типы \List{U}{1}{}
 реализация < {1}{s}>}

{%
Обратите внимание, что это не то же самое, что быть подтипом.
Например,  реализация ,
Но он не реализует .
Точно так же  реализация .
Также обратите внимание, что когда  реализация 
где  не является подтипом ,
 не может быть подтипом .%
?

%
Предположим, что  Тип и  Это сырой тип, такой, что  реализация .
Существуют уникальные типы {U}{1}{}
 реализация < {1}{s}>}
Затем мы говорим, что {U}{1}{s} являются
\IndexCustom{аргументы фактического типа  при $G$}{%
Введите аргументы! типа в сыром виде,
Мы говорим: «P000489». Это
\IndexCustom{$j$-фактический аргумент типа  при $G$}{%
Введите аргументы! тип в сыром виде, ,
Для любого .

{%
Например, аргумент типа  при 
.
Эта концепция особенно полезна, когда
цепочка прямых суперинтерфейсов от  до 
не просто передает все типы аргументов в неизменном виде, например,
с таким заявлением, как
 C<X, Y>\,\,\EXTENDS\,\,B<List<Y>\, Y, \,X> \{\}.%
?


 Наследование и преобладание


%%
 быть интерфейсом и  Будьте библиотекой.
Определить  быть набором подписей членов
\DefineSymbol{m)
Таким образом, все следующие положения:


  Доступен для  и
  Это прямой суперинтерфейс  и либо

  подпись участника  или
  является членом .

  Не переопределяется .


%
Кроме того, мы определяем  быть
Набор подписей членов \DefineSymbol{m'')
Таким образом, все следующие положения:


  Это интерфейс класса .
  декларирует подпись участника .
  Он имеет то же название, что и P000516.
  Доступен по адресу .
  Это прямой суперинтерфейс  и либо

  подпись участника  или
  Является членом .



%
 быть интерфейсом класса  Объявлено в библиотеке .
  Все участники 
  , если .

%
Все ошибки времени компиляции, связанные с преобладанием членов экземпляра
в разделе  Кроме того, вы можете переключаться между интерфейсами.

%
Если вышеприведенное правило вызовет несколько подписей членов
с таким же названием  Чтобы потом наследовать
Наследуется только один член, а именно:
Комбинированная подпись участника под названием ,
Прямой суперинтерфейс
% Это четко определено, потому что  Это интерфейс класса.
в текстовом порядке, в котором они декларируются,
В отношении 
(P000886).
Это ошибка времени компиляции
если вычисление упомянутой объединенной подписи участника не удается.


 Правильный член переопределяется


%
%
 и  подписи участников с таким же названием .
Далее   $m'$
Если удовлетворяются все следующие критерии:



 и  Это оба метода, оба геттера или оба сеттера.

Если  и  Оба являются геттерами:
Тип возврата  Должен быть подтип возвратного типа .

Если  и  Оба метода или оба сеттера:
 Тип функции 
За исключением того, что параметром типа является встроенный класс 
Для каждого параметра  который является ковариантным по декларации
(P000887).
 Тип функции .
 Это должно быть подтипом .

{%
Требование подтипа гарантирует, что список аргументов формирует
которые допустимы для вызова метода с подписью 
Допускается также вызов метода с подписью .
Например,  может принимать 2 или 3 позиционных аргумента;
 может принимать 1, 2, 3 или 4 позиционные аргументы, но не наоборот.
Это встроенное свойство правил подтипа функции.
?

% % TODO(Eernst): Приходите, это предупреждение удалено.
Если  и  Оба метода являются
 необязательный параметр ,
 является параметром  соответствует , имеет значение по умолчанию  и  имеет значение по умолчанию ,
 и  должны быть идентичными,
Происходит статическое предупреждение.


{%
Обратите внимание, что параметр, который является ковариантным по объявлению
должен иметь тип, удовлетворяющий еще одному требованию;
относительно соответствующих параметров во всех суперинтерфейсах,
прямые и косвенные
(P000888).
Мы не можем сделать это требование частью понятия правильных переопределений.
Потому что правильное переопределение касается только
Отношение к одному суперинтерфейсу.%
?


{Миксины}


%
Смесь описывает разницу между классом и его суперклассом.
Смесь либо получена из существующей классификации классов.
или вводится посредством смешанной декларации.
Это ошибка времени компиляции для получения смеси из
класс, который объявляет генеративного конструктора,
или из класса, который имеет суперкласс, отличный от P001005.

%
Применение миксина происходит, когда один или несколько миксинов смешиваются с
Классовая декларация через \WITH{} ().
Применение миксина может использоваться для расширения класса на раздел ;
Альтернативно, класс может быть определен как приложение для смешивания.
Как описано в следующем разделе.


{Смешанные классы}



<mixinApplicationClass> ::= \gnewline{}
<identifier> <typeParameters>? <mixinApplication> "

<mixinApplication> ::= <typeNotVoid> <mixins> <interfaces>


%
Ошибка времени компиляции, если элемент
Список типов  Положение о применении смеси является
переменная типа (),
тип функции (),
псевдоним типа, не обозначающий класс (),
Перечисленный тип (),
Отложенный тип (),
тип \DYNAMIC{} (),
Тип \VOID{} (),
или типа  >> для любого  (P000898).
Если  является типом в  пункт, {смешение}{тип! смешение
 является смесью, полученной из , если  обозначает класс,
или смеси, введенной , если  обозначает смешанную декларацию.

%
 быть заявлением класса смешанного применения формы



%
Ошибка времени компиляции  Перечисленный тип ().
Это ошибка времени компиляции, если какая-либо из  является перечисленным типом
().

%
Эффект  В библиотеке  Называется она «P000584». в
объем , связанный с классом () определяется пунктом
 \WITH{} , , 
с именем , как описано ниже.
Если  класс также реализует , , .
Если {}f декларация класса закрепляется встроенным идентификатором ,
Определяемый класс становится абстрактным классом.

%
Оговорка формы  {} , , 
с именем  определяет класс следующим образом:

%
Если существует только один миксин (), то \code {} 
определяет класс, полученный при применении смеси ()
Смешивание  () к классу, обозначаемому
 с именем .

%
Если существует более одного миксина (), то
 класс, определяемый по формуле { < , , 
с именем , где  Это новое имя, и сделать  Абстрактно. \WITH{} , ,  Определяет класс, полученный
посредством применения смеси смеси  к классу  с именем .

%
В любом случае, пусть  быть классной декларацией с
те же конструкторы, суперклассы, интерфейсы и экземпляры, что и
Определенный класс.
Это ошибка времени компиляции, если объявление  вызвать
Ошибка времени компиляции.
% TODO (Eernst): Не совсем!
% Мы не хотим супер-вызовов на ковариантных реализациях.
% - ошибки времени компиляции.

{%
Это ошибка, например, если  содержит декларацию члена 
который перекрывает подпись участника  в интерфейсе ,
но который не является правильным переопределением 
()%
?

 Смешанная декларация Срок


%
Смесь определяет ноль или более
{mixin Member Declaration}{mixin!member Declaration},
ноль или больше
{необходимый суперинтерфейс}{mixin!required superinterface},
один
{комбинированный} суперинтерфейс {mixin!combined superinterface}
ноль или больше
<{реализованные интерфейсы}{mixin!implemented interface}.

%
Смесь, полученная из классовой декларации:




 Требуемый суперинтерфейс
Комбинированный суперинтерфейс,
,   в виде реализованных интерфейсов
и, например, члены {члены} в виде смешанных деклараций.
Если  Это дженерик, как и смесь.

%
Декларация о смешивании вводит смешивание и обеспечивает объем
Статические декларации членов.


<mixinDeclaration> ::= \MIXIN{} <typeIdentifier> <typeParameters>
< ({} <typeNotVoidList>)? <Интерфейсы>.
\gnewline{}\{" (<metadata> <classMemberDeclaration>)*  "


%
Это ошибка времени компиляции, чтобы объявить конструктора в смешанной декларации.

%
Декларация о смешивании без  пункт является эквивалентным
в соответствии с пунктом \code{\ON{} Объект.

%
Пусть  будет \MIXIN Декларация формы



%
Ошибка времени компиляции, если какой-либо из типов  
или  это
переменная типа (),
тип функции (),
псевдоним типа, не обозначающий класс (),
Перечисленный тип (),
отложенный тип (),
тип < (),
тип < (),
или типа  для любого  ().

%
 быть интерфейсом, заявленным в декларации класса




где  - новое имя.
Это ошибка времени компиляции для декларации смешивания, если 
Классовая декларация вызовет ошибку компиляции,
{%
То есть, если какой-либо член объявлен более чем одним объявленным суперинтерфейсом,
И нет никакой конкретной подписи для этого члена среди супер.
интерфейсы %
}
Интерфейс  называется
 из заявления о смешивании .
{%
Если заявление о смешивании  имеет только один объявленный суперинтерфейс, 
интерфейс супервызова  имеют точно такие же члены
как интерфейс .%
?

%
 Интерфейс, который будет определяться декларацией класса




где  являются заявлениями членов
Декларация о смешивании  За исключением того, что все супервызовы лечатся
как если бы {} было действительным выражением со статическим типом .Ошибка времени компиляции для микширования  Если это  класс
Объявление вызовет ошибку компиляции времени,
{%
То есть, если необходимы суперинтерфейсы, реализованные интерфейсы
и декларации не определяют согласованный интерфейс,
если любое заявление члена содержит ошибку компиляции
кроме сверхпризыва,
или если супер-вызов не действителен против интерфейса .%
?
Интерфейс, представленный декларацией смешивания  имеет тот же член
Подписи и суперинтерфейсы .

%
Декларация о смешивании  Добавить микс
 требуется суперинтерфейс , , ,
{комбинированный суперинтерфейс} ,
\NoIndex{реализованный интерфейс}, , 
и члены экземпляра, объявленные в  как .


{Приложение миксина}


%
Для суперкласса может быть применена смесь, дающая новый класс.

%
 Будьте классом,
 Будьте смешаны с  требуется суперинтерфейс , , ,
{комбинированный суперинтерфейс} ,
{реализованные интерфейсы} , , 
\metavar{members} как \NoIndex{mixin members declarations},
Пусть P000666 будет именем.

%
Ошибка времени компиляции:    не осуществлять,
прямо или косвенно, все , , .
Это ошибка времени компиляции, если любой из 
супер-призыв члена 
{%
(например, , ,
или , %
?
 не имеет конкретного
Использование  который является действительным переопределением члена  в
интерфейс .
{%
Мы рассматриваем супер-призывы в миксах как
вызовы интерфейса на комбинированном суперинтерфейсе,
Поэтому мы требуем, чтобы суперкласс приложения для смешивания имел
действительные реализации этих членов интерфейса
На самом деле это супер-призыв.%
?

%
Применение микширования  до $S$ с именем  вводит новый
класс, , с именем , суперкласс ,
реализованный интерфейс 
и {members} в качестве членов экземпляра.
Класс  не имеет статических членов.
Если  объявляется любой генеративный конструктор, затем приложение
Внедрение генеративных конструкторов на  следующим образом:

%
Пусть  будет библиотекой, содержащей приложение микширования.
{%
То есть библиотека, содержащая пункт \code{$S$ < 
или пункт  << , \ldots, \ $M_k$,  >>
Появление приложения микширования.%
?

 Называется он «P000695».

Для каждого генеративного конструктора формы

 Это доступно для ,  иметь
неявно заявленный конструктор формы




где  Получено из  заменяя события ,
который обозначает суперкласс по  и  Получено из  через
Замена происшествий  который обозначает суперкласс по .
Если  Конструктор генеративных констов и  не объявляют ни одного
переменных экземпляров, P000714 также является конструктором.

%
Для каждого генеративного конструктора формы

 Это доступно для ,  иметьнеявно заявленный конструктор формы

<<<p000008>


где  Получено из  заменяя события ,
который обозначает суперкласс, по ,
 Получено из  заменяя события 
который обозначает суперкласс по ,
,  - постоянная оценка выражения
То же значение, что и P000738.
Если  Конструктор генеративных констов и  не объявлять ни одного
переменных экземпляров, P000741 также является конструктором.

%
Для каждого генеративного конструктора формы
 ,  , \{$T_{k+1}$$a_{k+1}$=$d_1$ \ldots,  $a_{k+n}$ = $d_n$\}]
 Это доступно для ,  иметь
неявно заявленный конструктор формы




где  Получено из  заменяя события 
который обозначает суперкласс по ,
 Получено из  заменяя события 
который обозначает суперкласс по ,
, , является постоянным выражением, оценивающим
То же значение, что и .
Если  Конструктор генеративных констов и  не объявляют ни одного
поля, P000768 также является конструктором.

%
Если  не декларирует никаких генеративных конструкторов, кроме приложения
не вызывает неявного стимулирования каких-либо экспедиторских конструкторов.
В частности, он не вызывает по умолчанию конструктора
().

{%
Конструктор по умолчанию будет иметь суперинициализатор
Это попытка вызвать конструктора, которого не существует.
Разработчики не смогут устранить эту ошибку.
?


 Расширения.


%
В этом разделе указаны расширения.
Этот механизм поддерживает декларирование функций.
которые аналогичны используемым методам,
аналогично неметодическим функциям в этом
Объявляются вне целевого класса,
Они решаются статически.
Резолюция основана на том, имеет ли соответствующее продление сферу действия.
и удовлетворяет ли призыв ряду других требований.


<extensionDeclaration> ::= <
{} <typeIdentifierNotType> <typeParameters> \ON{{<тип>>
\gnewline{}\{" (<metadata> <classMemberDeclaration>)*  "


%
Декларация, полученная из {продление Декларации}, известна как
{продление!декларация}.
Он вводит в действие

с данным {идентификатор}, если присутствует,
в пространство имен библиотеки, в которой находится
{(и, следовательно, в библиотечном объеме)},
и предоставляет возможность для объявления
членов-расширителей.
Дополнительно расширение вводится в пространство имен библиотеки
Свежее, личное имя.
Первый из них известен как
{объявленное имя}{расширение!объявленное имя}
от продления,
Последний известен как
{свежее имя}{расширение! свежее название
расширения.

{%
Также вводится новое название.
Пространство имен библиотеки текущей библиотеки
для каждого импортного расширения,
Даже если он импортируется с приставкой
(P000675).
%
Заявленное название расширения  это
введены в библиотечный объем текущей библиотеки
По тем же правилам, что и названия
Другие местные или импортные декларации, такие как классы.
%
Оригинальное название:  вводятся в этих случаях,
Но и в одном дополнительном случае:
когда происходит столкновение имени на объявленном имени .%
?

{%
Новое название позволяет использовать расширение
в неявном призыве
(),
Даже в том случае, если заявленное имя
или префикс импорта, обеспечивающий доступ к заявленному названию
оговаривается декларацией в промежуточном объеме,
Конфликт с именем, столкновение.%
?

{%
Тот факт, что расширение может быть использовано неявно даже в том случае, когда
не имеет объявленного имени
или объявленное имя затенено или противоречиво
отражает тот факт, что основное предполагаемое использование является неявным вызовом.
Хотя разработчик не может знать (и, следовательно, не может использовать) новое имя.
из данного расширения имплицитное обращение может использовать его.%
?

%
Ошибка времени компиляции, если библиотека имеет
Отсроченный импорт библиотеки 
Импортируемое пространство имен из  содержит
Имя, обозначающее расширение.

{%
Это означает, что импорт должен \HIDE{} или {} Для устранения
Наименования расширений от отложенного импорта.%
?

{%
Это ограничение гарантирует, что никакие расширения не будут введены.
использование отсроченного импорта,
что позволяет ввести семантику для таких расширений в будущем.
Без ущерба для существующего кода.%
?

%
{тип} в заявлении о продлении
называется расширением
\IndexCustom{\ON{} Тип {расширение!\ON{} Тип.
 Тип может быть любым действительным типом, включая переменную типа.

{%
Основная интуиция заключается в том, что расширение  может иметь
Тип {}  (определение типа получателя) и набора членов.
Если  выражение, статический тип которого 
 является членом, объявленным ,
 может ссылаться на указанный член со значением  Ссылка на .
Явно разрешенная форма  Доступен,
 Его можно использовать даже в случае
где  Выполняет какую-то другую функцию
Потому что некоторые другие расширения более специфичны.
Приведены подробности этих понятий, правил и механизмов.
в этом разделе и его подразделах.%
?

%
Заявленное название расширения не обозначает тип,
Можно использовать для обозначения самого расширения.
{%
(например, для доступа к статическим элементам расширения,или для того, чтобы разрешить вызов явно)%
}

%
Декларация о продлении вводит две сферы:



A {тип-параметр}{скоп!тип-параметр},
который является пустым, если расширение не является общим ().
Сфера охвата заявления о продлении по типу параметра включает
библиотечный объем текущей библиотеки.
Сфера действия параметра типа - это текущая сфера для
Параметры типа и для расширения  {тип}.

A {body scope}{scope!extension body}.
Включающая сфера охвата органа сфера действия заявления о продлении является
Сфера действия декларации о продлении типа.
В настоящее время сфера действия заявления о продлении членства является
сфера охвата прилагаемой декларации о продлении.


%
%
Пусть  будет заявлением о продлении с объявленным именем .
Декларация в  с модификатором {} обозначается как
{статический член}{расширение!статический член}
Декларация.
%
Декларация в  без модификатора \STATIC{} обозначается как
<{член организации}{расширение!член организации}
Декларация.
%
Ошибки в названии конфликтов могут возникать в 
В ситуациях, которые также происходят в классах и миксах
(P000678).
Кроме того, ошибка времени компиляции происходит в следующих ситуациях:


<<<p001102>  Объявляет участника, имя основания которого .
  Заявляет параметр типа .
  Объявляет участника, базовым именем которого является имя параметра типа
.
  объявляется экземпляром члена или статичным членом, имя основания которого
, , , ,
или {==}.
{%
То есть член, чье имя также является именем
Член экземпляра, который есть у каждого объекта.%
?
  объявляется застройщиком.
  объявляет переменную экземпляра.
  Объявляют абстрактным членом.
  Объявляет метод с формальным параметром
с модификатором .


{%
Абстрактные члены не допускаются, поскольку
Никакой поддержки для реализации.
%
Конструкторы не допускаются с момента продления
Не вводит никаких типов, которые могут быть построены.
%
Переменные не допускаются, поскольку память не выделяется.
для каждого объекта, доступного как {} в членах.
Разработчики могут эмулировать per-\THIS государство при необходимости,
Например, используя .
%
Члены с тем же именем базы, что и члены 
не допускаются, поскольку они могут быть вызваны только с использованием явного разрешения;
как в ,
Что может быть запутанным и подверженным ошибкам.%
?


 Явное прошение члена представительства о продлении


%
 E - простой или квалифицированный идентификатор
Это означает расширение.
{приложение расширения}{приложение расширения!
является выражением формы
 <typeArguments>? '(' <expression> ')'.
Член расширения с именем, доступным для текущей библиотеки
может быть явно использован на конкретном объекте
Выполняя призыв члена
()
где приемник является приложением расширения.

{%
Тип вывода еще не указан в этом документе,
Предполагается, что они уже имели место
(),
Ниже приводится описание предполагаемого лечения.
В этом разделе и его подразделах есть похожие комментарии о выводе типа.
Ниже, с пометкой «С выводом типа: .

 быть простым или квалифицированным идентификаторомРасширение под названием  и заявлены следующим образом:%
?




{%
Тип вывода для приложения расширения
в форме 
делается точно так же, как это было бы для
Тот же синтаксис, рассматриваемый как вызов конструктора
где  предполагается обозначать следующий класс,
и тип контекста пуст (без требований): %
?



{%
Это выведет аргументы типа для , и это будет
Введите тип контекста для выражения .
Например, если  объявляется как
\code{{) E<T> в Set<T> \{\,\ldots\,\}
 E(\{\})} будет содержать выражение \code{\{\}} с
тип контекста, который делает его буквальным.%
?

%
%
 быть расширением формы
< T}{1}{s}>($e$],
где  обозначает расширение, объявленное как


%
\EXTENSION{} E< <  \{\,\ldots\,.

%
Мы определяем
<обоснованно \ON{} {расширение!обосновано} \ON{} тип
 .
Мы определяем
{обоснование-к-связанности < Тип: {%}
Обсуждение:Instantiation-to-bound < тип
 
где {U}{1}{s} является результатом инстанциации в связанную
Типовые параметры 
(P000681).
%
Ошибка времени компиляции возникает, если
 T_j}{[T_1/X_1, \ldots, T_s/X_s]B_j},
<<<p000060>
{(то есть границы не могут быть нарушены)}.
%
Ошибка времени компиляции возникает, если статический тип  может быть присвоено
 Тип .
{%
Обратите внимание, что ошибка времени компиляции также происходит.
если статический тип  является 
()%
?

%
Это ошибка времени компиляции, если происходит расширение приложения.
в месте, где это {не} синтаксический приемник
простой или составной член вызов
(P000683).

{%
То есть, единственным допустимым использованием приложения расширения является
Призывать или отрывать от него членов.
Это похоже на то, как имена префиксов могут использоваться только в качестве
Цели призыва членов,
Кроме того, расширения могут декларировать операторы.
Например,  может быть действительным призывом к
Оператор {+} объявлено в расширении .%
?

%
Приложение расширения не имеет типа.

{%
Это согласуется с тем, что любое использование приложения расширения
где тип необходим - ошибка времени компиляции.
?

%
%
 Быть простым членом призыва
()
чей приемник  является расширением формы
%
< T}{1}{k}>($e$)}
\commentary{(который равен , когда  равен нулю)}
имя члена-корреспондента {n},
Предположим, что  Не имеет ошибок времени компиляции.
Ошибка времени компиляции возникает, если расширение не обозначено 
Объявляют члена по имени .
В противном случае пусть  X}{1}{k}}
Типовые параметры указанного расширения.
Пусть \DefineSymbol{s} будет подписью члена  Об этом сообщает .
Точно такие же ошибки времени компиляции происходят для  как
те, которые произошли бы для вызова члена 
который получен из  Заменить  через
переменной, тип которойКласс  в том же объеме, что и 
что объявляет члена  с подписью члена
:



%
Подпись участника  называется
\IndexCustom{подпись участника вызова}{%}
Расширение! подпись участника вызова.
.
Статический тип  является возвратным типом
подпись участника вызова .
 Например: }



{%
С выводом типа:
В случае, когда подпись участника вызова  является родовым,
Вывод типа происходит на  Точно так же, как это
Для вызова функции, тип которой является
Тип функции .%
?

%
%
Для динамической семантики,
 Быть простым, безусловным членом
чей приемник  является расширением формы
%
< T}{1}{k}> ()
где параметры типа  \DefineSymbol{\List{ X}{1}{k}}
и фактические значения \List{T}{1}{k} являются {{t}{1}{k}}
(),
%
и имя члена-корреспондента которого .
 быть членом  Называется он «P000102».
Оценка  доходов путем оценки
 для объекта ,
оценка и привязка фактических аргументов к формальным параметрам;
(),
и, наконец, выполнить 
в среде связывания, где \List{X}{1}{k} связаны с \List{t}{1}{k},
{} связано с ,
Каждый формальный параметр связан с соответствующим фактическим аргументом.
Значение  Это значение, возвращаемое вызовом .

%
Когда  является условным или составным призывом,
Статический анализ и динамическая семантика определяются
член призыва разбавлять
(P000687).

{%
Обратите внимание, что каскад ()
Приемник приложения расширения  представляет собой ошибку времени компиляции.
Это происходит потому, что подразумевается, что  обозначает объект, который не является истинным,
а также потому, что это заставит каждого 
для вызова члена того же расширения,
который вряд ли будет желательным.%
?


 Неявное прошение члена представительства о продлении


%
Вспомогательные элементы расширения могут быть вызваны или скрытно закупорены.
{(без упоминания названия расширения)},
Как если бы они были членами получателя
Призыв данного члена.

{%
Например, если  это
правильное явное обращение члена экземпляра  из
расширение ,
 может быть
Правильное неявное обращение с тем же значением.
Другими словами,
Получатель  вызова участника может быть неявно заменен
приложение расширения с приемником ,
если удовлетворяется ряд требований, описанных ниже.%
?

{%
Неявное вызов предназначен в качестве основного способа использования расширений.
с явным призывом в качестве запасного варианта
если неявное обращение является ошибкой, или
Неявное обращение решается в отношении члена
расширение, отличное от предполагаемого.%
?

%
%
Происходит неявное вызов члена-расширителя
Призыв к члену 
(перенаправлено с «P000689»)
с приемником  имя соответствующего члена  если {f}
(1)  не является буквальным,
(2) интерфейс статического типа  не имеет члена
имя основания которого является именем основания ,
3 существует уникальный, наиболее специфический
()Расширение, обозначаемое  который является доступным
()
и применимым
()
.

%
В случае, если ошибка компиляции не произошла,
 Обозначается как , который получен из  через
Замена ведущего  на .

{%
С выводом типа:
Когда  является родовым,
Типовой вывод, применяемый к  могут приводить реальные аргументы,
Получение  форма \code{$E$<\List{ T}{1}{k}>()}.
Если этот этап вывода типа не удается, то  не применяется
() %
?

{%
Неявное обращение члена инстанции о продлении
В каскаде также возможно
потому что каскад разбавляется до выражения, содержащего
одно или несколько обращений.%
?


 Доступность расширения.


%
Расширение  это
{доступно}{расширение!доступность}
в определенной сфере 
Если существует имя  Для лексического поиска  из 
()
доходность .

{%
Оригинальное название: P000149 может быть объявлено имя  или свежее название ,
Но так как свежее название всегда в объеме
Всякий раз, когда заявленное имя находится в пределах,
Достаточно вспомнить свежее название.
%
Когда новое имя  находится в библиотечном пространстве,
доступно в формате {любой},
потому что имя свежее и потому не может быть затенено
любым заявлением в любой промежуточной сфере.
%
Это означает, что если  Доступен в любом месте данной библиотеки 
Тогда он доступен везде в .%
?


 Применимость продления


%
%
Пусть  будет расширением.
 Быть членом призыва
()
с приемником  Статический тип 
и с именем соответствующего члена, базовое имя которого .
Мы говорим, что  это
{применимо}{расширение!применимо}
 Если соблюдены все следующие условия:



 Это не буквальный тип.

{%
Это означает, что призыв является вызовом члена экземпляра,
В том смысле, что  обозначает объект,
Таким образом, он может вызвать элемент экземпляра или быть ошибкой,
Это не может быть статический доступ.
Обратите внимание, что  Также не обозначает префикс или расширение,
И это не расширение,
Потому что у них нет типа.%
?

Тип  не имеет экземпляра с именем базы ,
 не является ни {}, ни .

{%
\DYNAMIC{{} и  Считается, что у них есть все члены.
Кроме того, это ошибка для доступа к члену на приемнике типа. <
(),
Расширения никогда не применяются к получателям
любой из типов ,  или .%
?

Для определения применимости расширения,
Типы функций и тип 
Считается, что у них есть член по имени .

<<%
Следовательно, расширения никогда не применимы к функциям.
когда базовое имя члена .
%
Участников, заявленных встроенным классом 
существуют на всех видах,
Для членов с такими именами расширение никогда не применяется.%
?
<<<p001202>
Рассмотрим приложение расширения  формы ,
где  Новая переменная со статическим типом .
Необходимо, чтобы возникновение в области, которая является текущей областью для 
Это не ошибка времени компиляции.

{%
Другими словами,  должен соответствовать  Тип .
С выводом типа могут быть добавлены выведенные фактические аргументы типа,
доходность \code{<\List S}{1}{k}>($v$],
Который не должен быть ошибкой.
Если этот шаг вывода не срабатывает, это не ошибка,
Это означает, что  Не применяется к .%
?

Расширение  Заявляет экземпляр члена с именем базы .


{%
С выводом типа:
Контекстный тип обращения не влияет
применимо ли продление,
и ни тип контекста, ни способ вызова не влияет
Типовой вывод ,
Но если сам метод расширения является общим,
тип контекста может повлиять на вызов члена.%
?


 Специфика расширения.


%
%
Когда  E}{1}{k} , являются расширениями
которые доступны и применимы к члену вызова 
(),
Мы определяем понятие
{специфичность}{расширение!специфичность}
который является частичным порядком на \List{E}{1}{k}.

{%
Специфика используется для определения того, какой метод расширения выполнить.
в ситуации, когда возможен более одного выбора.%
?

%
%
 быть участником вызова с получателем 
и соответствующее имя члена ,
%
и пусть  и  обозначает два различных
Доступные и применимые расширения для .
%
Пусть  быть инстанцирован  Тип  в отношении ,
 быть привязанным к инстанциации  Тип  в отношении ,

(P000698).
 Более конкретно, чем  В отношении 
если выполнено хотя бы одно из следующих условий:



 не указывается в системной библиотеке,
 Объявляется в системной библиотеке.

 и  Обоих заносят в системную библиотеку,
или ни один из них не объявлен в системной библиотеке, и


 \SubtypeNE{T_1}{T_2}, но не {T_2}{T_1}, или
 \SubtypeNE{T_1}{T_2} \SubtypeNE{T_2}{T_1} и
\SubtypeNE{S_1}{S_2}, но не \SubtypeNE{S_2}{S_1}.


{%
Другими словами, инстанциированный \ON{} тип определяет специфику,
и связанное с инстанциацией  Тип используется в качестве галстука
в случае, когда субтипирование не проводит различия между первым.%
?


{%
Следующие примеры иллюстрируют неявную резолюцию о продлении
когда доступно несколько расширений.%
?



<<%
Здесь применяются оба расширения,
 Более конкретно, чем  Потому что
\SubtypeNE{}{.%
?



{%
Здесь все три расширения относятся к обоим вызовам.
Для ,  Наиболее специфический,
 Является подтипом обоих
 и .
Следовательно, тип  является .

Для ,  является наиболее конкретным.
Презентация \ON{} Типы, которые сравниваются
  и
P000817 для двух других расширений.
Используя привязку к инстанциации \ON{} Типы как галстук,
Мы находим, что  менее точным, чем ,
Поэтому выбирается .Следовательно, тип  является .%
?

{%
Определение специфики направлено на выбор
Расширение, которое имеет более точную информацию о типе.
Это не обязательно приводит к наиболее точному типу результата.
(перенаправлено с «P000823») мог бы вернуться ,
Также важно, чтобы правило было простым.

На практике мы ожидаем, что непреднамеренные конфликты с именами участников расширения будут редкими.
Если автор предоставляет более специализированные версии
расширение для подтипов,
Выбор расширения, которое имеет наиболее точные типы
Скорее всего, это будет довольно неожиданное и полезное поведение.
?


 Статический анализ членов продления


%
Статический анализ деклараций членов в расширении 
опирается на объемы расширения
()
и соблюдает обычные правила, за исключением следующих:

%
При проведении статического анализа на теле экземпляра члена
расширение  с \ON{} тип ,
Статический тип .

%
Ошибка времени компиляции возникает, если тело члена расширения
содержит .

{%
Лексический поиск в расширении  может давать
Заявление метода экземпляра, заявленного в .
Как указано в другом месте
(),
Это означает, что члены экземпляра расширения
Теневые члены класса
при вызове из другого экземпляра члена внутри того же расширения
Использование вызова неквалифицированной функции
(перенаправлено с «P000825») и не ,
.
%
Это единственная ситуация, когда неявные призывы к
член расширения с именем базы 
может быть успешным, даже если интерфейс приемника имеет
Член с базовым именем .
%
С другой стороны, это согласуется с общим свойством Дарта, что
Лексически заключающие декларации теневые другие декларации, например,
Унаследованная декларация может быть скрыта глобальной декларацией.
Вот пример: %
?



{%
,
 принимает решение по декларации в ,
Если учесть, что  не имеет члена с базовым именем ,
Если нет других расширений, создающих конфликт.

Использование  в декларации 
не определяется в текущем лексическом разрезе,
Таким образом, он рассматривается как ,
потому что интерфейс  тип  иметь
 Геттер.

Использование  в \code{isOdd} Лексически разрешает
 выше него,
Таким образом, он рассматривается как ,
Даже если есть другие расширения в области
где  на .

Использование  в  Лексически разрешает
 выше него,
Несмотря на то, что есть член с тем же именем в
Интерфейс .
Геттер  нельзя назвать
в неявном призыве из любого места за пределами .
Это исключительный случай, упомянутый выше,
где член расширения тени
Постоянный член экземпляра на .
На практике продления очень редко вводят членов.
с тем же именем, что и член его < Тип интерфейса.

Неквалифицированный идентификатор  который не входит в сферу
Называется «P000849». Члены инспекции, как обычно
().
Если  не определяется статичным типом 
(тип {})
Это может быть заблуждением.
или может быть решена с использованием другого расширения.%
?


 Метод расширения Клозуризация


%
Метод расширенного экземпляра подлежит клозуризациианалогично методам классовых экземпляров
().

%
%
Ссылка на стр. Ссылка на стр. Заявка должна быть выполнена с помощью приложения
()
форма \code{$E$<\List{ S}{1}{m}>($e_1$)}.
%
 Y}{1}{m} - формальные параметры типа
Расширение .
%
%
Выражение  формы 
где  Если идентификатор известен как
{добыча имущества расширения}{%
Расширение! извлечение имущества.
Ошибка времени компиляции, если только  Заявляет о члене экземпляра по имени .
Если указанный экземпляр является методом, то
 имеет статический тип ,
где  Функциональный тип указанного способа декларирования.

{%
Если  является геттером  Это призыв Геттера,
который указан в другом месте
() %
?

%
Если  Это метод  Он известен как
<\IndexCustom{внутренний метод клозулизации}{%
Расширение! Пример метода клозурии
 на ,
и оценка 
\commentary{(который {$E$<\List{) S}{1}{m}>($e_1$).  ?
идет следующим образом:

%
Оценить  на объект .
 Будьте свежей конечной переменной, связанной с .
Далее  оценивает объект функции, эквивалентный:




где  объявляет параметры типа
,
требуемые параметры {p}{1}{n}
Названные параметры {n}{n}{n}{n}{n}{n}{n}}} с дефолтами  >>{1}{k}
использовать  для параметров, значение по умолчанию которых не указано.


где  объявляет параметры типа
,
требуемые параметры \List{p}{1}{n}
и факультативные позиционные параметры
\List{p}{n+1}{n+k) с дефолтами {d}{1}{k}
использовать  для параметров, значение по умолчанию которых не указано.


%
В функции буквы выше,
,
,
где  является типом соответствующего параметра в
Декларация .
Следует избегать захвата переменных типа в {S}{1}{m},
 должны быть переименованы, если {S}{1}{m} содержит любые случаи ,
Для всех .

{%
Другими словами, клозуляция является ценностью
функция, буквальная подпись которой совпадает с подписью ,
За исключением того, что фактические аргументы заменяются
Формальные параметры типа ,
Затем он просто передает призыв
 с захваченным объектом  как получатель.%
?

{%
Клосуризация двух расширений никогда не бывает равной
Если только они не идентичны.
Заметим, что это отличается от клозулизации классовых методов экземпляра,
которые равны, когда они отрывают один и тот же метод одного и того же приемника.
?

{%
Причина этого различия заключается в том, что даже если  и 
являются примерными методами клозулизации того же расширения 
применяется к тому же приемнику ,
могут иметь различные фактические типы аргументов, переданные ,
потому что аргументы этого типа определяются сайтом вызова
(и с выводом: по статическому типу выражения, дающего ),
И не только по свойствам  и вырванный метод.%
?

{%
Обратите внимание, что метод клозулизации экземпляра на расширении не является
Постоянное выражение,
Даже в том случае, когда приемник является постоянным выражением.
Это потому, что он создает новый объект функции каждый раз, когда он оценивается.
?

%Клозуризация методом расширения может происходить для неявного вызова
Член расширенного экземпляра.

{%
Это является следствием того факта, что неявное обращение
рассматривается как соответствующее явное обращение
().
Например,
 могут быть косвенно преобразованы в
,
которые затем обрабатываются, как указано выше.%
?

%
Клозуризация методом расширения подвергается инстанцированию общей функции
().
<<<p001280>> Например: }



{%
В этом примере \code{\{\}> напечатаны,
потому что функция объекта получена методом расширения клозуризации
подвергался инстанциации общей функции
В результате получился  значение ,
что делает оценку » ложной.%
?


  Участник продления.


%
Расширение может содержать {} метод, который используется неявно,
Аналогично вызову выражения функции
().

{%
Например,  
когда статический тип  Это нефункция, которая имеет метод под названием .
Вот пример, где  Метод происходит от расширения: %
?



<<<<<p001289>>
Это может показаться несколько удивительным,
Похожий подход с использованием :
\code{{} ( >>) i < 1[3])\,\,\{\,...\,\}}.
Мы ожидаем, что разработчики будут использовать эту мощность ответственно.
?

%
%
 быть приложением расширения
(),
 Выражение формы





{(где список аргументов типа опущен, когда  равен нулю)}.
 Затем рассматривается как
()


.

{%
Другими словами, вызов заявки на продление
немедленно рассматривается как вызов метода расширения, названного %
?

%
%
 выражение со статическим типом 
который не является выражением извлечения свойств
(),
 быть выражением формы





{(где список аргументов типа снова опущен, когда  равен нулю)}.
Если  , , или тип функции,
или интерфейс  имеет метод под названием ,
 указывается в другом месте
(P000712).
В противном случае, если  имеет неметодический экземпляр с базовым именем \CALL{}
 - ошибка времени компиляции.
В противном случае  Выражение  который является


.

{%
Обратите внимание, что  может быть неявным призывом к
метод расширения, названный ,
И это может быть ошибкой.
В последнем случае сообщения об ошибках должны быть сформулированы в терминах ,
не %
?

%
Это ошибка времени компиляции, если  является неявным призывом к
метод экземпляра расширения, названный .

{%
В частности,  Не может быть призывом к приемнику продления
Тип возврата - тип функции  или .%
?

{%
Обратите внимание, что нет поддержки для неявного извлечения имущества.
который отрывает метод расширения, названный .
Например, при условии расширения  в предыдущем примере: %
?



{%
Неявное извлечение собственности может быть разрешено,
Но это будет стоить читабельности.
Тип  Хорошо известен как не вызывающий,и неявный  Отрыв не будет иметь видимого синтаксиса.
Аргументы видны читателю,
Но для неявного отрыва от  функция,
Не существует видимого синтаксиса.%
?

{%
При желании извлечение имущества может быть выражено явно.
%
?


 Подсчеты


%
, или , используется для представления
фиксированное число постоянных значений.


<enumType> ::= \ENUM{} <идентификатор>
 >>{} <enumEntry> (',' <enumEntry>)* (',')? "

<enumEntry> ::= <metadata> <identifier>


%
Декларация перечня форм
 \ENUM{}  \{$m_0\,\,\id_0, \ldots,\ m_{n-1}\,\,\id_{n-1}$\}
имеет тот же эффект, что и классификация



{%
Это также ошибка времени компиляции для подкласса, смешивания или реализации перечня.
или в явном виде вводить в заблуждение.
Эти ограничения даются в нормативной форме.
в разделах , , 
 соответственно.%
?


 Дженерикс


%
Декларация класса (),
микширование (),
расширение (),
псевдоним (),
или функция ()  может быть ,
То есть  Могут быть объявлены формальные параметры типа.

%
Когда объект в этой спецификации описывается как общий,
и рассматривается особый случай, когда число аргументов типа равно нулю,
Список аргументов типа должен быть опущен.

{%
Это позволяет включать необычные случаи в качестве особых случаев.
Например,
вызов негенерической функции возникает как частный случай
где функция принимает аргументы нулевого типа,
Приводятся аргументы нулевого типа.
В этой ситуации некоторые операции также опускаются (не имеют эффекта).
операции, в которых формальные параметры типа заменяются
Типовые аргументы.%
?

%
A {общая декларация класса}{объявление класса}! общий
вводит общий класс в библиотечный объем текущей библиотеки.
A {общий класс}{класс!общий класс}
Это отображение, которое принимает список фактических аргументов типа.
и относит их к одному классу.
Рассмотрим типовую декларацию класса  Название: P000303
с официальными декларациями параметров типа
,
и параметризованный тип  формы .

%
Ошибка времени компиляции .
Ошибка времени компиляции  не является хорошо ограниченным
().

%
В противном случае указанный параметризованный тип  обозначает
Приложение общего класса, объявленное  Типовые аргументы
.
Это дает класс  члены которого эквивалентны членам
Классовая декларация, полученная из декларации  путем замены
каждое появление  по .

{%
Другие свойства  Например, отношения подтипа
указывается в других
()%
?

{%
Общие псевдонимы типа указаны в другом месте
() %
?

%
A {общий} Тип {type!generic} - тип, который вводится
декларация общего класса или псевдоним общего типа,
Или это тип .

%
A {объявление общей функции}{объявление функциональности!объявление общей функции}
Вводится общая функция () в существующий объем.

%
Рассмотрим функцию вызова выражения формы
,
где статический тип  Является общим типом функции
Формальные параметры типа
.
Ошибка времени компиляции .
Ошибка времени компиляции при наличии Например,  Не является подтипом .

{%
То есть, если количество аргументов неправильное,
или если аргумент  фактического типа не является подтипом
соответствующие границы,
где каждый формальный параметр типа заменяется
Соответствующий аргумент фактического типа.%
?


<typeParameter> ::= <metadata> <identifier> ({} <typeNotVoid>)?

<typeParameters>::= <<typeParameter> (>typeParameter>) "


%
Параметр типа  может быть суффиксирован с {} пункт
Это означает, что  для .
Если нет  Предложение присутствует, верхняя граница .
Это ошибка времени компиляции, если параметр типа является супертипом его верхней границы.
Когда верхняя граница сама по себе является переменной типа.

{%
Это не позволяет делать такие заявления, как
<< \EXTENDS{} X
и
{X} < Y, Y  X}%
?

%
Параметры типа объявлены в области параметров типа класса или функции.
Типовые параметры дженерика  являются по объему
- границы всех типовых параметров .
Типовые параметры типовой декларации классов  Кроме того, в сфере применения
\EXTENDS{} и {} Положения  (если они существуют)
и в теле .

{%
Однако параметр типа общего класса
Считается неполноценным типом
когда на него ссылается статический член
(P000726).
Объемы, связанные с параметрами типа общей функции
Описаны в (.%)
?

{%
Ограничение статических элементов необходимо, поскольку
переменная типа не имеет значения в контексте статического элемента;
Потому что статика делится между
Все генерические инстанциации генерического класса.
Однако переменная типа может ссылаться на инициализатор экземпляра,
Даже если  Недоступно.%
?

{%
Поскольку параметры типа находятся в пределах их объема,
Мы поддерживаем F-связанную количественную оценку.
Это позволяет код проверки типа, такой как:%
?



{%
Даже в тех случаях, когда параметры типа находятся в пределах
В настоящее время существует множество ограничений:


[]
Параметр типа не может использоваться для обозначения конструктора в
выражение создания экземпляра ().
[]
Параметр типа не может использоваться в качестве суперкласса или суперинтерфейса.
(, , .
[]
Параметр типа не может использоваться в качестве общего типа.


Приведены нормативные версии этих
в соответствующих разделах данной спецификации.
Некоторые из этих ограничений могут быть сняты в будущем.
?


 Разнообразие


%
Мы говорим, что тип   в виде  если {f}
 происходит в ковариантном положении в ,
но не в противоположном положении,
И не в инвариантном положении.

%
Мы говорим, что тип   в виде  если {f}
 происходит в противоположном положении в ,
но не в ковариантном положении,
И не в инвариантном положении.

%
Мы говорим, что тип   в виде  если {f}
 происходит в инвариантном положении в ,
 происходит в ковариантном положении, а также в противоположном положении.

%
Мы говорим, что тип  встречается в  в виде 
Если верно одно из следующих условий:


  

  Форма: 
где  обозначает общий класс происходит в ковариантном положении в  Для некоторых .

  имеет форму

где список параметров типа может быть опущен,
 происходит в ковариантном положении в .

  имеет форму



\code{($S_1x_1, \  S_k\ x_k
 S_{k+1}\ x_{k+1} = d_{k+1}, \ldots\  S_n\ x_n = d_n$]]


или в форме



\code{($S_1x_1, \  S_k\ x_k
 S_{k+1}\ x_{k+1} = d_{k+1}, \ldots\  S_n\ x_n = d_n$\}]


где список параметров типа и каждое значение по умолчанию могут быть опущены,
 происходит в противоположном положении в 
Для некоторых .

  Форма: 
где  обозначает общий тип псевдонимов, такой, что
,
Формальный параметр типа, соответствующий  является ковариантным,
 происходит в ковариантном положении в .

  Форма: 
где  обозначает общий тип псевдонимов, такой, что
,
Формальный параметр типа, соответствующий  является противоположным,
 Происходит в противоположном положении в .


<<<p001010>%
Мы говорим, что тип  встречается в  в виде 
Если верно одно из следующих условий:


  Форма: 
где  обозначает общий класс
 происходит в противоположном положении в 
Для некоторых .

  имеет форму

где список параметров типа может быть опущен,
 происходит в противоположном положении в .

  имеет форму



\code{($S_1x_1, , S_kx_k,
 S_{k+1}\ x_{k+1} = d_{k+1}, \ldots\  S_n\ x_n = d_n$]]


или в форме



\code{($S_1x_1, \  S_k\ x_k
 S_{k+1}\ x_{k+1} = d_{k+1}, \ldots\  S_n\ x_n = d_n$\}]


где список параметров типа и каждое значение по умолчанию могут быть опущены,
 происходит в ковариантном положении в 
Для некоторых .

  Форма: 
где  обозначает общий тип псевдонимов, такой, что
,
Формальный параметр типа, соответствующий  является ковариантным,
 происходит в противоположном положении в .

  Форма: 
где  обозначает общий тип псевдонимов, такой, что
,
Формальный параметр типа, соответствующий  является противоположным,
 происходит в ковариантном положении в .


%
Мы говорим, что тип  встречается в  в виде 
Если верно одно из следующих условий:


  Форма: 
где  обозначает общий класс или общий тип псевдонима,
 происходит в инвариантном положении в  Для некоторых .

  имеет форму



\code{($S_1x_1, \  S_k\ x_k S_{k+1}\ x_{k+1} = d_{k+1}, \ldots\  S_n\ x_n = d_n$]]


или в форме



\code{($S_1x_1, \  S_k\ x_k
 S_{k+1}\ x_{k+1} = d_{k+1}, \ldots\  S_n\ x_n = d_n$\}]


где список параметров типа и каждое значение по умолчанию могут быть опущены,
 происходит в инвариантном положении в 
для некоторых ,
 Происходит в 
Для некоторых .

  Форма: 
где  обозначает общий тип псевдонима,
,
Формальный параметр типа, соответствующий  является инвариантным,
 Происходит в .


%
Рассмотрим заявление о псевдонимах общего типа 
с официальными декларациями параметров типа
,
с правой стороны .
Пусть .
%
Мы говорим, что
 Формальный параметр типа  является инвариантным {%}
Типовой параметр! инвариантный
f  происходит инвариантно в , 
{является ковариантным}{параметр типа! ковариантный
f  происходит ковариантно в  и 
\IndexCustom{является контравариантным}{параметр типа!контравариантный}
f  Происходит противоположно в .

{%
Разница дает характеристику того, как тип варьируется
как значение субтермина варьируется, например, переменная типа:
Предположим, что  Тип, в котором переменная типа  происходит,
 и  Такие типы, как  Это подтип .
Если  происходит ковариантно в 
 Это подтип .
Точно так же, если  происходит противоположно в 
 Это подтип .
Если  происходит неизменно
 и  не гарантируется
быть подтипами друг друга в любом направлении.
Короче говоря: с ковариацией тип ковариаций;
с противопоставлением, типом противопоказаний;
с инвариантностью, все ставки выключены.%
?


 Суперсвязанные типы


%
В этом разделе описывается, как
Заявленные верхние границы формальных типовых параметров соблюдаются.
включая случаи, когда допускается ограниченная форма нарушения.

%
 Тип   Это подтип .
{%
Например, ,  и \VOID{} являются топовыми типами,
 и  Будущее или будущее Или <\DYNAMIC>{}>.%
?

% Мы определяем свойство быть регулярными для всех типов,
% сверхсвязанные для параметризованных типов и хорошо ограниченные
% для всех типов. Мы требуем, чтобы все типы были хорошо ограничены.
% охватывает каждый субтермин типа, который сам по себе является типом, и тогда мы
% требует, чтобы типы были регулярными при использовании в определенных
% ситуаций.

%
Каждый тип, который не является параметризованным, является .

{%
В частности, каждый необычный класс и каждый тип функций.
Регулярно ограниченный тип.%
?

%
 быть параметризованным типом формы

где  обозначает общий класс или псевдоним общего типа.
Давай.

Формальные декларации параметров типа .
  если {f}
 является подтипом
,
Для всех .

<<%
Это означает, что обычные типы являются такими типами.
которые не нарушают предельные значения параметров типа.%
?

%
 быть параметризованным типом формы
где  обозначает общий класс или псевдоним общего типа.
  Если оба условия истинны:



 не имеет регулярных ограничений.

 быть результатом замены каждого события в 
верхнего типа в ковариантном положении по ,
и каждое происшествие в 
 в противоположном положении по .
Затем требуется, чтобы  был регулярным.
%
Более того, если  обозначает общий тип псевдонима с телом ,
Необходимо, чтобы каждый тип, который возникает как субтермин

Он хорошо сложен (определяется ниже).


{%
Короче говоря, по крайней мере, один тип аргумента нарушает свои границы, но тип
регулярно ограниченный после замены всех случаев экстремального типа
противоположный экстремальный тип, в зависимости от их дисперсии.%
?

%
Тип   если {f}
Он либо регулярный, либо супер-ограниченный.

%
Любое использование типа  Что не очень хорошо ограничено, так это ошибка времени компиляции.

%
Ошибка времени компиляции при параметризированном типе  сверхограниченный
когда он используется любым из следующих способов:


  Непосредственная субтерминация нового выражения
()
Постоянное выражение объекта
(P000733).
  является немедленным субтермином перенаправляющего заводского конструктора
подпись
(P000734).
  является немедленным субтермином  пункт о классе
(),
или это происходит как элемент в списке типов  пункт
% % TODO(Eernst): Приходите с расширением, добавьте ref. Может быть, класс микширования?
(, , ,
или  пункт
(, ,
или это происходит как элемент в списке типов  пункт о смешивании
(P000741).


{%
Это ошибка , если возникает сверхограниченный тип.
как немедленный субтермин  пункт
который определяет границу переменной типа
()%
?

{%
Типы членов из сверхсвязанных классов вычисляются с использованием одного и того же
Правила как типы членов от других типов. Типы функциональных приложений
Сверхограниченные типы вычисляются по тем же правилам, что и типы
Функциональные приложения, включающие другие типы. Вот пример: %
?

<<<p000010>

{%
При этом  имеет статический тип ,
Даже если верхняя граница переменной типа  составляет .%
?

{%
Сверхограниченные типы позволяют выражать информативные общие супертипы.
некоторые типы, общие супертипы которых
В противном случае было бы гораздо менее информативным.
?

{%
Рассмотрим, например, следующий класс: %
?

<<<p000011>

{%
Без сверхограниченных типов,
Не существует типа  что делает  Общий супертип
Все виды формы 
(отмечая, что все виды должны быть регулярными)
Когда у нас нет понятия сверхограниченных типов.
Поэтому, если мы хотим, чтобы переменная содержала любой экземпляр типа  ''
Эта переменная должна использовать  или другой топ-тип
как своего рода аннотация,
что означает, что член, как  Не известно, существует ли
(что мы имеем в виду, говоря, что тип «менее информативный»).
?

<{%
Мы могли бы ввести понятие рекурсивных (бесконечных) типов и выразить
наименьшая верхняя граница всех типов формы  как
некоторый синтаксис, значение которого может быть приближено
.
%
Однако мы ожидаем, что любая такая концепция в Дарте
повлечет за собой значительные затраты
О разработчикам и реализацияхс точки зрения дополнительной сложности и тонкости,
Поэтому мы решили этого не делать.
Сверхограниченные типы конечны,
Но они предлагают полезное приближение, контролируемое разработчиком
для таких бесконечных типов.%
?

{%
Например,
 и
\code{C<C<\VOID>{}>
Это типы, которые разработчик может использовать в качестве аннотации типа.
Этот выбор служит обязательством
конечный уровень раскрытия бесконечного типа,
Это позволяет осуществлять определенный контроль
В точке, где разворачивание заканчивается:
%
Если  имеет тип \code{C<C\DYNAMIC>{}>>
 имеет тип <<
 не имеет ошибки времени компиляции,
Но если  имеет тип \code{C<C\VOID>{}>> Тогда уже
 - ошибка времени компиляции.
Таким образом, разработчики могут получить более или менее строгий режим.
выражения, тип которых выходит за пределы данного конечного развертывания.
?


 Обоснование границы.


%
В этом разделе описывается, как вычислять аргументы типа.
которые опущены из вида,
или от призыва к общей функции.

{%
% % TODO(Eernst): Когда мы добавим спецификацию типа вывода, мы будем корректировать
% - спецификация i2b, позволяющая принимать начальные значения для
% от фактического типа аргументов (то есть ) Дано "от звонящего"
% % для некоторых ; вероятно, «звонящий» всегда будет выводом типа; и
% Все остальные части алгоритма остаются неизменными.
% % Я думаю, что будет более запутанным, чем полезным начать писать это сейчас.
%% и это будет небольшая корректировка, когда мы добавим вывод типа. Так
На данный момент мы просто определяем i2b как автономный алгоритм.
Обратите внимание, что предполагается, что вывод типа уже имел место.
(),
Таким образом, аргументы типа не считаются опущенными, если они выведены.
Это означает, что инстанциация для связывания является механизмом резервного копирования,
который будет использоваться, когда нет информации для вывода.%
?

%
Рассмотрим ситуацию, когда термин  формы \synt
обозначает родовую декларацию типа,
Используется как тип или как выражение в прилагаемой программе.
{%
Это означает, что аргументы типа принимаются, но не предоставляются.
?
Мы используем фразу
 соответственно 
Для определения таких терминов.
В следующем мы упоминаем только сырые виды,
Но все говорили о сырых типах
Применяется к выражениям сырого типа очевидным образом.

{%
Например, с заявлением ,
Оценка экспрессии сырого типа  В теле уступает
Пример класса  ,
потому что  подлежит инстанциации в связанное.
Обратите внимание, что  не является синтаксически выражением,
Тем не менее, доступ к
 пример, подтверждающий 
без предварительного согласования,
потому что это может быть значение переменной типа.
?

{%
Мы можем однозначно определить сырые типы для обозначения
Результат применения генерического типа
к списку неявно представленных фактических аргументов типа,
Связывание — это механизм, который делает именно это.
Это связано с тем, что Дарт не поддерживает и не будет поддерживать более высокие типы.
Например, значение переменной типа  Это будет тип,
не может быть общим классом  как таковой,
и не может применяться к аргументам типа, например, .%
?

{%
В типичном случае, когда встречается только ковариация,
Инстанциация в связанное даст 
Все обычные типы, которые можно выразить.
Это позволяет разработчикам рассматривать сырой тип как тип.
который используется для указания того, что
«Актуальные аргументы типа не имеют значения» .%
?{%
Например, принимая декларацию
\code{{) C<X расширяет число \{\ldots\},
Включение в строку  доходность ,
Это означает, что  можно использовать для объявления переменной 
значение которого может быть  для \emph{все возможные} значения .%
?

{%
Рассмотрим ситуацию, когда
общий тип alias обозначает тип функции,
Он имеет один тип параметра, который является противоположным.
Обоснование, связанное с псевдонимом этого типа, затем даст 
всех типов, которые могут быть выражены
Разнообразие таких аргументов.
Это позволяет разработчикам рассматривать такой тип псевдонима, используемый в качестве сырого типа.
Тип функции, который позволяет передавать функцию клиентам.
где не имеет значения, какие ценности
для аргумента типа, который ожидает клиент .%
?
{%
Например, с
\code{\TYPEDEF{} F<X> = \FUNCTION(X)
Включение в строку  доходность ,
Это означает, что  можно использовать для объявления переменной 
чья ценность будет функцией, которая может быть передана клиентам
a  для \emph{все возможные} значения $T$.%
?


 Вспомогательные концепции для обоснования границ


%
Прежде чем мы определим моментализацию для связывания
Необходимо определить два вспомогательных понятия.
 Будьте сырым типом.
Тип  затем
{сырая зависимость на}{сырая зависимость на!тип}
 если выполняется одно или несколько из следующих условий:



 имеет форму \synt{typeName}, $S$ .
{%
Обратите внимание, что этот случай не применим
Если  является субтермином термина формы
 <typeArguments>,
То есть,
Если  Получает любые аргументы.
Также обратите внимание, что  не может быть переменной типа,
Потому что тогда ' не может удержаться.
См. обсуждение ниже и ссылку на ~
Подробнее о том, почему это так.%
?

 имеет форму {<typeName> <typeArguments>}
Один из типов аргументов зависит от .

 имеет форму \syntax{<typeName> <typeArguments>?} где
{typeName} обозначает псевдоним типа ,
и тело  Это зависит от .

 имеет форму
<<тип> \FUNCTION{} <typeParameters>? <parameterTypeList>
\syntax{<type>?}\ raw-depends on $T$,
или связанное в {<typeParameters>?}raw-depends на ,
или тип в {параметрTypeList} зависит от .


{%
Метаварианты
()
как  и  понимаются как обозначающие типы,
и они считаются равными (как в ' является $T$)
в том же смысле, что и в разделе о правилах подтипа
(P000746).
%
В частности,
Хотя два одинаковых синтаксиса могут обозначать два различных типа,
и две разные части синтаксиса могут обозначать один и тот же тип,
Свойство интереса здесь заключается в том, обозначают ли они один и тот же тип.
Не то, чтобы они были написаны одинаково.

Интуиция, стоящая за ситуацией, когда тип зависит от другого типа
Мы должны вычислить любые недостающие аргументы для последнего.
Чтобы понять, что означает первое.

В правиле о псевдонимах типа  может быть, а может и не быть,
Аргументы типа могут присутствовать или не присутствовать.Однако нет необходимости рассматривать результат замены.
фактические аргументы типа для формальных параметров типа в корпусе 
(или даже правильность передачи аргументов этого типа на ),
Потому что нужно только проверить
Все виды формы <<<Наименование типа>>
И на них не влияет такая замена.
Другими словами, сырая зависимость — это отношение.
Что легко и дешево вычислить.%
?

%
 быть общим классом или общим типом псевдонимов
 Официальные заявления о параметрах типа
содержащие формальные параметры типа  X}{1}{k} и границы  B}{1}{k}
Для любого 
Мы говорим, что формальный параметр типа  
когда выполняется одно из следующих требований:


  Опущено.

  включен, но не содержит ни одного из {X}{1}{k}.
Если  сырой зависит от сырого типа 
каждый параметр типа  Должен иметь простую границу.


%
Понятие простой связи должно толковаться индуктивно, а не
мондуктивно, т. е. если связан  общего класса или
общий псевдоним типа  достигается в ходе расследования, является ли
 Это простая граница, ответ - нет.

{%
Например, с
\code{{) C<X\EXTENDS{}C>\}
Параметр типа  не имеет простой границы:
Необработанный  используется в качестве границы для ,
 должны иметь простые границы,
Но одна из границ  является пределом ,
и эта граница , так что  должны иметь простые границы:
Это был цикл, поэтому ответ «нет».
 Не имеет простых границ.%
?

%
 быть общим классом или общим псевдонимом.
Мы говорим: 
{имеет простые границы}{тип! Общий, имеет простые границы.
f каждый параметр типа  Имеет простые границы.

{%
Теперь мы можем указать, в каком смысле инстанциация в связанное требует
Типы должны быть «достаточно простыми».
Мы накладываем следующее ограничение на все границы параметров типа:
потому что все параметры типа могут быть подвергнуты инстанциации в связанные.%
?

%
Это ошибка времени компиляции
Если формальный параметр типа связан  Содержит сырой тип ,
Если только  Имеет простые границы.

{%
Поэтому аргументы типа на границах могут быть только опущены.
Если они сами имеют простые границы.
В частности,
\code{\CLASS{) C<X\EXTENDS{}C>\}
Ошибка времени компиляции,
потому что связан  сырой,
и формальный параметр типа 
что соответствует опущенному типу аргумента
не имеет простой границы.%
?

%
 быть одним из видов формы <{название типа}
который обозначает общий класс или псевдоним общего типа
 Он сырой.
Тогда  эквивалентен параметризованному типу, который
результат, полученный путем применения инстанциации к привязке к .
Это ошибка времени компиляции, если инстанциация в связанный не удается.

{%
Это правило применимо ко всем случаям необработанных видов,
Например, когда это происходит как аннотация типа переменной или параметра,
в качестве функции возврата,
как тип, который тестируется в тесте типа,
как тип в  Часть.
и т.д.%
?


 Обоснование связанного алгоритма


%
Теперь уточним, как

Алгоритм продолжается.
 Будьте сырым типом.
Пусть {X}{1}{k} будут формальными параметрами типа в декларации ,
И пусть  B {1} {k} - их границы.Для каждого 
 обозначать результат инстанциации связанным на ;
В случае, когда граница  опущена, пусть  будет .

<<%
Если  Для некоторых  сырой (в общем: если сырой зависит от какого-то типа); 
Тогда все его (соответственно P000604) опущенные аргументы типа имеют простые границы.
Это ограничивает сложность инстанциации для ,
В частности, это не может быть связано с циклом зависимости.
где нам нужен результат от инстанциации, связанный с 
для того, чтобы вычислить инстанциацию, связанную с .%
?

%
Пусть  будет , для всех .
{%
Это "текущее значение" границы для переменной типа , на шаге 0;
В целом мы рассмотрим текущий шаг  и используем данные для этого шага.
Например, связанный , чтобы вычислить данные для шага , %
?

{%%} Область для определений отношений, используемых во время i2b

\def\Depends{\ensuremath{\rightarrow_m)
\def\TransitivelyDepends{{^{+}_m}}

%
Пусть {} будет отношением между переменными типа
\List{X}{1}{k) такой, что
  Происходит в .
{%
Таким образом, каждая переменная типа связана, то есть зависит от:
Каждая переменная типа в своей связи, которая может включать себя.
?
Пусть {} будет транзитивным
\commentary{(но не рефлексивно)}
Закрытие .
Для каждого , пусть $U_{i,m+1}$, для ,
определяется следующим итеративным процессом, где  обозначает
:


 [1]
Если существует  такой, что

{(то есть, если граф зависимостей имеет цикл)}
Пусть \List{M}{1}{p} будут сильно связанными компонентами (SCC)
в отношении .
{%
То есть максимальные подмножества  X}{1}{k)
где каждая пара переменных в каждом подмножестве
соотносятся в обоих направлениях по ;
Обратите внимание, что SCC попарно разъединены;
Кроме того, они однозначно определены вплоть до переупорядочения.
и порядок не имеет значения для этого алгоритма.
?
 быть объединением \List{M}{1}{p}
{(то есть все переменные, участвующие в цикле зависимости)}.
Пусть .
Если  не принадлежит   .
В противном случае существует  ;
 Получено из 
путем замены \DYNAMIC{} для каждого появления переменной в 
Находится в позиции  который не является контравариантным,
и заменить  для каждого появления переменной в 
Находится в противоположном положении в .

[2].
Иначе
{(при отсутствии циклов зависимости)},
<<<p000010> быть наименьшим числом, таким, что  встречается в  Для некоторых 
и  для всех  в 
{%
(то есть, граница ) не содержит переменных типа,
<<<p000018> возникает в связи с некоторой другой переменной типа )%
}
Для всех остальных (P000019).
 Получено из 
Заменить  Для каждого случая 
Находится в позиции  который не является контравариантным,
и замена  для каждого случая 
который находится в противоположном положении в .

[3].
Иначе
{(когда вообще нет зависимостей)},
Окончание с результатом .


{%
Этот процесс всегда заканчивается, так как общее количество
Возникновение переменных типа из  в
Текущие границы строго уменьшаются с каждым шагом, и мы заканчиваем.
когда это число достигает нуля.%
?

{%
Может показаться несколько произвольным рассматривать неиспользованные и инвариантные параметры.
аналогично ковариантным параметрам,
В частности, из-за неисправности инвариантных параметров
Чтобы оправдать ожидания, что
Необработанный тип обозначает супертип всех выразительных типов с регулярными границами.

Мы могли бы легко сделать все возможное, чтобы связать ошибку.
когда применяется к типу, где инвариантность возникает в любом месте
во время выполнения алгоритма.
Тем не менее, есть ряд случаев, когда этот выбор производит полезный тип.
И мы решили, что не стоит запрещать такие случаи.%
?



{%
Например, значение  может иметь динамический тип
.
Тем не менее, аргументы типа должны быть тщательно подобраны.
или результат не будет подтипом .
Например,  не может иметь динамический тип
{B<int, Inv<int>{}>}
 не является подтипом .%
?

%
Необработанный тип  является ошибкой времени компиляции, если инстанциация связана с 
Устанавливает тип, который не является хорошо ограниченным
(P000514).

<%
Такая ошибка может произойти, как показано на следующем примере:
?



{%
С этими заявлениями,
 В качестве аннотации типа используется ошибка времени компиляции:
Алгоритм дает  F<C\DYNAMIC{}>
И это не регулярный и не сверхограниченный тип.
%
Полученный тип может быть определен явно как
\code{C\DYNAMIC{} (C\DYNAMIC{})}.
Такой тип существует,
Мы просто не можем выразить это, передав аргумент типа ,
Таким образом, мы делаем ошибку, а не допускаем ее неявно.
?

{%
Основная причина, почему имеет смысл сделать такой сырой тип ошибкой
Дело в том, что не существует субтипа отношений.
между соответствующими параметризованными типами.%
?
{%
Например,  и  не связаны между собой,
Даже когда \SubtypeNE{\code{T1}}{\code{T2}> или наоборот.
На самом деле, нет типа  какой бы то ни было
переменная с заявленным типом 
может быть присвоен переменной типа
\code{C\DYNAMIC{} \FUNCTION(C\DYNAMIC{})}.
%
Таким образом, сырой , если это разрешено,
Это не было бы супертипом P000601. для всех возможных ,
Это будет тип, который не связан с 
{каждый сингл}  удовлетворяет пределу .
Это настолько бесполезно, что мы сделали ошибку.
?

%
Когда инстанциация в связь применяется к типу, она протекает рекурсивно:Для параметризированного типа < T}{1}{k}>
применяется к  T}{1}{k}
Для функционального типа
 T_0


и тип функции
 T_0
применяется к  T}{0}{n+k}.

{%
Это означает, что установление связи не влияет на
Тип, который не содержит сырых типов.
И наоборот, инстанциация к связанным действиям на типы, которые являются синтаксическими субтерминами.
когда они глубоко вложены.%
?

%
 Обоснование связи с общей функцией {%}
Общая функция! моментализация для связывания
использует алгоритм, описанный выше,
Взять формальные параметры  X}{1}{k) из декларации ,
с границами {B}{1}{k} и,
для каждого ,
Позволяет  обозначает результат инстанциации, связанный с ,
и позволить \DYNAMIC{} когда граница  опущена.

%
Пусть  будет родовой декларацией функции.
Если инстанция связана с  доходность
Список типов аргументов  T}{1}{k) такой,
для некоторых ,
 является или содержит тип, который не является хорошо ограниченным,
или если  T}{1}{k} не удовлетворяет границам
по формальным параметрам типа ,
Мы говорим, что
 не имеет аргументов типа по умолчанию {%}
Общая функция! не имеет аргументов типа по умолчанию.


{Метадата}


%
Dart поддерживает метаданные, которые используются для прикрепления
Пользователь определяет аннотации к структуре программы.


<metadata> ::= ('@' <metadatum>)*

<metadatum> ::= <
<identifier> | <qualifiedName> | <constructorDesignation> <аргументы>


%
Метаданные состоят из серии аннотаций,
Каждый из них начинается с персонажа. {@}
За ним следует постоянное выражение  извлекаемый из
{metadatum}.
Ошибка времени компиляции  не является одним из следующих:


 Ссылка на постоянную переменную.
 Призыв к постоянному конструктору.


{%
Выражение  происходит в постоянном контексте
(),
Это означает, что  Модификаторы не должны быть указаны явно.%
?

%
Метаданные, которые являются первым элементом
Правая сторона грамматического правила
связано с абстрактным синтаксическим деревом для нетерминального
на левой стороне грамматического правила
{(то есть он связан с его родителем)}.
В противном случае метаданные связаны с абстрактным деревом синтаксиса.
Конструкция программы  что сразу следует за метаданными
Правило грамматики.

{%
Эти правила призваны обеспечить минимальный уровень согласованности.
в ассоциации, которая связывает каждый метадатум с конструкцией программы.
Структура абстрактного дерева синтаксиса специфична для реализации.
И даже не гарантируется, что инструмент будет использовать что-либо.
Это называется абстрактным синтаксическим деревом.
В этом случае выражение "абстрактное синтаксическое дерево" следует толковать как
Субъекты представления программ, которые фактически используются.

Это означает, что тонкие детали связи между метаданными
Абстрактный узел дерева синтаксиса также является специфическим для реализации.
В частности, реализация может выбрать ассоциацию данного метадатума.
с более чем одним абстрактным узлом дерева синтаксиса.%
?

%
Метаданные могут быть получены во время выполнения с помощью отражающего вызова,
Предоставляется аннотированная конструкция программы  доступен через отражение.

{%
Метаданные также могут быть получены статически.
Анализ программы и оценка констант с помощью подходящего интерпретатора.
На самом деле, многие, если не большинство применений метаданных полностью статичны.В этом случае связывание каждого метадатума с абстрактным синтаксическим узлом дерева
определяется данным инструментом статического анализа,
которые, разумеется, не подпадают под какие-либо ограничения настоящего документа.
Конечно, будет полезно стремиться к согласованности всех инструментов.
относительно связывания метаданных с абстрактными узлами синтаксиса.%
?

<%
Важно, чтобы накладные расходы по времени выполнения не были понесены
Введение метаданных, которые фактически не используются.
Поскольку метаданные содержат только константы,
Время, в которое он вычисляется, не имеет значения.
Таким образом, реализация может пропустить метаданные.
при обычном разборе и исполнении,
и оценивать его лениво.%
?

{%
Метаданные можно связать с конструктами
Это может быть недоступно через рефлексию.
Например, локальные переменные
(хотя можно предположить, что в будущем,
Более богатые отражающие библиотеки также могут обеспечить доступ к ним.
Это не так бесполезно, как может показаться.
Как отмечалось выше, данные могут быть получены статически.
Если исходный код доступен.%
?

%
Метаданные могут появляться перед библиотекой, заголовком части, классом,
типдеф, параметр типа, конструктор, завод, функция,
параметр или переменная декларация,
перед импортом, экспортом или частичной директивой.

%
Постоянное выражение, данное в аннотации, проверяется и оценивается по типу.
в объеме, окружающем аннотируемое заявление.


 Выражения


%
 Это фрагмент кода Дарта.
Это можно оценить во время выполнения.

%
Каждое выражение имеет ассоциированный статический тип ().
может иметь ассоциированный статический тип контекста
% % TODO(Eernst): Это не определено до тех пор, пока CL 92782 не приземлится.
% (),
Это может повлиять на статический тип и оценку выражения.
Каждый объект имеет соответствующий динамический тип (P000518).

% % TODO(eernst): <expression> следует вывести <functionExpression> Кроме того,
% от того, что в настоящее время получено из <первичного> Побуждает к подлинному
%% и серьезный источник двусмысленности, что затрудняет разбор Дарта
% с использованием чего-либо, кроме модифицированного рекурсивного парсера спуска.


<expression> ::= <assignableExpression> <assignmentOperator> <expression>
 <условное выражение>
 <cascade>
 <rowExpression>

<expressionWithoutCascade> ::= <
<assignableExpression> <assignmentOperator> <expressionWithoutCascade>
 <условное выражение>
 <rowExpressionWithoutCascade>

<expressionList> ::= <expression> (',' <expression>)*

<primary> ::= <Это выражение>
 {} <unconditionalAssignableSelector>
 \SUPER{{}<аргумент
 <выражение функции>
 <буквально>
 <идентификатор>
 <newExpression>
 <constObjectExpression>
<constructorInvocation>
 <<выражение>> "

<literal> ::= <nullLiteral>
 <booleanLiteral>
 <нумерация
 <stringLiteral>
 <symbolLiter>
 <listLiteral>
 <setOrMapLiteral>


%
Выражение  всегда могут быть заключены в скобки,
Но это никогда не оказывает семантического влияния на .

{%>
Однако это может оказать влияние на окружающую экспрессию.
Например, для класса  со статическим методом
,  возвращается 42.
 - ошибка времени компиляции.
%
Дело в том, что значение 
определяется в нескольких частях,
Вместо того, чтобы быть строго композиционным.
Конкретно, значение  и \code{(C)} Как выражения одинаковы,
но значение  не определяется в терминахЗначение  как выражение,
и отличается от значения .%
• строго композиционная оценка всегда оценивала бы каждый
% субэкспрессия с использованием одних и тех же правил («оценка» всегда одно и то же);
%, а затем объединит результаты оценки в результат
% от общего выражения. Мы не будем расширять это здесь, и в частности мы не будем
% обсудить, имеет ли композиционная оценка смысл в контексте
% побочных эффектов. Но все же полезно помнить, что у нас есть эти
% «высоко некомпозиционные» элементы в семантике, такие как статические
% метод поиска.
?


 Оценка выражения


%
Оценка выражения либо
{производит объект}{выражение! производит объект
или это
<{броски}{выражение!броски}
Исключительный объект и связанный с ним след стека.
В первом случае мы также говорим, что выражение
{оценивает объект}.

%
При оценке одного выражения ,
определяется в терминах оценки другого выражения ,
обычно подвыражение ,
и оценки  Отбросьте исключение и след стопки,
Оценка  остановиться в этой точке
и бросает тот же самый объект исключения и след стека.


 Идентификация объекта


%
Предопределенная функция Дарта 
Определено так, что  Если [f]:


  Оценка нулевого объекта ()
или экземпляра  и , или
  и  являются экземплярами  и , или
  и  постоянные строки и , или
  и  Примеры 
и один из следующих пунктов:


  и  ненулевой и .
<<<p001010>  и  являются .
 Как , так и  являются .
  и  представляет собой значение NaN
с тем же базовым битовым рисунком.

или
  и  Постоянные списки, которые определяются как идентичные.
в спецификации буквальных выражений списка (), ИЛИ
  и  Постоянные карты, которые определены как идентичные
в спецификации буквальных выражений карты (), ИЛИ
  и  являются постоянными объектами одного класса 
и значение каждой переменной экземпляра  тождественны
значение соответствующей переменной экземпляра . или
  и  Это один и тот же объект.


{%
Определение  Двойники отличаются от равенства
в том, что NaN идентичен самому себе,
и что отрицательный и положительный ноль различны.%
?

{%
Определение равенства для двойников продиктовано стандартом IEEE 754.
Они утверждают, что NaN не подчиняются закону рефлексивности.
Учитывая, что оборудование реализует эти правила,
Их необходимо поддерживать по соображениям эффективности.

Определение идентичности не ограничено таким же образом.
Вместо этого предполагается, что бит-идентичные двойники идентичны.

Правила идентификации делают невозможным для программиста Дарта
наблюдать, является ли булевое или числовое значение упакованным или неупакованным.%
?


<\subsection{Константы)


{%
Все использования «постоянного» в Дарте связаны со временем компиляции.
Потенциально постоянное выражение — это выражение, которое обычно приводит кпостоянное значение при задании значений определенных параметров.
Постоянные выражения — это подмножество потенциально постоянных выражений.
что \emph{может} оцениваться в момент компиляции.%
?

{%
Постоянные выражения ограничиваются выражениями, которые
выполнять только простые арифметические операции, булевы условия;
Творчество и творчество.
Функциональное тело пользователя не выполняется
при постоянной оценке выражения,
Только члены системных классов
, , , ,
, ,  или .%
?

%
{потенциально постоянные выражения}{%
Потенциально постоянное выражение
и \IndexCustom{постоянные выражения}{постоянное выражение}
являются следующие:



Буквальный булев, \TRUE{} или \FALSE{} (),
Это потенциально постоянное и постоянное выражение.


Буквальное число () это
Потенциально постоянное и постоянное выражение
если он оценивает тип  или .
Слишком большое целое число буквально не оценивает объект.

 Буквальная строка () с интерполяцией строк
()
с выражениями , , <<p000096> потенциально постоянное выражение
Если , ,  Это потенциально постоянные выражения.
Буквальное — это постоянное выражение
Если , ,  Постоянные выражения
Оценивая экземпляры , ,
,  или .
{%
Эти требования носят тривиальный характер
Если в строке нет интерполяций.%
?
{%
Было бы заманчиво разрешить интерполяцию струн там, где
интерполированное значение - любая постоянная времени компиляции.
Однако для этого потребуется запустить  метод
для постоянных объектов,
может содержать произвольный код.%
?


Буквальный символ () это
потенциально постоянное и постоянное выражение.


Буквальный  () это
потенциально постоянное и постоянное выражение.


Идентификатор, обозначающий постоянную переменную, является
потенциально постоянное и постоянное выражение.


Квалифицированная ссылка на статическую постоянную переменную
()
который не соответствует отложенному префиксу,
Это потенциально постоянное и постоянное выражение.
{%
Например, если класс  объявляет постоянную статическую переменную ,
 Это постоянная.
То же самое верно, если  Доступ осуществляется через префикс ;
 является постоянной, если  Отложенный префикс.%
?


простой или квалифицированный идентификатор, обозначающий класс;
микширование или псевдоним типа, который не квалифицирован отложенным префиксом,
Это потенциально постоянное и постоянное выражение.
{%
Постоянное выражение всегда оценивается как  объект.
Например, если  Это имя класса или типа псевдонима,
Выражение  является постоянным,
Если  импортируется с приставкой ,  это
Константа  Пример, представляющий тип 
Если только  Отложенный префикс.%
?


 быть простым или квалифицированным идентификатором
Функция верхнего уровня ()
Статический метод ()
Это не квалифицируется отложенным префиксом.
Если  не подвержена генерической функции
()
P000121 является потенциально постоянным и постоянным выражением.
Если инстанциация общей функции применяется к и приведенные фактические аргументы типа  T}{1}{s}
 потенциально постоянное и постоянное выражение
если {}f каждый , является выражением постоянного типа
().


Выражение идентификатора, обозначающее параметр константного конструктора
()
что происходит в списке инициализаторов конструктора,
Это потенциально постоянное выражение.


Постоянное выражение объекта ()
Это потенциально постоянное и постоянное выражение.


Постоянный список буквальный ()
Это потенциально постоянное и постоянное выражение.


Постоянный набор буквальный ()
Это потенциально постоянное и постоянное выражение.


Постоянная карта ()
Это потенциально постоянное и постоянное выражение.


Выражение скобки  это
Потенциально постоянное выражение
Если  Это потенциально постоянное выражение.
Это также постоянное выражение, если  Это постоянное выражение.


Выражение формы  это
Потенциально постоянное выражение
Если  и  Потенциально постоянные выражения
 Статически связан с
Предопределенная функция dart  обсуждавшееся выше
(P000538).
Кроме того, это постоянное выражение
Если  и  Это постоянные выражения.


Выражение формы  это
эквивалентно  во всех отношениях,
В том числе, является ли она потенциально постоянной или постоянной.


Выражение формы  потенциально постоянный
Если  и  Оба эти выражения потенциально являются постоянными.
Он также постоянен, если оба  и  являются постоянными и
 Оценивается на пример 
Или пример, который имеет примитивное равенство.
(),
 Оценка нулевого объекта
().


Выражение формы  потенциально постоянный
Если  Потенциально постоянный.
Постоянна, если  Это постоянное выражение, которое оценивает
Пример типа .


Выражение формы  это
Потенциально постоянный, если  и 
Оба эти выражения потенциально являются постоянными.
Постоянна, если  является постоянным выражением и

 Для сравнения: , или
  \TRUE{} и  Это постоянное выражение
Оценивается на примере типа .



Выражение формы  это
Потенциально постоянный, если  и 
Оба эти выражения потенциально являются постоянными.
Постоянна, если  является постоянным выражением и

 (b) оценивает до , или
  \FALSE{} и  Это постоянное выражение
Оценивается на примере типа .



Выражение формы  это
Потенциально постоянное выражение
Если  Это потенциально постоянное выражение.
Это также постоянное выражение, если  это
Постоянное выражение, оценивающее пример типа 
Таким образом, {} обозначает вызов оператора экземпляра.

 выражение одной из форм ,
 или  потенциально постоянный
 и  Оба эти выражения потенциально являются постоянными.
Он также постоянен, если оба  и  Это постоянные выражения, которыеоба оцениваются по экземплярам ,
или оба экземпляра ,
Чтобы оператор символизировал
 \lit{ |}, соответственно {}
обозначает вызов оператора экземпляра.

 Выражение одной из форм
, ,
 потенциально постоянный
Если  и  Оба эти выражения потенциально являются постоянными.
Он также постоянен, если оба  и $e_2$ Это постоянные выражения, которые
оба оценивают на пример ,
Чтобы оператор символизировал
 \lit{}, соответственно {}
обозначает вызов оператора экземпляра.

 Выражение формы  это
Потенциально постоянное выражение, если  и 
Оба эти выражения потенциально являются постоянными.
Кроме того, это постоянное выражение
 и  Постоянные выражения
и либо оба оценивают на пример  или ,
или оба оценивают на пример ,
{+} обозначает вызов оператора экземпляра.


Выражение формы  потенциально постоянное выражение
Если  Это потенциально постоянное выражение.
Это также постоянное выражение, если  Это постоянное выражение, которое
Оценка для экземпляра типа  или ,
Таким образом, {-} обозначает вызов оператора экземпляра.

 Выражение формы , ,
, , ,
, , , или

потенциально постоянный
Если  и  Оба эти выражения потенциально являются постоянными.
Он также постоянен, если оба $e_1$ и $e_2$ Это постоянные выражения, которые
оценивать по экземплярам  или ,
что данный символ оператора обозначает
Воззвание оператора экземпляра.

 Выражение формы 
Потенциально постоянная, если ,  и 
Все это потенциально постоянные выражения.
Постоянна, если  является постоянным выражением и

  {} и  является постоянным выражением, или
  {} и  является постоянным выражением.


 Выражение формы  потенциально постоянный
Если  и  Оба эти выражения потенциально являются постоянными.
Он постоянен, если  является постоянным выражением и

  оценивает объект, который не является нулевым объектом, или
  оценивает нулевой объект и  Это постоянное выражение.


 Выражение формы  потенциально постоянный
Если  Это потенциально постоянное выражение.
Постоянна, если  Это постоянное выражение, которое
Для примера ,
 Обозначает вызов Геттера.

 Выражение формы  потенциально постоянный
Если  Это потенциально постоянное выражение.
Постоянна, если  Это постоянное выражение, которое
оценивает до нуля, или оценивает до экземпляра 
 Обозначает вызов Геттера.

 Выражение формы  потенциально постоянный
Если  потенциально постоянное выражение
 Постоянное выражение типа,и она является постоянной, если  Постоянная.
{%
Ошибка времени компиляции для оценки постоянного выражения
Если бы броська была, то есть
Если  Оценивает объект, который не является нулевым объектом
и не по типу %
?

 Выражение формы  потенциально постоянный
Если  потенциально постоянное выражение
 Постоянное выражение типа,
и она является постоянной, если  Постоянная.


Выражение формы 
эквивалентен  во всех отношениях,
В том числе, является ли она потенциально постоянной или постоянной.


%
А.

является одним из:



Простой или квалифицированный идентификатор
обозначение типовой декларации (типовой псевдоним, класс или смешанная декларация)
который не соответствует отложенному префиксу,
необязательно с последующими аргументами типа формы

где , ,  Это постоянные выражения типа.

Тип формы 
где  - выражение постоянного типа.

% % TODO(Eernst): Это не позволяет вводить переменные типа
% к самому типу. «Функция <X>(X)» может быть выражением постоянного типа,
%, но это не покрывается действующими правилами: "X" - переменная типа,
% и никогда не допускается.
Тип функции
 Функция <\metavar{typeParameters}>(\metavar{argumentTypes})}
(где $R$ и {<\metavar{typeParameters}>} можно опустить)
где , \metavar{typeParameters} и {argumentTypes}
(при наличии) содержат только постоянные выражения типа.

Тип .

Тип .


% Быть потенциально постоянным — это структурно, а не по типу.
%, но программа по-прежнему должна удовлетворять типизации в сильном режиме.

% Постоянные выражения (например, «const Foo(42)») всегда оцениваются как
% одно и то же значение, причем не более одного значения на местоположение источника.
% Потенциально постоянные выражения, которые не являются постоянными
% допускают простые операции по основным типам (нум, струна, бул, нуль). Они могут
% вычисляется статически без использования пользовательского кода.

Валидно типизированное потенциально постоянное выражение все еще может потерпеть неудачу при оценке.
% Если это происходит при вызывании const, это ошибка времени компиляции.

%
Ошибка времени компиляции, если требуется выражение
Постоянное выражение,
Но его оценка выльется в исключение.
Ошибка времени компиляции, если утверждение оценивается как часть
постоянная оценка выражения объекта,
И это утверждение станет исключением.

%
Ошибка времени компиляции, если значение постоянного выражения
зависит от самого себя.

{%
В качестве примера рассмотрим: %
?




 Дополнительные замечания о константах и потенциальных константах


{%
Нет требования, что
Каждое постоянное выражение оценивается правильно.
Когда требуется постоянное выражение
(например, для инициализации постоянной переменной,
или в качестве значения по умолчанию формального параметра,
или в виде метаданных
Мы настаиваем на том, что постоянное выражение
успешно оцениваться во время компиляции.%
?

{%
Вышеизложенное не зависит от программы управления потоком.
Простое наличие необходимой постоянной времени компиляции
Оценка, которая провалится в программе, является ошибкой.
Это также рекурсивно:
Поскольку константы соединений состоят из констант,
Если какая-либо часть константы выбрасывает исключение при оценке,
Это ошибка.
С другой стороны, поскольку реализации могут свободно компилировать код с опозданием,
Некоторые ошибки времени компиляции могут проявляться довольно поздно: %
?



<{%
Реализация свободна для немедленного выпускаОшибка компиляции ,
Но это не обязательно делать.
Он может отсрочить ошибки, если не будет немедленно компилироваться.
Деклараций, которые ссылаются на .
Например, он может задержать предоставление ошибки компиляции.
О методе  до первого вызова .
Однако он не мог выбрать выполнение ,
см. ветвь, которая относится к  не принимается,
2 Успешное возвращение.

Положение в связи с призывом  Это другое.
 не является константой времени компиляции (несмотря на то, что ее значение равно);
При составлении  не нужно давать ошибку времени компиляции.
Реализация может запускать код,
Это приведет к тому, что будет использован геттер для .
В этот момент инициализация  должны иметь место,
который требует компилирования инициализатора,
который вызывает ошибку компиляции.%
?

{%
Обработка  заслуживает некоторого обсуждения.
Рассмотрим \code{\NULL{} +2.
Это выражение всегда вызывает ошибку.
Мы могли бы не рассматривать его как постоянное выражение.
(и вообще не допускать ) как
Субэкспрессия числовых или булевых постоянных выражений.
Есть два аргумента в пользу его включения:
Во-первых, он постоянен, поэтому мы можем оценить его во время компиляции.
Во-вторых, представляется более полезным дать
ошибка, вытекающая из оценки явно.%
?

{%
Можно обоснованно спросить, почему
 \,\,$e_2$\,:$e_3$ и \,??? 
имеют постоянные формы.
Если  Статически известно, зачем его тестировать?
Ответ заключается в том, что существуют контексты, где  является переменной,
Например, в постоянных инициализаторах конструктора, таких как:
\code{\CONST{} C(foo):\ \THIS.foo = foo??? someDefault Стоимость: %
?

{%
Разница между
потенциально постоянное выражение и постоянное выражение
заслуживает некоторых объяснений.
Ключевой вопрос заключается в том, как относиться к формальным параметрам конструктора.

Если константный конструктор вызывается из постоянного выражения объекта,
Фактические аргументы должны быть постоянными выражениями.
Если бы мы были уверены, что
Константные конструкторы всегда вызывались из постоянных выражений объектов.
Можно предположить, что формальные параметры конструктора
Константы времени компиляции.

Постоянные конструкторы также могут быть использованы для
выражения для создания обычных экземпляров (),
Таким образом, приведенное выше предположение не является в целом обоснованным.

Тем не менее, использование формальных параметров константного конструктора
имеет значительную полезность.
Вводится понятие потенциально постоянных выражений для облегчения
ограниченное использование таких формальных параметров.
В частности, мы разрешаем использование
Формальные параметры константного конструктора
для выражений, в которых участвуют встроенные операторы,
Но не для постоянных объектов, списков и карт.
Например:%
?



{%
Задание на  допускается при допущении
 постоянный
(перенаправлено с «P000720») Как правило, это не постоянная времени компиляции.
Задание на  аналогично, но вызывает дополнительные вопросы.
В этом случае суперэкспрессия  является ,
Для этого необходимо, чтобы  быть числовым постоянным выражением
Все выражение должно считаться постоянным.
Формулировка спецификации позволяет предположить
 оценивает в целое число,
Для вызова этого конструктора в постоянном выражении.
Аналогичные аргументы приводятся для  и \code{q}
в задании .
Однако не допускаются следующие конструкторы: %
?



{%
Проблема в том, что все это противоречит правилам
постоянные списки (),
Карты (),
и объекты (),
все из которых независимо требуют, чтобы их подэкспрессии были
Постоянные выражения.%
?

{%
Все нелегальные конструкторы P000729 выше
Не может быть разумно вызван через ,
выражение, которое должно быть постоянным, не может зависеть от
Формальный параметр, который может быть или не быть постоянным.
В отличие от этого, юридические примеры имеют смысл независимо от того,
Призывается ли застройщик через \CONST{} или через .

Внимательные читатели, конечно, будут беспокоиться о случаях, когда
Фактическая аргументация в пользу  Они являются константами, но не числовыми.
Это исключается правилами о постоянных строителях.
(),
в сочетании с правилами оценки константных объектов ().%
?


 Постоянный контекст


%
 быть выражением;  происходит в

если применяется одно из следующих положений:

% Мы избегаем замкнутости "постоянный контекст зависит от постоянного списка буквальных,
% и т.д., что зависит от постоянного контекста, упомянув  модификатор
% прямо здесь. Таким образом, «постоянный контекст» является концепцией более низкого уровня.
% на основе синтаксиса и «постоянных выражений X» (например, «постоянный список буквальный»)
% построены на этом.


  является элементом списка или набора буквальных знаков, первый токен которого ,
 является ключом или ценностью входа
Буквальная карта, первым токеном которой является .
  Происходит как  в конструкцию, полученную из {metadata}.
  является фактическим аргументом в выражении, полученном из
{constObjectExpression}.
  является инициализирующим выражением постоянной переменной декларации
().
  является выражением случая переключения
(P000548).
  - немедленная субэкспрессия
Выражение  Это происходит в постоянном контексте.
где  это
% % Может быть добавлен позже:
% % не выражение \THROW{} () и
Не функция буквальная
().


{%
Постоянный контекст вводится в ситуациях, когда
Выражение должно быть постоянным.
Это позволяет использовать  Модификатор должен быть опущен
в случаях, когда он не предоставляет никакой новой информации.%
?


<{Нулевой}


%
Зарезервированное слово  Для сравнения: .


<nullLiteral> ::= <


%
Нулевой объект является единственным экземпляром встроенного класса .
% Нижеследующее может быть следствием объявления «нулевым»,
%, но мы не разъясняем, мы просто требуем, чтобы это было ошибкой.
Попытка мгновенно  вызывает ошибку компиляции времени.
Это ошибка времени компиляции для класса для расширения, смешивания или реализации.
.
 Класс расширяется до  класс
и не объявляет никаких методов, кроме тех, которые также объявлены .

{%
Нулевой объект имеет примитивное равенство.
()%
?

%
Статический тип \NULL{} является \code{Null} Тип.


 Номера.


%
A {цифровой буквальный}{буквальный!цифровой}
является либо десятичным, либо шестнадцатеричным числом, представляющим целое число,
или десятичное двойное представление.


<numericLiteral> ::= <Номер>
 <HEX Номер

<NUMBER> ::= <DIGIT>+ ('.' <DIGIT>+)? <Выражение>?
 <DIGIT>+ <EXPONENT>

<EXPONENT> ::= ('e' | 'E') ('+' | '-')? <DIGIT>+
<HEX Номер ::= "0x" <HEX DIGIT>+
 '0X' <HEX DIGIT>+

<HEX DIGIT> ::= 'a' ..'f "
 "А" .. "
 <DIGIT>


%
Буквальное число, начинающееся с "0x" или "0X" "
{hexadecimal integer literal}{literal!hexadecimal integer}.
Имеет числовое целое значение шестнадцатеричного числа
после "0x" (соответственно "0X").

%
Буквальное число, содержащее только десятичные цифры, является
{decimal integerliteral}{literal! десятичное целое число.
Он имеет числовое целое значение десятичного числа.

%
{integerliteral}{literal!integer}
является либо шестнадцатеричным целым буквальным, либо десятичным целым буквальным.

%
 быть целым буквальным числом, которое не является операндом
от оператора унарного минуса,
И пусть  Статический контекстный тип .
Если  Применяется к  и  не присваивается ,
Статический тип  является ;
Статический тип  является .

{%
Это означает, что целое число буквально обозначает 
когда он удовлетворяет требованиям типа и  Не стал бы.
В противном случае это , даже в ситуациях, когда это ошибка.
?

%
Буквальное число, не являющееся целым числом, является
{double literal}{literal!double}.
{%
Двойной буквальный всегда содержит либо десятичную точку, либо экспонентную часть.
?
Статический тип двойного литерала — .

%
Если  является целым буквальным числом с числовым значением  и статический тип ,
 не является операндой унарного оператора минус,
Далее следует оценка  идет следующим образом:


 Если  является шестнадцатеричным целым числом,
  Класс реализуется как
подписанные 64-битные два целых числа комплемента,
 Оценить на пример  класс
обозначает цифровое значение ,
< В противном случае  Оценивает экземпляр  класс
обозначает цифровое значение .
Это ошибка времени компиляции, если целое число  не могут быть представлены
Именно на примере .


{%
Целые числа в Дарте предназначены для реализации как
64-разрядные два дополняют целые представления.
На практике реализация может быть ограничена другими соображениями.
Например, Dart, скомпилированный для JavaScript, может использовать тип номера JavaScript.
эквивалентно Dart , чтобы представлять целые числа, и если да,
целые буквы с точностью более 53 бит не могут быть представлены
Точно.%
?

%
Двойная буквальная оценка для экземпляра  класс
64-битное число с плавающей точкой двойной точности
В соответствии со стандартом IEEE 754.

%
Буквальное целое число со статическим типом  и числовое значение 
Оценивается на пример  Класс, представляющий значение .
Это ошибка времени компиляции, если значение 
Не может быть представлен {точно} экземпляром .

{%
64-битное число с двойной точностью плавающей точки
Обычно используется для представления диапазона реальных чисел.
вокруг точного значения, обозначаемого числом
знак, мантисса и показатель.
Для целых букв, оценивающих до 
Значения, которые мы настаиваем на том, что целое число буквального значения
- точное значение  Пример:%
?

%
Это ошибка времени компиляции для класса для расширения, смешивания или реализации.
.
Это ошибка времени компиляции для класса для расширения, смешивания или реализации..
Ошибка времени компиляции для любого класса
кроме  и  для расширения, смешивания или реализации
.

%
Примеры  и  все вышеперечисленное
{==} оператор, унаследованный от  класс.


 Булева?


%
Зарезервированные слова \TRUE{} и {} оценивать объекты
{true}{true, the object} и
{false}{false, the object}
которые представляют булевы значения истинные и ложные соответственно.
Они являются {буквальные буквы}{буквальные!булин}.


<booleanLiteral> ::= <
 <


%
И \NoIndex{true} и {false} являются примерами
Встроенный класс ,
Нет других объектов, которые реализуют .
Ошибка времени компиляции для класса
расширять, смешивать или реализовывать .

{%
\TRUE{} и {} Объекты имеют примитивное равенство
()%
?

%
Ссылка на геттер  на булево значение возвращается
 Объект, который является значением выражения .
Статический тип булевого литерала — .


 Струны.


%
A  представляет собой последовательность кодовых единиц UTF-16.

{%
Это решение было принято для совместимости с веб-браузерами и Javascript.
Более ранние версии спецификации требовали, чтобы строка была
Последовательность действительных точек кода Unicode.
Программисты не должны зависеть от этого различия.
?


<stringLiteral> ::= (<multilineString> | <singleLineString>) +


%
Струна может быть последовательностью однострочных струн и многострочных струн.


<singleLineString> ::= <RAW Песня  LINE Стрельба
 < LINE СТРИНГ SQ BEGIN Окончание
 < LINE СТРИНГ SQ BEGIN\_MID> <выражение> \gnewline{}
(<Single) LINE СТРИНГ SQ  MID\_MID> <выражение>* <
< LINE СТРИНГ SQ MID Окончание
  LINE СТРИНГ\_ DQ BEGIN Окончание
 < LINE СТРИНГ DQ BEGIN\_MID> <выражение <
(<Single) LINE СТРИНГ  DQ MID\_MID> <выражение>* \gnewline{}
< LINE СТРИНГ DQ MID Окончание

<RAW Песня  LINE Строить>::=
 () \sq>\>\>)*\sq "
\alt 'r' ''' ((''' |'' |\')''

< Содержание  Коммон: :=  \\>>\sq>>>>\>
<<<p001209> <ESCAPE Последовательность
\alt\\\gtilde('n' | 'r' | 'b' | 't' | 'v' | 'x' | 'u' |>)
 < СТРИНГ Интерпол

< СОДЕРЖАНИЕ\_SQ> ::= <STRING\_ СОДЕРЖАНИЕ\_КОММОН> | ''

 LINE СТРИНГ SQ BEGIN\_END::= \gnewline{}
 < Содержание  SQ>* "
 LINE СТРИНГ SQ BEGIN\_MID> ::= \gnewline{}
 < СОДЕРЖАНИЕ\_SQ>*$ {''

 LINE СТРИНГ SQ MID\_MID> ::= \gnewline{}
'} <STRING СОДЕРЖАНИЕ SQ>*$ {''

< LINE СТРИНГ SQ MID\_END::= \gnewline{}
'} <STRING Содержание  SQ>* "

< СОДЕРЖАНИЕ\_DQ> ::= <STRING\_ Содержание  КОММОН > |  "

 LINE Строка  DQ BEGIN\_END::= \gnewline{}
<STRING СОДЕРЖАНИЕ\_DQ>* "

< LINE СТРИНГ DQ BEGIN\_MID ::= \gnewline{}
''' <STRING\_ СОДЕРЖАНИЕ\_DQ>*$ {''

< LINE СТРИНГ DQ MID\_MID ::= \gnewline{}
'} <STRING\_ СОДЕРЖАНИЕ\_DQ>*$ {''

< LINE СТРИНГ DQ MID\_END::= \gnewline{}
'} <STRING\_ СОДЕРЖАНИЕ\_DQ>* "


%
Одна строка ограничивается
Сопоставление однократных цитат или двойных цитат.

{%
Поэтому,  и  Оба они являются законными.
 и
.
 не является ни действительной строкой, ни .%
?

{%
Грамматика гарантирует, что одна строка не может превышать
одна строка исходного кода,
если оно не содержит интерполированного выражения, которое охватывает несколько линий.
?

%
Соседние струны неявно сцеплены, чтобы сформировать одну строку буквально.

{%
Вот пример: %
?



{%
Dart также поддерживает оператор + для конкатенации струн.

Оператор + на струнах требует аргумента струн.
Он не загоняет свой аргумент в нитку.
Это помогает избежать головоломок, таких как %.
?



{%
Что означает «P000780» вместо того, чтобы
.
По соображениям эффективности,
Рекомендуемая идиома Дарта - использование интерполяции строк.%
?



{%
Струнная интерполяция хорошо работает в большинстве случаев.
Главная ситуация, когда она не вполне удовлетворительна
Это для струнных букв, которые слишком велики, чтобы поместиться на линии.
Многолинейные струны могут быть полезны, но в некоторых случаях
Мы хотим визуально выровнять код.
Это можно выразить, написав
меньшие струны, разделенные белым пространством: %
?



%
вспомогательный

поддерживается вне парсера,
Для того, чтобы обеспечить правильное сопоставление струнных интерполяций.

{%
Это необходимо, поскольку выражение непростых интерполяций струн
может содержать строковые буквы
с их собственными непростыми интерполяциями.%
?

%
Правила с именами BEGIN Мид.
маркер нажимается на вспомогательный стек, чтобы указать, что
Началась струнная интерполяция данного вида,
где вид \lit{\sq}, \lit{}, \lit{\sqsqsq}, или \lit{"""].
Правила с именами \_MID Мид.
Можно использовать только правило с видом сверху вспомогательного стека.
Правила с именами  MID Конец.
можно использовать только правило с видом сверху вспомогательного стека,
После этого маркер выскакивает.


<multilineString> ::= <RAW MULTI LINE Стрельба <MULTI LINE Строка  SQ BEGIN Окончание
 <MULTI LINE СТРИНГ SQ BEGIN\_MID> <выражение <
(<MULTI) LINE СТРИНГ SQ MID\_MID> <выражение>* \gnewline{}
<MULTI LINE СТРИНГ SQ MID Окончание
 <MULTI LINE СТРИНГ DQ BEGIN Окончание
 <MULTI LINE СТРИНГ DQ BEGIN\_MID> <выражение> \gnewline{}
(<MULTI) LINE СТРИНГ DQ MID\_MID> <выражение>* \gnewline{}
<MULTI LINE СТРИНГ DQ MID Окончание

<RAW MULTI LINE STRING>::= 'r' '\sqsqsq < "
\alt 'r' ''''.*? ''''

<QUOTES\_SQ> ::= |\sq >> "

< СОДЕРЖАНИЕ\_TSQ> ::= \gnewline{}
<QUOTESSQ> (<STRING) СОДЕРЖАНИЕ КОММОН> | ''' | '\" |\"

<MULTI LINE СТРИНГ SQ BEGIN\_END::= <
 < СОДЕРЖАНИЕ TSQ>  "

<MULTI LINE СТРИНГ SQ BEGIN\_MID::= <
 < СОДЕРЖАНИЕ TSQ>* <QUOTES\_SQ> '$ {''

<MULTI LINE\_ СТРИНГ SQ MID\_MID ::= \gnewline{}
'} <STRING\_ СОДЕРЖАНИЕ\_TSQ>* <QUOTES\_SQ>$ {''

<MULTI LINE СТРИНГ SQ MID\_END::= <
'} <STRING\_ СОДЕРЖАНИЕ\_TSQ>*\sqsqsq "

<QUOTES\_DQ> ::= | '''

< СОДЕРЖАНИЕ\_TDQ> ::= \gnewline{}
<QUOTES\_DQ> (<STRING) СОДЕРЖАНИЕ КОММОН> |\sq" |\" |\"

<MULTI LINE СТРИНГ\_ DQ BEGIN\_END::= <
"""" <STRING СОДЕРЖАНИЕ\_TDQ>*""

<MULTI LINE СТРИНГ DQ BEGIN\_MID::= <
""" <STRING СОДЕРЖАНИЕ TDQ> <QUOTESDQ>$ {''

<MULTI LINE СТРИНГ DQ MID\_MID ::= <
'} <STRING СОДЕРЖАНИЕ TDQ> <QUOTESDQ>$ {''

<MULTI LINE СТРИНГ DQ MID\_END::= \gnewline{}
'} <STRING\_ СОДЕРЖАНИЕ\_TDQ>*""

<ESCAPE ПОСЛЕДОВАНИЕ > ::=\>>\<\r>>\>\>\ "
\ <HEX DIGIT> <HEX Диджит>
\alt\ <HEX DIGIT > <HEX DIGIT> <HEX DIGIT> <HEX Диджит>
\alt\{) <HEX DIGIT Последовательность "

<HEX DIGIT Последовательность>::= <
<HEX DIGIT > <HEX Цифра? <HEX Цифра?
\gnewline{} <HEX Цифра? <HEX Цифра? <HEX Цифра?


%
Многолинейные струны разграничиваются либосовпадающие тройки одиночных котировок или
Сопоставление тройных двойных цитат.
Если первая строка многострочной строки состоит исключительно из
символы белого пространства, определяемые производством <<<p001293>> Белый мир
(),
может быть приставлено \syntax{,
И эта линия игнорируется,
включая разрыв линии в конце.

{%
Идея состоит в том, чтобы игнорировать первую линию многострочной строки только в белом пространстве.
где белое пространство определяется как вкладки, пространства и окончательный разрыв линии.
Они могут быть представлены непосредственно,
Но так как для большинства персонажей префиксирование с помощью backslash
Идентичность в несырьевой струне,
Мы допускаем и такие формы.%
?

%
В правиле для \synt RAW MULTI LINE Стринг,
Два случая {.*?}
обозначает шаг нежелательного распознавания токенов:
Он заканчивается, как только Lookahead является указанным следующим токеном.
\commentary{(то есть {} или {"""))}.

{%
Обратите внимание, что многострочная интерполяция опирается на
Стек состояния интерполяции вспомогательной струны,
как однострочная интерполяция.%
?

%
Струны поддерживают последовательности побега для специальных персонажей.
Побеги бывают:



\syntax{\'} для новой линии, эквивалентной {\0A'}.

\syntax{\'} для возврата перевозки, эквивалентного {\0D'}.

\syntax{\'} для подачи формы, эквивалентной {\0C'}.

\syntax{\'} для заднего пространства, эквивалентно {\08'}.

\syntax{\'} для вкладки, эквивалентной {\09'}.

\syntax{\'} для вертикальной вкладки, эквивалентной {\0B'}.

< <HEX DIGIT> <HEX DIGIT>, эквивалентный

<< <HEX DIGIT> <HEX DIGIT>

< <HEX DIGIT> <HEX DIGIT> <HEX DIGIT> <HEX DIGIT>,
эквивалентный


< <HEX DIGIT> <HEX DIGIT> <HEX DIGIT> <HEX DIGIT>

<< <HEX DIGIT Последовательность > }
Точка кода Unicode, представленная
<<HEX DIGIT Последовательность.
Ошибка времени компиляции, если значение
<<HEX DIGIT Последовательность
не является действительным кодом Unicode.
 Например, }
{\color{комментарийЦвет}{<<<< P001351>>>{'\{0A}'}}\color{нормативный цвет}}
{является кодовой точкой U+000A. ?

\$ Это указывает на начало интерполированного выражения.

{% Нам нужно определение для  Чтобы иметь возможность использовать его в .
\def\k$k$
В противном случае \syntax{\\\k> указывает характер < для
любой  не в {'n', 'r', 'f', 'b', 't', 'v', 'x', 'u'}.
?


%
Любая строка может быть приставлена к символу {r},
Указывая, что это ,В этом случае не распознаются побеги или интерполяции.

%
Линейные разрывы в многострочной строке представлены
 LINE Производство.
Перерыв линии вводит один новый символ линии (U+000A)
в струнное значение.

%
Это ошибка времени компиляции, если буквальная строка не содержит
символьная последовательность формы \syntax{\'} за которым не следует
последовательность из двух шестидесятых цифр.
Это ошибка времени компиляции, если буквальная строка не содержит
символьная последовательность формы \syntax{\> за которым не следует
либо последовательность из четырех шестидесятых цифр,
или путем кудрявой скобки разграниченной последовательности шестидесятых цифр.


<LINE BREAK > ::=\ "
 "
 "


%
Все строковые буквы оцениваются на экземпляры встроенного класса .
Ошибка времени компиляции для класса
расширять, смешивать или реализовывать .
 класс переопределяет оператор {==}
 класс.
Статический тип строкового литерала — .


 Струнная интерполяция


%
Можно встраивать выражения в несырьевые струнные буквы,
Чтобы эти выражения были оценены,
Полученные объекты преобразуются в струны и
соединяются с ограждающей струной.
Этот процесс называется «P000577».


<stringInterpolation> ::= < СТРИНГ Интерпол
 <${<выражение>> "

<< СТРИНГ Интерполация>::= \gnewline{}
"$" (<IDENTIFIER) Нет DOLLAR > | <BUILT В  Идентификатор> | \THIS


{%
Читатель заметит, что выражение внутри интерполяции
может включать в себя струны,
которые могут быть интерполированы рекурсивно.%
?

%
Неудавшийся \lit{\$} символ в строке означает
Начало интерполированного выражения.
 Знак может сопровождаться либо:


 Единый идентификатор  не содержит \lit{\$} характер
(но это может быть встроенный идентификатор),
или зарезервированным словом .
 Выражение  Ограничен кудрявыми брекетами.


%
Форма \code{\\{\id\}}.
интерполированная строка , с содержанием
\code$s_0$\e_1>>>_1\ldots{}s_{n-1}<<< P000331>>>\{$e_n$\}$s_{n}$ "
(где любой из ) может быть пустой
Оценивается путем оценки каждого выражения  ()
в строку  в порядке, в котором они встречаются в исходном тексте, следующим образом:


 Оценка  Объект .
 Нажмите  Метод  без аргументов,
Пусть  быть возвращенным объектом.
 Если  является нулевым объектом, возникает динамическая ошибка.


%
Наконец, результат оценки  это
сцепление струн , , ,  и .


{Символы}


%
A {символ буквальный}{буквальный!символ}
обозначает имя, которое будет
действительное название декларации, действительное название библиотеки или .


<symbolLiteral> ::=
'#' (<identifier> ('.' <identifier>)* | <operator> | )


%Статический тип буквального символа .

%
 Идентификатор, который не начинается с подчеркивания
(«P000789»).
Буквальный символ 
Оценить на пример 
идентификатор .

%
Буквальный символ 
где <<<<<<<<<<<<<<<<<<<<<<<<<<<<<<<<<<<<<<<<<<<<<<<<<<<<<<<<<<<<<<<<<<<<<<<<<<<<<<<<<<<<<<<<<<<<<<<<<<<<<<<<<<<<<<<<<<<<<<<<<<<<<<<<<<<<<<<<<<<<<<<<<<<<<<<<<<<<<<<<<<<<<<<<<<<<<<<<<<<<<<<<<<<<<<<<<<<<<<<<<<<<<<<<<<<<<<<<<<<<<<<<<<<<<<<<<<<<<<<<<<<<<<<<<<<<<<<
Для сравнения: 
представление этой конкретной последовательности идентификаторов.
{%
Этот тип символа буквально обозначает название библиотечной декларации.
Как указано в {название библиотеки}.
Названия библиотек не подлежат конфиденциальности, даже
если некоторые из его идентификаторов начинаются с подчеркивания.
?

%
Буквальный символ \code{\#\metavar{op}}
где \metavar{op} происходит от \synt{operator}
Оценить на пример 
Укажите конкретное имя оператора.

%
Буквальный символ 
Для сравнения: 
Представляет собой зарезервированное слово .

%
За значение  Буквальный символ, представляющий термин исходного кода
Как указано в предыдущих параграфах, мы говорим, что  является
{нечастный символ на основе}{символ!нечастный, на основе}
строка, содержание которой является символами этого термина, без белого пространства.

{%
Обратите внимание, что это не относится к частным символам.
которые обсуждаются ниже.
Частный символ не основан на какой-либо строке.
?

%
Если  является значением вызова  конструктор
в форме
, ,
или ,
где  является выражением
\commentary{(постоянно, если необходимо)}
который оценивает строку ,
Мы говорим, что  {неприватный символ на основе} .

<{%
Обратите внимание, что  является частным символом,
и отличается от , как описано ниже.%
?

%
Предположим, что ,
 значение постоянного выражения
Символ, основанный на строке .
Если , то  и  Это один и тот же объект.
 То есть символьные экземпляры канонизированы. ?

%
Если  и  Нечастные символы
{(не обязательно постоянная)}
основанный на струнах  и 
 и  равны в соответствии с оператором \lit{===
если и только если 
(P000554).

%
Буквальный символ , где  Является идентификатором
Оценить на пример 
представляет собой частный идентификатор  библиотеки-приложения.
Все случаи  {в той же библиотеке} оценить
Тот же объект,
и никаких других символов, буквальных или  призыв конструктора
оценивает один и тот же объект,
и не для  Пример, равный этому объекту
в соответствии с оператором {==}.

{%
Можно задаться вопросом, какова мотивация введения буквальных символов?
В некоторых языках символы канонизированы, а струны нет.
Однако буквальные струны уже канонизированы в Дарте.
Символы немного легче печатать по сравнению со строками.
Их употребление может стать странным,
Но это не является достаточным основанием для добавления
Буквальная форма языка.
Основная мотивация связана с использованием рефлексии.
Веб-специфическая практика, известная как минификация.

Минификация последовательно сжимает идентификаторы во всей программе.
Чтобы уменьшить размер загрузки.
Эта практика создает трудности для рефлексивных программ.
Это относится к программным декларациям через строки.
строка будет ссылаться на идентификатор в источнике,Но идентификатор больше не будет использоваться в минированном коде.
и отражающий код, использующий их, не сработает.
Поэтому в отражении Дарта используются объекты типа 
Вместо струн.
Случаи  гарантируется стабильная
Что касается минирования.
Предоставляя буквальную форму для символов делает
Отражающий код легче читать и писать.
Символы легко печатаются и часто могут действовать как
Удобные заменители для энумов являются вторичными преимуществами.
?


<\subsection{Сборник литературы}


%
В этом разделе указаны несколько буквальных выражений, обозначающих коллекции.
Некоторые синтаксические формы могут обозначать более одного вида коллекции.
В этом случае выполняется этап раздвоения
Для того чтобы определить вид
(P000555).

%
Подразделы настоящего раздела посвящены механизмам, которые
Общие для всех видов коллекции буквы
(,
,
Далее следует спецификация букварей списка
(, ,
Затем следует спецификация того, как различать и выводить типы
для наборов и карт
(, ,
И, наконец, спецификация наборов
()
и карты
(P000563).


<listLiteral> ::=  <typeArguments>? <элементы>? "

<setOrMapLiteral> ::=  <typeArguments>? <элементы>? "

<элементы>::= <элемент> (',' <элемент>)*,'?

<element> ::= <expressionElement>
 <mapElement>
 <spreadElement>
 <ifElement>
 <для элемента>

"Экспрессии" :: <Экспрессии>

<mapElement> ::= <expression> <expression>

<spreadElement> ::= ('...' | '...?') <выражение>

<ifElement> ::= \IF{} '(' <выражение> ')' <элемент> ({} <элемент>)?

<forElement> ::=  \FOR{} '(' <forLoopParts> ')' <элемент>


%
Синтаксически, a  может быть
а \synt{listLiteral} или {setOrMapLiteral}.
Содержание коллекции определяется как последовательность
{сбор буквальных элементов}{сбор буквальных элементов}! элементы},
коротко .
Каждый элемент может быть декларативной спецификацией одного объекта,
Например, \synt{expressionElement} или \synt{mapElement},
может указывать коллекцию, которая должна быть включена,
в форме <{spreadElement}
или это может быть вычислительный элемент, определяющий
Как получить ноль или больше объектов
в форме \synt{ifElement} или {forElement}.

{%
Термины, полученные из {элемента},
Способность создавать коллекции из них;
Также известен как
.%
?

%
 элемент  полученный из
\synt{expressionElement} или {mapElement}
.
Листовые элементы элемента формы
 или
\code\FOR\,({forLoopParts})$\ell$}
Это элементы листьев .
Листовые элементы элемента формы

Это соединение листовых элементов  и .
Элементами листа {spreadElement} является пустой набор.

{%
Листовые элементы сборника в буквальном смысле всегда
Набор элементов выражения и/или элементов карты.%
?

%
Для того, чтобы позволить сбору букв в качестве постоянных выражений,
Уточним, что это означает для элемента 
быть
{constant}{сборник буквальный элемент!constant}
или
{потенциально постоянная}{%
Буквальный элемент коллекции! Потенциально постоянные:



Когда  {выражениеЭлемент} формы :

 потенциально постоянный элемент
Если  является потенциально постоянным выражением,
 является постоянным элементомЕсли  Это постоянное выражение.

Когда  {mapElement}
в форме :

 потенциально постоянный элемент
 и  потенциально постоянные выражения,
Это постоянный элемент, если они являются постоянными выражениями.

Когда  {распространенный элемент}
в форме \code{...} или \code{...?}:

 потенциально постоянный элемент
Если  Это потенциально постоянное выражение.

 является постоянным элементом
Если  Постоянное выражение, оценивающее
,  или  пример
Первоначально созданный списком, набором или картой буквально.

Кроме того,  является постоянным элементом, если это ,
где  Постоянное выражение, оценивающее
к нулевому объекту.

Когда  {ifElement}
в форме

или форма
:

 потенциально постоянный элемент
Если  является потенциально постоянным выражением,
 потенциально постоянный,
и так далее , если присутствует.

 является постоянным элементом, если  Это постоянное выражение
и:



  и
 {} и  является постоянным,
 {} и  Потенциально постоянный.

  и
 Показатель ,
 является постоянным,
 потенциально постоянный;
 Показатель ,
 потенциально постоянный,
 является постоянным.



{%
A {forElement} никогда не может происходить в постоянном сборе букв.%
?

% % TODO(Eernst): Мы можем изменить этот текст на комментарий или удалить его,
% состоит из описаний ошибок, но они уже охвачены
% в других местах.
% %
% ошибок компиляции.
% %
% % Следующие ошибки времени компиляции указаны в функции
Спецификация %%, но они уже подразумеваются указанными ошибками
% в других местах.
% %
%% - "Неосведомленный элемент спреда имеет статический тип ".
%% Для списка буквально: ListLiteralInference делает эту ошибку случайностью
Процентный анализ на элементе «Spread».
% % Set/map буквально: setAndMapLiteralInference, ditto.
% %
"Спредный элемент в списке или наборе буквальных элементов имеет статический тип, который является
% % не  и не является подтипом ".
% % Список буквальный: анализ случаев на «элементе спреда».
% % Набор буквальный: анализ кейсов по «элементу спреда»: преуспевает только в том случае, если
% %  реализует «Карту» (а не «Итерируемую»), поэтому имеет ключ/значение
Тип %% и тип элемента без ошибок (P000564).
% %
%% - "Элемент спреда в списке или наборе имеет статический тип, который
% Реализация   и
%%  не присваивается типу элемента списка".
% % Перечислите буквально: ср. «Элемент распространения» означает, что элемент распространения
% % имеет тип элемента , а неназначение является ошибкой .
% % Набор буквальный: cf. 'Spread element' for set/maps: ditto.
% %
%% - "Спредный элемент на карте в буквальном смысле имеет статический тип, который является
% % не  и не является подтипом ".
% % Анализ кейсов «Spread element» подразумевает, что он реализует
% % «Изменяемый», а не «Карта», в этом случае у него нет пары типа ключ/значение,
%, что является ошибкой, .
% %
%% - "Если статичный тип элемента карты распространяется 
% для некоторых  и  и  не присваивается ключевому типуКарта %% или  не присваивается типу значения карты".
% % Это ошибка, .
% %
%% - "Переменная в {для элемента} {(либо для входа)
%% элемент или C-стиль)} объявляется вне элемента окончательным или
% % не иметь сеттера».
% % Эта ошибка указана в  и
%% .
% %
%% - "Тип выражения итератора в синхронном элементе включения
% не может быть присвоено  Для любого типа .
В противном случае  итератора составляет .
%, покрытых тем же текстом, что и предыдущая ошибка.
% %
%% - "Переходный тип итератора в синхронном элементе включения
% не может быть отнесено к типу переменной "for-in".
%, покрываемых одним и тем же текстом.
% %
%% - "Тип выражения потока в асинхронном \AWAIT{} вход
Элемент %% не может быть присвоен  Для любого типа .
% % В противном случае  Поток имеет значение ".
%% Один и тот же текст (текст «для элемента» включает в себя как ожидание, так и ожидание).
% %
%% - "Тип потока итератора в асинхронном \AWAIT{} элемент for-in
% не может быть отнесено к типу переменной "for-in".
% опять тот же текст.
% %
%% - "\AWAIT{} << Применяется, когда функция непосредственно заключает
Буквальный набор %% не является асинхронным".
% опять тот же текст.
% %
% % — «\AWAIT Используется перед C-стилем {для элемента}.
%%  Его можно использовать только в петлях ".
% опять тот же текст.
% %
%% - "Тип выражения состояния \commentary{(второй пункт)}
% в стиле C \FOR{} Элемент не может быть присвоен ".
% опять тот же текст.


 Тип поощрения


%
{ifElement} взаимодействует с продвижением типа
Таким же образом, как  Заявления делают.
 быть {ifElement} формы
 или
.
Если  показывает локальную переменную  имеет тип 
Тип  Известен как  в ,
Если ни одно из следующих утверждений не верно:


  потенциально мутирует в ,
  Потенциально мутирует в функции
кроме той, где  объявляется, или
  Доступ осуществляется с помощью функции, определенной в  и
 Потенциально мутирует в любом месте в пределах .


% % TODO(Eernst): Приходите, обновите это.
{%
Продвижение типа, вероятно, станет более сложным в будущей версии Dart.
Когда это происходит,
< Элементы будут продолжать соответствовать  заявления
()%
?


 Литературная оценка элементов


%
Оценка последовательности сбора буквальных элементов
()
доходность a
\IndexCustom{сбор буквальной последовательности объектов}{%
Коллекция буквально! объектная последовательность.
Также называется {объектная последовательность}, когда не может возникнуть двусмысленности.

%
Мы используем нотацию
%  Не может быть использован после '@', поэтому мы используем приближение.
\IndexCustom{\LiteralSequence{\ldots}}{[[...]]@%
<< \hspace{-0.6mm}[{-0.6mm}]}}
обозначать объектную последовательность с явно перечисленными элементами,
и мы используем '' для вычисления сцепления последовательностей объектов
\commentary{(как в $s_1 + s_2$)},
Это операция, которая будет успешной и не будет иметь побочных эффектов.
Каждый элемент в последовательности является объектом.  или парой .
Не существует понятия типа элемента для последовательности объекта,
Не существует динамических ошибок, возникающих из
Тип несоответствия во время конкатенации.

{%
Последовательности объектов можно безопасно рассматривать как механизм низкого уровня.
которые могут пропустить другие необходимые действия, такие как динамические проверки типа
Каждый доступ к объектной последовательности происходит
в коде, созданном языком, определяемым desugaring
на статически проверенных конструкциях.
Остается неустановленным, как реализуется объектная последовательность,
Требуется только, чтобы он содержал указанные объекты или пары.
в данном порядке.
Для каждой коллекции,
последовательность используется в данном порядке для заполнения коллекции,
в том виде, который специфичен для данного вида,
который указывается отдельно
(, , %

Может существовать реальная структура данных.
представление последовательности объекта во время выполнения,
но последовательность объектов также может быть устранена, например,
потому что каждый элемент вставляется непосредственно в целевой набор
Как только она будет вычислена.
Обратите внимание, что каждая последовательность объектов будет содержать исключительно объекты.
Или же он будет содержать только пары,
Любая попытка создать смешанную последовательность вызовет
Ошибка во время компиляции или во время выполнения
(последний может иметь место для элемента с разбросом со статическим типом \DYNAMIC{}.%)
?

%
Допустим, что дается дословный сборник \metavar{target},
и будет использоваться последовательность объектов, полученная, как описано ниже
населять {цель}.
 обозначает динамический тип {цель}.

{%
Доступ к типу {цель} необходим ниже
Для увеличения динамических ошибок в определенных точках во время
Оценка объектной последовательности.
Обратите внимание, что динамический тип {цель} статически известен.
за исключением связывания переменных любого типа в его {тип Аргументы}.Это означает, что на некоторые вопросы можно ответить во время компиляции, например:
Или нет (P001045). происходит как суперинтерфейс
.
В любом случае,  гарантируется осуществление
 (когда {цель} - список или набор)
или  (когда {цель} является картой),
но никогда оба.%
?

%
Предположим, что дается место в коде и динамический контекст,
Таким образом, возможна обычная оценка выражения.
 Оценка последовательности буквальных элементов коллекции {%}
Буквальный элемент коллекции! оценка последовательности
В этой связи следует отметить, что в этой связи следует принять во внимание, что в данном случае

%
Пусть $s_{\metavar{syntax}}$ формы \List{\ell}{1}{k} быть
Последовательность сбора буквальных элементов.
Последовательность объектов 
полученных путем оценки 
является сцепление последовательностей объектов
Полученные путем оценки каждого элемента , :
,
где \EvaluateElement{_j} обозначает объектную последовательность, полученную
Оценка одного литературного элемента коллекции .

%
При псевдозаявлении формы

используется в нормативном кодексе ниже,
обозначает расширение  с последовательностью объекта
по результатам оценки ,
но также обозначает спецификацию действий, предпринятых для
Последовательность объектов,
и для получения побочных эффектов, связанных с этим вычислением,
как подразумевается при оценке выражений и исполнении заявлений
Как указано ниже для оценки \EvaluateElement{\ell}.

%
При псевдозаявлении формы

происходит в нормативном кодексе ниже,
Это происходит в точке, где вычисление завершено.
и указывает, что значение \EvaluateElement{\ell} составляет .
 Оценка литературного элемента коллекции {%}
Буквальный элемент коллекции! оценка
в данном контексте
к объектной последовательности
\IndexCustom{\EvaluateElement{\ell>>>}}
Оценка Element(l)@\emph Элемент\code{()}
Затем указывается следующее:

%
 Элемент выражения
В данном случае  выражение ;
 оценивается на объект 
.


%
 Элемент карты
В данном случае  Это пара выражений
;
 оценивается на объект ,
 оценивается на объект ,
.


%
 Элемент распространения.
Элемент  Он имеет форму  или .
Оценить  на объект .



Когда  означает :
% % TODO(Eernst): Приходите в NNBD, эта ошибка больше не может произойти: удалите.
Если  Это нулевой объект, после чего происходит динамическая ошибка.
В противном случае оценка осуществляется на этапе 2.

Когда  означает :
Если  Тогда нулевой объект
.
В противном случае оценка осуществляется на этапе 2.

 Динамический тип .
 Статический тип .
Когда  Не является высшим типом
(),
Пусть  будет .
Когда <<<p000079> Это топовый тип:
Если {цель} - это список или набор
Пусть  будет ;
иначе
\commentary{(где \metavar{цель} - карта)},
Пусть  будет .



Когда {цель} - список или набор
 реализация
()
,
следующий код выполняется в контексте, где  происходит,
где , ,  и  свежие переменные, является переменной нового типа, связанной с
Аргумент фактического типа  в 
():

\vspace{-2ex}[t]{]



{%
Код использует псевдовариабельный  обозначает последовательность объектов.
Мы не указываем тип ,
Эта переменная используется только для указания требуемого значения.
Семантические действия, предпринятые для сбора полученной объектной последовательности.
В случае, если реализация не
Изображение  Во всяком случае,
Действие может заключаться в немедленном продлении {target}.
Аналогичный подход используется в последующих случаях.%
?

Когда {цель} - карта и  реализация
,
следующий код выполняется в контексте, где  происходит,
где , ,   и 
свежие переменные,
и  и  свежие переменные типа, связанные с
первый, соответственно второй фактический тип аргумента
из  по адресу :

\vspace{-2ex}[t]{]



% не изменится при nnbd: аргументы типа «распространение» могут быть .
Реализация может задерживать динамические ошибки
Это происходит, если данный  не имеет типа ,
или отданное  не имеет своего типа ,
Но это не может произойти после того, как пара была добавлена к .

В противном случае возникает динамическая ошибка.

{%
Это происходит, когда цель является повторяемой соответственно картой.
Распространения нет, что возможно для
спред, статический тип которого .%
?



{%
Это может быть не самый эффективный способ перемещения предметов в коллекции.
Реализация может, конечно, использовать любой другой подход.
с таким же наблюдаемым поведением.
Тем не менее, чтобы дать реализациям больше возможностей для оптимизации.
Допустим также следующее.%
?

%
Если  Объект, динамический тип которого реализует
()
,  или ,
Реализация может выбрать вызов  на объекте.
Если  является объектом
чей динамический тип ,
реализация может выбрать вызов оператора 
Для доступа к элементам из списка.
Если он это сделает, он будет только передавать индексы.
которые являются неотрицательными и меньше значения, возвращенного .

{%
Это позволяет использовать более эффективный код для
Распределение коллекции и доступ к ее частям.
Эти классы должны иметь
Эффективная и бесперебойная реализация
 и оператора {[]}.
Реализация Dart может определить, применяются ли эти варианты.
в момент компиляции на основе статического типа ,
или во время выполнения на основе фактического значения.%
?


%
 Если элемент
Когда  {ifElement} формы
 или
,
Состояние  оценивается по значению .
Если   затем
.
Если  {} и  Присутствует тогда
,
Если  Не присутствует тогда
.
%  может иметь тип .
Если  не является ни \TRUE{}, ни \FALSE{} Затем происходит динамическая ошибка.


%
 Для элемента
 быть получено из {<forLoopParts>} и
Пусть  будет {для элемента} формы
,Где «P001482»? Это означает, что  может присутствовать или отсутствовать.
Для оценки ,
следующий код выполняется в контексте, где  происходит,
где < Присутствует тогда и только тогда, когда присутствует в :

\vspace{-2ex}[t]{]





<{Список литературы}


%
В этом разделе указывается, как буквальный список \metavar{list} пройден и
{предполагаемый тип элемента}{список буквальный! Тип элемента
Для {список} определяется.
Сначала мы указываем, как определить тип элемента одного элемента.
Как использовать этот результат для вывода
Тип элемента {список} в целом.

%
Тип контекста 
()
Для каждого элемента 
Получено из контекстного типа {список}.
Если вывод вниз ограничивает тип 
 или  Для некоторых 
 .
В противном случае  \FreeContext{}
(P000969).

%
 Термин, полученный из {элемент}.
Вывод типа элемента  Тип контекста 
Происходит следующим образом,
где тип контекста для вывода типа элемента всегда ,
Если ничего не сказано об обратном:

%
 Элемент выражения
В данном случае  Выражение .
Тип элемента  это
% Тип  Не является типом элемента, поэтому мы упоминаем .
Предполагаемый тип  в контексте .


%
 Элемент карты
% Мы не можем точно сказать, является ли это ошибкой в  возникает первым
% (скажем, когда буквальный список имеет явный тип аргумента, он не является
% гарантирует, что реализация выполняет вывод типа на ней вообще.
% Таким образом, мы _must_ укажем ошибку в , и это местоположение должно
% — это комментарий, потому что у нас уже есть ошибка в другом месте.
{%
Этого не может произойти:
Это ошибка времени компиляции
Когда листовой элемент списка буквально является элементом карты
()%
?


%
 Элемент распространения.
 Выражение .
Если  означает ,
 Выберите тип  в контексте .
Иначе
{(когда  является )},
%% TODO (eernst): Приходите в NNBD, добавьте  в спецификацию
% — «неотменяемый тип».
 быть ненулевым типом
% % TODO(Eernst): Уточнить, будет ли вывод действительно иметь такой контекст;
% ясно, что мы должны устранить «Null» из , а также что $e$
% допускается иметь потенциально нулевой тип, и это кажется неудобным.
% %, если мы используем 'Iterable< в качестве контекстного типа и отказа в случае, когда ,
% % говорят: «Список <?» Для некоторых .
Предполагаемый тип  в контексте .



Если  реализация ,
Тип элемента  это
Аргумент типа  при .

Если  ,
Предполагаемый тип элемента  является .

Если   и оператор спреда {...?},
Тип элемента  является .

В противном случае возникает ошибка времени компиляции.

\vspace{-5 мм)


%
 Если элемент
В данном случае  имеет форму
 или
.
Состояние  Всегда выводится контекстный тип .
Допустим, что "" отсутствует.
Затем, если указанный тип элемента  является ,
Тип элемента  является .

В противном случае присутствует «\code{\ELSE\,\,}».
Если выбранный тип элемента  является  и
Тип элемента  является ,
Тип элемента  это
наименьшая верхняя граница  и .


%
 Для элемента
В данном случае  имеет форму

где  является производным от {forLoopParts} и
 Это означает, что  может присутствовать или отсутствовать.

Ошибки времени компиляции для  как
Ошибки, которые могут возникнуть с соответствующим  заявление
 \,\,\FOR\,\,$P$\,\{\}>,
находится в том же объеме, что и .
Кроме того, ошибки и анализ типов  выполняется
как если бы это произошло в пределах объема тела указанного {} заявления.
{%
Например, если  имеет форму
<<<p001102>
переменная  Находится в пределах .%
?

Вывод для частей
{%
(таких, как выразительное выражение фор-ин,
или \synt{forInitializerStatement} для петли)%
?
выполнено как для соответствующего \FOR{} заявление,
включая < если и только если элемент содержит .
Затем, если указанный тип элемента  является ,
Предполагаемый тип элемента  является .

{%
Другими словами, вывод течет вверх от элемента тела.
?
\vspace{3 мм)


%
Наконец, мы определяем
\IndexCustom{тип вывода в списке буквальный}{%
Типовой вывод! список буквальный
в целом.
Предположим, что \metavar{list} является производным от \synt{listLiteral}
и содержит элементы \List{\ell}{1}{n},
и тип контекста для \metavar{list} является .



Если   затем
Предполагаемый тип элемента для \metavar{list} является ,
где  является наименьшей верхней границей
Выводимые типы элементов \List{\ell}{1}{n}.


% % TODO(Eernst): Спецификация функции говорит , но как мы знаем  Это тип?
В противном случае
Предполагаемый тип элемента для {list} является ,
где  определяется путем нисходящего вывода.


%
В обоих случаях статический тип {список} является .


<{Список}


%
A {список буквальный}{буквальный}! список
обозначает объект списка, который является целым индексируемым набором объектов.
Правило грамматики для {listLiteral} указано в другом месте.
(P000974).

%
При заданном списке буквальный  Нет никаких аргументов,
Тип аргумента  Выбирается так, как указано в другом месте
(),
и  отныне рассматривается как
()
.

{%
Статический тип списка буквальный формы 

()%
?

%
 быть буквальным списком формы
< \List{}{1}{m}]}.
Это ошибка времени компиляции, если элемент листа  является
{mapElement}.
Это ошибка времени компиляции, если для некоторых 
 не имеет элемента типа,
Тип элемента  Не может быть назначено .

%
Список может содержать ноль или более объектов.
Количество объектов в списке – это их размер.
Список имеет связанный набор индексов.Пустой список имеет пустой набор индексов.
Непустый список имеет индекс 
где  Это размер списка.
Динамическая ошибка при попытке получить доступ к списку
использовать индекс, который не входит в его набор индексов.

{%
Системные библиотеки определяют множество членов для типа ,
Но мы указываем только минимальный набор требований.
Используется самим языком.%
?

%
Если список буквальный  начинается с зарезервированного слова 
 происходит в постоянном контексте
(),
Это а
{постоянный список буквально}{буквальный!список!постоянный},
Что является постоянным выражением
()
и поэтому оценивается во время компиляции.
В противном случае это a
{список времени выполнения буквальный}{буквальный!список!срок выполнения}
и оценивается во время выполнения.
Только список времени выполнения буквы могут быть мутированы
После того, как они были созданы.
% Эта ошибка может произойти потому, что постоянное является динамическим свойством.
Попытка мутировать постоянный список буквально приведет к динамической ошибке.

{%
% Прямо или косвенно верно следующее: Существует 
% модификатор в списке буквальный, или список буквальный в целом
% в постоянном контексте.
Обратите внимание, что в сборнике буквальные элементы постоянного списка буквальные
происходит в постоянном контексте
(),
что означает  Модификаторы не должны быть указаны явно.%
?

%
Это ошибка времени компиляции
Если элемент постоянного списка буквален, то он не постоянен.
Это ошибка времени компиляции, если аргумент типа постоянного списка буквален.
{(независимо от того, является ли это явным или предполагаемым)}
Не является постоянным выражением типа
(P000981).

{%
Связывание формального параметра типа замкнутого класса или функции
не известны во время составления,
Поэтому мы не можем использовать такие параметры типа в постоянных выражениях.
?

%
Значение постоянного списка буквально
 < \List{}{1}{m}}
является объектом  чей класс реализует встроенный класс

где  - фактическое значение 
(),
и содержимое которого представляет собой объектную последовательность {o}{1}{n}
Оценка \List{}{1}{m}
(P000983).
  (по индексу ) В этом случае .

%


Будьте двумя постоянными буквами.
Пусть  с содержимым  Реальный тип аргумента 
соответственно
 Содержание  Реальный тип аргумента 
быть результатом их оценки.
Далее  Показатель {} если {f}
 и 
{} для всех .

{%
Другими словами, постоянные буквальные списки канонизированы.
Нет необходимости рассматривать канонизацию.
для других экземпляров типа ,
Поскольку такие случаи не могут быть
результат оценки постоянного выражения.%
?

%
Список времени выполнения буквальный
< \List{}{1}{m}}
Оценка проводится следующим образом:



Оцениваются элементы \List{\ell}{1}{m}
(),
в объектную последовательность \LiteralSequence{\List{o}{1}{n}}.

Свежий пример () , размер ,
чей класс реализует встроенный класс 
выделяется,
где  - фактическое значение 
().

Оператор \lit{[]=} вызывается на $o$ с
Первый аргумент  Второй аргумент
.

Результатом оценки является .


%Объекты, созданные буквами списка, не переопределяются
{==} оператор, унаследованный от P001117 класс.

{%
Обратите внимание, что в этом документе не указывается приказ
в котором установлены элементы.
Это позволяет выполнять параллельные задания в списке.
Если этого желает реализация.
Порядок может соблюдаться только следующим образом (и на него нельзя полагаться):
Если элемент  не является подтипом элемента типа списка,
Ошибка динамического типа возникает, когда  присваивается .%
?


<{Сет и карта Буквальная неопределенность}


%
Некоторые термины, такие как  и  \id\,\}> являются двусмысленными:
Они могут быть либо буквальными, либо буквальными.
Эта двусмысленность устраняется в два этапа.
На первом этапе используется только тип синтаксиса и контекста.
% % TODO (Eernst): Включите эту ссылку, когда будет определен «тип контекста».
% % ()
и описано в этом разделе.
Второй этап использует типы выражений и описывается следующим образом.
(P000988).

%
Пусть  будет {setOrMapLiteral}
с помощью элементов «>>>>>>>>>>>>>>>>>>>>>>>>>>>>>>>>>>>>>>>>>>>>>>>>>>>>>>>>>>>>>>>>>>>>>>>>>>>>>>>>>> и типами «Публикация»
Если   Тогда пусть  быть неопределенным.
% % TODO(Eernst): Определить «наибольшее закрытие» типа контекста
% при определении «контекстного типа».
В противном случае пусть  быть наибольшей закрытостью  С)
(P000989).

% % TODO(Eernst): Исключить слова "тип контекста", "наибольшее закрытие".
{%
В будущей версии этого документа будут указаны типы контекста.
Основная интуиция заключается в том, что

Это тип, объявленный для
Получающая организация, такая как
формальный параметр  или объявленной переменной .
Этот тип будет контекстным для
фактический аргумент, переданный ,
соответственно начальное выражение для .
В некоторых ситуациях контекст не имеет ограничений.
Например, когда переменная объявлена с помощью {} Вместо аннотации типа.
Это приводит к возникновению
{неограниченный тип контекста}{контекстный тип!неограниченный},
\IndexCustom{}{[@],
Это может также происходить в сложном термине, например, .
% % TODO(Eernst): Уточните, почему мы не используем i2b
% - введение понятия наибольшего (и наименьшего) закрытия.
Наибольшее закрытие контекстного типа  это
Наименее распространенный супертип всех типов
Получается путем замены \FreeContext{} по типу.%
?

%
Шаг неоднозначности этого раздела является
первая применимая запись в следующем списке:


 Когда  Типовые аргументы  T}{1}{k} :
Если , то  Буквальный набор со статическим типом
.
Если , то  Буквальная карта со статическим типом
.
В противном случае возникает ошибка времени компиляции.

Когда  реализация
()
 но не ,
 Это набор буквальный.
Когда  Реализация  но не ,
 - буквальная карта.

Когда 
\commentary{(то есть ) имеет элементы листьев)}:
Если  Содержит {mapElement}
а также {выражениеЭлемента},
Происходит ошибка времени компиляции.
В противном случае, если  содержит {выражениеЭлемента},
 - это буквальный набор.
В противном случае  содержит \synt{mapElement} и $e$ Это буквальная карта.

Когда  имеет форму \code{\{\}} и $S$ неопределенным,
 - буквальная карта.
{%Нет более глубокой причины для такого выбора.
Но факт, что \code{\}} Это карта по умолчанию
Он был полезен, когда были введены буквы,
Потому что это было бы переломным изменением, чтобы сделать его установленным.
?

В противном случае  остается неоднозначным.
{%
В данном случае  не является пустым, но содержит только спреды
0 или более раз в  Элементы или {для элементов}s.
Затем во время вывода будет происходить раздвоение
()%
?


{%
Когда этот шаг не определяет статический тип,
будет определяться по типу вывода
()%
?

%
Если этот процесс успешно дезавуирует буквальный
Мы говорим, что  это
{неоднозначно множество}{set!unambiguously}
или
{неоднозначно карта}{map!unambiguously},
по мере необходимости.


<\subsubsection{Сет и карта буквального вывода}


%
В этом разделе указывается, как  Литературный  пройденный
и связанных с ним
{предполагаемый тип элемента}{набор или карта буквально! Тип элемента
и/или связанный
{пара типа ключа и значения}{%
Настройка или карта буквально! пара типа ключа и значения
определяется.

{%
Если  имеет тип элемента, то он может быть набором,
Если у него есть пара типа ключа и значения, то это может быть карта.
%
Однако, если буквальный  содержит разветвленный элемент типа ,
Этот элемент не может быть использован для определения  Это набор или карта.
Неоднозначность представлена как имеющая \emph{оба}
Тип элемента и пара типа ключа и значения.

Ошибка, если неясность не определена другими элементами,
Но если это решено, то требуется динамический элемент распространения.
Чтобы оценить подходящий пример
(реализация , когда ) Это набор,
и реализации , когда  Это карта,
Это означает, что это динамическая ошибка, если есть несоответствие.
В других ситуациях ошибка компиляции состоит в том, что оба
тип элемента и пара типа ключа и значения,
 Это должен быть и набор, и карта.
Вот пример: %
?



%
Пусть <{коллекция} будет сборником буквально
полученный из {setOrMapLiteral}.
Предполагаемый тип {элемент} является типом элемента ,
пара типа ключа и значения  или оба.
Вычисляется относительно типа контекста 
(),
который определяется следующим образом:



Если {коллекция} является однозначно установленным
()
 ,
% % TODO(Eernst): Добавьте ссылку, когда мы укажем вывод.
где  определяется нисходящим выводом,
и может быть 
% % TODO(Eernst): Исправьте эту ссылку, когда мы указываем типы контекста.
()
Если нисходящий вывод не ограничивает его.

% % TODO(Eernst): Удалите это, когда мы укажем вывод.
{%
В будущей версии этого документа будет указан вывод,
Понятие нисходящего вывода,
и ограничения.
Краткая интуиция заключается в том, что вывод выбирает значения для
Типовые параметры в генерических конструкциях
где не приводятся аргументы,
нацеливание на тип, который соответствует данному типу контекста;
Нисходящий вывод делает это путем передачи информации.
из данного выражения в его подвыражения,
Вывод вверх распространяет информацию в противоположном направлении.
Ограничения выражаются в терминах контекстных типов;
Быть неограниченным означает иметь {} в качестве контекстного типа.
Имеет тип контекста, который 
один или несколько случаев 
обеспечивает частичное ограничение на предполагаемый тип.%?

Если {коллекция} является однозначной картой
 
где  и  определяется нисходящим выводом,
и может быть 
Если нисходящий контекст не ограничивает одно или оба.

В противном случае {коллекция} неоднозначна,
и нисходящий контекст для элементов 
.


%
Мы говорим, что коллекция буквальный элемент
{может быть множество}{сборник буквальный элемент! может быть набором.
если он имеет тип элемента;
это
{может быть карта}{сборник буквальный элемент! Это может быть карта
если он имеет пару типа ключа и значения;
это
{должен быть набор}{сборник буквальный элемент! Должен быть набор.
если она может быть установленной и не иметь пары типа ключа и значения;
и это
{должна быть карта}{сборник буквальный элемент! Должна быть карта
может быть картой и не имеет типа элемента.

%
 Термин, полученный из {элемент}.
 Вывод типа {%}
Типовой вывод! Буквальный элемент коллекции
 Тип контекста  После этого доходы
следующим образом:

%
 Элемент выражения
В данном случае  Выражение .
%
Если  ,
Тип элемента  это
Предполагаемый тип  в контексте .
%
Если  ,
Тип элемента  это
Предполагаемый тип  в контексте .


%
 Элемент карты
В данном случае  Это пара выражений .
%
Если  ,
Выводимая пара типа ключа и значения  является ,
где  и  это
Предполагаемый тип  соответственно ,
в контексте .
%
Если  ,
Выводимая пара типа ключа и значения  является ,
где  это
Предполагаемый тип  в контексте  и
 это
Вычисленный тип  в контексте .


%
 Элемент распространения.
В данном случае  имеет форму
 или .
Если    быть
Предполагаемый тип  в контексте .
Тогда:



Если  реализация ,
Тип элемента  это
Аргумент типа  при .

{%
Это результат сопоставления ограничений для 
Используя ограничение  X.
Обратите внимание, когда  Реализует такой класс, как  или ,
не может быть подтипом 
()%
?

Если  реализация ,
Выводимая пара типа ключа и значения  является ,
где  является первым и  Второй тип аргумента
 в .

{%
Это результат сопоставления ограничений для  и  использовать
Ограничение  X Y$>}$%

Обратите внимание, что этот случай и предыдущий случай
может совпадать на одном и том же элементе одновременно
когда  Реализует оба  и .
Та же ситуация возникает несколько раз ниже.
В таких случаях мы полагаемся на другие элементы, чтобы разъяснить.
?

Если   затем
Предполагаемый тип элемента  является ,
и пара типа ключа и значения  это
.

{%
Здесь мы производим как тип элемента, так и пару типа ключа и значения.и полагаться на другие элементы для разбора.%
?

Если   и оператор спреда {...?} затем
Тип элемента  является ,
и пара типа ключа и значения .

В противном случае возникает ошибка времени компиляции.



% % TODO(Eernst): Уточните, почему мы не можем использовать контекстный тип
% % «Iterable» <$P_e$> Точно так же: вывод  в контексте «Итерируемый» <$P_e$>
%% также.
В противном случае, если   Тогда пусть  быть
Предполагаемый тип  в контексте , а затем:



Если  реализация ,
Тип элемента  это
Аргумент типа  в .
{%
Это результат сопоставления ограничений для  использовать
Ограничение  X$>}$%
?

Если  ,
Предполагаемый тип элемента  является .

Если   и оператор спреда {...?},
Предполагаемый тип элемента  является .

В противном случае возникает ошибка времени компиляции.



В противном случае, если   Тогда пусть  быть
Предполагаемый тип  в контексте , а затем:



Если  Для реализации ,
Выводимая пара типа ключа и значения  является ,
где  является первым и  Второй тип аргументации
 в .
{%
Это результат сопоставления ограничений для  и  использовать
Ограничение .%
?

Если  ,
Выводимая пара типа ключа и значения  это


.

Если   и оператор спреда {...?},
пара типа ключа и значения .

В противном случае возникает ошибка времени компиляции.

{-5 мм}


%
 Если элемент
В данном случае  имеет форму
 или
.
Условие  Всегда выводится контекстный тип .

Предположим, что "" отсутствует. Тогда:



Если выбранный тип элемента  является ,
Предполагаемый тип элемента  является .

Если пара типа ключа и значения  является ,
Выводимая пара типа ключа и значения  является .


В противном случае Присутствует .
Ошибка компиляции  Должен быть набор и  Должна быть карта,
или наоборот.

{%
Это означает, что один не может распространить карту на одну ветвь и набор на другую.
С тех пор \DYNAMIC{} обеспечивает как тип элемента, так и пару типа ключа и значения,
a \DYNAMIC{} Распространение в любой ветви не приводит к возникновению ошибки.%
?

Тогда:



Если выбранный тип элемента  является  и
Предполагаемый тип элемента  является ,
Тип элемента  это
наименьшая верхняя граница  и .
<<<p001703>
Если выбрана пара типа ключа и значения  это

и пара типа ключа и значения  это
,
Пара типа ключа и значения  это
,
где  является наименьшей верхней границей  и , и является наименьшей верхней границей  и .

{-5 мм}



 Для элемента
В данном случае  имеет форму

где  является производным от {forLoopParts} и
 Это означает, что  может присутствовать или отсутствовать.

Ошибки времени компиляции для  как
Ошибки, которые могут возникнуть при соответствующем  заявление
 \,\,\FOR\,\,$P$\,\{\}>,
Находится в том же объеме, что и .
Кроме того, ошибки и анализ типов  выполняется
как если бы это произошло в пределах объема тела указанного .

{%
Например, если  имеет форму

переменная  находится в пределах .%
?

Вывод для частей
<{%
(таких, как выразительное выражение фор-ин,
или \synt{forInitializerStatement} для петли)%
?
выполнено как для соответствующего \FOR{} заявление,
включая << если и только если элемент содержит .



Если выбранный тип элемента  является  затем
Предполагаемый тип элемента  является .

Если пара типа ключа и значения  является ,
Выводимая пара типа ключа и значения  является .


{%
Другими словами, вывод течет вверх от элемента тела.
?
\vspace{3 мм)


%
Наконец, мы определяем
{тип вывода на наборе или карте буквально}{%
Типовой вывод! Набор или карта буквально
в целом.
Предположим, что \metavar{collection} происходит от \synt{setOrMapLiteral},
и тип контекста для \metavar{collection} является .



Если <{коллекция} является однозначно установленным:



Если  \FreeContext затем
Статический тип  
где  является наименьшей верхней границей
Выводимые типы элементов.

% % TODO(Eernst): Спецификация функции говорит , но как мы знаем  Это тип?
В противном случае статический тип {сборник} является 
где  определяется путем нисходящего вывода.

{%
Обратите внимание, что вывод никогда не приведет к паре типа ключа и значения.
с данным типом контекста.%
?


Статический тип \metavar{сборник} является .

Если {коллекция} является однозначной картой
где   или  \FreeContext
и парами типа ключа и значения являются
{K}{V}{1}{n):

Если  {} или  ,
Тип статического ключа  
где  является наименьшей верхней границей {K}{1}{n}.
В противном случае статическим ключевым типом \metavar{collection} является $K$
где  определяется выводом вниз.

Если  {} или  ,
Тип статического значения \metavar{collection} $V$
где  является наименьшей верхней границей {V}{1}{n}.
В противном случае тип статического значения \metavar{collection} является $V$
где  определяется путем нисходящего вывода.

{%
Обратите внимание, что вывод никогда не приведет к типу элемента.
С учетом этого контекста.%
?

Статический тип {сборник} является .
В противном случае {коллекция} все еще неоднозначна,
нисходящий контекст для элементов 

Диагностика производится с использованием
Непосредственные элементы {сборник}:



Если все элементы могут быть набором,
По меньшей мере один элемент должен быть набором,
<{коллекция} - набор буквальных слов с
статический тип , где  это
наименьшая верхняя граница элементов типов элементов.

Если все элементы могут быть картой,
и хотя бы один элемент должен быть картой, то  это
Буквальная карта со статическим типом   это
наименьшая верхняя граница ключевых типов элементов и  это
наименьшая верхняя граница типов значений.

В противном случае возникает ошибка времени компиляции.
<<<p001764>> В этом случае буквальный не может быть расшифрован. ?



{%
Эта последняя ошибка может произойти, если буквальный {должен быть} и набором, и картой.
Вот пример: %
?




<{%
Или, если ничего не указывает на то, что это {или} набор или карта:%
?




<\subsubsection{Сеть)


%
A {set literal}{literal!set} обозначает заданный объект.
Правило грамматики для {setOrMapLiteral}
Набор буквальных и картографических букв происходит в другом месте.
().
Буквальный набор состоит из нулевых или более элементов.
(P000998).
Термин получен из {setOrMapLiteral}
может быть набор буквальный или карта буквальный,
и определяется с помощью шага раздвоения
Буквально ли это набор или карта
(, .

%
При заданном наборе буквальный  Нет никаких аргументов,
Тип аргумента  Выбирается так, как указано в другом месте
(),
и  отныне рассматривается как
()
.

{%
Статический тип набора буквальных форм 

()%
?

%
 быть буквальным набором формы
\code{$T$>\{\List{\ell}{1}{m}}.
Это ошибка времени компиляции, если элемент листа  является
{mapElement}.
Это ошибка времени компиляции, если для некоторых 
 не имеет элемента типа,
Тип элемента  Не может быть назначено .

%
Набор может содержать ноль или более объектов.
Наборы имеют метод, который можно использовать для вставки объектов;
Это приведет к динамической ошибке, если набор не может быть изменен.
В противном случае при вставке объекта  в набор ,
Если объект  Существует в  такой, что
 Показатель 
не вносит никаких изменений в ;
Если такого объекта не существует,
 Добавлено в ;
В обоих случаях вставка завершается успешно.

%
Набор упорядочен: итерация над элементами набора
происходит в том порядке, в котором элементы были добавлены в набор.

{%
Системные библиотеки определяют множество членов для типа ,
Но мы указываем только минимальный набор требований.
Используется самим языком.%

Обратите внимание, что реализация может потребовать
последовательные определения нескольких членов
класса, реализующего  Чтобы работать правильно.
Например, может быть геттер  который требуется
поведение, которое в некотором смысле согласуется с оператором; {==}.
Такие ограничения документируются в системных библиотеках.
?

%
Если набор буквальный  начинается с зарезервированного слова 
 происходит в постоянном контексте
(),
Это а{constant set literal}{literal!set!constant}
Что является постоянным выражением
()
и поэтому оценивается во время компиляции.
В противном случае это a
{run-time set literal}{literal!set!run-time}
и оценивается во время выполнения.
Только буквальные значения времени выполнения могут быть мутированы после их создания.
% Эта ошибка может произойти, потому что постоянное является динамическим свойством.
Попытка мутировать постоянный набор букв приведет к динамической ошибке.

{%
% Прямо или косвенно верно следующее: Существует 
Модификатор % на буквальном наборе, или мы используем правило «немедленной субэкспрессии».
% — постоянные контексты.
Обратите внимание, что выражения элементов постоянного множества буквальны
происходит в постоянном контексте
(),
что означает  Модификаторы не должны быть указаны явно.%
?

%
Ошибка времени компиляции, если
Литературный элемент в постоянном наборе
Это не постоянное выражение.
Ошибка времени компиляции, если
элемент в постоянном наборе буквальный
Не существует примитивного равенства.
().
Ошибка времени компиляции, если два элемента постоянного набора букв равны
В соответствии с их \lit{======================================================================================================================================================================================================================================================= оператор
().
Это ошибка времени компиляции, если аргумент типа постоянного множества буквален.
{(независимо от того, является ли это явным или предполагаемым)}
Не является постоянным выражением типа
().

{%
Связывание формального параметра типа замкнутого класса или функции
не известны во время составления,
Поэтому мы не можем использовать такие параметры типа в постоянных выражениях.
?

%
Значение постоянного множества буквально
 <\,\,$T$>\{\List{\ell}{1}{m}<<< P002058>>>
является объектом  чей класс реализует встроенный класс

где  - фактическое значение 
(),
и чье содержание представляет собой совокупность объектов в
объектная последовательность {o}{1}{n} полученный
Оценка \List{\ell}{1}{m}
().
Элементы  происходит в том же порядке, что и
объекты в указанной последовательности объектов
{(что можно наблюдать итерацией)}.

%
Пусть {? < \{\,$\ell_{11}, \ldots, \ell_{1m_1}$\,\}
 \,\,$T_2$>\{\,$\ell_{21}, \ldots, \ell_{2m_2}$\,\}
Будьте двумя постоянными буквами.
 Содержание  Аргумент типа 
соответственно
 Содержание  Аргумент типа 
быть результатом их оценки.
Далее  Показатель  если {f}
% % TODO(Eernst): См. понятие nnbd «один и тот же тип».
 и 
{} для всех .

{%
Другими словами, буквы с постоянным набором канонизируются, если они имеют
Один и тот же аргумент и одни и те же ценности в одном и том же порядке.
Два постоянных набора букв никогда не идентичны.
Если они имеют различное количество элементов.
Нет необходимости рассматривать канонизацию.
для других экземпляров типа ,
Поскольку такие случаи не
результат оценки постоянного выражения.%
?

%
Набор времени выполнения \code{$T$>\{\List{\ell}{1}{n}}
Оценка проводится следующим образом:



Оцениваются элементы \List{\ell}{1}{m}
(),
к объектной последовательности \LiteralSequence{\List{o}{1}{n}}.

Свежий объект () 
внедрение встроенного класса  Он был создан,
где  Это действительное значение 
().
Для каждого объекта  в \List{o}{1}{n}, по порядку
 Вводится в .
{%
Обратите внимание, что это оставляет  без изменений, когда  уже содержит
и объект  что равно  По словам оператора \lit{===.%
?

Результатом оценки является .


%
Предметы, созданные установленными буквами, не переопределяются
{==} оператор, унаследованный от P001183 класс.


\subsubsection{Карты}


%
A \IndexCustom{mapliteral}{буквально} обозначает объект карты,
Это отображение от ключей к значениям.
Правило грамматики для {setOrMapLiteral}, которое охватывает оба
Буквы карты и набор букв происходит в другом месте
().
Буквальная карта состоит из нулевых или более элементов.
().
Термин получен из \synt{setOrMapLiteral}
может быть набор буквальный или карта буквальный,
и определяется с помощью шага раздвоения
Буквально ли это набор или карта
(, .

%
При заданной карте буквальный  Нет никаких аргументов,
Тип аргументов  и  Выбирается так, как указано в другом месте
(),
и  отныне рассматривается как
()
.

{%
Статический тип карты буквальный формы 

(\ref{setAndMapLiteralInference})%
?

%
 Будьте картой буквальной формы
\code{$K$,\,$V$>\{\List{\ell}{1}{m}}.
Это ошибка времени компиляции, если элемент листа  является
{выражениеЭлемент}.
Это ошибка времени компиляции, если для некоторых 
 не имеет пары ключа и значения;
или парой типа ключа и значения  является ,
 не может быть присвоено  или
 Не может быть назначено .

%
Объект карты состоит из нуля или более записей карты.
Каждая запись имеет  и ,
Мы говорим, что карта
{binds}{map!binds} или
{maps}{map!maps}
Ключ к ценности.
Парой ключей и ценностей является
Добавлено на карту с помощью оператора \lit{[=],
и значение для данного ключа извлекается из карты с помощью оператора {[]}.
Ключи карты обрабатываются аналогично набору
():
При связывании ключа  значение  на карте 
\commentary{(как в \code{[]\,=\,})},
Если  Имеется ключ  такой, что
 Показатель ,
 Привязывается  к ;
иначе
{(когда такой ключ  не существует)},
Обязательный от  до  Добавлено в .

%
Упорядочена карта: итерация по ключам, значениям или парам ключ/значение
происходит в том порядке, в котором клавиши были добавлены в набор.

{%
Системные библиотеки поддерживают множество операций на экземпляре
чей тип реализуется ,
Но мы указываем только минимальный набор требований.
которые используются самим языком.

Обратите внимание, что реализация может потребовать
последовательные определения нескольких членов
класса, реализующего  Чтобы работать правильно.
Например, может быть геттер  который требуется
поведение, которое в некотором смысле согласуется с оператором; {==}.
Такие ограничения документируются в системных библиотеках.
?

%
Если карта буквальная  начинается с зарезервированного слова ,
или если  происходит в постоянном контексте
(),
Это а
{постоянная карта буквально}{буквально!map!constant}
Что является постоянным выражением
()и поэтому оценивается во время компиляции.
В противном случае это a
<{run-time mapliteral}{literal!map!run-time}
и оценивается во время выполнения.
Только буквальные карты времени выполнения могут быть мутированы после их создания.
% Эта ошибка может произойти, потому что постоянное является динамическим свойством.
Попытка мутировать постоянную карту буквально приведет к динамической ошибке.

{%
% Прямо или косвенно верно следующее: Существует 
Модификатор % на буквальной карте, или мы используем правило «немедленной субэкспрессии»
% — постоянные контексты.
Обратите внимание, что ключевые и ценностные выражения постоянной карты буквальны.
происходит в постоянном контексте
(),
что означает  Модификаторы не должны быть указаны явно.%
?

%
Это ошибка времени компиляции
Литературный элемент в постоянной карте не является постоянным.
Ошибка времени компиляции, если
Ключ в постоянной карте
Не существует примитивного равенства.
().
Это ошибка времени компиляции, если два ключа постоянной карты равны.
В соответствии с их {=== оператор
().
Это ошибка времени компиляции, если аргумент типа постоянной карты буквален.
{(независимо от того, является ли это явным или предполагаемым)}
Не является постоянным выражением типа
().

{%
Связывание формального параметра типа замкнутого класса или функции
не известны во время составления,
Поэтому мы не можем использовать такие параметры типа в постоянных выражениях.
?

%
Значение постоянной карты буквально
 \,\,$T_1$\,$T_2$>\{\List{\ell}{1}{m}
является объектом  чей класс реализует встроенный класс
,
где  и  - фактическое значение  соответственно 
().
Пары ключей и значений  это
пары объектной последовательности {k}{v}{1}{n} полученный
Оценка \List{}{1}{m}
(),
в таком порядке.

%
Пусть ? \,\,$U_1$\,$V_1$>\{\List{\ell}{1}{m_1}
 <\,<\,<$U_2$\,$V_2$>\{\List{\ell}{1}{m_2}
Будьте двумя постоянными буквами карты.
Пусть  с содержимым \KeyValueList{k_1}{v_1}{1}{n)
Реальные типы аргументов , 
соответственно
 с содержанием \KeyValueList{k_2}{v_2}{1}{n)
Реальный тип аргумента , 
быть результатом их оценки.
Далее  Показатель  если {f}
, ,


Для всех (P000653).

{%
Другими словами, постоянные буквальные обозначения карт канонизированы.
Нет необходимости рассматривать канонизацию.
Для других типов ,
Поскольку такие случаи не могут быть
результат оценки постоянного выражения.%
?

%
Карта времени выполнения буквальная
<\code{$T_1, T_2$>\{\List{\ell}{1}{m}}}
Оценка проводится следующим образом:



Оцениваются элементы \List{\ell}{1}{m}
(),
в объектную последовательность \LiteralSequence{\KeyValueList{k}{v}{1}{n}}.

Свежий пример () 
чей класс реализует встроенный класс 
выделяется,
где  и  являются действительными значениями  соответственно 
().

Оператор \lit{[]=} вызывается на $o$
с первым аргументом  и второй аргумент ,
Для каждого , в таком порядке.
Результатом оценки является .


%
Объекты, созданные буквальными картами, не переопределяются
тот {=== оператор, унаследованный от  класс.


 Бросок.


%
 Используется для исключения.


<rowExpression> ::= << <выражение>

<rowExpressionWithoutCascade> ::= \THROW{} <expressionWithoutCascade>


%
Оценка броскового выражения формы
\code{{} ;
идет следующим образом:

%
Выражение  оценивается на предмет 
().

{%
Нет требования, чтобы выражение  должны оценивать
Любые специальные объекты.%
?

%
Если  является нулевым объектом (), затем бросается .
В противном случае пусть  быть следом стека, соответствующим текущему состоянию исполнения,
и  Заявление с броском  как объект исключения
 След стека ().

%
Если  является примером класса  или его подкласса,
Это первый случай, когда  объект выбрасывается,
След стека  Сохраняется на  Чтобы он был возвращен
 Геттер унаследован от .

{%
Если то же самое  объект выбрасывается более одного раза,
 Геттер вернет стековый след
< >> {первое время было выброшено.%
?

%
Статический тип выражения броска .


<{Функциональные выражения}


%
A {функция дословная}{дословная!функция}
является анонимной декларацией и выражением
который инкапсулирует исполняемый блок кода.

% % TODO(Eernst): Это очень неоднозначно, потому что <первичный> получает
% <выражение функции>. Dart.g извлекает <functionExpression> из
% % <expression>, но допускает только блок в <functionPrimary>, который
% является производным от <primary>. Функция «ExpressionWithoutCascade»
%%, что позволяет выполнять функцию «=» буквально в каскадном задании.
% % Тем не менее, мы должны очень тщательно оценить поломку, прежде чем принимать
% от этого подхода, поскольку он предотвращает разбор буквальных функций
%% как условное выражение, если нулевое выражение, ... постфиксное выражение.

<functionExpression> ::= <formalParameterPart> <functionExpressionBody>

<functionExpressionBody> ::= ? <=> <выражение>
 (\ASYNC{} '*'? | \SYNC{} '*')) <блокировать>


%
Грамматика не позволяет функции буквально объявить тип возврата,
Но возможно, что функция, дословно имеющая
<{объявлено о возвращении Тип {буквальный!функциональный!декларированный тип возврата},
потому что она может быть получена посредством вывода типа.
Такой тип возврата включает
Когда мы говорим об объявленном типе функции возврата.

{%
Тип вывода будет указан в будущей версии этого документа.
В настоящее время мы рассматриваем вывод типа как завершенный этап,
Этот документ определяет значение программ Дарта.
где уже добавлены предполагаемые типы.%
?

%
Мы говорим, что тип 
{derives a future type}{type!derives a future type}
 в следующих случаях, используя первый применимый случай:

% % TODO(Eernst): Обратите внимание, что «X расширяет X?» может создать бесконечный цикл.
% % Возможно, мы сделаем такую ошибку.


% % TODO(Eernst): Приходите классы микширования и типы расширений: добавьте их.
Если  Это тип, который вводится
класс, микширование или декларирование enum,
Если  Прямой или косвенный суперинтерфейс
()
  для некоторых , Получает будущий тип .

Если  Тип  для некоторых ,
 Получает будущий тип .
<<<p001900>
Если   Для некоторых 
 Получает будущий тип ,
 Получает будущий тип .
<<<p001901>
Если  является переменной типа со связанным , и
 Получает будущий тип ,
 Получает будущий тип .
<<<p001902>
{%
Не существует правила для случая, когда  имеет форму 
Потому что этого никогда не произойдет
(это понятие используется только в , который определен ниже.)
?


%
Если ни один из этих случаев не применим,
Мы говорим, что  Он не имеет будущего типа.

{%
Обратите внимание, что если  Получает будущий тип  < T}{F},
 всегда имеет форму  или ,
где   или \code{FutureOr}. Доказательство является
по индукции на структуру :


<<<p001907>
% % TODO(Eernst): Приходите классы микширования и типы расширений: добавьте их.
Если  Это тип, который вводится
класс, микширование или декларирование enum,
Если  Прямой или косвенный суперинтерфейс
()
  Для некоторых , тогда
 и , .
<<<p001908>
Если  Тип  Для некоторых , то по рефлексивности,
.  и
, следует, что .
<<<p001909>
Если   Для некоторых 
 Получает будущий тип ,
С помощью индукционной гипотезы, ,
где  имеет форму
 или  и   или
. Следовательно, , и путем замены,
. С тех пор  для всех 
. Итак, позволив  и ,
Отсюда следует, что .

Если  является переменной типа со связанным , и
 Получает будущий тип ,
С помощью индукционной гипотезы, ,
где  имеет форму
 или  и   или
. Кроме того, с тех пор  является пределом , 
транзитивность, . Таким образом, позволяя  и , это
Далее следует .


Также обратите внимание, что «производные» в этом контексте относятся к вычислениям.
где тип  Приведены супертипы  обыскиваются,
и тип  Выбирается одна из этих форм.
Нет никакой связи с понятием «производный класс», означающим «подкласс». "
Некоторые языковые сообщества используют.%
?

%
Определяем вспомогательную функцию
\IndexCustom{{T}}{flatten(t)@{flatten}}
следующим образом, используя первый применимый случай:


 Если  
переменная типа  Тип  затем


, если  Получает будущий тип 
затем <\DefEquals{\flatten{T}}{\code{\flatten !
 В противном случае
\DefEquals{{T}}{{ Икс.


 Если  Получает будущий тип 
 T}}{S}

 Если  Получает будущий тип 
\  \DefEquals{\flatten{T}}{\code{?}}.

<<<p001931> В противном случае \DefEquals{\flatten{ T}}{T}.


{%
Это определение гарантирует, что для любого типа 
. Доказательством является индукция на
Структура :


<<<p001935> Если   затем


, если  Получает будущий тип ,
Затем <- и , так .
По индукционной гипотезе .
С тех пор  В данном случае следует, что
<
.
<<<p001937> В противном случае .
По индукционной гипотезе, .
С тех пор  В данном случае следует, что
<
.


 Если  Получает будущий тип 
или , то, поскольку ,
Отсюда следует, что . С тех пор 
Из этого следует, что .

 Если  Получает будущий тип  или
, с тех пор ,
Отсюда следует, что .
 для любого типа 
(это может быть показано с использованием правил подтипа союзного типа и
 по ковариантности, т. е. по транзитивности,
. С тех пор  в данном случае,
Отсюда следует, что .

 В противном случае, 
. С тех пор
< Из этого следует, что
.

?

%
 Позиция, стрелка.
Статический тип функции буквальный формы








это
{T_0}

<<<p001946>
% % [вывод]: Статический тип функции буквально может быть выведен.
где  Статический тип .


%
 Позиция, стрелка, будущее.
Статический тип функции буквальный формы





 () [] \ASYNC{} => $e$>


это
 Будущее<\flatten{T_0}>


где  Статический тип .


%
 Названа, стрелка.
Статический тип функции буквальный формы





 \{$T_{n+1}\ x_{n+1} = d_1, \ldots,\ T_{n+k}\ x_{n+k} = d_k$\}> 

<<<p001963>
это
{T_0}


где  Статический тип .
<<<p001966>

%
{Название, стрелка, будущее}
Статический тип функции буквальный формы





 () \{$T_{n+1}\ x_{n+1} = d_1, \ldots,\ T_{n+k}\ x_{n+k} = d_k$\} \ASYNC{} => $e$


это
 Будущее<\flatten{T_0}>


где  Статический тип .


%
 Позиционно, блок.
Статический тип функции буквальный формы

<<<p001979>
\code{\TypeParameters{X}{B}{S}>


 () []   


это

<<<p001987>

%
 Позиция, блок, будущее.
Статический тип функции буквальный формы





 () [] < <<<p002101>  

<<<p001993>
это
% % TODO(Eernst): Настройка для учета выводов типа..


%
<<<p001996> Позиция, блок, поток.
Статический тип функции буквальный формы





 () [] <   


это
% % TODO(Eernst): Настройка для учета выводов типа.
\FunctionTypePositionalStdCr{\code{Stream}}.


%
{Позиционное, блоковое, итерабельное}
Статический тип функции буквальный формы





 () [] <   


это
% % TODO(Eernst): Настройка для учета выводов типа.
.


%
<{Название, блок)
Статический тип функции буквальный формы

<<<p002013>>



 () [] <<<p002107>  


это
% % TODO(Eernst): Настройка для учета выводов типа.
.


%
{Название, блок, будущее}
Статический тип функции буквальный формы





 () \{$T_{n+1}\ x_{n+1} = d_1, \ldots,\ T_{n+k}\ x_{n+k} = d_k$\} < >>  $s$\}


это
% % TODO(Eernst): Настройка для учета выводов типа.
.


%
{Название, блок, поток}
Статический тип функции буквальный формы





 \{$T_{n+1}\ x_{n+1} = d_1, \ldots,\ T_{n+k}\ x_{n+k} = d_k$\} <   


это
% % TODO(Eernst): Настройка для учета выводов типа.
.


%
{Название, блок, итерабельный}
Статический тип функции буквальный формы





 \{$T_{n+1}\ x_{n+1} = d_1, \ldots,\ T_{n+k}\ x_{n+k} = d_k$\} <  <<<p000016> 


это
% % TODO(Eernst): Настройка для учета выводов типа.
.


%
Во всех вышеперечисленных случаях,
Списки аргументов типа опускаются, когда 
и всякий раз  Не уточняется, ,
Считается, что он был указан как .

%
Оценка функции буквально приводит к объекту функции .

{%
Тип времени выполнения  определяется на основе
Статический тип  Функция буквальная
и связывания переменных типа, происходящих в 
в случае, когда оценка
()%
?


 Это.
<<<p001231>

%
Зарезервировано слово  обозначает
Цель вызова нынешнего члена инстанции.


<ThisExpression> ::= \THIS{}


%
Статический тип {} является интерфейсом
Немедленно включающий класс, энум или микширование, если таковые имеются.
Статический тип  это
 типом ограждающего расширения, если таковое имеется
(P000866).

{%
Если ни одна из этих деклараций не существует,
Возникновение {} является ошибкой времени компиляции
()%
?

%
Это ошибка времени компиляции, если {} появляется, неявно или явно,
в функции верхнего уровня или переменном инициализаторе, в заводском конструкторе,
или в статическом способе или переменном инициализаторе,
или в инициализации выражения переменной непоздней инстанции.


<<<p001280>> Создание инстанций.


%
Выражения создания инстанций обычно дают примеры
И призывают строителей их инициализировать.

{%
Исключением является то, что
Призыв заводского конструктора работает как обычный вызов функции.
Он может, конечно, оценить выражение создания экземпляра и, таким образом,
Приведите свежий пример,
Но никакие новые примеры не создаются как прямое следствие
Производитель Invoction.%
?

%
Ошибка времени компиляции, если
тип  в случае создания выражения одной из форм


\code{{} . (: : 


\code{{) (: : 


\code{{} $T$\id: : 


\code{\CONST{}  (): : 


Перечисленный тип ().


 Новый.


%
 Строитель (P000869).


<newExpress> ::= \NEW{} <конструкторПроектирование> <аргументы>


%
 Быть новым выражением формы


\code{{) $T$\id: : 
или форма


\code{{) (: : .
%
Ошибка времени компиляции  не является
класс или параметризированный тип, доступный в текущем объеме;
или если  Это параметризированный тип, который не является классом.
{%
Например, \code{{} F<int>()} является ошибкой, если  Типовой псевдоним
не обозначает класс.%
?

%
Если  параметризованный тип ()
,
 быть общим классом ,
И пусть

Формальные параметры типа .
Если  не является параметризованным типом, пусть  будет .



Если  имеет форму
\code{{) $T$\id (: : 
Ошибка времени компиляции, если  Это не имя
Конструктор, объявленный , или  недоступна.

Если  имеет форму
\code{{) (: : 
Ошибка времени компиляции  Это не имя
Конструктор, объявленный .


%
Пусть  будет вышеупомянутым конструктором, названным  или .

%
Ошибка времени компиляции  абстрактный
 Это не заводской конструктор.
Ошибка времени компиляции  является негенерическим классом
 - параметризованный тип.
% % Мы предполагаем, что вывод имел место, поэтому фактические аргументы типа
% всегда указывается явно.
Ошибка времени компиляции  Является общим классом
 Это не параметризированный тип.
Ошибка времени компиляции  Это общий класс,
 Параметризованный тип и .
 То есть количество аргументов типа неверно. ?
Ошибка времени компиляции  Это общий класс,
 является параметризованным типом,
 не является регулярным
(P000871).

%
Если  является перенаправляющим заводским конструктором,
Ошибка времени компиляции  В некоторых количествах
Перенаправление заводских перенаправлений перенаправляется на себя.
{%
Возможно и разрешено перенаправлять завод 
Чтобы войти в бесконечный цикл, например,
 Перенаправление на неперенаправленный заводской конструктор
 Тело которого использует  в выражении создания экземпляра.
Только петли, которые состоят исключительно из перенаправления заводских перенаправлений
Обнаружено в момент компиляции.%
?

%
 быть статичным типом
Формальный параметр конструктора  (перенаправлено с «P000099»)
соответствует фактическому аргументу , .
Ошибка времени компиляции, если статический тип

не присваивается .
{%
Необычный случай покрыт .%
?

%
Статический тип  является .

%
Оценка  идет следующим образом:

%
Во-первых, аргументная часть





Оценивается, получая оцененную фактическую часть аргумента


.


 Обратите внимание, что необычный случай охватывается разрешением .
% Эта ошибка может произойти из-за неявного слепка.
Если для любого

Тип времени выполнения  не является подтипом
,
Возникает динамическая ошибка.

%
\Case{Незагруженные отложенные конструкторы}
% Эта ошибка может произойти, потому что загрузка является динамическим свойством.
Если  отложенный тип с префиксом ,
Если  не был успешно загружен,
Возникает динамическая ошибка.


%
 Генеративные конструкторы.Когда  является генеративным конструктором
()
оценка приступает к выделению нового экземпляра
(), , класса .
% Мы приводим аргументы типа как часть класса экземпляра.
%, потому что  Включает тип аргументов, но мы также предоставляем их.
% как связующее звено, доступное конструктору: Иначе мы не смогли бы
% доступа к параметрам типа при инициализации выражений
Список инициализаторов %, где нет доступа к .
 Выполняется для инициализации  в отношении
(b) привязки, возникшие в результате оценки списка аргументов, и
Если  Это общий класс,
с параметрами типа, связанными с .

%
Если исполнение  обычно (),
 Показатель .
В противном случае выполняется  Выбрасывает объект исключения  и стековый след ,
и затем оценка  также бросает объект исключения  и след стека 
(P000875).


%
 Перенаправление заводских строителей
Когда  является перенаправляющим заводским конструктором
()
Форма  или
в форме 
где  Это означает, что  может присутствовать или отсутствовать,
Оставшаяся оценка  является эквивалентным
оценка
\code{{)  (: : 
в среде, где
 является свежей переменной, связанной с  для , и
  для .
{%
Нам нужен доступ к переменным типа, потому что  Может содержать их.%
?


%
 Неперенаправляющие заводские конструкторы
Когда  является неперенаправляемым заводским конструктором,
Тело  исполняется в отношении
(b) привязки, возникшие в результате оценки списка аргументов,
и с параметрами типа, если таковые имеются,  связанный
Аргументы реального типа .
Если это исполнение возвращает объект
()
Затем  оценивает возвращаемый объект.
В противном случае, если исполнение завершается нормально или возвращается без объекта,
 оценивает нулевой объект ().
В противном случае исполнение бросает исключение  и след стека ,
и затем оценка  Также бросок  и 
(P000879).

{%
Заводской конструктор может быть объявлен в абстрактном классе и использоваться безопасно.
как он будет производить либо действительный экземпляр или бросок.%
?



 Конст!


%
 Призывает постоянного строителя
(P000880).


<constObjectExpression>:= \CONST{} <конструкторПроектирование> <аргументы>


%
 быть постоянным выражением объекта формы


\code{{) . : : 
или форма


\code{{)  (: : .

%
Ошибка времени компиляции  не является
класс или параметризированный тип, доступный в текущем объеме;
или если  Это параметризированный тип, который не является классом.
Ошибка времени компиляции  является отложенным типом
(P000881).
 В частности,  не должна быть переменной типа.

%
Ошибка времени компиляции  Не является постоянным выражением
для некоторых .

%
Если  параметризованный тип ()
,
 быть общим классом ,
И пусть

Формальные параметры типа .Если  не является параметризованным типом, пусть  будет .

%
Если  является параметризованным типом,
Ошибка времени компиляции  не является постоянным выражением для любого
.



Если  имеет форму
\code{{) . : : 
Ошибка времени компиляции  Это не имя
Константный конструктор, объявленный , или  недоступна.

Если  имеет форму
\code{{) : : 
Ошибка времени компиляции  Это не имя
Константный конструктор, объявленный .


%
 быть упомянутым выше константным конструктором под названием  или .

% % TODO(Eernst): Эти ошибки такие же, как и в случае «новых». Можно ли избежать
Проценты говорят дважды. Нам нужно обратиться к неуклюжей части
% текста в предыдущем подразделе, или просто сказать "точно
% одинаковых ошибок.
%
Ошибка времени компиляции  абстрактный
 Это не заводской конструктор.
Ошибка времени компиляции  является негенерическим классом
 - параметризованный тип.
% % Мы предполагаем, что вывод имел место, поэтому фактические аргументы типа
% всегда указывается явно.
Ошибка времени компиляции  Является общим классом
 Это не параметризированный тип.
Ошибка времени компиляции  Это общий класс,
 является параметризованным типом и .
 То есть количество аргументов типа неверно. ?
Ошибка времени компиляции  Это общий класс,
 является параметризованным типом,
 не является регулярным
(P000883).

%
 быть статичным типом
Формальный параметр конструктора  (перенаправлено с «P000229»)
в соответствии с фактическим аргументом , .
Ошибка времени компиляции, если статический тип

не присваивается .
 Необычный случай покрыт .

%
Статический тип  является .

%
Оценка  идет следующим образом:

%
Если  имеет форму
\code{{) .  (: : 
 значение выражения :


\code{{) . : : .

{%
 быть результатом оценки ,
в определенный момент времени выполнения программы
в библиотеке , где $e$ происходит.
Результат оценки  в 
в другое время и/или в каком-либо другом исполнении
получить результат , такой, что  Заменить на 
Канонизация, как описано ниже.
Это означает, что значение хорошо определено.%
?

%
Если  имеет форму
\code{{)  (: : 
 быть ценностью
\code{{) : : .
 Что вполне понятно по той же причине.


 Если во время выполнения программы,
Постоянное выражение объекта уже оценивается как
Пример  Класс  Типовые аргументы :


 Для каждого экземпляра переменная  , быть значением переменной экземпляра  в , и
 быть значением переменной экземпляра  в .
Если 
Для всех переменных  
Тогда значение  составляет ,
В противном случае значение  составляет .

 В противном случае значение  составляет .


{%
Другими словами, постоянные объекты канонизированы.
Чтобы определить, является ли объект на самом деле новым, нужно его вычислить.
Затем его можно сравнить с любыми кэшированными экземплярами.
Если в кэше есть эквивалентный объект,
Мы выбрасываем вновь созданный объект и используем кэшированный.
Объекты эквивалентны, если
Они имеют одинаковые аргументы и одинаковые переменные экземпляров.
Поскольку конструктор не может вызвать никаких побочных эффектов,
Исполнение конструктора не наблюдается.
Конструктор должен быть выполнен только один раз на сайте вызова, во время компиляции.
?

%
Ошибка времени компиляции при оценке постоянного объекта
Это приводит к тому, что выбрасывается непойманное исключение.

{%
Чтобы увидеть, как такие ситуации могут возникнуть, рассмотрим следующие примеры:
?

% % TODO(Eernst): Удалите некоторые {} при интеграции имплицитного создания. мд


{%
Из-за правил, регулирующих постоянных строителей,
Оценка конструктора 
с аргументом  или аргумент \code{{} IntPair(1, 2)}
Это может привести к тому, что он бросит исключение, что приведет к ошибке времени компиляции.
В последнем случае ошибка вызвана тем, что
\code{{} Их можно использовать только с несколькими «известными» типами.
Для того, чтобы избежать произвольного кодирования во время
Оценка постоянных выражений.%
?


 Изолирующий термин.


%
Изолировать Isolate Срок достигается через то, что синтаксически
Обычный метод вызова,
использование одного из статических методов  или  определяемый
 класс в библиотеке .
Однако такие призывы имеют смысловой эффект создания
Новый Изолятор с собственной памятью и нитью управления.

%
Память изолированного термина конечна, как и пространство, доступное для
Стек вызовов его нити.
% Эта ошибка может произойти, потому что использование памяти является динамическим свойством.
Возможно, что запущенный изолят исчерпает свою память или стек.
В результате возникает динамическая ошибка, которая не может быть эффективно обнаружена.
Это приведет к тому, что изолятор будет приостановлен.

{%
Как описано в разделе ,
Обработка приостановленного изолят является ответственностью времени выполнения.%
?


 Функция вызова.


%
Призыв к функции происходит в следующих случаях:
выражение функции ()
(перенаправлено с «P000886»)
когда метод (),
Геттер (), 
или сеттер ()
Призывается,
или когда застройщик
(либо путем создания экземпляра (),
перенаправление конструктора (),
или супер инициализации.
Различные виды функции призвания различаются по
Как определяется функция, на которую следует ссылаться, ,
Как и в случае с  () является обязательным.
Однажды  был определен,
Формальные параметры типа  связаны с
соответствующие фактические типы аргументов,
и формальные параметры  связаны с соответствующими фактическими аргументами.
Когда тело  исполняется, будет исполняться
с вышеупомянутыми переплетами.

%
Выполнение корпуса формы  эквивалентно исполнению
а тело формы \{ возвращение ; .Исполнение тела формы \code{{=>  эквивалентно исполнению
тело формы \code{\ASYNC{}  возвращение ; .

%
Если  является синхронным и не является генератором () затем
Выполнение корпуса  Начинается немедленно.
Если казнь тела  возвращает объект 
(),
Вызов оценивается в .
Если исполнение завершается нормально или возвращается без предмета,
вызов оценивает нулевой объект ().
Если исполнение бросает объект исключения и стекает след,
Вызов бросает тот же самый объект исключения и след стека
(P000897).

{%
Полная функция тела никогда не может сломаться или продолжиться.
()
потому что \BREAK{} или \CONTINUE{} Заявление всегда должно происходить внутри
Заявление, являющееся целью {} или .
Это означает, что функциональное тело может
Либо завершить нормально, бросить, либо вернуться.
Обычное завершение или возвращение без предмета обрабатываются.
то же самое, что возврат с нулевым объектом (),
Таким образом, результат выполнения функции тела всегда может быть использован как
результат оценки выражения,
либо оценивая объект, либо оценивая бросание.%
?

%
Если  Помечено  (),
Свежий пример () 
осуществление  немедленно возвращается,
где  Это реальный тип
()
соответствует типу элемента 
().

{%
Реализация Дарта должна обеспечить
Конкретная реализация 
Возвращается по адресу  методы.
Типичной стратегией было бы привести пример
Подкласс класса  Определено в .
Единственный метод, который нужно добавить
по реализации Dart в этом случае .%
?

%
Реализация должна соответствовать
Контракт  и не должны
принимать любые меры, признанные исключительно эффективными в этом договоре.

{%
В договоре прямо упоминается ряд ситуаций
где определенные итерации могут быть более эффективными, чем обычно.
Например, путем предварительного вычисления их длины.
Нормальные итерабельные элементы должны повторяться над своими элементами, чтобы определить их длину.
Это верно в случае синхронного генератора.
где каждый элемент вычисляется функцией.
Было бы неприемлемо предварительно вычислять результаты генератора.
и кэшировать их, например.%
?

%
Когда итерация над итерируемым начинается, получая
Итератор  от итерируемого и вызывающего ,
Выполнение корпуса  Начнем.
При исполнении корпуса  Полностью (),


 Если он возвращается без предмета или завершается нормально.
(),
 расположен после последнего элемента,
Таким образом, его текущее значение является нулевым объектом.
()
и текущий вызов на  на  возвращает ложные
Как и все остальные призывы.
 Если он бросает объект исключения  и след стека 
Тогда текущее значение  является нулевым объектом
()
и текущий призыв к  броски  и  Тоже самое.
Дополнительные вызовы  Вернуть ложное.


%
Каждый итератор запускает отдельное вычисление.
Если  функция нечистая,
Последовательность объектов, полученных каждым итератором, может отличаться.

{%
Можно получить более одного итератора из данного итератора.
Обратите внимание, что операции на самом итераторе могут создавать различные итераторы.Примером может служить .
Можно предположить, что различные итераторы могут давать
Последовательности разной длины.
Такую же осторожность нужно соблюдать при написании  функционирует как когда
Написать  класс.
В частности, он должен изящно обрабатывать несколько одновременных итераторов.
Если итератор зависит от внешнего состояния, которое может измениться,
Он должен проверить, что государство по-прежнему действует после каждой доходности.
(перенаправлено с «P001033») Если это не так.%
?

%
Каждый итератор работает со своими небольшими копиями всех локальных переменных.
В частности, каждый итератор имеет одинаковые исходные аргументы.
Даже если их связывания изменяются функцией.
<<%
Два исполнения итератора взаимодействуют только через состояние вне функции.
?
% Альтернативой будет кэширование результатов итератора в итерабельном,
% и проверьте кэш на каждом . Это также будет иметь странные проблемы.
% Полученная стоимость может отличаться от выражения в доходности. И это -
% потенциальная утечка памяти, поскольку кэш поддерживается любым итератором.

%
Если  Помечено \ASYNC{} (),
Свежий пример () 
Он связан с призывом,
где динамический тип  реализация ,
где  Это реальный тип
()
соответствует типу будущего значения .
Затем тело  выполняется до тех пор, пока не будет приостановлено или завершено,
После чего возвращается .
{%
Тело  может быть приостановлено во время оценки  выражение
или выполнение асинхронного \FOR{} петля.%
?
Будущее (P000326) завершается при исполнении корпуса  полный
().
Если выполнение тела возвращает объект,  завершается этим объектом.
Если она завершается нормально или возвращается без предмета,
 завершается нулевым объектом (),
и если он делает исключение  и след стека ,
 Завершается ошибка  и след стека .
Если выполнение тела бросает перед тем, как тело приостанавливается в первый раз,
Обсуждение  Происходит это в какой-то момент после того, как призыв вернулся.
{%
Звонителю нужно время, чтобы настроить обработку ошибок для возвращенного будущего.
Будущее не завершается ошибкой
<<<<<<<<<<<<<<<<<<<<<<<<<<<<<<<<<<<<<<<<<<<<<<<<<<<<<<<<<<<<<<<<<<<<<<<<<<<<<<<<<<<<<<<<<<<<<<<<<<<<<<<<<<<<<<<<<<<<<<<<<<<<<<<<<<<<<<<<<<<<<<<<<<<<<<<<<<<<<<<<<<<<<<<<<<<<<<<<<<<<<<<<<<<<<<<<<<<<<<<<<<<<<<<<<<<<<<<<<<<<<<<<<<<<<<<<<<<<<<<<<<<<<<<<
?

%
Если   (),
Свежий пример () 
осуществление  немедленно возвращается,
где  Это реальный тип
()
соответствует типу элемента 
().
Когда  Прослушивается, исполнение тела  Начнем.
При исполнении корпуса  завершается:


 Если он завершается нормально или возвращается без объекта
(),
Если  Был отменен
Тогда его отмена в будущем завершается нулевым объектом ().
 Если он бросает объект исключения  и след стека :


 Если  Он был отменен, а его будущее аннулировано.
с ошибкой  и след стека .
 Ошибка  и след стека  Излучается с помощью .

  Закрыто.


{%
Тело асинхронного генератора функции
Не может сломаться, продолжиться или вернуться с объектом.
(P000919).
Первые два допускаются только в тех случаях, когда
Выполнит перерыв или продолжит,
Возвращать заявления с выражением не допускается
Функции генератора.%
?
{%
Когда поток асинхронного генератора был отменен,
Очистка будет происходить в  (P000920) внутри генератора.
Мы выбираем любые исключения, которые происходят в это время.
Будущее отмены, а не их потеря.%
?


<{Списки фактических аргументов}


%
Фактические списки аргументов имеют следующий синтаксис:


<argument> ::=('' (<argumentList> '','')?'' "

<argumentList> ::= <namedArgument> (',' <namedArgument>)*
\alt <expressionList> (',' <namedArgument>)*

<namedArgument> ::= <label> <expression>


%
Списки аргументов позволяют необязательную запятую после последнего аргумента
({','?}).
Список аргументов с такой задней запятой эквивалентен во всех отношениях.
Тот же список параметров без задней запятой.
Все списки аргументов в этой спецификации показаны без запятой,
Правила и семантика в равной степени применимы к
Соответствующий список аргументов с запятой.

%
 быть перечнем аргументов формы

При этом следует исходить из того, что вычисляемый тип является узким, а узкий — узким. <<<<p><p><p><p><p><p><p>>.
  тогда
.

%
 Статический тип списка аргументов





 Будьте формальным списком параметров





где каждый параметр может быть отмечен 
\commentary{(не показано, но разрешено)}.

%
Мы говорим, что  это
а 
f ,  и  это
Подтип  для всех .
Мы говорим, что  это
 
f ,  и  это
Применяется к  для всех .

%
 Статический тип списка аргументов





 Будьте формальным списком параметров


\code{(\{$T_{n+1}\ x_{n+1} = d_1, \ldots,\ T_{n+k}\ x_{n+k} = d_k$\})


где каждый параметр может быть отмечен 
{(не показано, но разрешено)}.

%
Мы говорим:  это
а  
f ,
,
 является подтипом  для всех ,
 является подтипом   и
 Для всех
.
Мы говорим, что 
 для 
f ,
,
 присваивается  для всех ,
 Применяется для   и
 Для всех
.

{%
Короче говоря, фактический список аргументов соответствует официальному списку параметров.
когда первый можно безопасно передать второму.
?


\subsubsection{Оценка списка фактических аргументов}


%
Функциональное призвание предполагает оценку
список фактических аргументов для функции,
Связывание результатов с формальными параметрами функции.

%
При анализе списка аргументов может возникнуть двусмысленность, потому что
Один и тот же исходный код может быть одним общим вызовом функции.
Это могут быть два или более реляционных и/или сменных выражений.
В этой ситуации выражение всегда разбирается
как родовое призвание функции.

% Должны ли мы указать точное правило неоднозначности здесь?:
% Мы видели 'a', '<', соответствующий '>' и '(', где
% «а» сложно, потому что у нас могут быть такие вещи, как «новый Фу».
% 'x..y(...).m<...>(...) и т.д. в основном все, что может предшествовать
% аргументация Часть в грамматике.

{%
Примером может служить ,
который может быть вызовом  проход
два действительных аргумента типа , илиСсылка на  результат, возвращенный путем
a вызов общей функции .
Обратите внимание, что двусмысленность можно устранить, опустив
скобки вокруг выражения ,
или добавление скобок вокруг одного из реляционных выражений.%
?

{%
Когда намерением является прохождение
несколько реляционных или сдвиговых выражений в качестве фактических аргументов;
и есть двусмысленность, исходный код можно легко отрегулировать
форма, которая является однозначной.
Кроме того, мы ожидаем, что будет более распространенным иметь
Общие вызовы функций как реальные аргументы
чем наличие реляционных или сменных выражений, которые совпадают
и имеют скобки в конце, так что возникает двусмысленность.
?

%
Оценка фактической аргументной части формы



идет следующим образом:

%
Аргументы типа  оцениваются
в том порядке, в котором они появляются в программе,
производящие виды .
Аргументы  оцениваются
в том порядке, в котором они появляются в программе,
Производить объекты .

{%
Проще говоря, аргументная часть, состоящая из  Типы аргументов,
 Позиционные аргументы и  Названы аргументы
оценивается слева направо.
Обратите внимание, что список аргументов типа опускается, когда 
()%
?


 Обязательный акт для формальностей


{%
В нижеследующем случае необычный случай охватывается неявно:
Когда число аргументов фактического типа равно нулю
Весь список аргументов типа \code{<\ldots{}>} опущен,
и аналогично для пустых списков параметров типа ().%
?

%
Подумайте о призыве  Функция  с
Фактическая аргументация часть формы
.

{%
Обратите внимание, что  обозначает функцию в семантическом смысле,
Вместо синтаксической конструкции.
Ссылка на этот раздел используется в других разделах.
когда указан статический анализ вызова,
и статический тип  Это было определено.
Сама функция может быть получена из декларации функции.
от экземпляра, связанного с {} и декларацией метода экземпляра,
или как объект функции, полученный путем оценки выражения.
Из-за этого мы не можем указать здесь, какая синтаксическая конструкция
Соответствует .
%
Ссылка на этот раздел также используется в других разделах.
когда фактические аргументы должны быть связаны с соответствующими формальными параметрами,
 Для определения динамической семантики будет использовано слово.%
?

{%
Мы не звоним  а «функциональный объект» здесь, поскольку мы не хотим подразумевать
Каждый вызов функции должен включать отдельную оценку.
выражение, которое дает объект функции,
Затем следует вызов этого функционального объекта.
Например, реализация должна быть разрешена для компиляции вызова.
Функция верхнего уровня как серия шагов, посредством которых кадр стека
Создан, за ним следует низкоуровневый скачок к генерируемому коду для организма.
В этом разделе,
слово "функция" является более низкоуровневым, чем "функциональный объект",
но «функция» по-прежнему означает семантический объект
который связан с декларацией функции,
даже при том, что не может быть соответствующей сущности в куче во время выполнения.%
?

%
% Мы не можем пройти один и тот же параметр дважды.
Ошибка времени компиляции  для любого .

%
Для данного типа  вводим понятие
 Ограниченный тип. T0 ограниченный.
 Сам по себе  ограниченный;
Если   ограниченный и
 является переменной типа со связанным 
  ограниченный;Наконец, если   ограниченный и
 является переменной типа
  ограничен.
В частности, а
\IndexCustom{\DYNAMIC{} Ограниченный тип {type!dynamic bounded}
является либо  само собой
или переменной типа, граница которой  ограниченный,
или пересечение, второй операнд которого ограничен .
Аналогично для a
\IndexCustom{\FUNCTION{} ограниченный тип}{type!function limited}.

%
А.
{Ограниченный тип функции}{Ограниченный тип функции!
Тип   где  является функциональным типом
().
Ограниченный тип функции  имеет

Уникальный тип функции    ограничен.

%
Если статический тип  является  ограниченный или  ограниченный,
Никаких дальнейших статических проверок при вызове 
{(кроме отдельных статических проверок на подтерминалы, как аргументы)},
Статический тип  является .
В противном случае это ошибка времени компиляции, если статический тип  не является
Функциональный тип ограничен.

%
Если ошибки не произошло и статический анализ  не является полной
Затем статический тип  является ограниченным по типу функции;
Пусть  будет ассоциированным типом функции .

%
 Тип возврата ,
давать 
быть формальными параметрами типа,
 количество требуемых параметров,
Пусть  будет позиционными параметрами,
и пусть  Необязательные параметры .
 Статический тип формальных параметров ,
и для каждого  пусть  тип параметра, названного ,
где каждый тип параметра получается путем замены  на ,
в аннотации типа данного параметра.
Наконец, позвольте  Статический тип .

{%
У нас есть фактический список аргументов, состоящий из  Типы аргументов,
 позиционные аргументы и  Названы аргументы.
У нас есть функция с  Типовые параметры,
 Необходимые параметры и дополнительные параметры .
Рисунок  Показывает, как это происходит.%
?

% Мнение по поводу декларации:
%
% | ** ... * | * * ... *
% <-s type par.s-> <-h required par.s--> <-k optional par.s->
% <--n pos, par.s, h=n--> <--k именуется par.s--> Нареченный
% <-------------------------------------------- n=h+k--------- POS
%
% Фактическая аргументация часть:
%
% | ** ... * | * * ... *
% <-r тип arg.s-> <---m позиционный arg.s--> <-l названный arg.s->
[h]
 1{\mbox{\commentary{{#1} Аргументы}}
 1{\mbox{\commentary{{#1} Параметры}}}
%
 Реальные аргументы:

%
 Декларация Термин с названными параметрами: 

%
 Декларация Термин с факультативными позиционными параметрами: 

%
 Возможные фактические части аргументов и формальные части параметров.



%
% Вывод типа считается полным, поэтому мы должны
% правильное количество аргументов.
Ошибка времени компиляции .
Ошибка времени компиляции  и для некоторых ,
.
Это ошибка времени компиляции, если только .
Если ,
Ошибка времени компиляции, если только  Назовите параметры и
.

{%То есть количество типов аргументов должно совпадать
количество типовых параметров,
И границы должны соблюдаться.
Мы должны получить как минимум необходимое количество позиционных аргументов.
и не более общего числа позиционных параметров.
Для каждого названного аргумента должен быть названный параметр с тем же именем.
?

%
Статический тип  является .

%
Ошибка времени компиляции  не может быть назначен
.
Ошибка времени компиляции  не может быть назначен
.

{%
Рассмотрим случай, когда вызов функции в фокусе здесь
Примерный метод вызова.
Для каждого конкретного аргумента,
Соответствующий параметр может быть ковариантным.
Однако вышеупомянутые требования к уступке применяются в равной степени.
как когда параметр ковариантен, так и когда он не ковариантен.
?

{%
Ковариантность параметров в вызове метода экземпляра может быть введена
подтип статически известного типа приемника,
Это означает, что любая попытка пометить данный фактический аргумент как опасный.
Благодаря динамической проверке типа, которой он будет подвергнут
будет неполным:
Некоторые фактические аргументы могут быть подвергнуты такой динамической проверке типа.
Хотя это не известно статически на сайте вызова.
Это неудивительно для такого механизма, как ковариация параметров.
предназначен для того, чтобы разработчики могли явно запрашивать
что этот конкретный вид безопасности компиляции нарушается.
Дело в том, что этот механизм откладывает исполнение
лежащая в основе инвариантная для бега время,
Взамен они позволяют создавать полезные программы.
которые в противном случае были бы отклонены во время компиляции.%
?

%
Для динамической семантики,
 Выберите функцию  Типовые параметры  требуемые параметры;
 позиционные параметры ;
 быть необязательными параметрами, объявленными .

%
Оцениваемая фактическая аргументация





производной от фактического аргумента части формы





Связано с формальными параметрами типа и формальными параметрами  следующим образом:

%
Прохождение неправильного количества аргументов фактического типа.
Если  и  затем
Если  Не имеет аргументов типа по умолчанию
()
Затем происходит динамическая ошибка.
В противном случае замените фактический список аргументов типа:
Пусть  будет  и пусть   быть результатом
инстанциация для связывания
()
по формальным параметрам типа ,
Замена действительных значений любых переменных свободного типа
().
В противном случае, если ,  бросается.

%
% Передача именованных аргументов функции с необязательными позиционными параметрами.
Если  и , выкидывают .
% Слишком мало или слишком много позиционных аргументов.
Если , или , бросается .
% Когда l>0, h=n и есть k названных параметров p_{h+1}..p_{h+k}.
Кроме того, каждый
,
должен иметь соответствующий названный параметр в наборе
,
<.
% Наконец-то!
Далее  привязанный к
<p000539>,
 Связано с .
Все остальные формальные параметры  Они привязаны к своим значениям по умолчанию.

{%
Все остальные параметры обязательно являются необязательными.
и, следовательно, значения по умолчанию.%
?

%
% Проверьте тип аргументов.
% Эта ошибка может произойти из-за ковариации и из-за динамического вызова.
Ошибка динамического типа, если  не является подтипом фактической границы
()
 аргумент типа , для фактических аргументов типа .% Проверьте типы позиционных аргументов.
% Эта ошибка может возникать из-за неявных слепков, ковариации и динамических вызовов.
Ошибка динамического типа, если  не является нулевым объектом ()
и фактический тип
()
 Не является супертипом динамического типа .
% Проверьте типы названных аргументов.
% Эта ошибка может возникнуть из-за неявных слепков, ковариации и динамических вызовов.
Ошибка динамического типа, если  это
не нулевой объект и фактический тип
()
 Не является супертипом динамического типа .


 Неквалифицированный призыв


%
Неквалифицированный вызов функции  имеет форму


,


где  - идентификатор.

{%
Обратите внимание, что список аргументов типа опускается, когда  ()%
?

%
Выполните лексический поиск 
()
Местонахождение: .

%
 Лексический поиск дает декларацию.
Пусть  будет декларацией, полученной в результате лексического поиска .

«P000843»

Когда  является типовой декларацией, то есть
объявление класса, смеси, перечня, псевдонима типа или параметра типа,
Применяется следующее:
Если  является декларацией класса 
Конструктор имеет имя 
Тогда значение  зависит от контекста:
Если  происходит в постоянном контексте
(),
Тогда  рассматривается как
()
;
Если  Не происходит в постоянном контексте.
 Обозначается как .
Если  не является классовой декларацией,
или класс, названный  не имеет конструктора с именем ,
Происходит ошибка времени компиляции.

В противном случае, если  является декларацией
локальная функция,
функции библиотеки, или
библиотека или статический геттер, или переменная,
 рассматривается как
()
Выражение функции Invocation
(P000937).

В противном случае, если  это
Статический метод или геттер
{(которые могут быть индуцированы статической переменной)}
в замкнутом классе или смеси ,
 рассматривается как
()

(P000939).

Если  является экземпляром расширения 
с параметрами типа  X}{1}{k}
 Обрабатывается как \code{$E$<\List{X}{1}{k}>(\THIS). .
Статический анализ и оценка
Происходит с преобразованным выражением,
Поэтому нет необходимости в дальнейшем уточнять лечение .

{%
Иными словами, внутри  Примеры членов  Тень будет
Примеры членов  Тип, то есть
Расширение имеет более высокий приоритет, чем интерфейс объекта.
Противоположное верно для призывов повсюду за пределами [P000590].

Нет необходимости рассматривать пример члена класса,
Потому что лексический поиск в классе никогда не даст декларации.
Является членом организации.%
?

{-2ex}


 Лексический поиск дает префикс импорта
При лексическом поиске  дает префикс импорта,
Происходит ошибка времени компиляции.


 Лексический поиск ничего не дает
При лексическом поиске  ничего не дает,
 рассматривается как
()
Обычный метод вызова

().

{%
Это происходит, когда лексический поиск определяет, что
 должен ссылаться на экземпляр члена класса или расширения,Местонахождение:  Вы можете получить доступ к 
Интерфейс ограждающего класса имеет член ,
или имеется соответствующее расширение с таким членом.
Статический анализ и оценка продолжаются
,
поэтому нет необходимости дополнительно указывать лечение .%
?


{%
Обратите внимание, что неквалифицированный вызов не определяет семантику оценки.
Это происходит потому, что каждый случай, который не является ошибкой, заканчивается выводом, что
Неквалифицированный призыв следует рассматривать как какую-либо другую конструкцию.
которые указаны в другом месте.%
?


 Функция экспрессии вызова.


%
Выражение функции  имеет форму





где  - выражение.

{%
Обратите внимание, что список аргументов типа опускается, когда 
()%
?

%
Рассмотрим ситуацию, когда  обозначает класс 
содержит декларацию конструктора по имени ,
или имеет форму  где
 обозначает класс  который содержит декларацию
Конструктор по имени .
Если  происходит в постоянном контексте
()
Затем  рассматривается как ,
Если  Не происходит в постоянном контексте.
 Обозначается как .

{%
Когда  рассматривается как другая конструкция ,
Статический анализ и динамическая семантика
указаны в разделе о 
() %
?

%
В противном случае это ошибка времени компиляции, если  Тип буквальный.

{%
Эта ошибка уже была указана в другом месте.
()
Для случая, когда  является идентификатором,
 могут иметь и другие формы, например, %
?

%
В противном случае, если  является идентификатором , затем  обязательно должен обозначать
локальная функция, библиотечная функция, библиотека или статический геттер,
или переменной, как описано выше,
 не рассматривались бы как вызов функции выражения.

%
Если  является выражением извлечения свойств
()
 рассматривается как обычный метод вызова
().

{%
 рассматривается как метод вызова метода
 на объекте ,
не как призыв к геттеру  на 
Затем следует вызов функции .
Если метод или приемник  Они существуют, они будут эквивалентны.
Однако, если  не определено в ,
В результате возникает вызов  будет отличаться.
 Перейти к  описать
Призыв к методу  с аргументом  В первом случае,
и вызов геттера  (без аргументов) в последнем.%
?

%
 Статический тип .
Если  - тип интерфейса, который имеет метод под названием ,
 рассматривается как
()
Обычный призыв


.

%
В противном случае, статический анализ  выполняется как указано
В разделе ,
Использование  как статический тип вызываемой функции,
Статический тип  как там указано.

%
Оценка вызова выражения функции





Прибыли на оценку , дающие объект .
 быть свежей переменной, связанной с .
Если  является объектом функции, то функция вызывает





Оценивается путем связывания фактического с формальным.Как указано в разделе ,
и выполнение корпуса  с этими привязками;
Возвращенный результат является результатом оценки .

%
В противном случае  не является функциональным объектом.
Если  имеет метод, названный 
Оценивается следующий обычный метод вызова,
и его результат является результатом оценки :


.

%
В противном случае  не имеет метода, названного .
Новый экземпляр  заданного класса  Он был создан,
такой, что:


 Показатель .
  Оценивает символ .
  Оценка в неизменяемый список
чей динамический тип ,
содержащие объекты, полученные в результате оценки
.
  Оценка на неизменяемую карту
чей динамический тип ,
с ключами и значениями, полученными в результате оценки

\code{<Symbol, Object>: : .
  Оценка в неизменяемый список
чей динамический тип ,
значения, полученные в результате оценки

.


%
Затем метод вызова  оценивается,
и его результат является результатом оценки .

{%
Ситуация, при которой  Призыв может только возникнуть
когда статический тип  является .
Семантика времени выполнения гарантирует, что
Вызов функции может быть приравнен к вызову
Метод экземпляра .
Тем не менее, тип интерфейса с методом под названием \CALL{}
не является подтипом какого-либо типа функции.
()%
?


<{Клозуризационный процесс)


%
 быть выражением, обозначающим
объявление функции верхнего уровня, локальной функции,
Статический метод класса, смеси или расширения
();
Или пусть  Быть функцией буквально
().
Оценка  Получает объект функции
Что является результатом 
применяется к декларации, обозначенной 
соответственно функции буквальной  рассматривается как функциональная декларация.
% Обратите внимание, что мы не обещаем, что это будет новый функциональный объект.
%, так что постоянные выражения также охватываются, и другие могут быть
% идентичны, даже если не требуются, например, локальные функции без свободных переменных.
\IndexCustom{Клосуризация}{Клосуризация}
обозначает примерный метод клозулизации
()
а также функции клозура,
Он также используется в качестве краткого описания для любого из них.
Когда нет двусмысленности.

%
В то же время, по сравнению с предыдущими исследованиями >>>>>
сводится к созданию функционального объекта 
что является примером класса  чей интерфейс является
подтип фактического типа 
()
соответствующий подписи в функциональной декларации ,
используя текущие связывания переменных типа, если таковые имеются.
Не существует функционального типа  который является подтипом 
 Это подтип .

{%
Если  Статический метод или функция верхнего уровня
 примитивное равенство
() %
?

{%
Другими словами,  имеет свободу быть подобающим подтипом
тип функции, который мы можем прочитать из декларации 
Поскольку это может быть определенная внутренняя платформа определенного класса,
 не имеет свободы быть подтипом
Другой и более особый тип функции, и он не может быть .%
?

%Призыв  с определенным списком аргументов свяжет фактические данные с формальными
Таким же образом, как и призыв к 
(),
и затем выполнить тело 
в захваченной области применения, измененной сферой применения связанного параметра,
Получить такое же завершение
()
Ссылка на «P000714» Он бы сдался.

%
 и  Это два постоянных выражения, которые
оценивает объект функции, полученный путем клозуризации функции
той же функциональной декларации.
В данном случае  будет оцениваться как истинное.

{%
То есть, постоянные выражения, оценка которых является функцией клосуризации.
Канонические.%
?


 Обоснование общей функции


% % TODO(Eernst): Спецификация дженерикальной инстанциации функции полагается
% на вывод типа иначе, чем другие механизмы, которые мы имеем
%, поскольку невозможно утверждать, что вывод типа
%, которые, как предполагается, уже имели место, а затем мы просто рассматриваем
% - программа, в которой отсутствуют аргументы типа или аннотации типа.
% добавляются синтаксически. Следовательно, этот раздел, вероятно, должен быть
% % переписываются довольно широко, когда мы добавляем спецификацию вывода.

%
Инстанциация общей функции - это механизм, который дает
объект негенерической функции, основанный на данной общей функции.

{%
Сущность универсальной инстанциации функции
Для того, чтобы получить «проклятые» призывы,
В том смысле, что общая функция может получить свою фактическую
Введите аргументы отдельно
(Тогда он должен получить аргументы типа , а не только некоторые из них),
Это приводит к появлению объекта негенерической функции.
Аргументы типа передаются неявно, на основе вывода типа;
% % TODO(Eernst): Приходите, конструкторы, пересмотрите это.
Будущая версия Dart может позволить передавать их явно.
?
 Вот пример:




{%
Каждый объект функции хранится в 
имеет динамический тип ,
и получено имплицитно
Прохождение фактического типа аргумента  ''
для соответствующей общей функции.%
?

%

 выражение, статический тип которого  это
%
Тип общей функции формы
 T_0}{X}{B}{}{{p}}
где  Получается из {formalParameterList}.
\commentary{Примечание:   является родовым.
Предположим, что тип контекста является негенерическим типом функции. .
В этой ситуации происходит ошибка компиляции времени.
(,
,
,
,
,
,
<<<<<<<<<<<<<<<<<<<<<<<<<<<<<<<<<<<<<<<<<<<<<<<<<<<<<<<<<<<<<<<<<<<<<<<<<<<<<<<<<<<<<<<<<<<<<<<<<<<<<<<<<<<<<<<<<<<<<<<<<<<<<<<<<<<<<<<<<<<<<<<<<<<<<<<<<<<<<<<<<<<<<<<<<<<<<<<<<<<<<<<<<<<<<<<<<<<<<<<<<<<<<<<<<<<<<<<<<<<<<<<<<<<<<<<<<<<<<<<<<<<<<<<<<<<<<<<<
,
,
,
,
,
,
Если только следующий шаг не увенчается успехом:

%
 Инстанциация типа общей функции {%}
Общий тип функции:
Вывод типа применяется к  Тип контекста ,
и приводит фактический список аргументов типа
%
 T}{1}{s}.

{%
Генерическая инстанциация типа функции не работает
в случае, когда вывод типа не удается,
В этом случае возникает вышеупомянутая ошибка времени компиляции.
Он будет указан в будущей версии этого документа.
как вычисляется тип вывода  T}{1}{s}
()%
?

%
Допустим, что генерический тип функции удался.
 F'} обозначает тип
p.
{%
Обратите внимание, что гарантируется, что  присваивается ,
или вывод был бы неудачным.%
?
Отныне в статическом анализе
Это явление  Считается, что имеет статический тип. .

%Выполнение  действует следующим образом.
Оценить  на объект .
% % TODO(Eernst): Приезжайте, снимите следующую строку.
Динамическая ошибка возникает, если  Это нулевой объект.
 S_0}{Y}{B'}{} быть динамическим типом 
{(по звучанию этот тип является подтипом )}.
 Затем оценивает объект функции  Динамический тип
q,
где  Это действительное значение , для .
Призыв  с реальными аргументами {args} имеет
Тот же эффект, что и вызов 
с аргументами реального типа {t}{1}{} и фактические аргументы {args}.

%
 и  быть двумя постоянными выражениями, подлежащими
Генетическая инстанциация функций.
Предположим, что  и  без оценки типа контекста 
 Таким образом, P001106 является истинным.
Предположим, что данные типы контекста приводят к успешному
Генетическая функция типа инстанциации с
те же самые аргументы фактического типа для  и ,
Получение объектов негенерической функции  соответственно .
В данном случае  будет оцениваться как истинное.

{%
Постоянные выражения, оценка которых
Генетическая инстанциация функций канонизирована.
на основе основной функции и фактических аргументов типа.
Как следствие, они также равны по оператору. \lit{===.%
?

%
 и  быть двумя выражениями
{(которые могут быть или не быть постоянными)}
которые подлежат обобщению функции.
Предположим, что  и  без оценки типа контекста 
соответственно  так, что  верно.
Предположим, что данные типы контекста приводят к успешному
Генетическая функция типа инстанциации с
те же аргументы фактического типа для  и ,
Получение объектов негенерической функции  соответственно .
В данном случае  будет оцениваться как истинное.

{%
Если одно или оба выражения не являются постоянными,
не уточняется, является ли
 {} или ,
но оператор \lit{========================================================================================================== истинно для объектов равной функции
с теми же аргументами реального типа.%
?

{%
Никакое понятие равенства не подходит, когда аргументы типа отличаются.
Даже если результирующая функция объекта
Получается, что у них точно такой же тип во время бега,
потому что выполнение двух функций объектов, которые отличаются в этих способах
могут иметь различные побочные эффекты и возвращать разные результаты.
при исполнении, начиная с точно такого же состояния.%
?


<{Смотреть}


%
 Это процедура, которая выбирает
a конкретного заявления члена, основанного на прохождении
последовательность классов, начиная с данного класса 
Продолжение суперкласса текущего класса на каждом шагу.
Поиск может быть частью статического анализа и может быть выполнен.
Во время бега. Он может преуспеть или потерпеть неудачу.

{%
Мы определяем несколько видов поиска с очень похожей структурой.
Мы прописываем каждую из них, несмотря на избыточность.
Во избежание внедрения механизмов абстракции на метауровне
только для этой цели.
Дело в том, что мы должны указать для каждого поиска
Какого человека он ищет,
потому что, например, используются «методический поиск» и «высший поиск».
в различных ситуациях.%
?

{% Сфера применения определения "lookup".

#1
%
Результатом a
  в  В отношении  в классе является результатом поиска  в  в отношении .
Результат {  в  в отношении  Это:
Если  объявляет конкретный экземпляр {#1} с именем 
Это доступно для ,
Затем, что {#1} заявление является результатом поиска {#1},
и мы говорим, что {#1} было { Посмотреть в 
В противном случае, если  имеет суперкласс ,
Результатом поиска {#1} является
Результат поиска по запросу  в  в отношении .
В противном случае, мы говорим, что поиск не удался.

?

%
 быть идентификатором  Объект,  библиотека,
 Класс, который является классом  или их суперкласса.

\LookupDefinitionWithStart{метод}
\LookupDefinitionWithStart{getter}}
\LookupDefinitionWithStart{setter}

#1
Результатом a
  в  в отношении 
является результатом a
 в $o$ В отношении 
Начиная с класса .
?

%
 Идентификатор:  Объект и  Библиотека.
{метод}
\LookupDefinition{getter}}
{setter}

} % Конец области применения определений поиска.

{%
Обратите внимание, что для поиска Геттера (Сеттера) результат может быть
Геттер (настройщик), который был вызван переменной экземпляра
Декларация.%
?

{%
Обратите внимание, что мы иногда используем такие фразы, как метод поиска  "
Чтобы указать, что метод поиска выполняется,
и аналогично для сеттеров и геттеров.%
?

{%
Мотивация игнорировать абстрактных членов во время поиска
В основном, чтобы обеспечить более гладкую композицию смеси.%
?


 Вызов Геттера высшего уровня.


%
Оценка вызова геттера верхнего уровня  формы ,
где  является идентификатором,
идет следующим образом:

%
Применяется функция геттера .
Значение  Это результат, возвращаемый вызовом к функции геттера.
{%
Обратите внимание, что призыв всегда определен.
Согласно правилам идентификации ссылок,
Идентификатор не будет рассматриваться как вызов геттера высшего уровня.
Если только геттер не  определяется.%
?

<<<<<<<<<<<<<<<<<<<<<<<<<<<<<<<<<<<<<<<<<<<<<<<<<<<<<<<<<<<<<<<<<<<<<<<<<<<<<<<<<<<<<<<<<<<<<<<<<<<<<<<<<<<<<<<<<<<<<<<<<<<<<<<<<<<<<<<<<<<<<<<<<<<<<<<<<<<<<<<<<<<<<<<<<<<<<<<<<<<<<<<<<<<<<<<<<<<<<<<<<<<<<<<<<<<<<<<<<<<<<<<<<<<<<<<<<<<<<<<<<<<<<<<<<<<<<<<<
Статический тип  Объявленный тип возврата .


{Вызовы участников}


[h]
\begin{minipage}[t]{]
<<<p000010>


\vspace{4 мм)
\begin{minipage}[t]{]

В таблицах выше,
, , ,  являются выражениями;
\metavar{args} происходит от \synt{arguments};
\metavar{types} происходит от \synt{typeArguments};
 является \synt{оператор}, который не является <{==}
% % TODO(Eernst): Добавить исключение для &&=, | |= при добавлении поддержки.
и '\code{=}' является \synt{compoundAssignmentOperator}.

%
 Призывы членов с синтаксическим приемником .



%
 Это выражение с определенной синтаксической формой.
чья динамическая семантика включает призыв
одного или двух членов данного получателя,
или призыва членов продления.
В этом разделе указывается, какие синтаксические формы являются призывами членов,
Определение некоторой терминологии
который необходим для обозначения
Конкретные части нескольких синтаксических форм вместе.

%
Статический анализ и динамическая семантика
каждой из синтаксических форм, являющихся призывами членов
указывается отдельно,Этот раздел касается только
Синтаксическая классификация и терминология.

{%
Например, один вид вызова члена - это обычное вызов метода.
()%
?

%
А.
<{простой призыв члена}{призыв члена} простой
соответственно
<{составной член призыва}{член призыва! составной
на
{синтаксический приемник}{призыв к члену! синтаксический приемник
Выражение  является выражением
Одна из форм, показанных на рис.~.
Каждое обращение имеет
{соответствующее название члена}{%
Призыв членов! имя соответствующего члена
Как показано на рисунке.

%
Каждый член вызова на рис.~ который содержит
<? является
{условное членское призвание}{членское призвание!условное}.
Ан
{безусловный вызов члена}{призыв члена!безусловный}
Это призыв к члену, который не является условным.

{%
Для простого вызова члена имя соответствующего члена является
Имя члена, на которое ссылается
в случае, когда член вызова вызывает члена экземпляра.
Для композитного члена вызов это имя геттера
и имя основания как геттера, так и сеттера.%
?

{%
Обратите внимание, что  Не может быть 
Несмотря на то, что  Используйте метод экземпляра.
Это связано с тем, что семантика суперинвокации отличается от
других призывов.
Среди бинарных операторов {==} не входит.
Это связано с оценкой 
Он включает в себя больше шагов, чем призыв члена инстанции.
Точно так же \lit{\&\&} и \lit{ | |} не включены
потому что их оценка не включает вызов метода.%
?

[t]
[h] []


%
 Отказ от призывов членов.
 является {compoundAssignmentOperator}.
Используется первое применимое правило.
В частности, "\code{=}" Не может быть  "
когда более раннее правило с "" совпадает.
 называется {заменитель приемника}
Композитный член вызывает,
и указывается в основном тексте
().



%
Составной член вызов является сокращенной формой
значение которого сводится к простым призывам членов
Как показано на рис.~.
Этот шаг называется  трансформация,
И мы говорим, что полученное выражение было .
Фиг.~ содержит
Несколько случаев 
который является
{заменитель приемника}{метод вызова! замещающий приемник
Композитный метод вызова.
Смысл каждого появления  определяется следующим образом:

%
Когда получатель  является приложением расширения
()
в форме < T}{1}{k}>()
\commentary{(где ) означает, что список аргументов типа отсутствует)}:
 быть свежей переменной, связанной со значением 
и с тем же статичным типом, что и ,
затем  \code{\List{T}{1}{k}>() Когда это происходит
в качестве получателя приглашения,
иным образом  .

%
Когда  не является приложением расширения,
 является свежей переменной, связанной со значением ,
с тем же статичным типом, что и .

{%
Это соответствует внешнему виду \LET{} В каждом правиле
На рисунке.
где  происходит,
и четкое различие между двумя формами ,
Но цифра была бы значительно более многословной, если бы она была
указанным образом.%
?


 Метод вызова.


%
Метод вызова может принимать несколько форм, как указано ниже.


 Обычный вызов


%
Обычный метод вызова может быть условным или безусловным.

%

Подумайте
{условный вызов обычного метода}{%
Метод призыва! условный обыкновенный
 в форме
.

{%
Обратите внимание, что необычные призывы возникают как особый случай, когда
Число типов аргументов равно нулю,
В этом случае список аргументов типа опущен,и аналогично для формальных списков параметров типа ().%
?

%
Статический тип  то же, что и статический тип


.


Ошибки времени компиляции, которые могут быть вызваны





Они также генерируются в случае .

%
Оценка  идет следующим образом:

%
Если  является буквальным типом или обозначает расширение,  рассматривается как


.

%
В противном случае оценивать  Объект .
Если  является нулевым объектом,  оценивает нулевой объект ().
В противном случае пусть  быть свежей переменной, связанной с  оценивать

к объекту .
 Показатель .


%

A {статический вызов члена}{призыв члена!статический}
 является обращением к форме
,
где  Буквальный тип, или  обозначает расширение.

{%
Необычные призывы возникают как особый случай
где число аргументов типа равно нулю ().%
?

%
Происходит ошибка времени компиляции
Если только  не обозначает класс, смесь или расширение, которое объявляет
статический член, названный ,
которые мы будем называть
{обозначенный член}{статический член призыв! обозначенный член
.
Когда обозначенный элемент является статичным методом, пусть  быть его функциональным типом;
когда обозначенный элемент является статическим геттером, пусть  быть его возвратным типом;
Когда обозначенный член не является ни тем, ни другим, возникает ошибка времени компиляции.

%
Статический анализ  Затем выполняется
Как указано в разделе ,
считая  статическим типом функции вызова,
Статический тип  как там указано.

%
Оценка вызова статического метода  в форме





идет следующим образом:

%
Если обозначенный член  Это статический метод,
Пусть  будет функцией, объявленной этим членом.
Осуществляется привязка фактических аргументов к формальным параметрам
Как указано в разделе .
Тело  Затем выполняется в отношении привязки
Это вытекает из оценки аргументной части.
Значение  Это объект, возвращаемый в результате выполнения тела .

%
Если обозначенный член  Статический геттер,
Вызовите сказанное Getter и дайте  быть
Новая переменная, привязанная к возвращенному объекту.
Тогда значение  является ценность
.


%

 Безусловный обычный метод вызова {%}
Метод призыва! Безусловный обыкновенный
 имеет форму

где  Это выражение, которое не является буквальным,
и не обозначает расширение.

{%
Необычные призывы возникают как особый случай
где число аргументов типа равно нулю ().%
?

%
 Статический тип .

%
Если  \DYNAMIC ограниченный
()
 является одним из
, ,  или 
Мы говорим, что  является
 членский призыв {
динамический Воззвание члена-объекта.
Статический анализ которого приводится отдельно ниже.

%
В противном случае это ошибка времени компиляции, если  не имеет доступной
()
член экземпляра по имени , если только:



 \DYNAMIC{} ограниченный;В этом случае дальнейшие статические проверки не выполняются на 
{(кроме отдельных статических проверок терминов, таких как аргументы)}
Статический тип  является .
Или

 \FUNCTION{} ограниченный
()
 ;
В этом случае дальнейшие статические проверки не выполняются на 
{(кроме отдельных статических проверок терминов, таких как аргументы)}
Статический тип  является .
{%
Это означает, что для вызовов метода экземпляра, названного ,
Получатель типа {} трактуется как приемник типа .
Ожидается, что любой конкретный подкласс 
будет осуществлять ,
Но нет подписи метода
который может быть принят для \CALL{} в \FUNCTION{}
Каждая подпись будет противоречить
некоторые потенциальные приоритетные декларации.%
?


%
Если  Не имел доступного участника по имени 
Статический тип  является ,
и никаких дальнейших статических проверок на 
{%
(за исключением подвыражений ) подвергаются собственному статическому анализу.
}

%
Если  является динамическим  Призыв члена
()
Затем статический тип элемента указывается в
Раздел .
В этом случае, если   или 
Тогда пусть  быть возвратным типом указанного геттера;
если   или 
 Будьте типом указанного метода.

{%
Обратите внимание, что это всегда ошибка времени компиляции, если 
 или .
?

%
В противном случае  обозначает экземпляр члена.
Пусть  Библиотека, содержащая .
 быть результатом поиска по методу  в  в отношении ,
и если метод поиска удался, то пусть  Статический тип .

%
В противном случае пусть  быть результатом Getter Lookup
  в отношении ,
 Тип возврата .
{%
С тех пор  Мы не можем иметь сбой в обоих поисках.%
?
Если геттер возвращает тип  является типом интерфейса
метод, который называется ,
 рассматривается как
()
Обычный призыв


,


который определяет дальнейший статический анализ.

%
В противном случае статический анализ  выполняется
Как указано в разделе ,
считая  статическим типом функции вызова,
Статический тип  Как там указано,
Кроме того, что вызовы методов, названных 
 на подтипах  не являются подтипами 
имеют специальные правила, аналогичные правилам для добавок ()
мультиплиативные (P000967) операторы.

%
 быть обращением к форме 
 Контекстный тип .
Тип контекста  Затем определяется следующим образом:
Если \SubtypeNE{T}{ и не \SubtypeNE{T}{\code{Never}}, то:


\item{} Если \SubtypeNE{}{C}, а не {}{C},
и \SubtypeNE{T}{
Контекстный тип  является .
\item{} Если \SubtypeNE{}{C}, а не {}{C},
и не \SubtypeNE{T}{
Контекстный тип  является .< В противном случае тип контекста  является .


%
Далее  Статический тип .
Если {T}{ и не \SubtypeNE{T}{
 присваивается ,
Статический тип  определяется следующим образом:


< Если {T}{
Статический тип  является .
< В противном случае, если \SubtypeNE{S}{\code{double}}
и не \SubtypeNE{S}{}
Статический тип  является .
\item{} В противном случае, если \SubtypeNE{T}{\code{int}},
\SubtypeNE{S}{}} и не \SubtypeNE{S}{}
Статический тип  является .
\item{} В противном случае статический тип  является .


%
Пусть  быть обращением к форме ,
где \SubtypeNE{T}{ и не \SubtypeNE{T}{}
 Контекстный тип .
Контекстный тип  и  Затем определяется следующим образом:


< Если \SubtypeNE{T}{}
\SubtypeNE{}{S} и не {}{S}
Таким образом, контекстный тип  и  является .
\item{} Если \SubtypeNE{T}{}
\SubtypeNE{}{S} и не {}{S}
Таким образом, контекстный тип  и  является .
< В противном случае тип контекста $e_2$ и $e_3$ является \code{num}.


%
Далее  Статический тип  и $T_3$ быть статичным типом
.


< Если все ,  и  являются подтипами , но
Не подтипы , а статический тип  .
\item{} Если все ,  и  являются подтипами , но
Не подтипы , а статический тип  .
\item{} В противном случае статический тип  является .


%
Это ошибка времени компиляции, чтобы вызвать метод экземпляра на типе буквальный
Затем следует токен «.» (период).
 Например,  Это ошибка.

{%
Причина этого правила заключается в том, что доступ участника к типу буквальный.
Зарезервировано для вызова статических членов.
Призыв статического члена класса, смеси, перечня или расширения
использует указанную сущность в качестве {namespace},
не в качестве фактического класса, смеси или расширения.
В частности, синтаксический приемник не оценивается.
к объекту, который даже не был бы возможен для расширения.%
?

{%
Доступ пользователя к типу буквальный
(например, ,  или ,
всегда трактует декларацию, обозначаемую буквальным как
Пространство имен для доступа к статическим элементам или конструкторам.
Например,  Это ошибка
 не декларирует статический элемент, названный .
Он не будет оценивать  до  объект
Затем назовите его  Примерный метод.
Для этого можно использовать .
Обратите внимание, что каскады бывают разными:
Они всегда сначала оценивают свой приемник на объект.

Как естественное следствие,
Буквальный тип не может быть приемником в
неявное обращение к методу расширения
()%
?

%
Оценка вызова безусловного обычного метода  в форме

идет следующим образом:

%Во-первых, выражение  оценивается на предмет .
Пусть  быть результатом поиска
()
Способ   в отношении текущей библиотеки .

%
Если метод поиска удался,
Осуществляется привязка фактических аргументов к формальным параметрам
Как указано в разделе .
Тело  Затем выполняется в отношении привязки
В результате оценки списка аргументов,
и с < Ссылка: .
Значение  Это объект, возвращаемый в результате выполнения тела .

%
Если метод поиска не сработал,
Тогда пусть  быть результатом поиска геттера
()
  в отношении .

%
Если поиск геттера удался, нажмите на геттер 
 быть возвращенным объектом.
Тогда значение  является ценность


.

%
Если поиск Геттера также не удался,
затем новый экземпляр  заданного класса  Он был создан,
такой, что:


<<<<<<<<<<<<<<<<<<<<<<<<<<<<<<<<<<<<<<<<<<<<<<<<<<<<<<<<<<<<<<<<<<<<<<<<<<<<<<<<<<<<<<<<<<<<<<<<<<<<<<<<<<<<<<<<<<<<<<<<<<<<<<<<<<<<<<<<<<<<<<<<<<<<<<<<<<<<<<<<<<<<<<<<<<<<<<<<<<<<<<<<<<<<<<<<<<<<<<<<<<<<<<<<<<<<<<<<<<<<<<<<<<<<<<<<<<<<<<<<<<<<<<<<<<<<<<<< <<<<<<<<<<<<<<<<<<<<<<<<<<<<<<<<<<<<<<<<<<<<<<<<<<<<<<<<<<<<<<<<<<<<<<<<<<<<<<<<<<<<<<<<<<<<<<<<<<<<<<<<<<<<<<<<<<<<<<<<<<<<<<<<<<<<<<<<<<<<<<<<<<<<<<<<<<<<<<<<<<<<<<<<<<<<<<<<<<<<<<<<<<<<<<<<<<<<<<<<<<<<<<<<<<<<<<<<<<<<<<<<<<<<<<<<<<<<<<<<<<<<<<<<<<<<<<<< Оценивает до <<<<<<<<<<<<<<<<<<<<<<<<<<<<<<<<<<<<<<<<<<<<<<<<<<<<<<<<<<<<<<<<<<<<<<<<<<<<<<<<<<<<<<<<<<<<<<<<<<<<<<<<<<<<<<<<<<<<<<<<<<<<<<<<<<<<<<<<<<<<<<<<<<<<<<<<<<<<<<<<<<<<<<<<<<<<<<<<<<<<<<<<<<<<<<<<<<<<<<<<<<<<<<<<<<<<<<<<<<<<<<<<<<<<<<<<<<<<<<<<<<<<<<<<<<<
 Оценивает символ .
  Оценка в неизменяемый список
чей динамический тип ,
содержащие объекты, полученные в результате оценки
.
  Оценка на неизменяемую карту
чей динамический тип ,
с ключами и значениями, полученными в результате оценки

\code{<Symbol, Object>: : .
 Оценка в неизменяемый список
, динамический тип которого реализуется ,
содержащие объекты, полученные в результате оценки
.


%
Тогда метод  Вы можете посмотреть в  и
Ссылаясь на аргументы ,
и результатом этого вызова является результат оценки .

{%
Ситуация, при которой  Ссылка может возникнуть только
когда статический тип  является .
Обратите внимание, что формулировка не позволяет повторно оценить приемник  и
аргументы .%
?



 Каскады


%
 Это своего рода выражение, которое позволяет выполнять
несколько операций на данном объекте
без сохранения в переменной и доступа к ней через имя.
В общем, каскад можно распознать по использованию ровно двух периодов,
{..}


<cascade> ::= <cascade> '..' <cascadeSection>
\alt <conditionalExpression> (?..' | '..') <cascadeSection>

<cascadeSection> ::= <cascadeSelector> <cascadeSectionTail>

<cascadeSelector> ::=[<выражение>] "
 <идентификатор>

<cascadeSection Хвост ::= <cascadeAssignment>
 <selector>* (<assignableSelector> <cascadeAssignment>)?

<cascadeAssignment> ::= <assignmentOperator> <expressionWithoutCascade>


%
A {каскадный доступ к члену}{каскадный доступ к члену}
Это выражение происходит от {каскад}.

{%
A {cascadeSection} позволяет получить доступ к участникам, включая установщики.
Мотивация для каскадного доступа к члену заключается в том, что он позволяет
выполнение цепочки операций на основе объекта
при сохранении ссылки на этот объект для дальнейшей обработки.%
?

{%
 быть приложением расширения
().
Обратите внимание, что это ошибка времени компиляции
иметь {каскад} формы
 или ,
где  {каскадная секция}.%
?

{%Например,  позволяет получить ссылку на
Объект  который является результатом оценки ,
и в то же время использовать  Для вызова геттера  и
Сеттер  на величину, возвращаемую этим Getter.%
?

%
Выражение формы  где
 - \synt{условный} Выражение
 является {каскадная секция}
\IndexCustom{первоначально условный каскадный доступ}{%
Каскадный доступ! Первоначально условный.
Более того, если  это
первоначально условный каскадный доступ и
 Происходит от {каскадная секция} затем
 Кроме того, он изначально является условным.

{%
Короче говоря, каскад изначально условен, если
"Внутренние точки" - это {?..} Вместо {..}
Обратите внимание, что только самые внутренние точки могут иметь значение {?}.
Все невнутренние неявно пропускаются, если приемник недействителен,
Так что любой \lit{?} на невнутреннем \lit{..} будет бесполезным.%
?

%
Статический анализ и динамическая семантика каскадного доступа
определяется в терминах следующего шага разбавления.

%
 быть каскадным доступом, который изначально не является условным.
Это означает, что существуют термины \List{c}{1}{k)
полученный из <\synt{каскадная секция}
и a {условный} Выражение 
которая не является расширением,
 .
В данном случае  Обезжиренный для


<<<< P000305>>>}{$v_1$}{$v$. $c_1$\ $\cdots$}{$v_k$}} P000311. $c_k$}{}.

%
 Это каскадный доступ, который изначально является условным.
Это означает, что существуют термины \List{c}{1}{k)
полученный из <\synt{каскадная секция}
и {условный} Выражение 
которая не является расширением,
 .
В данном случае  Обезжиренный для


\Let{}{}}
$v$\,==\,\NULL? \ \NULL\ :\ \LetMany{$v_1$}{$v$. }{}{. .

{%
Обратите внимание, что грамматика такова, что  это
Синтаксически правильное выражение для всех .%
?


 Сверхвызовы


% Условные суперинвокации бессмысленны: {} не является нулевым.

%
  имеет форму
%


.

{%
Обратите внимание, что необычные призывы возникают как
случай, когда число аргументов типа равно нулю,
В этом случае список аргументов типа опущен,
и аналогично для формальных списков параметров типа ().%
?

%
Это ошибка времени компиляции, если происходит суперинвокация метода.
функция верхнего уровня или переменный инициализатор,
в списке инициализаторов переменных или инициализаторов,
в классе ,
на заводе-конструкторе,
или в статическом способе или переменном инициализаторе.

{% Сфера применения суперкласса. }

\def\SuperClass{{S_{{{super}}}}

%
%
Пусть P001632 будет суперклассом (P000974).
непосредственного класса для ,
%
Пусть  Библиотека, содержащая .
Пусть декларация  быть
Результат поиска метода  в \SuperClass{}
В отношении  (),
 Статический тип .В противном случае, если метод поиска не сработал,
Пусть заявление  быть результатом поиска
Геттер  в отношении  в \SuperClass
(),
 Тип возврата .
Если оба поиска не удались, возникает ошибка времени компиляции.

%
Иначе
{(когда один из поисков удался)},
Статический анализ  выполняется
как указано в разделе ~,
учитывая, что функция имеет статический тип ,
Статический тип  как там указано.

{%
Обратите внимание, что поисковые системы игнорируют абстрактные декларации.
Это означает, что будет ошибка времени компиляции.
если целевой член  является абстрактным,
когда его вообще не существует.%
?

%
Ан
\IndexCustom{неявно {} супервызов {%}
Метод суперинвокации! подразумеваемый 
имеет форму


,


и рассматривается как


.

 Список аргументов типа снова опускается, когда .

%
%
Оценка  идет следующим образом:
 быть текущим связыванием ,
 быть классом для ,
И пусть {} быть суперклассом () .
 Библиотека, содержащая .
%
Пусть декларация  быть результатом поиска
Способ   в  Начнем с 
().
Если поиск удался,
 обозначает функцию, связанную с .
%
Иначе
{(когда метод поиска не удался)},
Пусть декларация  быть результатом поиска
Геттер   в  Начнем с 
().
Если поиск геттера удался,
Ссылка на сказанное Геттер с \THIS{} Ссылаясь на ,
 Обозначает возвращенный объект.

{%
Не может случиться так, что оба дисплея провалятся.
потому что соответствующие поиски тогда не удались бы во время компиляции,
В этом случае программа имеет ошибку времени компиляции.%
?

%
В противном случае выполняется привязка фактических аргументов к формальным параметрам для

Как указано в разделе ,
и выполнить корпус  с указанными привязками
плюс связывание \THIS{} с .
Результат вернулся по  Это результат оценки .

} % Конец диапазона для названия суперкласса.


 Отправка сообщений.


%
Сообщения являются единственным средством общения между изолятами.
Сообщения отправляются с использованием определенных методов в библиотеках Дарта; для отправки сообщения не существует определенного синтаксиса.

{%
Другими словами, методы, поддерживающие отправку сообщений
воплощают примитивы Дарта, недоступные обычному коду;
Очень похоже на методы, которые изолируют нерест.%
?


 Добыча имущества


%
 Извлечение собственности.
Это позволяет получить доступ к члену в качестве свойства, а не функции.
Добыча имущества может быть любой:


 Пример метода клозулизации,
Преобразует метод в объект функции
(P000982).
Или
 Призыв Геттера, который возвращается
результат вызова метода Геттера
(P000983).


{%
Функциональные объекты, полученные от членов посредством клозуризации
Они известны как слезы.%
?

%
Добыча имущества осуществляется в нескольких формах, как описано ниже.

%
<{условный}
Рассмотрим {выражение условного извлечения свойств}{%
Добыча имущества! условный формы .

%
Если   трактуется как .

%
В противном случае, статический тип  Это то же самое, что
Статический тип .
 Статический тип ,
Пусть  Будьте свежей переменной типа .
За исключением ошибок внутри  и ссылки на имя ,
точно такие же ошибки времени компиляции, которые были бы вызваны 
Они также генерируются в случае .

%
Оценка условного выражения извлечения имущества 
Форма  идет следующим образом:

%
Если  Буквальный или  обозначает расширение,
Оценка  Оценка составляет .

%
В противном случае оценка  к объекту .
Если  является нулевым объектом,  оценивает нулевой объект ().
В противном случае пусть  быть свежей переменной, связанной с 
и оценить  Объект .
Далее  Показатель .


%
<
 быть идентификатором;
а \IndexCustom{выделение статического свойства}{%}
Добыча имущества! статический
 является выражением формы ,
где  Тип буквальный или  обозначает расширение.

%
Происходит ошибка времени компиляции
Если только  обозначает класс, смесь или расширение, которое объявляет
статический элемент, названный ,
которые мы будем называть
{обозначенный член}{статическое извлечение свойств} обозначенный член
.
Если обозначенный член является статическим геттером,
Статический тип  является возвратным типом указанного геттера;
если обозначенный элемент является статичным методом,
Статический тип  является функциональным типом указанного метода;
Если обозначенный член не является ни тем, ни другим, возникает ошибка времени компиляции.

%
Оценка извлечения статического свойства  формы 
идет следующим образом:
Если обозначенный член  Статический геттер,
Сказано, что Геттер вызывается, уступая объект ,
 В результате получается .
Если обозначенный член  Это статический метод,
 оценивает объект функции по функции клозуризации
()
Применяется к указанному статическому методу.


%
 Безусловный
 быть идентификатором;
a {безусловная добыча имущества}{%
Недвижимость!Безусловно
является выражением формы 
где  Это выражение, которое не является буквальным,
не обозначает расширение
();
или является выражением формы 
(P000987).


%
<<<p001670>> Неявный
 Выражение, статический тип которого
Тип интерфейса, который имеет метод под названием .
В случае, когда тип контекста для  является функциональным типом
или типа ,
 Обозначается как .

{%
Это означает, что «призывной объект» может рассматриваться как функция.
Поддерживает механизм, похожий на функцию клозуризации.
()
путем десугарирования его методом клозулизации на .
Это происходит только тогда, когда статически известно, что это вызывающий объект.
когда тип контекста требует функции.
?



 Получение доступа и извлечение метода


%
Рассмотрим безусловную добычу имущества 
()
Форма .
Ошибка времени компиляции  Это имя
экземпляр члена встроенного класса 
 - буквальный тип.

{%Это означает, что мы не можем использовать 
Получить объект функции для  Методика
 Об этом сообщает «P001191».
Но мы можем использовать :
 Тогда это не буквальный тип, а выражение в скобках.%
?

{%
Это прагматичный компромисс.
Способность отрывать методы экземпляров на экземплярах 
считался менее полезным,
Было сочтено более полезным настаивать на простом правиле, которое
метод отрыва по типу буквальный является {всегда} отрыв
Статический метод по обозначенному классу.%
?

%
 Статический тип .

%
Если   ограниченный
()
 является одним из
, ,  или 
Мы говорим:  является
 извлечение имущества {
динамический Извлечение имущества объекта.
Статический анализ которого приводится отдельно ниже.

%
В противном случае это ошибка времени компиляции, если  не иметь
Способ или приемник, названный 
Если  не является  ограниченный,
или  \FUNCTION ограниченный
()
 .
Статический тип  Это:


 Тип возврата, указанный в разделе 
Если  является динамическим  вывоз имущества
  или .
 Статический тип, указанный в разделе ~
Если  является динамическим  вывоз имущества
  или .
 Объявленный тип возврата ,
Если  Имеет доступный экземпляр Getter, названный .
 Тип функции подписи метода ,
Если  Имеет доступный метод экземпляра под названием .
 \FUNCTION{} если  \FUNCTION{}  .
 Тип {} в противном случае.
 Это происходит только тогда, когда   ограниченный.


{%
Обратите внимание, что тип метода отрыва игнорирует
является ли любой данный параметр ковариантным.
Динамический тип функционального объекта
Полученные таким образом данные учитывают ковариацию параметров.%
?

%
Оценка извлечения имущества  формы 
идет следующим образом:

%
Во-первых, выражение  оценивается на предмет .
 быть результатом поиска () метод
()
  в отношении текущей библиотеки .
Если метод поиска удался, то  оценивать на
Метод клозулизации  на объекте 
().

<%
Обратите внимание, что  Это никогда не абстрактный метод.
Потому что метод поиска пропускает абстрактные методы.
Если метод поиска не сработал, например,
потому что есть абстрактная декларация , но нет конкретной декларации,
Мы продолжим следующий шаг.
Однако, поскольку методы и геттеры никогда не перевешивают друг друга,
Взгляд на Геттера тоже обязательно провалится,
и  в конечном итоге будет использован.
Прискорбный вывод заключается в том, что ошибка будет относиться к отсутствующему геттеру.
Вместо того, чтобы пытаться скрыть абстрактный метод.
?

%
В противном случае  - это вызов геттера.
Пусть  будет результатом поиска () геттер
()
  в отношении .
В противном случае тело  Выполняется с помощью \THIS{} Ссылка на .
Значение  Это результат, возвращаемый вызовом к функции геттера.

%Если поиск Геттера провалился,
затем новый экземпляр  заданного класса  Он был создан,
такой, что:


 <<<p001210>> Показатель .
  Оценивает символ .
  Оценивать объект
чей динамический тип ,
который является пустым и неизменяемым.
<<<p001707>  Оценивать объект
чей динамический тип ,
который является пустым и неизменяемым.
  Оценивать объект
чей динамический тип ,
который является пустым и неизменяемым.


%
Тогда метод  смотреть вверх
 и с аргументом ,
и результатом этого вызова является результат оценки .

{%
Ситуация, при которой  Призыв может только возникнуть
когда статический тип  составляет .%
?


 Супер Геттер Доступ и Клозуризация Методов


%
Подумайте об извлечении имущества  формы .

%
Пусть P000497 будет суперклассом немедленно закрывающего класса.
Ошибка времени компиляции  не иметь
доступный способ экземпляра или геттер с именем .
Статический тип  Это:


 Объявленный тип возврата ,
Если  Имеет доступный экземпляр Getter по имени .
 Тип функции подписи метода ,
Если  Имеет доступный метод экземпляра под названием .
 Тип {} в противном случае.
 Это происходит только тогда, когда  \DYNAMIC{} или \FUNCTION.


{%
Обратите внимание, что тип метода отрыва игнорирует
является ли любой данный параметр ковариантным.
Динамический тип функционального объекта
Полученные таким образом данные учитывают ковариацию параметров.%
?

%
Оценка извлечения имущества  из формы 
идет следующим образом:

%
 быть реализацией метода, выполняемого в настоящее время,
 класс, в котором объявлен .
 Это суперкласс P000511.
 быть результатом метода поиска  в 
Что касается нынешней библиотеки .
Если метод поиска удался, то  оценивать на
Метод клозулизации 
в отношении суперкласса 
().

%
В противном случае  - это вызов геттера.
 быть результатом поиска
getter  >> в  в отношении .
Тело  Выполнено с помощью \THIS{} Связано с текущим значением .
Значение  Это результат, возвращаемый вызовом к функции геттера.

{%
Поиск Getter не потерпит неудачу, потому что это ошибка времени компиляции.
извлечение супер свойств члена  когда суперкласс 
не имеет конкретного члена под названием .%
?


 Метод клозулизации инстанций


%
В этом разделе указана динамическая семантика клозулизации методом экземпляра.

{%
Обратите внимание, что необычный случай охватывается неявно с использованием ,
В этом случае в декларации о параметрах типа указывается
и фактические списки аргументов типа, переданные в призывах
Опущены ().%
?

%

является клосуризацией какого-либо метода на каком-либо объекте, определяемом ниже;
или суперклозуляция ().

%
 Будьте объектом, и пусть  быть свежей конечной переменной, связанной с .  на объекте 
определяется как эквивалентный
{(кроме равенства, как указано ниже)}
к:


%    , если  Оригинальное название: P000536
% и  является одним из \code{<<, >, <=, >=, ==, -, +, /, /, *, , ,
% << , ,  (это исключает клозуризацию унарных -).
%  << , если  Он называется P001735.
%  , если  Называется «P001224».
%  , если  Называется «P001225».


где  Метод экземпляра, названный 
с указанием типовых параметров
,
требуемые параметры {p}{1}{n}
именованные параметры \List{p}{n+1}{n+k) с дефолтами {d}{1}{k}
использование  для параметров, значение по умолчанию которых не указано.


где  Метод экземпляра, названный 
с указанием типовых параметров
,
требуемые параметры {p}{1}{n}
и факультативные позиционные параметры
{p}{n+1}{n+k) с дефолтами {d}{1}{k}
использование  для параметров, значение по умолчанию которых не указано.


%
, определяется следующим образом:
Если  является примером недженерического класса, .
В противном случае пусть  быть
формальные параметры типа класса ,
и  являются аргументами фактического типа.
Далее .

{%
То есть мы заменяем формальные параметры типа закрывающего класса,
если есть,
по соответствующим фактическим типам аргументов.%
?

%
Типы параметров  определяются следующим образом:
Пусть декларация метода  Внедрение 
которое вызывается выражением в теле.
 Класс, который содержит .

{%
Обратите внимание, что  является динамическим типом , или его суперклассом.
?

%
Для каждого параметра , , если  ковариантный
()
 Встроенный класс .

{%
Это связано с динамическим типом объекта функции, полученного посредством
Членская клозуляция.
Статический тип выражения, который приводит к клозуризации члена
указывается в другом месте
(,
.
Обратите внимание, что для статического типа игнорируется, является ли параметр ковариантным.
?

%
Если  необычный класс для ,
% % TODO(Eernst): См. понятие nnbd «один и тот же тип».
 аннотация типа, обозначающая тот же тип, что и
который обозначается аннотацией типа на
соответствующее объявление параметров в .
Если это объявление параметров не имеет аннотации типа, то  .

%
В противном случае  представляет собой обобщенное представление общего класса .
 быть формальными параметрами типа ,
 быть фактическим типом аргументов  при .
Далее  аннотация типа, обозначающая
,
где  аннотация типа соответствующего параметра в .
Если это параметрическое объявление не имеет аннотации типа, то  .

%
Есть один способ, которым
Функциональный объект, полученный методом примера, отличается от
Функциональный объект, полученный путем клозулизации функции на
Вышеупомянутая функция дословно:Предположим, что  и  являются объектами,  является идентификатором,
 и  Функциональные объекты
получено путем клозулизации  на  соответственно .
Далее  Оценка истинности
если и только если  и  Это один и тот же объект.

{%
% Опишите последствия для "==" и для "идентичности", для получателей
% и для клозураций.
В частности, две клозулизации метода 
от одного и того же объекта равны,
и две клозулизации метода 
Неидентичные объекты не равны.
Предположим, что  является свежей переменной, связанной с объектом, ,
Из этого следует, что  должно быть ложным
когда  и  Они не связаны с одним и тем же объектом.
Однако реализация Dart не требуется.
канонизировать функциональные объекты,
что означает 
Не гарантируется, что это правда,
Даже если известно, что  и  Они связаны с одним и тем же объектом.%
?

{%
Особый режим равенства в этом случае облегчает
Использование извлеченных функций свойств в API, где обратные вызовы
Слушатели часто должны быть зарегистрированы, а затем незарегистрированы.
Распространенным примером является DOM API в веб-браузерах.%
?


 Суперклозургия


%
Этот раздел определяет динамическую семантику
суперклозуляция.

{%
Обратите внимание, что необычный случай охватывается неявно с использованием ,
в этом случае объявления параметров типа опущены ().%
?

%
Рассмотрим выражения в теле класса  который является
подкласс данного класса ,
где декларация метода, реализующая  существует в ,
Не существует класса  который является
Подкласс  и суперкласса  который реализует .

{%
Рассмотрим ситуацию, когда
Сверхвызов  Выполнить  как указано в .%
?

%

Это клозуляция метода по отношению к классу, как определено далее.
  в отношении класса 
определяется как эквивалентный
\commentary{(кроме равенства, как указано ниже)}
к:


%  <<    Оригинальное название: P000629
% и  является одним из \code{<<, >, <=, >=, ==, -, +, /, \gtilde/, *, , ,
% <<  , .
%  \gtilde\SUPER;\}, если  Называется он «P001766».
%  , если  Называется «P001232».
%    Называется «P001233».


где  Метод экземпляра, названный 
с указанием типовых параметров
,
требуемые параметры \List{p}{1}{n}
именованные параметры \List{p}{n+1}{n+k) с дефолтами \List{d}{1}{k}.


где  Метод экземпляра, названный 
с указанием типовых параметров
,
требуемые параметры {p}{1}{n}
и факультативные позиционные параметры
\List{p}{n+1}{n+k) с дефолтами \List{d}{1}{k}.


{%
Обратите внимание, что суперклозуризация является методом примера клозуризации.
как определено в ()%
?

%
, определяется следующим образом:
Если  Это необычный класс .
В противном случае пусть  быть формальными параметрами типа , быть фактическим типом аргументов \THIS{} при .
Далее .

{%
То есть, мы заменяем формальные параметры типа закрывающего класса, если таковые имеются,
соответствующих фактических аргументов.
Мы должны рассмотреть тип аргументов в отношении конкретного класса.
потому что класс может передавать аргументы другого типа
к своему суперклассу
чем те, которые он получает сам.%
?

%
Типы параметров  определяются следующим образом:
Пусть декларация метода  быть реализацией  в .

%
Для каждого параметра , , если  ковариантный
()
 Встроенный класс .

{%
Это связано с динамическим типом объекта функции, полученного посредством
Суперклозуляция.
Статический тип выражения, порождающий суперклозуризацию
указывается в другом месте
(,
.
Обратите внимание, что для статического типа игнорируется, является ли параметр ковариантным.
?

%
Если  является негенерическим классом для ,
% % TODO(Eernst): См. понятие nnbd «один и тот же тип».
 аннотация типа, обозначающая тот же тип, что и
который обозначается аннотацией типа на
соответствующее объявление параметров в .
Если это параметрическое объявление не имеет аннотации типа, то  .

%
В противном случае  представляет собой обобщенное представление общего класса .
Пусть  будут формальными параметрами типа ,
 быть фактическим типом аргументов  при .
Далее  аннотация типа, обозначающая
,
где  аннотация типа соответствующего параметра в .
Если это параметрическое объявление не имеет аннотации типа, то  является .

%
Есть один способ, которым
Функциональный объект, полученный суперклозурацией, отличается от
Функциональный объект, полученный путем клозулизации функции на
Вышеупомянутая функция дословно:
Предположим, что выражение .  в данном классе
Оценивается в двух случаях, когда 
Связано с  соответственно ,
и результирующие объекты функции  соответственно :
 тогда истинно
если и только если  и  Это один и тот же объект.


 Общий метод обоснования


% % TODO(Eernst): Как и обобщение функций, общий метод
Инстантификация % % зависит от типа вывода. Смотрите комментарий в
%%  для более подробной информации.

%
Общий метод инстанциации является механизмом, который дает
объект негенерической функции,
на основе извлечения имущества, которое обозначает примерный метод клозулизации
(, .

{%
Это механизм, очень похожий на метод клозулизации экземпляра.
Это происходит только в ситуациях, когда
В противном случае произошла бы ошибка времени компиляции.%
?

{%
Сущность обобщенного метода инстанциации
Для того, чтобы получить «проклятые» призывы,
в том смысле, что общий метод экземпляра может получить свой фактический
Аргументы по отдельности при клозурии
(Тогда он должен получить аргументы типа , а не только некоторые из них),
Это приводит к появлению объекта негенерической функции.
Аргументы типа передаются неявно, на основе вывода типа;
Будущая версия Dart может позволить передавать их явно.
?
 Вот пример:




{%
Объект функции, который хранится в  в конце 
динамический тип ,
и получено имплицитноПрохождение фактического типа аргумента  ''
к обозначенному универсальному методу экземпляра,
Таким образом получают негенерическую функцию объекта указанного типа.
%
Обратите внимание, что этот функциональный объект принимает необязательный позиционный аргумент.
Хотя это не является частью
статически известный тип соответствующего способа экземпляра,
ни контекстного типа.%
?

{%
Другими словами, обобщенный метод инстанциации дает функцию
Подпись которого максимально соответствует типу контекста,
но в отношении формы перечня параметров
(то есть количество позиционных параметров и их факультативность);
или набор имен названных параметров,
будет определяться методом подписи фактического метода экземпляра
от данного получателя.
Разница может быть только в том, что фактический тип
% Речь идет о динамическом типе, поэтому исключения для  нет.
подтип данного контекстного типа,
в противном случае заявление о таком методе
Ошибка времени компиляции.%
?

%
 быть выражением извлечения свойств формы
,  или 
(, 
который статически решен для обозначения способа экземпляра, названного ,
 Статический тип .
Рассмотрим ситуацию, когда  является функциональным типом формы
 T_0}{X}{B}{s}{{параметры}}

{(то есть  является общим типом функции)},
Контекстный тип является негенерическим типом функции. .
В этой ситуации происходит ошибка компиляции времени.
(,
,
,
,
,
,
,
,
,
,
,
,

За исключением случаев, когда универсальный тип функции
()
преуспевает, то есть:

%
Вывод типа применяется к  Тип контекста ,
и это увенчалось успехом, приведя фактический список аргументов типа
 T}{1}{s}.

{% Сфера применения .

\def\gmiName{\metavar{gmiName\ensuremath{_{\id}}}}

% для любого заданного названия метода , < - это имя, связанное с
% используется для неявно индуцированного метода, который мы используем для получения функции
% объект, являющийся результатом обобщенного метода инстанциации. Позвоним.
%, что связывает метод «общего метода инстанциации».
%
% Основная идея заключается в том, что мы рассматриваем все библиотеки, включенные в
% компиляции полной программы Dart, а затем мы выбираем суффикс,
% говорят «_*», что является уникальным во всем мире (поскольку нет идентификатора, написанного пользователем);
% может иметь этот суффикс, и компилятор не использует ничего подобного.
% для других целей.
%
% Когда мы нашли такой «глобально свежий суффикс», его легко увидеть.
%, что мы можем использовать его для определения неявно индуцированных методов
%, которые сделают все правильно:
%
% Допустим, что класс "C" имеет общий метод "foo", и он является предметом
% к обобщенному методу.
%
% Затем мы создаем общий метод инстанциации для «foo» в «C», с помощью
% название «foo_*». Если есть подкласс «D» от «C», который имеет приоритет
Процентная декларация «foo», затем мы также генерируем «foo_*» обобщенное представление
Метод % в «D», и что можно вернуть закрытие с другим параметром
Форма списка %, потому что это способ объявления «foo» в «D».
%
% Затем мы можем назвать e.foo_*<...>() где e имеет статический тип C,
% в месте, где исходный код имел «e.foo» и контекст;
Тип % был необычным (поэтому был введен общий метод).
% Это может вызвать «foo_*» в «C», или «D», или в каком-либо другом классе.
% динамический тип результата оценки «е».
%
% В целом это гарантирует получение функционального объекта, параметр которогоФорма списка % точно соответствует форме списка параметров «foo» в
% от динамического типа приемника, как это и должно быть. Не будет никакого
% имя конфликтует с пользовательскими декларациями и основными декларациями
% на самом деле будут переопределять, как они должны, потому что они все используют то же самое
% «глобально свежий суффикс».
%
% Текст ниже пытается сообщить, как это может работать, не будучи слишком
%, что очень похоже на конкретную реализацию.

%
Рассмотрим ситуацию, в которой инстанциация типов функций увенчалась успехом.
Пусть < быть свежим именем, которое связано с ,
Что является частным тогда и только тогда, когда  является частным.
{%
Реализация может использовать, скажем, , когда  является ,
Как известно, свежий, потому что
Пользовательские идентификаторы не могут содержать .%
?
Затем программа изменяется следующим образом:


 Когда  :
Заменить  на \code{?. \gmiName<\List T}{1}{s}>()}.
 Когда  :
Заменить $i$ на \code{$e$.\gmiName<\List{ T}{1}{s}>()}.
 Когда  :
Заменить  на \code{. \gmiName<\List T}{1}{s}>()}.


%
Вставленные выражения не имеют ошибки времени компиляции и могут быть выполнены.
потому что соответствующий общий метод индуцируется неявно.
Мы используем фразу

чтобы обозначить эти неявно индуцированные методы,
и определить метод, который его вызвал, как его
{цель}{общий метод инстанциации! Цель.

%
Предположим, что класс  объявляет общий метод экземпляра, названный ,
с сигнатурой метода, соответствующей типу общей функции ,
Формальные параметры типа ,
и формальные декларации о параметрах {параметры}.
Пусть {аргументы} обозначают соответствующий фактический список аргументов,
прохождения этих параметров.

<%
Например, {параметры} может быть


\code$\PairList{T}{p}{1}{n},\ $\{$T_{n+1}\ p_{n+1} = d_1, \ldots,\ T_{n+k}\ p_{n+k} = d_k$\}


в каком случае <{аргументы}
\code{{p}{1}{n},\ $p_{n+1}$:\ \ $p_{n+k}$:\ $p_{n+k}$.%
?

%
 быть того же типа функции, что и ,
за исключением того, что в нем отсутствуют формальные объявления параметров типа.
{%
Например, если  это
\FunctionTypeSimpleGeneric{\VOID}{,  \EXTENDS\ num}{X x, List<Y>ys}
 это
\FunctionTypeSimple{}{X x, List<Y> ys}.
Обратите внимание, что  Обычно содержит переменные свободного типа.%
?

%
Метод экземпляра с именем \gmiName{} После этого он становится неявным,
с тем же поведением, что и следующее заявление
(за исключением равенства возвращаемого объекта функции,
которые указаны ниже:

% Использовать ...>, поскольку может существовать параметр .


%
 Пример класса, который содержит
Неявно индуцированная декларация 
Как описано выше.
%
Рассмотрим ситуацию, когда программа оценивает
два вызова этого метода с одним и тем же приемником ,
с аргументами типа, действительные значения которых
% % TODO(Eernst): См. понятие nnbd «один и тот же тип».
одинаковых типов {t}{1}{s} для обоих вызовов,
И предположим, что призывы вернулись.
Примеры  соответственно .
%
При этом гарантируется, что  и  равными
По словам оператора {==}.
Не уточняется, является ли

Оценивается как {} или .
} % Конец диапазона для .

{%
Никакое понятие равенства не подходит для разных получателей.
и когда приводятся различные аргументы,
потому что выполнение двух функций объектов, которые отличаются в этих способах
могут иметь различные побочные эффекты и возвращать разные результаты.
при исполнении, начиная с точно такого же состояния.%
?


<\subsection{Назначение}


%
Задание изменяет значение, связанное с переменной,
Или взывает к сеттеру.


<назначение оператора> ::== "
 <compoundAssignmentOperator>

<compoundAssignmentOperator> ::= '*= "
</= "
 <~/=
\%=
τ+=τ
 "
= "
= "
= "
\alt\&=
 "
 "
<<??= "


%
\Case{{=$e$>
Подумайте о назначении  формы \code{ = $e$},
где  - идентификатор.
Выполните лексический поиск  Местонахождение: .



Когда лексический поиск дает декларацию  локальной переменной 
{(который может быть формальным параметром)},
Ошибка времени компиляции возникает, если  является окончательным
или если статический тип  не присваивается заявленному типу .

Когда лексический поиск дает декларацию 
не является локальной переменной,
Он гарантированно будет сеттером
{(что может быть явно или косвенно вызвано переменной)}
Потому что другие декларации не имеют названия.
Форма .

Если  является декларацией статического сеттера в классе или смеси 
 рассматривается как
()
присваивание \code{$C$.\id{} = $e$}.

{%
Дальнейший анализ, а также оценка \code{$C$.\id{} = $e$}
Доходы, как указано в другом месте.%
?

В противном случае возникает ошибка времени компиляции.
Если только статический тип  присваивается параметру типа .

Когда лексический взгляд ничего не дает,
 рассматривается как
()
\code{\THIS.\id{}=$e$}.

{%
В данном случае известно, что  Имеет доступ к 
(),
Интерфейс ограждающего класса имеет член, названный ,
или имеется соответствующее расширение с членом, названным .
Статический анализ и оценка продолжаются
\code{\THIS.\id{}=$e$,
поэтому нет необходимости дополнительно указывать лечение .%
?

{%
Лексический поиск никогда не может дать префикс импорта.
Потому что у них никогда нет названия формы. .%
?


%
Во всех случаях
{(независимо от того,  является локальной переменной и т.д.
Статический тип  Статический тип .

%
Оценка присвоения $a$ формы \code{\id{} = $e$}
действует следующим образом.
Выполните лексический поиск  Местонахождение: .



В случае, когда лексический поиск дает
Декларация  локальной переменной ,
{(который может быть формальным параметром)},
Выражение  оценивается на предмет ,
и переменной  Обозначается как .
 оценивает объект 
().
В случае, когда лексический поиск 
Местонахождение: 
дает декларацию ,
 обязательно является установщиком верхнего уровня 
{(возможно, неявно вызванное переменной)}.

Выражение  оценивается на предмет .
Далее сеттер  на него ссылаются
с его формальным параметром, связанным с .
Далее  оценивает объект .

{%
 не может быть статическим установщиком в классе ,
 Затем рассматривается как
\code{.=$e$.
которые указаны в другом месте.%
?
<<<p001901>
{%
Случай, когда лексический поиск 
ничего не может произойти,
Поскольку этот случай рассматривается как
\code{.=$e$,
чья оценка указана в другом месте.%
?



%
\Case{\code{\id{}=$e$>
Подумайте о назначении  формы \code{$p$.\id{} = $e$},
где  является префиксом импорта и  является идентификатором.

%
происходит ошибка времени компиляции,
Если только  имеет член, который является сеттером  Оригинальное название: P001258
{(которые могут быть косвенно вызваны переменной декларацией)}
Статический тип 
присваивается параметру типа .

%
Статический тип  Статический тип .

%
Оценка назначения  по форме \code{.\id{} = $e$}
идет следующим образом:
Выражение  оценивается на предмет .
Затем сеттер обозначается  на него ссылаются
с его формальным параметром, связанным с .
 оценивает объект .
<<<p001911>

%
\Case{
Рассмотреть задание  из формы .
Ошибки времени компиляции, которые могут быть вызваны
 Они также генерируются в случае .
Статический тип  Статический тип .

%
Оценка присвоения  формы 
идет следующим образом:
Если  Тип буквальный,  Это эквивалентно P001264.
В противном случае оценка  к объекту .
Если  является нулевым объектом,  оценивает нулевой объект ().
В противном случае пусть  быть свежей переменной, связанной с 
оценивать  Объект .
Далее  Показатель .


%

Подумайте о назначении  формы .
 Статический тип .
Если  , дальнейшие проверки не проводятся.
В противном случае это ошибка времени компиляции, если
 Имеет доступный набор экземпляров под названием .
Это ошибка времени компиляции, за исключением статического типа 
может быть отнесен к заявленному типу формального параметра указанного установщика.
Является ли  
Статический тип  Статический тип .

%
Оценка назначения формы 
идет следующим образом:
Выражение  Оценивается на объект .
Выражение  Оценивается на объект .
Затем сеттер  Проверено ()
в  в отношении текущей библиотеки.
% Эта ошибка может возникнуть из-за неявных слепков и динамических вызовов.Ошибка динамического типа, если динамический тип 
не является подтипом фактического параметра типа упомянутого затворника
().
В противном случае тело сеттера выполняется с
его формальный параметр, связанный с  и  Ссылка на .

%
Если поиск сеттера не удался, то новый экземпляр  из
Предопределенный класс  Он создан таким образом, что:


  Оценка до .
  Оценивает символ .
<<<p001921>  Оценивать объект
чей динамический тип ,
которая неизменяема и содержит те же объекты, что и
.
  Оценивать объект
чей динамический тип ,
который является пустым и неизменяемым.
  Оценивать объект
чей динамический тип ,
который является пустым и неизменяемым.


%
Тогда метод  Об этом сообщает 
Призыв к спору, .

{%
Ситуация, при которой  Призыв может только возникнуть
когда статический тип  .%
?

%
Значение выражения присваивания 
Независимо от того, был ли поиск сеттера неудачным или успешным.


%

Подумайте о назначении  формы .
 Будьте суперклассом немедленного замкнутого класса.
Ошибка времени компиляции  не иметь
конкретный доступный набор экземпляров, названный .
В противном случае, это ошибка времени компиляции, если статический тип 
не может быть присвоен статическому типу формального параметра указанного сеттера.
Статический тип  Статический тип .

%
Оценка назначения формы 
идет следующим образом:
 Будьте текущим методом выполнения и позвольте  быть
Класс, в котором  Его посмотрели вверх.
 Это суперкласс P000893.
Выражение  оценивается на предмет .
Затем сеттер  Проверено ()
в  в отношении текущей библиотеки.
Тело  выполняется с формальным параметром, связанным с 
и  Связано с текущим значением .

<{%
Поиск сеттера не провалится, потому что это ошибка времени компиляции.
когда ни один конкретный сеттер не назван  существует в .%
?

%
Значение выражения присваивания .

%
% Эта ошибка может произойти из-за неявных отливок и ковариации смешивания.
Ошибка динамического типа  не является нулевым объектом ()
и динамический тип  это
не является подтипом фактического типа формального параметра 
() в .
<<<p001931>

%

Подумайте о назначении  из формы .
 Статический тип .
Если  , дальнейшие проверки не проводятся.
В противном случае это ошибка времени компиляции, если
<<<p000011> Метод получил название .
 быть статичным типом
первый формальный параметр метода ,
 Статический тип второго.
Это ошибка времени компиляции, за исключением статического типа  
может быть присвоено  соответственно .
Является ли  
Статический тип  Статический тип .

%
Оценка уступки  
идет следующим образом:
Оценить  на объект ,
затем оценить  для объекта ,
И, наконец, оценить  для объекта .
Назовите метод  на 
 В качестве первого аргумента:  как второй аргумент.
Далее  Показатель .
% Стоит ли добавлять: Это динамическая ошибка, если  оценивать на
Постоянный список или карта?


%

Подумайте о назначении  формы .
 Будьте суперклассом немедленного замкнутого класса.
Ошибка времени компиляции  не имеет метода .
В противном случае пусть  быть статичным типом
первый формальный параметр метода ,
 Статический тип второго.
Это ошибка времени компиляции, если статический тип  соответственно 
не может быть присвоено  соответственно .
Статический тип  Статический тип .

%
Для оценки присваивается форма 
Это эквивалентно выражению .



 Назначение соединения


%

Рассмотрим сложное назначение  
где  является идентификатором или идентификатором, квалифицированным префиксом импорта.
Ошибки времени компиляции, которые могут быть вызваны

Они также генерируются в случае .
Статический тип  это
наименьшая верхняя граница статического типа  и статический тип .

%
Оценка комплексного назначения  
идет следующим образом:
Оценить  на объект .
Если  не является нулевым объектом (),  Показатель .
В противном случае оценка  к объекту ,
А потом <<<p000079> Оценка до .


%
\Case{\code{. ??= }>
Рассмотрим сложное назначение  формы 
где  Тип буквальный
которые могут быть или не быть квалифицированы префиксом импорта.
Ошибки времени компиляции, которые могут быть вызваны

Они также генерируются в случае .
Статический тип  является наименьшей верхней границей
Статический тип  и статический тип .

%
Оценка комплексного назначения  формы \code{. ??= }
где  является типом буквальных поступлений следующим образом:
Оценка  Объект .
Если  не является нулевым объектом (),  Показатель .
В противном случае оценка  для объекта ,
 Показатель .


%
\Case\code{. ??= }
Рассмотрим сложное назначение  из формы . Статический тип   Будьте свежей переменной типа .
За исключением ошибок внутри  и ссылки на имя ,
Ошибки времени компиляции, которые могут быть вызваны

Они также генерируются в случае .
Более того, это ошибка времени компиляции, если  не имеет геттера с именем .
Статический тип  является наименьшей верхней границей
Статический тип  и статический тип .

%
Оценка комплексного назначения  формы 
идет следующим образом:
Оценить  на объект .
 быть свежей переменной, связанной с .
Оценить  на объект .
Если  не является нулевым объектом (),  Показатель .
В противном случае оценка  к объекту ,
 Показатель .


%
\Case{
Рассмотрим сложное назначение  формы .
Ошибки времени компиляции, которые могут быть вызваны

Они также генерируются в случае .
Кроме того, это ошибка времени компиляции
если статический тип  Не имеет значения .
Статический тип  является наименьшей верхней границей
Статический тип  и статический тип .

%
Оценка комплексного назначения  в форме

идет следующим образом:
Оценить  на объект  а затем оценить  Объект .
Позвоните по телефону  Метод  с аргументом ,
 быть возвращенным объектом.
Если  не является нулевым объектом (),  Показатель .
В противном случае оценка  для объекта 
Затем нажмите  Метод 
 В качестве первого аргумента:  как второй аргумент.
 Показатель .


%
\Case{
Рассмотрим сложное назначение  из формы .
Ошибки времени компиляции, которые могут быть вызваны

Они также генерируются в случае .
Кроме того, точно такие же ошибки времени компиляции, которые могут быть вызваны
Оценка выражения 
Они также генерируются в случае .
Статический тип  является наименьшей верхней границей
Статический тип  и статический тип .

%
Оценка комплексного назначения  
идет следующим образом:
Оценить  на объект .
Если  не является нулевым объектом (), затем  Показатель .
В противном случае оценка  ,
 Показатель .


%
\Case{
Рассмотрим сложное назначение  формы .
Ошибки времени компиляции, которые могут быть вызваны

Они также генерируются в случае .
% Примечание: Мы используем статический тип  вместо того, чтобы
%  Хотя последнее было бы проще. Это потому, что
% первое останется верным, если будет введено NNBD, и поскольку оно уменьшает
% от количества синтетического синтаксиса.
Статический тип  является наименьшей верхней границей
Статический тип  и статический тип .

%
Оценка комплексного назначения  формы \code{?. ??= }
идет следующим образом:
Оценить  на объект .Если  является нулевым объектом (), затем  Оценивает нулевой объект.
В противном случае пусть  быть свежей переменной, связанной с .
Оценка  Объект .
Если  не является нулевым объектом (), затем  Показатель .
В противном случае оценка  к объекту ,
 Показатель .
<<<p001807>

%

Составное назначение формы 
где  Тип буквальный
который может или не может быть квалифицирован импортным префиксом
Это эквивалентно выражению .


%
\Case{
Для любого другого действующего оператора ,
a сложное присвоение формы 
эквивалентен ,
где  является идентификатором или идентификатором, квалифицированным префиксом импорта.


%
\Case{
Составное назначение формы 
где  Тип буквальный
который может или не может быть квалифицирован импортным префиксом
Это эквивалентно P001369.


%
\Case{
Рассмотрим сложное назначение  формы .
 Свежие переменные, статическим типом которых является статический тип .
За исключением ошибок внутри  и ссылки на имя ,
Ошибки времени компиляции, которые могут быть вызваны

Они также генерируются в случае .
Статический тип  Статический тип .

%
Оценка комплексного назначения  формы 
идет следующим образом:
Оценить  на объект   быть свежей переменной, связанной с .
Оценить  на объект $r$
 Показатель .


%
\Case{
Рассмотрим сложное назначение  формы .
Пусть  и  Будьте свежими переменными
где статическим типом первого является статический тип 
статическим типом последнего является статический тип .
За исключением ошибок внутри  и  и ссылки на
имена  и ,
Ошибки времени компиляции, которые могут быть вызваны

Они также генерируются в случае .
Статический тип  Статический тип .

%
Оценка комплексного назначения  Форма

идет следующим образом:
Оценить  на объект  и оценить  Объект .
 и  быть свежими переменными, связанными с  и  соответственно.
Оценить  на объект ,
А потом  Показатель .


%
\Case\code{?. = }
Рассмотрим сложное назначение  формы .
Ошибки времени компиляции, которые могут быть вызваны

Они также генерируются в случае .
Статический тип  Статический тип .

%
Оценка комплексного назначения  в форме

идет следующим образом:
Оценить  на объект .
Если  является нулевым объектом, то  оценивает нулевой объект ().
В противном случае пусть  быть свежей переменной, связанной с .
Оценить  на объект .
Далее  Показатель .


%

Составное назначение формы 
где  Тип буквальныйЭто эквивалентно выражению .



<{условный}


%
 Оценка одного из двух выражений
основанный на боолеиновом состоянии.


<условный Выражение>::= <ifNullExpression>
\gnewline{} (?<expressionWithoutCascade>: <expressionWithoutCascade>)?
<<<p001202>

%
Оценка условного выражения  формы 
идет следующим образом:

%
Во-первых,  оценивается на предмет .
% Эта ошибка может произойти из-за неявных отливок и нулевых.
Это динамическая ошибка, если тип времени выполнения  Это не P001391.
Если  , то значение  это
Результат оценки выражения .
В противном случае значение  является результатом оценки выражения .

%
Если  показывает локальную переменную  имеет тип ,
Тип  Известно, что  в ,
Если ни одно из следующих утверждений не верно:


  потенциально мутирует в ,
  потенциально мутирует в пределах функции
, где  объявляется, или
  доступ к функции, определенной в  и
 потенциально мутирует в любом месте в области .


%
Ошибка времени компиляции, если
Статический тип  Не может быть присвоено .
Статический тип  наименьшая верхняя граница () из
Статический тип  и статический тип .


 Если нулевые выражения


%
 оценивает выражение и
Если результатом является нулевой объект (), оценивает другой.


<ifNullExpression> ::= <logicalOrExpression> (??? <logicalOrExpression>)


%
Оценка выражения  формы 
идет следующим образом:

%
Оценить  на объект .
Если  не является нулевым объектом (), затем  Показатель .
В противном случае оценка  к объекту ,
А потом  Показатель .

%
Статический тип  наименьшая верхняя граница () из
Статический тип  и статический тип .


 Логические булевы выражения


%
Логические булевы выражения объединяют булевы объекты с помощью
Булевы операторы соединения и разъединения.


<logicalOrExpression> ::= \gnewline{}
<logicalAndExpression> (<logicalAndExpression>)

<логический AndExpression> ::= <equalityExpression> ( <equalityExpression>)


%
 является выражением равенства
(),
или вызов логического булевого оператора на
Выражение  с аргументом .

%
Оценка логического булевого выражения  формы  причины
Оценка  Объект .
% Эта ошибка может произойти из-за неявных отливок и нулевых.
Это динамическая ошибка, если тип времени выполнения  Это не P001394.
Если  , то результат оценки  ,
иначе  оценивается на предмет .
% Эта ошибка может произойти из-за неявных отливок и нулевых.
Это динамическая ошибка, если тип времени выполнения  Это не P001395.
В противном случае результат оценки  является .

%
Оценка логического булевого выражения  формы Причин оценки  Производить объект .
% Эта ошибка может произойти из-за неявных отливок и нулевых.
Это динамическая ошибка, если тип времени выполнения  Это не P001396.
Если  , то результат оценки  ,
иначе  Оценивается на объект .
% Эта ошибка может произойти из-за неявных отливок и нулевых.
Это динамическая ошибка, если тип времени выполнения  Это не P001397.
В противном случае результат оценки  является .

%
Логическое булево выражение  формы 
показывает локальную переменную  имеет тип 
если выполняются оба следующих условия:


 Либо  Показывает, что  имеет тип 
 Показывает, что  имеет тип .
  не мутирует в  или в пределах функции
За исключением того, где объявлен .


%
Если  показывает локальную переменную  имеет тип ,
Тип  Известен как  в ,
Если ни одно из следующих утверждений не верно:


% % Первый пункт здесь не нужен для звука,
Сохраняется в основном для обратной совместимости.
% % Если  Показывает, что  имеет тип , затем любое назначение
% % в  Это не является недействительным.
Удаление линии видно в семантике, потому что:
% % от числа x;
%% (x??=42) null & x is int & x.toRadixString(16)!=""
% допускается без линии и не допускается.
% % На момент написания анализатор запрещает код.
  потенциально мутирует в ,
  потенциально мутирует в ,
  потенциально мутирует в пределах функции
где  объявляется, или
  доступ к функции, определенной в  и
 Потенциально мутирует в любом месте в области .


%
Ошибка времени компиляции, если
Статический тип  не может быть присвоено 
или если статический тип  не может быть назначено на .
Статический тип логического булевого выражения .


{Равенство}


%
Выражения равенства проверяют объекты на равенство.


<equalityExpress> ::= <
<relationalExpression> (<equalityOperator> <relationalExpression>)
 \SUPER{}<equalityOperator> <Реляционное выражение>

<equalityOperator> ::=== "
 >>!= "


%
 является либо реляционным выражением
(),
или вызов оператора равенства на любой 
или выражение , с аргументом .

%
Оценка выражения равенства  формы 
идет следующим образом:


 Выражение  Оценивается на объект .
 Выражение  оценивается на предмет .
 Если  или  является нулевым объектом (),
 Показатель {}  и  является нулевым объектом
и {} в противном случае.
В противном случае
 Оценка  эквивалентно методу вызова
.
<<<p001209>

%
Оценка выражения равенства  в форме
\code{{== 
идет следующим образом:


 Выражение  оценивается на предмет . Если \THIS{{} или  является нулевым объектом (),
 Оценки для оценок \TRUE{}
Если оба \THIS{{} и $o$ является нулевым объектом
и {} в противном случае.
В противном случае
 Оценка  эквивалентно методу вызова
.
<<<p001210>

{%
В результате приведенного выше определения,
Пользовательский {=== Методы могут предполагать, что их аргумент не является нулевым.
Избегайте стандартного прелюдия котельной:



Еще один вывод заключается в том, что необходимости никогда не существует.
Использование  Для проверки на соответствие ,
Никто не должен беспокоиться о том, писать ли
\NULL{{= или == %
?

%
Выражение равенства формы  является эквивалентным
Выражение .
Выражение равенства формы \code{\SUPER{}!=  является эквивалентным
выражение <! (\SUPER{}) .

% выражение  оценивается на объект ;
% от выражения  оценивается на предмет .
% Далее, если $o_1$ и $o_2$ Это один и тот же объект,
% далее  , в противном случае  Показатель .

%
Статический тип выражения равенства .


{Реляционные выражения}


%
Реляционные выражения вызывают реляционных операторов на объектах.


<реляционный Выражение>::= <bitwiseOrExpression> <<
(<typeTest> | <typeCast> | <relationalOperator> <bitwiseOrExpression>)?
 \SUPER{} <relationalOperator> <bitwiseOrExpression>

<relationalOperator> ::=>= "
 "
 ><=
< "


%
 Это либо немногословное выражение
(),
или вызов реляционного оператора на любой 
или выражение , с аргументом .

%
Реляционное выражение формы    является эквивалентным
Метод вызова .
Реляционное выражение формы    является эквивалентным
Метод вызова .


<{Битвы выражений}


%
Битовые выражения вызывают битовые операторы на объектах.


<bitwiseOrExpression> ::= <<
<bitwiseXorExpression> (<bitwiseXorExpression>)*
 \SUPER{} (' |' <bitwiseXorExpression>) +

<bitwiseXorExpress> ::= <<
<bitwiseAndExpression> ('^' <bitwiseAndExpression>)
\alt \SUPER{} ('^' <bitwiseAndExpression>)+

<bitwiseAndExpression> ::= <shiftExpression> ( <shiftExpression>) *
 < ( <shiftExpression>)+

<bitwiseOperator> ::= "
 "
 "
<<<p001212>

%
 является выражением сдвига (),
или обращение к оператору bitwise
В-третьих, на любой В.М.: В.М.: В.М.: В.М. В.М.: В.М., в частности, в связи с тем, что
с аргументом .

%
Пошаговое выражение формы  является эквивалентным
Метод вызова .
Побитовое выражение формы \code{\SUPER{}   является эквивалентным
Метод вызова .

{%
Должно быть очевидно, что правила статического типа для этих выражений
Определяется эквивалентностью выше --ergo,
Правила типа для вызова метода и
подписи операторов по типу .
То же самое относится и к аналогичным ситуациям в рамках данной спецификации.%
?


 Сдвиг!
<<p001773>

%Выражения сдвига вызывают операторов сдвига на объектах.


<shiftExpression> ::= \gnewline{}
<additiveExpression> (<shiftOperator> <additiveExpression>)*
 \SUPER{} (<shiftOperator> <additiveExpression>)+

<shiftOperator> ::= "
 "
 "


%
 является либо аддитивным выражением
(),
или вызов оператора смены
либо << или выражение ,
с аргументом .

%
Сдвиг выражения формы    является эквивалентным
Метод вызова .
Сдвиг выражения формы    является эквивалентным
Метод вызова .

{%
Обратите внимание, что это определение подразумевает порядок оценки слева направо
Среди сменных выражений:

оценивается как
 который эквивалентен
.
То же самое относится к аддитивным и мультипликативным выражениям.
?


<<<p001909>> Аддитивные выражения


%
Аддитивные выражения вызывают операторов сложения на объектах.

<<<<<<<<<<<<<<<<<<<<<<<<<<<<<<<<<<<<<<<<<<<<<<<<<<<<<<<<<<<<<<<<<<<<<<<<<<<<<<<<<<<<<<<<<<<<<<<<<<<<<<<<<<<<<<<<<<<<<<<<<<<<<<<<<<<<<<<<<<<<<<<<<<<<<<<<<<<<<<<<<<<<<<<<<<<<<<<<<<<<<<<<<<<<<<<<<<<<<<<<<<<<<<<<<<<<<<<<<<<<<<<<<<<<<<<<<<<<<<<<<<<<<<<<<<<<<<<<
<additiveExpression> ::= <multiplicativeExpression>
\gnewline{} (<additiveOperator> <multiplicativeExpression>)*
<<<p001911> \SUPER{} (<additiveOperator> <multiplicativeExpression>)+

<additiveOperator> ::=> "
 "


%
 является мультипликативным выражением
(),
или вызов аддитивного оператора
либо << или выражение ,
с аргументом .

%
Добавочное выражение формы   является эквивалентным
Метод вызова .
Добавочное выражение формы \SUPER{}   является эквивалентным
Метод вызова .

%
Статический тип аддитивной экспрессии обычно определяется
по подписи, указанной в декларации используемого оператора.
Призывы операторов  и  из
класс ,  и 
Обрабатываются специально чекером.

%
 быть аддитивным выражением формы ,
 Статический тип ,
 Контекстный тип .

Если \SubtypeNE{T}{ и не \SubtypeNE{T}{}
Тип контекста  определяется следующим образом:


<< Если \SubtypeNE{}{C}, а не {}{C},
и \SubtypeNE{T}{}}
Контекстный тип  является .
<< Если \SubtypeNE{}{C}, а не {}{C},
и не \SubtypeNE{T}{
Тогда контекстный тип  является .
<< В противном случае тип контекста  является .


%
Далее  Статический тип .
Если {T}{ и не \SubtypeNE{T}{\code{Never}>
 присваивается ,
Статический тип  определяется следующим образом:


< Если \SubtypeNE{T}{
Статический тип  является .
<< В противном случае, если \SubtypeNE{S}{}
и не  >>{S}{,
Статический тип  является .
< В противном случае, если \SubtypeNE{T}{\code{int}},
\SubtypeNE{S}{}} и не \SubtypeNE{S}{,
Статический тип  является .<< В противном случае статический тип  является .


 Множественные выражения


%
Умножительные выражения вызывают операторы умножения на объекты.


<multiplicativeExpress> ::= <
<unaryExpression> (<multiplicativeOperator> <unaryExpression>)*
<<<p001941> \SUPER{} (<multiplicativeOperator> <unaryExpression>)+

<multiplicativeOperator> ::= '* "
 "
 "
</> "


%
 является унарным выражением
(),
или вызов мультипликативного оператора
либо << или выражение ,
с аргументом .

%
Умножительное выражение формы    эквивалентно
Метод вызова .
Множественное выражение формы    является эквивалентным
Метод вызова .

%
Статический тип мультипликативного выражения обычно определяется
по подписи, указанной в декларации используемого оператора.
Призывы операторов  и  из
класс ,  и 
Обрабатываются специально чекером.

%
 быть мультипликативным выражением формы 
где  является одним из  или ,
 Статический тип ,
 Контекстный тип .

Если {T}{ и не \SubtypeNE{T}{}
Тип контекста  определяется следующим образом:


<< Если \SubtypeNE{}{C}, а не {}{C},
и \SubtypeNE{T}{}}
Контекстный тип  является .
<< Если \SubtypeNE{}{C}, а не {}{C},
и не {T}{
Контекстный тип  является .
<< В противном случае контекстный тип  является .


%
Пусть далее  будет статичным типом .
Если {T}{ и не \SubtypeNE{T}{\code{Never}>
 присваивается ,
Статический тип  определяется следующим образом:


<< Если {T}{
Статический тип  является .
< В противном случае, если {S}{}
и не  >>{S}{,
Статический тип  является .
< В противном случае, если \SubtypeNE{T}{\code{int}},
\SubtypeNE{S}{}} и не \SubtypeNE{S}{\code{Never},
Статический тип  является .
<< В противном случае статический тип  является .



 Унарные выражения


%
Унарные выражения вызывают унарных операторов на объектах.


<unaryExpression> ::= <prefixOperator> <unaryExpression>
 <awaitExpression>
 <postfixExpression>
 (<minusOperator> | <tildeOperator>) <
 <incrementOperator> <assignableExpression>

<prefixOperator> ::= <minusOperator>
< <отрицатель
<<<p001978> <tildeOperator>

<minusOperator> ::= "

<negationOperator> ::= '! "

<tildeOperator> ::=~ "


%
 Это либо постфиксное выражение
(),
Выражение ожидания ()
или вызов оператора префикса на выражениеили вызов унарного оператора на любой  Выражение .

<<<p001700>%
Выражение  рассматривается как
()
.

%
Выражение формы  рассматривается как
.
Выражение формы \code{-{}-} трактуется как
.

%
 быть выражением формы 
где  является буквальным целым числом () с числовым целым числом ,
со статическим типом контекста .
Если  Применяется к  и  не присваивается ,
Статический тип  является ;
Статический тип  является .

%
Если статический тип  , то  оценивать на
Пример  класс, представляющий числовое значение .
Если  Это ноль и P001493. класс может представлять собой отрицательное нулевое значение,
Вместо этого полученный экземпляр представляет собой отрицательное нулевое значение.
Ошибка времени компиляции, если целое число  не могут быть представлены
Именно на примере .

%
Если статический тип  , то  оценивать на
Для примера  класс, представляющий числовое значение .
Если  равна нулю, полученный экземпляр представляет собой
{отрицательное} нулевое двойное значение, .
Это ошибка времени компиляции, если целое число  не могут быть представлены
Именно на примере .
{%
Мы рассматриваем  {как будто} это одно целое число
с отрицательным числовым значением.
Мы не оцениваем  индивидуально, как выражение,
или заботится о своем статическом типе.
?

%
Любое другое выражение формы  является эквивалентным
Метод вызова .
Выражение формы  является эквивалентным
метод вызова () .


 Подождите выражения.


%
 позволяет кодировать
Управление выходом до асинхронной операции
()
Готово.


<awaitExpress> ::= <unaryExpression>


%
%
 Выражение формы .
 Статический тип .
Статический тип  Затем < С.
().

%
Оценка  идет следующим образом:
Во-первых, выражение  оценивается на предмет .
Пусть \DefineSymbol{T} будет \flatten С.
Если тип времени выполнения  является подтипом ,
Тогда пусть \DefineSymbol{f} be $o$
В противном случае пусть  быть результатом создания
новый объект с использованием конструктора 
 В качестве аргумента.

%
Далее, поток, связанный с
Внутренняя замкнутая асинхронная < петля
(),
Если есть, то приостанавливается.
Нынешнее призвание функционального тела, непосредственно заключающего в себе 
Приостанавливается до тех пор, пока  Готово.
Спустя некоторое время после завершения операции  управление возвращается к текущему вызову.
Если  Завершилась ошибка  и след стека ,
 броски  и 
().
Если  завершается объектом ,  Показатель .

<{%
Использование \flattenName{} для поиска 
и, следовательно, определить динамический тип испытания
Это означает, что мы ожидаем будущего в каждом случае, когда этот выбор будет правильным.
?

<{%
Интересный случай на краю этого компромисса - когда Имеет статический тип .
Можно сказать, что за этим типом стоит намерение
Значение  является ,
Или это  Который не является будущим,
Или это .
Так что, предположительно, мы должны ждать первого рода,
Второе и третье должны быть без изменений.
Однако второй вид может быть .
Этот объект не является , и это не ,
Таким образом, он {должен} считаться во второй группе.
Тем не менее, \code{FutureOr<Object>?}>
Таким образом, мы ожидаем  .
Мы выбрали эту семантику, потому что это было наименьшее изменение.
относительно семантики в более ранних версиях Дарта,
А также потому, что позволяет простое правило:
Тип  Используется для того, чтобы решить, есть ли
Будущее (если оно есть) ждет, и исключений нет.
В таких случаях, как этот пример, тип подразумевает, что
 Не следует ждать.
Таким образом, мы ожидаем каждое будущее, которое мы можем разумно ожидать.
?

<{%
Выражение ожидания может происходить только в объявленной функции.
асинхронно.  идентификатор не имеет особого значения
в контексте нормальной функции, поэтому случаи 
В этих функциях не вводится выражение ожидания.
Тем не менее,  может быть действительным вызовом функции в
неасинхронные функции.%
?

<<%
Ожидаемое выражение не могло осмысленно возникнуть в синхронной функции.
Если бы такая функция была приостановлена в ожидании будущего,
Он больше не будет синхронным.
?

<%
Это не ошибка времени компиляции, если тип  не является
супертип или подтип .
Инструменты могут дать подсказку в таких случаях.%
?


 Постфиксные выражения


%
Выражения Postfix вызывают операторов postfix на объектах.


<postfixExpression> ::= <assignableExpression> <postfixOperator>
 <primary> <selector> *

<postfixOperator> ::= <IncrementOperator>

<constructorInvocation> ::= <
<typeName> <typeArguments> <identifier> <аргументы>

<selector> ::= '! "
 <assignableSelector>
 <Аргументы>

<argumentpart> ::=
<тип Аргументы>. <аргументы>

<incrementOperator> ::=++ "
<- "


%
 является первичным выражением;
функция, метод или вызов Геттера;
Призыв названного конструктора;
или вызов оператора постфикса на выражение .
Все, кроме последних двух, указаны в другом месте.

%
 Призывы конструктора
Рассмотрим <  формы
<\code{$n$<<\metavar{тип Аргументы}>. \id(\metavar{аргументы})}.
Если  не обозначает класс 
который объявляет конструктора по имени ,
Происходит ошибка времени компиляции.

%
В противном случае, если  происходит в постоянном контексте
()
Затем  рассматривается как ,
Если  Не происходит в постоянном контексте.
Затем  рассматривается как .

% Мы можем добавить поддержку передачи аргументов типа статичным методам.
%, но это не та функция, которая появится в ближайшее время.
{%
Обратите внимание, что  Не может быть ничего, кроме создания экземпляра.
(постоянный или нет)
 приводит фактические аргументы типа ,
который не поддерживается, если  обозначает префикс библиотеки,
и если  Статический метод вызова.%
?


%
 \code{-{}-}}
Рассмотрим выражение postfix  формы ,где  является идентификатором и  является либо \lit{++} или \lit{-{}-}.
Ошибка времени компиляции возникает, если только  обозначает переменную,
 обозначает геттер и имеется связанный с ним сеттер .
 Статический тип переменной  Возвратный тип геттера.
Ошибка времени компиляции возникает, если  не является 
 У вас нет оператора {+} (когда \op{} является \lit{++})
или оператор \lit{-} (когда \op{} является \lit{-{}-}),
или если тип возврата этого оператора не присваивается
переменная соответственно типу аргумента множителя.
Ошибка времени компиляции возникает, если  не может быть присвоено
Тип параметра указанного оператора.
Статический тип  является .

%
Оценка выражения postfix 
Форма  соответственно \code{-{}-}
где  является идентификатором, действует следующим образом:
Оценить  на объект  Пусть  быть свежей переменной, связанной с .
% % TODO(Eernst): Для того, чтобы исключить синтаксический сахар, мы должны
%% переписывают это, чтобы указать эффект непосредственно (случаи, чтобы помнить: 'v' - это
% % локальная переменная/параметр, экземпляр/статический/глобальный/импортный геттер.
Оценка  соответственно .
 Показатель .

{%
Вышесказанное гарантирует, что если оценка включает в себя геттер,
Он называется ровно один раз.
Аналогично в случаях ниже.%
?


%
 . -{}-}}
Рассмотрим постфиксное выражение  формы ,
где  является буквальным и  является либо \lit{++}, либо {-{}-}.
Ошибка времени компиляции возникает, если только  обозначает статический геттер
и имеется связанный статический установщик 
{(возможно, неявно вызванная статической переменной)}.
 Возвратный тип указанного геттера.
Ошибка времени компиляции возникает, если  не является 
 У вас нет оператора {+} (когда \op{} является \lit{++})
или оператор {-} (когда \op{} является \lit{-{}-}),
или если тип возврата этого оператора не присваивается
Тип аргументации сеттера.
Ошибка времени компиляции возникает, если  не может быть присвоено
Тип параметра указанного оператора.
Статический тип  является .

%
Оценка выражения postfix 
Форма  соответственно . -{}-}
где  является типом буквальных поступлений следующим образом:
Оценить  на объект 
 быть свежей переменной, связанной с .
Оценка  соответственно .
 Показатель .


%
\code{.++} . -{}-}}
Рассмотрим выражение postfix  формы 
где  {++} или {-{}-}.
 Статический тип .
Ошибка времени компиляции возникает, если  иметь
Геттер по имени  и сеттер по имени 
{(возможно, неявно вызванная переменной экземпляра)}.
 Возвратный тип указанного геттера.
Ошибка времени компиляции возникает, если  не является 
 У вас нет оператора {+} (когда \op{} \lit{++})
или оператор {-} (когда \op{} является \lit{-{}-}),
или если тип возврата этого оператора не присваиваетсяТип аргументации сеттера.
Ошибка времени компиляции возникает, если  не может быть присвоено
Тип параметра указанного оператора.
Статический тип  является .

<%
Оценка выражения postfix 
Форма  соответственно . -{}-}
идет следующим образом:
Оценить  на объект $u$  быть свежей переменной, связанной с .
Оценить  на объект 
 быть свежей переменной, связанной с .
Оценка  соответственно .
Далее  Показатель .


%
 \code{[-{}-}}
Рассмотрим выражение postfix  формы 
где  {++} или {-{}-}.
 Статический тип 
 Статический тип .
Ошибка времени компиляции возникает, если  иметь
оператор \lit{[]} и оператор \lit{[]=}.
Пусть  будет возвратным типом первого.
Ошибка времени компиляции возникает, если только  может быть присвоено
Первый тип параметра указанного оператора {[]=}.
Ошибка времени компиляции возникает, если  не является 
 У вас нет оператора \lit{+} (когда \op{} является \lit{++})
или оператор {-} (когда \op{} является {-{}-}),
или если тип возврата этого оператора не присваивается
Второй тип аргумента указанного оператора {[]=}.
% Разрешается «e1[e2]++»; также, когда запись имеет двойной статический тип,
%, поэтому мы не можем просто сказать «назначаемый».
Ошибка времени компиляции возникает при прохождении целого числа 
как аргумент в пользу оператора \lit{+} или {-} Это было бы ошибкой.
Статический тип  является .

%
Оценка выражения postfix 
Форма  соответственно  -{}-}
идет следующим образом:
Оценить  на объект $u$ и  Объект .
 и  быть свежими переменными, связанными с  и  соответственно.
Оценить  на объект 
 быть свежей переменной, связанной с .
Оценка  соответственно .
Далее  Показатель .


%
 ? -{}-}}
Рассмотрим выражение postfix  формы 
где  {++} или {-{}-}.
Ошибки времени компиляции, которые могут быть вызваны

Они также генерируются в случае .
Статический тип  Статический тип .

%
Оценка выражения postfix 
Форма  соответственно ? -{}-}
идет следующим образом:
Если  Буквальный тип, оценка  является эквивалентным
Оценка  соответственно . -{}-}.
В противном случае оценка  к объекту .
Если  является нулевым объектом,  оценивает нулевой объект ().
В противном случае пусть  быть свежей переменной, связанной с .
Оценка  соответственно . -{}-} к объекту .
Далее  Показатель .



 Соответствующие выражения


%
 Назначаемые выражения - это терминыОн может появиться на левой стороне задания.
В этом разделе описывается, как оценивать термины этих терминов, когда это необходимо.
Семантика назначения в целом описывается в другом месте.
().

{%
Грамматика присваиваемых выражений включает очень общие формы.
как идентификатор  Квалифицированный идентификатор .
Следовательно, присваиваемое выражение может иметь множество различных значений,
В зависимости от связывания этих идентификаторов в контексте.

Например, термин  является присваиваемым выражением.
 может относиться к геттеру  на
объект, на который ссылается переменная ,
В этом случае  будет оцениваться на предмет
Прежде чем перейти к  члена этого объекта.
Термин «P001568» может также обозначать класс  ссылаться через
Префикс импорта ,
 Статическая переменная этого класса.
В данном случае  не будет оцениваться в стоимостном выражении.%
?


<assignableExpression> ::= <primary> <assignableSelectorPart
 \SUPER{} <безусловный Подписываемый Выборщик>
 <идентификатор>

<assignableSelectorPart> ::= <selector>* <assignableSelector>

<unconditionalAssignableSelector> ::= ‘[<выражение>] "
 <идентификатор>

<assignableSelector> ::= <unconditionalAssignableSelector>
< <определитель>>
 >> '?' <выражение> '' "



%
Раздел о назначениях
()
определяет статический анализ и динамическую семантику
различные формы присвоения.
Каждый из этих случаев применим, когда указанные термины удовлетворяют
Заданные боковые условия
{%
(например, один случай формы  
быть выражением, тогда как 
требует, чтобы  был буквальным типом)%
}
Случаи, требующие, чтобы термины были выражениями, считаются наименее конкретными.
Они используются только в том случае, если нет других совпадений.
{%
(содержащий ) используется
если соответствующий термин является буквальным)%
}

{%
Синтаксически эти выражения не являются производными от {выражение},
но они могут быть получены из {выражение}, например, потому что
<< Выражение может содержать {первичное}
Выводится из {выражение}.
Мы используем следующее правило, чтобы найти такие выражения:
?

%
%
Пусть  будет {назначаемое выражение}.
Предположим, что  Это такой термин, как  могут быть получены из
\syntax{$t$ <assignableSelector>}.
В данном случае мы говорим, что  Это
<{приемный термин}{назначаемое выражение! термин "получатель"
.
Когда  Это выражение, мы говорим, что  Это
{выражение приемника}{назначаемое выражение! выражение приемника
.

{%
Короче говоря, мы получаем  путем отсечения назначенного селектора
с конца .
Легко заметить, что только некоторые 
иметь срок приема.
Например, простой  нет.%
?

%
Пусть  будет присваиваемым выражением.
Предположим, что  имеет выражение приемника .
Оценка  происходит таким же образом, как и
Оценка любого другого выражения.


<{Лексический поиск}


%
В этом разделе указано, как искать имя
основанные на замкнутых лексических рамках.
Это называется P001298.
Когда  является идентификатором,
Вы можете найти имя  Форма  Как и форма .

%
Лексический поиск может дать декларацию или префикс импорта.
И он ничего не может дать.{%
Это не ошибка времени компиляции, когда лексический поиск ничего не дает.
В этой ситуации данное имя  преобразуется в
,
Статический анализ полученного выражения
может иметь или не иметь ошибок времени компиляции.%
?

%
Лексический поиск может привести к ошибке времени компиляции.
как указано ниже.
Однако это отличается от результата, полученного в результате лексического поиска.
Потому что эта спецификация никогда не указывает на распространение ошибок.

<{%
Другими словами, когда другие части этой спецификации указывают, что
проводится лексический осмотр,
Они должны рассмотреть дальнейшие шаги, предпринятые
Когда всплывающее окно дает декларацию,
Когда он дает префикс импорта,
И когда он ничего не дает.
Но они не упоминают об этом, например,  Это ошибка
потому что лексический поиск для  Погрешность.%
?

{%
Лексический поиск отличается от поиска члена экземпляра
()
потому что эта операция ищет через последовательность суперклассов,
В то время как лексический поиск ищет через последовательность замкнутых областей.
Лексический поиск отличается от простого поиска в закрытых областях
потому что лексический поиск «связывает» геттеров и сеттеров,
Подробнее ниже.%
?

%
%
Рассмотрим ситуацию, когда имя 
Базовое имя 
()
где  является идентификатором,
и лексический поиск  выполняется из определенного места
\DefineSymbol{\ell}.

{%
Мы указываем имя и место, откуда выполняется лексический поиск.
Месторасположение не всегда избыточное:
В некоторых ситуациях мы выполняем поиск сеттера по имени ,
Но токен  В программе этого не происходит.
Чтобы справиться с такими ситуациями, мы должны указать оба.
имя, которое высматривается,
и местоположение, которое определяет, какие области являются закрытыми.%
?

%
Когда мы говорим, что лексический поиск идентификатора  выполняется,
Понятно, что поиск осуществляется из
место данного происшествия .

%
 быть самой внутренней лексической областью, содержащей 
который имеет декларацию с именем базы .
В случае, когда  иметь
Декларация под названием  а также заявление под названием ,
давать  D) быть декларацией под названием .
В ситуации, когда  иметь
Одно заявление с именем basename ,
Пусть  будет такой декларацией.

{%
Нелокальная вариабельная декларация с именем  будет косвенно стимулировать
Геттер  и, возможно, сеттер  в текущем объеме.
Это означает, что  может обозначать неявно индуцированный геттер или сеттер
Вместо базовой переменной декларации.
Это важно в том случае, когда возникает ошибка.
потому что поиск был для одного вида, но только другой вид существует.
?

{%
Если мы ищем имя  с базовым именем ,
Мы прекращаем поиск, если находим какое-либо заявление под названием  или .
Если в этой области имеются заявления как для , так и для ,
Мы возвращаем тот, который имеет запрашиваемое имя .
В случае, если присутствует только одно заявление, мы возвращаем его.
Даже если он может иметь имя , когда  является , или наоборот.
Эта ситуация может привести к ошибке, как указано ниже.
?

%
На первом этапе мы проверяем несколько возможных ошибок.

% Даже если мы нашли заявление, оно может иметь неправильное название.
%
 существует
В этом случае, по крайней мере, одно заявление с базовым именем 
находится в зоне действия по адресу .Это ошибка времени компиляции, если имя  не является ,
Если только  является членом экземпляра или локальной переменной
{(который может быть формальным параметром)}.

{%
То есть, это ошибка, если мы ищем сеттера и находим геттера,
или наоборот,
Но не ошибка, если мы ищем сеттера и находим локальную переменную.
Если мы ищем сеттера и находим экземпляр геттера, или наоборот,
Это не ошибка,
Потому что сеттер может быть унаследован.
Это проверяется после того, как ничего не дает.
(что означает, что {} будет подготовлено.%)
?

%
Если  является членом инстанции,
Ошибка времени компиляции  Не имеет доступа к .


<%
 не существует
Ошибка времени компиляции  не имеет доступа к 
().


{%
Мы всегда смотрим вверх   и ,
Независимо от того, является ли   или .
Такой подход создает более тесную связь между парой деклараций.
где один является геттером по имени 
А другой — сеттер по имени .
Это позволяет разработчикам думать о
геттер и сеттер, которые объявляются вместе как единое целое,
Вместо двух независимых деклараций.
?

{%
Например, если термин относится к  Нужен сеттер,
и самой внутренней декларации, названной  или  является геттером 
и нет соответствующего сеттера,
Это ошибка времени компиляции.
Эта ошибка возникает даже в том случае, когда более удаленная область охвата имеет
Декларация сеттера  
потому что мы уже взяли на себя обязательство использовать 
(так что это на самом деле «пара сеттера / геттера, где сеттер отсутствует»)
Можно сказать, что эта «парная» тень :%
?



%
На втором и последнем шаге, если ошибки не произошло,
действовать, как описано в первом применимом случае из следующего перечня:



Когда  не существует,
Лексический поиск ничего не дает.
{%
В этом случае гарантируется, что  имеет доступ к ,
 будет рассматриваться как .
Ошибки могут возникать, например,
потому что нет члена интерфейса  ,
а также нет доступного и применимого метода расширения.%
?

Рассмотрим случай, когда  Формальное заявление о параметрах типа
класса или смеси.
Ошибка времени компиляции  происходит внутри
статический метод, статический геттер или статический сеттер,
внутри статического переменного инициализатора.
% NB: существует ошибка _no_, когда она возникает в инициализаторе переменной экземпляра;
%, что означает, что в данном конкретном случае разрешается «доступ к этому».
% Это может показаться непоследовательным, но это не должно быть вредным.
В противном случае лексический поиск дает .

Рассмотрим случай, когда  является заявлением государства-члена в
Класс или смесь .
Лексический поиск ничего не дает.
{%
В этом случае гарантируется, что  имеет доступ к .%
?

% Случаи, рассмотренные здесь:  это
% - импорт с префиксом ;
% - псевдоним класса или типа; %, чье имя должно быть , не нужно говорить, что.
% — библиотечная переменная, геттер или сеттер;
% - статический метод/настройщик/настройщик класса, смеси, перечня или расширения;
% - примерный способ расширения;
% — локальная переменная (которая может быть формальным параметром);
% — локальная функция.
В противном случае лексический поиск дает .


{%Обратите внимание, что лексический поиск никогда не даст объявление о том, что
Член инстанции.
В каждом случае, когда установлено, что ошибки нет,
и один из членов экземпляра является результатом поиска,
Лексический поиск ничего не дает.

Причина этого заключается в том, что не может быть декларации в объеме,
но интерфейс класса может иметь требуемый член,
Возможно применение метода расширения.
Для единообразия,
В таких ситуациях мы всегда сообщаем, что ничего не нашли.
Это означает, что  Должен быть добавлен.%
?


 Ссылка на идентификатор


%
 состоит из одного идентификатора;
Он обеспечивает доступ к объекту через неквалифицированное имя.
{typeIdentifier} - идентификатор, который может быть использован
как название вида декларации.

{%
A {qualifiedName} не является выражением идентификатора,
Но мы определяем его синтаксис здесь, потому что он используется в нескольких различных контекстах.
и более тесно связана с простым идентификатором
чем для любого из тех правил грамматики, где он используется.
?


<identifier> ::= <Идентификатор>
 <BUILT В  Идентификатор
 <Другой Идентификатор

<typeIdentifierNotType> ::= <IDENTIFIER>
 <Другой Идентификатор   Тип
 \DYNAMIC{}

<typeIdentifier> ::= <typeIdentifierNotType>
 \TYPE{}

<qualifiedName> ::= <typeIdentifier> <identifier>
 <typeIdentifier> <typeIdentifier> <identifier>

<BUILT В  Идентификатор: := \ABSTRACT{} | {} | {} | {}
\alt\hspace{-3mm}{} |{} |{} |{} !
\FACTORY{} |{} |{}
\alt\hspace{-3mm}{} |{} |{} |{} !
\LIBRARY{} |{} |{}
\alt\hspace{-3mm} \PART{} | \REQUIRED{} | \SET{} | \STATIC{} | \TYPEDEF{}

<Другой Идентификатор  Не  Тип ::= \gnewline{}
\ASYNC{} | {} | {} | {} | {} | {} !
\AWAIT{} | \YIELD{}

<Другой Идентификатор: := \gnewline{}
<Другой Идентификатор   TYPE> | \TYPE{}

<IDENTFIER NO  ДОЛЛАР: := <IDENTFIER Начало No ДОЛЛАР
\gnewline{} <IDENTFIER Часть  No Доллар*

<IDENTFIER Начало  No DOLLAR > ::= <LETTER> "

<IDENTFIER Часть  No ДОЛЛАР: := \gnewline{}
<IDENTFIER Начало  No ДОЛЛАР > | <ПРИЛОЖЕНИЕ

<IDENTIFIER> ::= <IDENTFIER Начало> <IDENTIFIER Часть*

<IDENTFIER Начало>::= <IDENTFIER Начало  NO ДОЛЛАР "

<IDENTFIER Часть ::= <IDENTFIER Начало>> <Цифры>

<LETTER> ::= 'a' .. 'z' | 'A' .. 'Z'

<DIGIT> ::= 0.. "9"

<WHITESPACE> ::= ( | <LINE BREAK > +


%
Упорядочение лексических правил выше гарантирует, что 
 Идентификатор  NO DOLLAR не выводят никаких встроенных идентификаторов.
Аналогично, лексическое правило для зарезервированных слов (\ref{reservedWords})
должно считаться предшествующим правилу для  BUILT В  Идентификатор.
такой, что \synt{IDENTIFIER}} и \synt{ Идентификатор  No ДОЛЛАР
Также не следует выводить какие-либо зарезервированные слова.

%
 является одним изидентификаторы, полученные в результате производства  BUILT В  Идентификатор.

{%
Обратите внимание, что это синтаксическая ошибка, если встроенный идентификатор
используется как заявленное название
% % TODO(Eernst): Приходите с расширениями, добавьте их.
префикс, класс, микширование, число, параметр типа, псевдоним типа или расширение.
Аналогично, это синтаксическая ошибка для использования встроенного идентификатора.
кроме других \DYNAMIC{} или {}
как идентификатор в аннотации типа или связанный параметр типа.%
?

{%
Встроенные идентификаторы — это идентификаторы, которые используются в качестве ключевых слов в Dart.
Но это не зарезервированные слова.
Встроенный идентификатор не может использоваться для обозначения класса или типа.
Другими словами, они рассматриваются как зарезервированные слова при использовании в качестве типов.
Это устраняет многие запутанные ситуации.
как для читателей, так и во время парсинга.%
?

%
Ошибка времени компиляции, если один из идентификаторов \AWAIT{} или {}
используется в качестве {идентификатора} в функциональном теле
отмеченный либо ,  или .

<<<<<<<<<<<<<<<<<<<<<<<<<<<<<<<<<<<<<<<<<<<<<<<<<<<<<<<<<<<<<<<<<<<<<<<<<<<<<<<<<<<<<<<<<<<<<<<<<<<<<<<<<<<<<<<<<<<<<<<<<<<<<<<<<<<<<<<<<<<<<<<<<<<<<<<<<<<<<<<<<<<<<<<<<<<<<<<<<<<<<<<<<<<<<<<<<<<<<<<<<<<<<<<<<<<<<<<<<<<<<<<<<<<<<<<<<<<<<<<<<<<<<<<<<<<<<<<<
Это делает идентификаторы \AWAIT{} и \YIELD{} вести себя как зарезервированные слова
в ограниченном контексте.
Этот подход был выбран потому, что он был менее разрушительным, чем это было бы.
делать \AWAIT{} и \YIELD{} зарезервированные слова или встроенные идентификаторы;
В то время, когда эти функции были добавлены к языку.
?

%
 Это два или три идентификатора, разделенные {.}.
Все, кроме последнего, должно быть {typeIdentifier}.
Он используется для обозначения декларации, которая импортируется с префиксом.
или заявление {} в виде класса, смеси, перечня или расширения, или обоих.

%
Статический тип выражения идентификатора  который является идентификатором 
определяется следующим образом.
Выполните лексический поиск 
(перенаправлено с «P001269»)
Местонахождение .

%
 Лексический поиск дает декларацию.
 быть декларацией, полученной в результате лексического поиска .



Если  объявляет класс, миксинг, энум, псевдоним типа, перечисленный тип,
или параметра типа,
Статический тип  является .

Если  Декларация библиотечного геттера
{(которые могут быть неявно вызваны библиотечной переменной)},
Статический тип  является статичным типом
библиотечный геттер вызов 
().

Если  Статический метод, библиотечная функция или локальная функция.
Статический тип  Тип функции .

{%
Обратите внимание, что  может быть впоследствии подвергнут
обобщение функции
()%
?

Если  Объявление статического геттера
{(которые могут быть индуцированы статической переменной)}
 встречается в классе ,
Статический тип  Тип возврата Геттера
.

Если  локальная переменная декларация
{(который может быть формальным параметром)}
Статический тип  тип переменной  Заявлено ,
Если только  Известно, что имеет некоторый тип ,
где  является подтипом любого другого типа 
Таким образом,  Известно, что имеет тип ,
В этом случае статический тип  является .

% Продление геттера вызова по объему.
Если  Декларация экземпляра Getter
в заявлении о продлении ,
Статический тип  Тип возврата .

% Отрыв метода расширения по объему.Если  является декларацией метода экземпляра
в декларации о продлении ,
Статический тип  Тип функции .

{%
Лексический взгляд никогда не даст декларации.
который является членом класса.%
?

\vspace{-1ex)


%
 Лексический поиск дает префикс импорта
В этом случае лексический поиск
()
 Дает префикс импорта .
Префикс никогда не может быть использован в качестве отдельного выражения.
В этом случае возникает ошибка времени компиляции,
Если токен сразу после  {.}.
Статический тип не связан с  в данном случае.

% Тип указывается в , поэтому было бы последовательно сообщать
%. Однако из этого типа не будет использоваться, поскольку статический тип
% от всех построений формы .Что-то указано без ссылки
% к статическому типу . Поэтому мы не упоминаем этот тип здесь.
{%
Такой тип не нужен, потому что каждая конструкция
Префикс импорта  используется и сопровождается {.} указывается
таким образом, что тип  Не используется.%
?


%
 Лексический поиск ничего не дает
Когда лексический взгляд
()
 ничего не дает,
 рассматривается как
()
.

{%
В данном случае известно, что  Имеет доступ к 
().
Статический анализ и оценка продолжаются
,
поэтому нет необходимости дополнительно указывать лечение .%
?


%
Оценка выражения идентификатора  формы 
идет следующим образом:

%
 Лексический поиск дает декларацию.
В этом случае лексический поиск
(перенаправлено с «P001277»)
 выдает декларацию .
Оценка  идет следующим образом:



Если  является классом, микшированием, enum или типом alias,
Значение  является объектом, реализующим класс 
который соответствует соответствующему типу.

Если  Типовой параметр  Тогда значение  это
значение фактического типа аргумента, соответствующего 
Это было передано созидателю, который создал
Текущее связывание .

Если  Декларация библиотечного геттера
{(которые могут быть неявно вызваны библиотечной переменной)},
Оценка  равносильно оценке призыва к
Библиотечный геттер 
().

% Мы могли бы сказать, что «выталкивание вызвано постоянной переменной»; но всегда будет
% - переменная, потому что декларация Геттера не может дать константу, поэтому мы
% используйте ярлык и не упоминайте геттер полностью здесь.
Если  библиотека, класс или локальная постоянная переменная одной из форм
\code{{) $v$ = $e'$;} или \code{\CONST{}   = ;
Тогда значение  значение постоянного выражения .

Если  является декларацией
функция верхнего уровня, статический метод или локальная функция;
 оценивает функцию объекта, полученную путем клозулизации
()
.

Если  является заявлением на получение экземпляра в заявлении о продлении ,
 оценивает результат вызова указанного геттера
с текущим связыванием , и
текущие связывания параметров типа, объявленные по .
Если <<p001106> является заявлением метода экземпляра в заявлении о расширении 
с параметрами типа {X}{1}{}
<<<p001108> оценивает результат клозуризации метода расширения
\code{\List{X}{1}{s}>(). 
().

Если <<p001110> локальная переменная 
{(который может быть формальным параметром)}
 оценивает текущее связывание .


{%
Обратите внимание, что  не может быть декларацией
статическая переменная, статический геттер или статический сеттер, объявленный в классе ,
Потому что в этом случае  рассматривается как
()
Добыча имущества
()
,
который также определяет оценку .%
?


%
 Лексический поиск дает префикс импорта
Такая ситуация не может возникнуть,
Потому что это ошибка времени компиляции
для оценки префикса импорта как выражения,
и никаких конструкций с импортным префиксом
\commentary{(например, извлечение свойств )}
Оцените префикс импорта.


%
 Лексический поиск ничего не дает
Такая ситуация не может возникнуть,
Это происходит только тогда, когда  рассматривается как
,
чья оценка указана в другом месте
().



 Тип испытания


%
 Тестирование, если объект является членом типа.


<typeTest> ::= <isOperator> <typeNotVoid>

<isperator> ::= {}!


%
Оценка выражения  <  идет следующим образом:

%
Выражение  оценивается на объект .
Если динамический тип  является подтипом ,
Показатель выражается в .
В противном случае он оценивает до .

{%
Из этого следует, что {$e$ \IS{} Объект всегда верен.
Это имеет смысл в языке, где все является объектом.

Также обратите внимание, что \code{\NULL{} \IS{}  ложный
Если только   или .
Первые два бесполезны, как и все остальное.
форма \code{$e$ \IS{} Объект или  \IS{} .
Пользователи должны проверить нулевой объект () непосредственно
а не с помощью типовых тестов.%
?

%
Выражение  {}!  является эквивалентным
<! () \IS{} .

%
 Быть локальной переменной
{(который может быть формальным параметром)}.
Выражение формы \code{$v$ \IS{} 
Показывает, что  имеет тип 
Если  является подтипом типа выражения .
%
В противном случае
если заявленный тип  переменная типа ,
 является подтипом границы ,
 является подтипом типа выражения ,
 Показывает, что  Имеет тип .
%
В противном случае  Не показывает, что  имеет тип  для любого .

{%
Мотивацией для «показов, что v имеет отношение типа T» является
Чтобы уменьшить ложные ошибки, тем самым обеспечивая более естественный стиль кодирования.
Правила в текущей спецификации намеренно держатся простыми.
В будущем было бы целесообразно усовершенствовать эти правила;
Такое усовершенствование позволило бы принимать больше кода без ошибок.
Но не отказывайтесь от кода без ошибок.

Правило касается только местных жителей и параметров.Нелокальные переменные могут быть изменены
с помощью функций или методов, которые не доступны
для локального анализа.

Бессмысленно выводить более слабый тип, чем тот, который уже известен.
Кроме того, это приведет к ситуации, когда несколько типов
с переменной в заданной точке,
Это усложняет спецификацию.
Отсюда требование, что продвигаемый тип является подтипом текущего типа.

В любом случае, это не ошибка, когда тип теста не показывает
что данная переменная не имеет «лучшего» типа, чем ранее известный,
Но инструменты могут дать подсказку в таких случаях.
если подходящие эвристики указывают, что продвижение, вероятно, будет предназначено.%
?

%
Статический тип является выражением .


 Тип Каста


%
 Это гарантирует, что объект является членом типа.


<typeCast> ::= <asOperator> <typeNotVoid>

<asOperator> ::= \AS{}


%
Оценка литого выражения \code{$e$ \AS{}  идет следующим образом:

%
Выражение  Оценивается на объект .
% Эта ошибка может произойти по замыслу «как».
Ошибка динамического типа, если  не является нулевым объектом (),
и динамический тип  не является подтипом .
В противном случае  Показатель .

%
Статический тип литого выражения \code{$e$ \AS{}  .


 Заявления


%
 Это фрагмент кода Dart, который может быть выполнен во время выполнения.
Заявления, в отличие от выражений, не оцениваются на предмет,
Но вместо этого они выполняются для их влияния на состояние программы и поток управления.


<statements> ::= <statement> *

<statement> ::= <label>* <nonLabelled Statement>

<nonLabelled Statement> ::= <блокировать>
 <localVariableDeclaration>
 <for Statement>
 <в то время как заявление>
 <doStatement>
 >> <переключатель
 <если заявление>
 <rethrow Statement>
 <try Statement>
 <break Statement>
 <продолжение заявления>
 <возврат заявления>
 <yieldStatement>
 <yieldEachStatement>
 <выражение>
 <assert Statement>
 <localFunctionDeclaration>



 Заявление о завершении


%
Исполнение заявления \IndexCustom{завершение}{завершение}
одним из пяти способов:
либо
{завершается нормально}{завершается нормально},
<{breaks}{завершение!breaks}
Или это <\IndexCustom{продолжение}{завершение!продолжение}
(либо на этикетке, либо без этикетки),
{возвращается}{завершение!возвращается} (с объектом или без него),
Или это {броски}{завершение!броски}
Исключительный объект и связанный с ним след стека.

%
В описаниях выполнения заявления по умолчанию является то, что
исполнение обычно завершается, если не указано иное.

%
Если исполнение заявления, ,
определяется с точки зрения выполнения другого заявления,
и исполнение этого другого заявления обычно не завершается,
если не указано иное, исполнение  остановки
В этот момент и завершается одинаково.
{%
Например, если исполнение тела a  петля возвращает объект,
Так же выполняется выполнение {} петлевое высказывание.%
?

%
Если исполнение заявления определяется в терминах оценки выражения
и оценка этого выражения бросает,
Если не указано иное, исполнение заявления прекращается.
в этот момент и бросает тот же самый объект исключения и след стека.
{%
Например,
при оценке выражения состояния {} заявление бросает,
Затем выполняется утверждение {}.
Аналогично, если оценка выражения  заявление бросает,
Так же выполняется выполнение {} заявление.%
?

% % TODO(Eernst): Использовать предложение Лассе для добавления определения понятия
% — утверждение, для которого «статически известно, что оно не будет завершено».
%% обычно': Cf. https://github.com/dart-lang/language/issues/139.
% % Это определение, вероятно, должно быть здесь.


 Блоки.


%
 Поддерживает секвенирование кода.

%
Выполнение заявления о блокировке  идет следующим образом:

%
Для  выполняется.

%
Блок-заявление вводит новый объем,
сфера охвата которого является текущим объемом заявления о блоке.


 Высказывания


%
 Состоит из выражения, которое не
Начните с \lit{\{} характер.


<выражение> ::= <выражение>? "


%
Выражение высказывания не допускается
Начнем с {\{}.
{%
Это означает, что если исходный текст может быть
распаривается как выражение, за которым следует {;},
Тогда это производство грамматики не применяется.
когда выражение начинается с \lit{\{}.%
?
{%
Ограничение устраняет двусмысленность при разборе, где
 Можно запустить любой блок () или
Буквальная карта ().
Запрещая последнему начинать выражение,
Парсеру не нужно смотреть дальше
перед тем, как принять решение о том, что он анализирует утверждение блока.%
?

%
Исполнение заявления о выражении  Для этого необходимо оценить .
Если выражение оценивает объект,
Затем объект игнорируется, и выполнение выполняется нормально.


 Местная переменная декларация


%
A ,
Также известен как P000836.
имеет следующую форму:


<localVariableDeclaration> ::= <metadata> <initialized Переменная декларация>; "


%
Каждая локальная переменная вносит
а {местная переменная}{вариабельная!локальная}
в текущем объеме.

{%
Локальные переменные не вызывают геттеров и сеттеров.
Заметим, что официальное заявление о параметрах также вводит
локальная переменная в ассоциированной области формальных параметров
() %
?

%
Свойства быть
{initialized}{variable!initialized}
{constant}{variable!constant}
относится к локальным переменным с теми же определениями, что и для других переменных.
(P000721).

% % TODO(Eernst): Может потребоваться обновление / удаление при интегрированном анализе потока.
%
Мы говорим, что локальная переменная  является 
в некоторой сфере 
Если  не является окончательным, и назначение на  Происходит в .

%
Локальная переменная декларация формы \code{\VAR{} ; является эквивалентным
\code{{)  = ;}.
Если  является нулевым типом
()
затем локальное переменное объявление формы 
Это эквивалентно P000862.

% % TODO(Eernst): Пересмотр при добавлении анализа потока.
{%
Если  Потенциально неотменяемый типзатем локальная переменная декларация формы  допускается,
но выражение, которое дает оценку 
Ошибка времени компиляции, если анализ потока
()
показывает, что переменная гарантированно была инициализирована.
?

%
Локальная переменная имеет ассоциированную
{объявленный тип}{местная переменная! объявленный тип
который определяется его декларацией.
Локальная переменная также имеет ассоциированную
{тип}{местная переменная! тип
который определяется анализом потока
(перенаправлено с «P000724»)
Процесс, известный как продвижение типа
(P000725).

<<<p001010>%
Объявленный тип локальной переменной с объявлением одной из форм



это .

%
Объявленный тип локальной переменной с объявлением одной из форм



определяется следующим образом:



Если статический тип   затем
Заявленным типом  является .

Если статический тип  Форма: 
где  является переменной типа
(),
Заявленным типом  является .
{%
В данном случае  Он был немедленно повышен до 
()%
?

В противном случае, заявленный тип  Статический тип .


%
 быть локальной переменной, объявленной инициализацией переменной декларации;
Пусть  будет ассоциированным выражением инициализации.
Это а
{ошибка времени компиляции}, если статический тип  не может быть присвоено
Заявленный тип .

% % TODO(Eernst): Пересмотр при добавлении анализа потока.
{%
Если локальная переменная  является \FINAL{} и не ,
не является ошибкой компиляции, если объявление  это
не является инициализацией переменной декларации,
но выражение, которое дает оценку 
Ошибка времени компиляции, если анализ потока не показывает, что
Переменная гарантированно была инициализирована.
Аналогично, выражение, которое приводит к заданию 
Ошибка времени компиляции, если анализ потока не показывает, что
гарантируется, что переменная {not} была инициализирована.%

В любой ситуации, не охватываемой предыдущим пунктом,
Ошибка времени компиляции для присвоения локальной переменной
Что такое  и не 
() %
?

%
Предположим, что  локальная переменная с модификатором {}
что объявляет переменную ,
который имеет начальное выражение .
Это \Error{ошибка времени компиляции}, если $e$ Содержит {} выражение 
(),
Если нет функции <<<p000079> который является
Непосредственно прилагаемая функция для ,
 не является непосредственной функцией для .

{%
Другими словами,
выражение инициализации не может ждать выражения непосредственно,
любое ожидающее выражение должно быть синтаксически вложено в какую-либо другую функцию;
То есть функция буквальная.%
?

%
{ошибка времени компиляции}
локальная переменная упоминается в месте исходного кода, которое находится перед
окончание его инициализации выражения, если таковое имеется,
иным образом до объявления о наступлении
Идентификатор, который называет переменную.

<<%
Приведенный ниже пример иллюстрирует ожидаемое поведение.
На библиотечном уровне объявляется переменная ,
и еще один  объявлен внутри функции .%
?



{%
Декларация внутри «P000877» скрывает прилагающую.Все ссылки на «P000878»  "
Внутренняя декларация «P000880».
Однако многие из этих ссылок являются незаконными.
Потому что они появляются перед объявлением.
Одним из таких случаев является присвоение .
Задание на «\code{x}» в \IF{} заявление
страдает от многочисленных проблем.
На правой стороне перед декларацией написано «\code{x}»,
и левая сторона присваивает "" до его объявления.
Каждая из них является самостоятельной ошибкой времени компиляции.
Печатное заявление внутри P001213 также является незаконным.

Внутренняя декларация  сам является ошибочным
Правая сторона пытается
читать "" до прекращения действия декларации.
Возникновение , которое объявляет и называет переменную
(т.е. слева от «P000888» во внутренней декларации)
Это не ссылка, и это законно.
Последнее печатное заявление также совершенно законно.
?

{%
В качестве другого примера \code{\VAR{} x = 3, y = x;} является законным,
потому что  ссылается на его инициализатор.

Особенно извращенный пример включает
Локальное переменное имя, скрывающее тип.
Это возможно, потому что Дарт
Единое пространство имен для типов, функций и переменных.%
?



{%
Внутри , «P000891» означает локальную переменную.
Тип  скрыт переменной с тем же названием.
«Попытка представить»  вызывает ошибку времени компиляции
потому что она ссылается на локальную переменную до ее объявления.
Аналогично, для декларации .%
?

%
Выполнение заявления с переменной декларацией одной из форм
\code{\VAR{}  = ;

\code{{)  = ;
\code{\CONST{} <<<p000090>  = ;
\code{\FINAL{}  = ;
\code{{)   = ;
идет следующим образом:

%
Выражение  оценивается на предмет .
% Следующая ошибка может возникнуть из-за неявных слепков.
происходит динамическая ошибка
если динамический тип  не является подтипом фактического заявленного типа
()
.
В противном случае переменная  Связано с .

{%
Обратите внимание, что  Они могли быть изменены в результате неявного принуждения.
Например,  могут быть преобразованы в
 за счет инстанциации общей функции
(P000731).
Предполагается, что такие преобразования уже имели место.
в декларации выше.%
?

%
Пусть  будет \LATE{} и \FINAL{} Местные переменные
Объявляет переменную .
Если объект  Применяется для 
В ситуации, когда  несвязанный
 Связано с .
Если объект  Применяется для 
В ситуации, когда  Привязан к объекту 
Затем происходит динамическая ошибка.
\commentary{(не имеет значения, ) Это тот же объект, что и P000117.


 Декларация о местной функции Срок


%
Декларация функции объявляет новую локальную функцию.
(P000732).


<localFunctionDeclaration> ::= <metadata> <functionBody>


%
Декларация функции одной из форм
 {подпись}  <{заявления} 
или
  {подпись}  <{заявления} Это приводит к добавлению в текущий объем новой функции под названием .
Ошибка времени компиляции для ссылки на локальную функцию перед ее декларацией.

{%
Это означает, что локальные функции могут быть рекурсивными.
но не взаимно рекурсивными.
Рассмотрим следующие примеры: %
?



{%
Невозможно написать пару взаимно рекурсивных локальных функций.
Потому что один всегда должен прийти до того, как будет объявлен другой.
Эти случаи довольно редки,
И всегда можно управлять, сначала определив пару переменных.
затем присваивая им соответствующие функции букв:%
?



{%
Правила локальных функций немного отличаются от правил локальных переменных.
в том, что к функции можно получить доступ в рамках ее декларации
Но к переменной можно получить доступ только после ее объявления.
Это связано с тем, что рекурсивные функции полезны.
Рекурсивно определенные переменные почти всегда являются ошибками.
Поэтому имеет смысл согласовать правила для местных функций.
Для функций в целом, а не
с правилами для локальных переменных.%
?


 Если


%
 допускает условное исполнение заявлений.


<if Statement> ::= \IF{} '(' <выражение> ')' <statement> ({} <statement>)?


%
Если заявление формы
\code{{) ()  < 
где  Не является ли блок-заявление эквивалентным заявлению
\code{\IF{) ()  \ELSE{} .
Если заявление формы
\code{{} ()  < 
где  Не является ли блок-заявление эквивалентным заявлению
\code{\IF{} ()  < \{$s_2$\}}.

{%
Причиной такой эквивалентности является улавливание ошибок, таких как %.
?



{%
В разумных пределах нормы такого кодекса являются проблематичными.
Если предположить, что  объявляется
в рамках метода ,
когда  является ложным,
 будет неинициализирована при доступе.
Самым чистым подходом было бы требование блока после испытания.
Вместо произвольного заявления.
Это противоречит давнему обычаю,
подрывает цель Дарта - знакомство.
Вместо этого мы выбираем вставить блок, вводя область,
вокруг заявления, следующего за предикатом
(и аналогично для {} и петель).
Это приведет к ошибке времени компиляции в приведенном выше случае.
Конечно, если есть декларация  в окружающем пространстве,
Программисты могут быть удивлены.
Мы ожидаем, что инструменты будут освещать случаи тени.
Чтобы избежать подобных ситуаций.%
?

%
Исполнение заявления формы if
\code{{} ()  < 
где  и  Блоковые заявления,
идет следующим образом:

%
Во-первых, выражение  Оценивается на объект .
% Эта ошибка может произойти из-за неявных отливок и нулевых.
Это динамическая ошибка, если тип времени выполнения  Это не P000903.
Если  , то утверждение блока  казнен,
В противном случае блокировка  выполнено.

%
Ошибка времени компиляции, если тип выражения 
Не может быть присвоено .

%
Если  показывает локальную переменную  имеет тип ,
Тип  Известно, что  в ,
Если только одно из следующих не является истинным

  потенциально мутирует в ,
  потенциально мутирует в пределах функции
где  объявляется, или
  доступ к функции, определенной в  и
 Потенциально мутирует в любом месте в области .


%
Если заявление формы \code{{) ()  является эквивалентным
Если заявление \code{{) ()  \ELSE{} .


 Для


%
 Поддерживает итерацию.


<for Statement> ::=  \FOR{} '(' <forLoopParts> ')

<forLoopParts> ::= <forInitializerStatement> <expression>? <expressionList>?
 <forInLoopPrefix> < <выражение>

<forInLoopPrefix> ::= <metadata> <declaredIdentifier
 <идентификатор>

<forInitializerЗаявление> ::= <localVariableDeclaration>
 <выражение>? "


%
Заявление имеет три формы - традиционная для петли
и две формы вводного утверждения — синхронная и асинхронная.


 Для петли.


%
Исполнение заявления о форме
\code{{} (\VAR{})  = ; ;   идет следующим образом:

%
Если  Пусть  будет \TRUE В противном случае пусть .

%
Сначала заявление с переменной декларацией \VAR{}  выполнено.
Тогда:




Если это первая итерация цикла, пусть  будет .
В противном случае пусть  Переменная  Созданный в
предыдущее выполнение шага .

Выражение  оценивается на предмет .
% Эта ошибка может произойти из-за неявных отливок и нулевых.
Это динамическая ошибка, если тип времени выполнения  Это не P000905.
Если  , цикл завершается нормально.
В противном случае исполнение продолжается на этапе .


Заявление  выполнено.

Если это исполнение завершается нормально, продолжается без ярлыка,
или продолжает маркироваться ()
Для этого используется приставка {} Заявление (),
Затем исполнение заявления рассматривается так, как если бы оно было выполнено нормально.


 Будьте свежей переменной.
 Связано со значением .

Выражение  оценивается, и
Процесс повторяется на этапе .


{%
Вышеприведенное определение предназначено для предотвращения общей ошибки.
где пользователи создают объект функции внутри цикла,
намереваясь закрыть поверх текущего связывания петлевой переменной, и найти
(обычно после болезненного процесса отладки и обучения)
Все созданные объекты функции захватили
то же значение - один ток в последней выполненной итерации.

Вместо этого каждая итерация имеет свою собственную переменную.
Первая итерация использует переменную, созданную начальной декларацией.
Выражение, выполненное в конце каждой итерации, использует
свежая переменная ,
привязанный к значению текущей переменной итерации,
и затем изменяет  Для следующей итерации.%
?

%
Это ошибка времени компиляции, если статический тип 
Не может быть присвоено .

%A для указания формы \code{\FOR{} (; ;  
% эквивалентно следующему коду:

% {%\{$d$;
%\WHILE{} () 
% 
% ;
%
%

% Если  является пустым, он интерпретируется как .


 Для входа


%
Пусть  Получается из {<finalConstVarOrType>?}.
Заявление формы \code{\FOR{} ()  \IN{} $e$\,\,$S$]
Затем он рассматривается как следующий код.
где  и  Новые идентификаторы:



%
Если статический тип  является высшим типом
()
 
иначе  Статический тип .
Ошибка времени компиляции  не может быть присвоено
.

{%
Отсюда следует, что это ошибка времени компиляции.
% Следующая ошибка существует также в случае, когда  Определенно
% без назначения перед циклом: Цикл может работать >1 раз.
Если  Пусто и  является конечной переменной.
Кроме того, это динамическая ошибка, если  имеет тип ,
 Оценка для экземпляра типа
не является подтипом .%
?


 Асинхронный For-in


%
Заявление для входа может быть асинхронным.
Асинхронная форма предназначена для итерации по потокам.
Асинхронный для петли отличается
ключевое слово \AWAIT{} непосредственно перед ключевым словом .

%
 Получается из {<finalConstVarOrType>?}.
Выполнение заочного заявления, , форма
\code{{) \FOR{} ()  \IN{}  
идет следующим образом:

%
Выражение  оценивается на предмет .
% Эта ошибка может произойти из-за неявных отливок и нулевых.
Ошибка динамического типа, если  не является примером
Класс, который реализует .
Ошибка времени компиляции  пустой
 является конечной или постоянной переменной.

%
Поток, связанный с самой внутренней замкнутой асинхронной для петли,
Если есть, то приостанавливается.
Поток  Прослушивается, производя потоковую подписку ,
и выполнение асинхронного цикла for-in приостанавливается
До тех пор, пока не появится трансляция.
{%
Это позволяет выполнять другие асинхронные события.
Пока этот цикл ждет потоковых событий.%
?

Приостановка асинхронного для цикла означает приостановку
Подписка на соответствующий поток.
Потоковая подписка приостанавливается, называя ее  метод.
Если подписка уже приостановлена, реализация может пропустить
Далее следует запрос .

{%
 Звонок может бросить, хотя этого никогда не должно произойти.
для правильно реализованного потока.%
?

%
Для каждого  из ,
заявление  Выполняется с помощью  связанный
значение текущего события данных.

{%
Либо исполнение  является полностью синхронным или содержит
асинхронная конструкция (, , \YIELD, или 
который приостановит подписку потока на окружающий его асинхронный цикл.
Это гарантирует, что ни одно другое событие  происходит
перед исполнением  Полный,
Если  Это правильно реализованный поток.
Если  не действует как поток,
Например, не соблюдая требования паузы,
Поведение асинхронного цикла может стать непредсказуемым.
?

%
В случае исполнения  Продолжается без этикетки,
или на этикетке ()
префиксирует асинхронное утверждение (),Затем выполнение  рассматривается как завершенное в обычном режиме.

%
Если исполнение  в противном случае не завершается нормально,
подписка  Отменяется путем оценки \code{\AWAIT{} $v$.cancel()
где  Это новая переменная, ссылающаяся на подписку на поток .
Если эта оценка бросает,
Выполнение  Отбрасывает одно и то же исключение и стекает след.
В противном случае выполнение  Завершается так же, как и исполнение .
% Уведомление: В предыдущей спецификации было неясно, что произошло, когда
% отменяется подписка. Данный документ является открытым, и существует
% реализации не могут должным образом ждать отмены вызова.
В противном случае выполнение  вновь приостановлено, ожидая
Следующее мероприятие по подписке,
 возобновляется, если он был приостановлен.
{%
 Звонок может бросать, в этом случае асинхронный для
Петля тоже бросает.
Это никогда не должно произойти для правильно реализованного потока.
?

%
На  от ,
Объект ошибки  и след стека ,
подписка  Отменяется путем оценки \code{\AWAIT{} v.cancel()}
где  Это новая переменная, ссылающаяся на подписку на поток .
Если эта оценка бросает,
Выполнение  бросает тот же самый объект исключения и след стека.
В противном случае выполнение  бросать
 в качестве объекта исключения и  в качестве следа стека.

%
Когда  выполняется, выполняется  обычно завершается.

%
Это ошибка времени компиляции, если появляется асинхронное заявление для входа.
внутри синхронной функции ().
Это ошибка времени компиляции, если она является традиционной для цикла (). это
Приставка по ключевому слову \AWAIT{}.

{%
Асинхронный цикл не имеет смысла в синхронной функции.
По тем же причинам, по которым ожидание выражения не имеет смысла.
в синхронной функции.%
?


 Пока


%
Хотя утверждение поддерживает условную итерацию,
где состояние оценивается перед циклом.


<while Statement> ::= \WHILE{} '(' <выражение> ''


%
Исполнение некоторого заявления формы \code{{) () ;
идет следующим образом:

%
Выражение  Оценивается на объект .
% Эта ошибка может произойти из-за неявных отливок и нулевых.
Это динамическая ошибка, если тип времени выполнения  Это не P000916.

%
Если  , то выполнение заявления в то время обычно завершается.
(P000743).

%
В противном случае  \TRUE{} Затем выполняется утверждение .
Если исполнение выполняется нормально или продолжается без ярлыка
или на этикетке () Приставка  заявление
(),
Затем заявление повторно исполняется.
Если казнь прерывается без ярлыка,
Выполнение заявления при этом завершается нормально.
{%
Если исполнение прерывается ярлыком с префиксом {} заявление,
Это завершает выполнение цикла,
Но сам разрыв управляется
Окружающая маркированная запись ().%
?

%
Ошибка времени компиляции, если
Статический тип  Не может быть присвоено .


 Правда?


%
Заявление do поддерживает условную итерацию.
где состояние оценивается после петли.


<do Statement> ::= \DO{} \WHILE{} '(' <выражение> '''); "


%Исполнение заявления о выполнении формы \code{\DO{}  < ()
идет следующим образом:

%
Заявление  выполнено.
Если эта казнь продолжится без ярлыка,
или на этикетке () что префиксирует заявление do
(),
Затем выполнение  рассматривается как завершенное в обычном режиме.

%
Выражение  оценивается на предмет .
% Эта ошибка может произойти из-за неявных отливок и нулевых.
Это динамическая ошибка, если тип времени выполнения  Это не P000918.
Если  , выполнение заявления do завершается нормально
(P000749).
Если  , то заявление do выполняется повторно.

%
Это ошибка времени компиляции, если статический тип 
Не может быть присвоено .


{Переключатель}


%
 поддерживает диспетчерский контроль между
большое количество случаев.


<switchStatement> ::= <
\SWITCH{} '(' <выражение> '''''''{<switchCase>* <defaultCase>?''' "

<switchCase> ::= <label>* \CASE{} <expression>> <statements>

<defaultCase> ::= <label>*{} <statements>


%
Рассмотрим коммутационное заявление формы




или форма



{%
Обратите внимание, что каждое выражение  происходит в постоянном контексте
(),
что означает  Модификаторы не должны быть указаны явно.%
?

%
Ошибка времени компиляции, если каждое выражение не
 Постоянная.
Ошибка времени компиляции, если значение выражений
 Не являются также:


 экземпляры одного и того же класса , для всех , или
 Класс, который реализует ,
Для всех 
 Класс, который реализует ,
Для всех .


{%
Другими словами, все выражения в случаях оцениваются константами
один и тот же класс пользователей или определенных известных типов.
% % TODO(Eernst): Обновление, когда мы указываем вывод: const List<int> xs = []
% может быть контрпримером: Значение — это список, который «знает», что это
% % «Список» независимо от типа аннотации, но это не будет
% % «Список <int>», если такого рода аннотации не было.
Обратите внимание, что значения выражений известны во время компиляции.
независимо от аннотаций любого типа.%
?

%
Это ошибка времени компиляции, если один или несколько указанных экземпляров класса 
Нет примитивного равенства
(P000751).

{%
Запрет на равенство пользователей позволяет нам
эффективно реализовать коммутатор для определенных пользователем типов.
Вместо этого мы могли бы сформулировать соответствие с точки зрения идентичности.
с той же эффективностью.
Однако, если тип определяет оператора равенства,
Вероятно, программисты сочтут это довольно удивительным.
Если равные объекты не совпадают.%
?

{%
 Заявление должно использоваться только в
очень ограниченные ситуации (например, переводчики или сканеры).
?

%
Выполнение заявления о переключателе формы




или форма

<<<p000010>


идет следующим образом:

%
Заявление \code{\VAR{}  = ; оценивается,τ
где  - свежая переменная.
% Эта ошибка может произойти из-за неявных слепков и стандартного подчинения.
% % TODO(Eernst): Но почему не может  Быть примером подтипа?!
Это динамическая ошибка, если значение  это
не относится к тому же классу, что и константы .

{%
Обратите внимание, что если нет кейсов (),Тип  Не имеет значения.%
?

%
Далее, пункт \CASE{} :  Сопоставляется с ,
если .
В противном случае, если есть  пункт,
Заявления по делу  выполняется ().

%
Соответствие {} Статья \CASE{}  Заявление о переключателе

<<<p000011>


против величины переменной  идет следующим образом:

%
Выражение  оценивается на предмет .
% Эта ошибка может произойти из-за неявных отливок и нулевых.
Это динамическая ошибка, если тип времени выполнения  не является .
Если   В следующем случае,
<  противопоставляется , если ,
Если , то  Заявления пункта исполняются
(P000753).
Если  , то пусть  быть наименьшим числом
Таким образом,  и  не являются пустыми.
Если нет, то  Существует, пусть .
Заявления по делу  Затем выполняется ().

%
Соответствие {} Статья \CASE{}  Заявление о переключателе




против величины переменной  идет следующим образом:

%
Выражение  Оценивается на объект .
% Эта ошибка может произойти из-за неявных отливок и нулевых.
Это динамическая ошибка, если тип времени выполнения  Это не P000925.
Если   В следующем случае,
<  Сопоставляется с , если .
Если  , то пусть $h$ быть наименьшим целым числом
 и  не являются пустыми.
Если это так, P000317 Существует, случай заявления  исполняется
(P000755).
В противном случае заявление о переключателе обычно завершается
(P000756).

%
Это ошибка времени компиляции, если тип 
не может быть присвоен типу .
 быть последним утверждением последовательности высказываний .
Если  Это непустая блокировка, пусть  вместо того чтобы быть
Последнее заявление блока.
Ошибка времени компиляции  не является
а , , , или  заявление,
или выражение, где выражение является выражением {}.

{%
Поведение переключателей намеренно отличается от традиции С.
Неявное падение является известной причиной ошибок программирования.
и, следовательно, запрещены.
Почему бы просто не разорвать поток в конце каждого дела?
Вместо того, чтобы требовать явного кода для этого?
Это действительно было бы чище.
Кроме того, было бы чище настаивать на том, что в каждом случае
Одно (возможно, сложное) утверждение.
Мы решили не делать этого, чтобы
облегчает перенос коммутационных выражений с других языков.
Неявно нарушая контрольный поток в конце дела молча
изменить значение портированного кода, который полагался на промах,
Потенциально вынуждая программиста иметь дело с тонкими ошибками.
Наш дизайн гарантирует, что разница будет немедленно
Внимание кодера.
Программист будет уведомлен во время компиляции, если он забудет закончить дело.
с заявлением, которое прекращает прямолинейный контрольный поток.

Изощренность анализа сквозного провала является еще одним вопросом.
На данный момент мы выбрали очень простое синтаксическое требование.
Очевидно, что существуют ситуации, когда код не проходит.
и все же не соответствует этим простым правилам, например:
?



{%
Очень сложный код в кейсе, вероятно, плохой стиль в любом случае.
Такой код всегда можно рефакторировать.
?%
Это статическое предупреждение, если соблюдены все следующие условия:


 Заявление о переключателе не имеет пункта по умолчанию.
 Статический тип  является перечисленным
Тип с элементами .
 Наборы $\{e_1, \ldots, e_k\} $ и $\{\id_1, \ldots, \id_n\}$
Они не одинаковые.


{%
Другими словами, будет выпущено статическое предупреждение.
если выписка переключателя по перечню не является исчерпывающей.%
?


 Переключите показания по делу.


%
Исполнение заявлений по делу  из заявления переключателя




или заявление о переключении




идет следующим образом:

%
Выполнить .
Если эта процедура выполняется нормально,
Если  не являются показаниями последнего случая переключения
() Если нет  пункт,
 если имеется пункт {},
Затем выполнение корпуса переключателя выбрасывает ошибку.
В противном случае  являются последними показаниями дела о переключении,
И выполнение дела переключателя завершается нормально.

{%
Иными словами, между непустыми случаями нет неявного разрыва.
Последний случай в переключателе (по умолчанию или иным образом) может «пройти» "
В конце статьи.%
?

Если выполняется  без ярлыков (),
Затем выполнение заявления переключателя завершается нормально.

Если исполнение  Продолжение этикетки
(),
и этикетка ,
где  Если  Заявление имеет ,
 Если нет ,
где ,
Тогда пусть  быть наименьшим числом, таким, чтобы  и  были непустыми.
Если нет, то  существует,
 Если  Заявление имеет ,
В противном случае пусть .
Заявление по делу  Затем они выполняются ().

Если исполнение  завершается любым другим способом,
Выполнение заявления {} завершается таким же образом.


 Сброс


%
 используется для
Перебросьте исключение и связанный с ним след стека.


<rethrow Statement> ::= \RETHROW{{}}} "


%
Рассмотрим  заявление .
%  Как определено ниже, гарантируется существование: Каждое высказывание происходит в
% <block> (заявление может быть частью другого заявления , но
Гипотеза % индукции  происходит в <block>, а <block> может быть только
% происходит как утверждение или в заявлении (например, в <конечной части>), и это
Индукция % может закончиться только в «функциональном теле».
 быть функцией непосредственного включения .
Ошибка времени компиляции происходит, если  находится в \ON-\CATCH{} пункт
Функция немедленного включения — .

%
Исполнение a  заявление продолжается следующим образом:

%
 быть непосредственной функцией,
И пусть 
Включая пункт «Улов» ().

<<%
A < >> Заявление всегда появляется внутри  пункт,
и каждый \CATCH{} пункт рассматривается как некоторый \CATCH{} пункт о форме
.
Таким образом, мы можем рассмотреть  Заявление должно быть заключено в
a \CATCH{} в этой форме.%
?

%
 Заявление затем бросает ()
Со значением  в качестве объекта исключения,
и значение  Как след стека.


 Попробуй!


%
Заявление о попытке поддерживает определение кода обработки исключений
структурированным способом.

<try Statement> ::= \TRY{} <block> (<onPart>+ <finallyPart>? | <finallyPart>)

<onPart> ::= <catchPart> <block>
 << <typeNotVoid> <catchPart> <блокировать>

<catchPart> ::= \CATCH{} '(' <identifier> (',' <identifier>)?' "

<finallyPart> ::= << <блокировать>


%
Пробное утверждение состоит из блок-заявления, за которым следует по меньшей мере одно из:



Набор \ON{}-{} Положения, каждый из которых определяет
(явно или неявно)
тип объекта исключения, подлежащего обработке,
один или два исключительных параметра,
и блок-заявление.

A  >> пункт, который состоит из блок-заявления.


{%
Синтаксис предназначен для повышения совместимости с
Существующие программы Javascript.
 пункт можно опустить,
оставляя то, что выглядит как оговорка об улове JavaScript.%
?

%
Пробное заявление формы
\code{\TRY{}  ;
является эквивалентом заявления
\code{\TRY{)   << .

%
An \ON{}-\CATCH{} пункт о форме
\code{\ON{}  << (перенаправлено с «P000371») 
эквивалентен \ON{}-\CATCH{} пункт
\code{\ON{)  << (,  
где  - новый идентификатор.

%
An \ON{}-{} пункт о форме
\code{\ON{)  
является эквивалентом \ON{}-\CATCH{} пункт
\code{{)  \CATCH{} (,  
где  и  Это свежие идентификаторы.

%
An \ON{}-{} пункт о форме
\code{{} () 
является эквивалентом \ON{}-\CATCH{} пункт
\code{{) << \CATCH{} (,  
где  - новый идентификатор.

An \ON{}-{} пункт о форме
\code{{) (,  
является эквивалентом \ON{}-\CATCH{} пункт
\code{{) < \CATCH{} (,  .

%
An \ON{}-{} пункт о форме
\code{{)  {} (,  
Введен новый объем  где конечные локальные переменные
 и  определены.
Заявление  Он находится в пределах .
Статический тип  является 
Статический тип  является .

%
Исполнение a  Заявление  По форме:




идет следующим образом:

%
Сначала выполняется .
Если выполняется  броски ()
За исключением объектов  и след стека ,
 и  Сопоставлены с \ON{}-\CATCH{} положения
Получить новое завершение ().

Тогда, даже если исполнение [P000417] обычно не завершается
или сопоставление с \ON{}-\CATCH{} пункты не были выполнены нормально,
 Блок выполнен.

Если исполнение  обычно не завершается,
Выполнение заявления {} завершается таким же образом.
В противном случае, если выполнение  бросили (),
 Заявление завершается так же, как
Соответствие с \ON{}-\CATCH{} пункты.В противном случае  Заявление завершается так же, как
Выполнение .

%
Ошибка времени компиляции, если   - отложенный тип.


\subsubsection{{}-{} пункты }


%
Соответствие объекта исключения  и след стека  против
a (потенциально пустая) последовательность \ON{}-{} пункты формы

<<<p000016>
идет следующим образом:

%
Если нет {}-{} оговорок (), соответствующие броски
Объект исключения  и след стека 
(P000765).

%
В противном случае исключение совпадает с первым пунктом.

%
В противном случае, если тип  является подтипом ,
Первый из них – это первый,
 привязывается к объекту исключения 
 привязан к следу стека ,
и  выполняется в этой области.
Соответствие завершается так же, как и это исполнение.

%
В противном случае, если первый пункт не соответствует ,
 и  Рекурсивно сопоставляются с
Оставшийся \ON{}-\CATCH{} пункты:




 Возвращение!


%

возвращает результат абоненту синхронной функции,
завершает будущее, связанное с асинхронной функцией;
или прекращает поток или итерируемый, связанный с генератором
(P000766).


<return Statement> ::= \RETURN{} <выражение>? "


%
Рассмотрение заявления о возвращении  формы \code{\RETURN{} .
 Статический тип , если  присутствует,
 быть непосредственной функцией,
 Объявленный тип возврата .

%
 Синхронные негенераторные функции
Рассмотрим случай, когда  Синхронная негенераторная функция
(P000767).
%
% Возвращение без объекта — это нормально только для «пустых» типов возврата.
Это ошибка времени компиляции, если  ,
Если только  не ,  или .
%
% Возвращение с объектом в функции пустоты
% — это нормально только при возврате типа «пустота».
Это ошибка времени компиляции, если  {\RETURN{} ;}
 ,
 не является ни ,  или .
%
% Возвращение пустоты в «непустой» функции является ошибкой.
Это ошибка времени компиляции, если  \code{\RETURN{} ;}
 не является ни ,  или ,
 .
%
В противном случае возврат неназначенного значения является ошибкой.
Это ошибка времени компиляции, если  {\RETURN{} ;}
 не является ,
 не присваивается .

{%
Обратите внимание, что  Не может быть ,  или  В последнем случае,
Потому что все типы относятся к этим типам.
Ошибка не возникает, если  не имеет декларируемого типа возврата,
Поскольку тип возврата будет , к которому можно отнести каждый тип.
Синхронная негенераторная функция
который объявляет тип возврата, который не является «пустым» " "
Должен вернуть выражение явно.%
?
{%
Это помогает уловить ситуации, когда пользователи забывают.
вернуть объект в отчете о возврате.%
?


%
 Асинхронные негенераторные функции
Рассмотрим случай, когда  асинхронная негенераторная функция
(P000768).
%% Возвращение без объекта - это нормально только для типов возврата с асинхронным "пустым".
Это ошибка времени компиляции, если  ,
Если только  T
()
,  или .
%
{%
Асинхронный негенератор всегда возвращает будущее.
Если не будет дано никакого выражения, будущее будет завершено нулевым объектом.
()
Что мотивирует это правило.%
?
% Возвращение с объектом в пустоте функции асинхронизации только в порядке
%, когда это значение асинхронно- «пустота».
Это ошибка времени компиляции, если  \code{\RETURN{} ;}
{T} - ,
и {S} не является ни ,  или .
%
% Возвращение асинхронизации в функции «не-асинхронизации» является ошибкой.
Это ошибка времени компиляции, если  {\RETURN{} ;}
{T} не является ни , ни  или ,
{S} — .
%
В противном случае возврат не поддающегося анализу значения является ошибкой.
Это ошибка времени компиляции, если  \code{\RETURN{} ;}
{S} не является ,
 Будущее <{S}> не может быть присвоено .

{%
Обратите внимание, что {T} не может быть ,  или 
В последнем случае,
Потому что тогда  присваивается  для \emph{all} .
В частности, когда  
(что эквивалентно ),
 Будущее <{S}>> присваивается  для всех .
Это означает, что ошибка времени компиляции не возникает.
Но {только} нулевой объект ()
или экземпляр 
Его можно успешно вернуть во время бега.
Это не аномалия,
соответствует лечению синхронной функции
Тип возврата ;
Но инструменты могут дать намек на то, что понижение вряд ли удастся.

Ошибка не будет поднята, если  не имеет декларируемого типа возврата,
Поскольку тип возврата будет ,
 Будущее <\flatten{S}>} присваивается \DYNAMIC{} для всех .
Асинхронная негенераторная функция
который объявляет тип возврата, который не является «пустым» " "
Должен вернуть выражение явно.%
?
{%
Это помогает уловить ситуации, когда пользователи забывают.
возвращать объект в обратной записи асинхронной функции.%
?


%
 Функции генератора.
Ошибка времени компиляции, если заявление о возврате
форма \code{{} ; появляется в функции генератора.

{%
В случае функции генератора объект, возвращаемый функцией, является
Иль с ним, или с ним,
и отдельные элементы добавляются к этому итерабельному с использованием отчетов о выходе,
Возвращение объекта не имеет смысла.%
?


%
 Генеративные конструкторы.
Ошибка времени компиляции, если заявление о возврате
форма \code{{) ; появляется в генеративном конструкторе
(P000772).

{%
Легко забыть добавить  Модификатор для конструктора,
случайное преобразование фабрики в генераторный конструктор.
Статическая шашка может обнаружить несоответствие типа
В некоторых, но не во всех случаях.
Вышеприведенное правило помогает уловить такие ошибки.
которые в противном случае очень трудно распознать.
В этом нет никакого реального недостатка,
Возвращение объекта от генеративного конструктора бессмысленно.
?


%
Выполнение заявления о возвращении \code{{) ; идет следующим образом:

%Сначала выражение  оценивается, образуя объект .
 Выберите тип времени выполнения  и
 Тип возврата 
(P000773).
Если тело  Помечено \ASYNC{} ()
 является подтипом  Будущее<\flatten{T}>
 быть результатом оценки 
где  является свежей переменной, связанной с .
В противном случае пусть .
Затем обратная запись возвращает объект 
(P000775).

% % TODO(Eernst): У нас есть несколько особых случаев с динамической семантикой.
%, указанные выше, которые мы можем рассмотреть:
% %
% % Future<int> foo() async
% возврата нового Future<Object>.value(42);
% % }
% %
% % Статически мы позволим это, потому что Будущее <flatten (Будущее <Объект>)>
%, т.е. «Будущее<Объект>», присваивается «Будущее<int>». Но во время бега
% % мы получаем ошибку, потому что «Будущее» не является подтипом
% % «Future<flatten(Future<int>)» = «Future<int>». Мы могли бы ждать
% от стоимости только потому, что мы собираемся вернуть «Будущее», и это будет
Процентная доля тогда удалась. Мы поступаем правильно? Может быть и хорошо, но
% % кажется удивительным, что компилятор иногда вставляет «ождать».
% и здесь мы должны сделать это явно, чтобы добиться успеха.
% %
% % Объект /* или динамический и т.д. */ foo() async {
% % Объект О;
% %, если (некоторые условия) o = 42; другие o = новое будущее<S>.value(42);
% возврата o;
% % }
% %
% % Здесь мы будем вставлять «ожидание» на возвращенную стоимость всякий раз, когда
% Мы возвращаем будущее (для всех S). Это означает, что мы «ждем»
% в некоторых ситуациях, а не в других. Это может удивить разработчиков.
%%, что мы можем иметь это динамическое изменение в месте в коде
%, где нет статического обоснования для ожидания будущего.

%
 Это тип времени выполнения .



Если тело  Помечено \ASYNC{} (перенаправлено с «P000776»)
% Эта ошибка может произойти из-за неявных слепков.
Ошибка динамического типа, если  Не является подтипом .

% Эта ошибка может произойти из-за неявных слепков.
В противном случае это ошибка динамического типа, если  Не является подтипом .


%
Выполнение заявления о возврате без выражения,

Возвращение без объекта
(P000777).


<{Методы}


%
 Это идентификатор, за которым следует толстая кишка.
 Заявление с приставкой .
 является клаузулой случая в заявлении о переключении
(перенаправлено с «P000778»)
Префиксируется этикеткой .

{%
Единственная роль этикеток состоит в том, чтобы обеспечить цели для
брейк () и продолжение () заявлений.%
?


<label> ::= <identifier> "


%
Исполнение маркированного заявления , , состоит из выполнения .
Если выполняется  Перерывы на этикетке  (),
Выполнение  обычно завершается,
иное исполнение  Завершается таким же образом, как и исполнение .

%
Пространство имен этикеток отличается от используемого для
Типы, функции и переменные.

%
Область применения ярлыка, обозначающего утверждение , составляет .
Область применения ярлыка, обозначающего кейсовую оговорку выключателя , составляет .

{%
Программисты должны избегать ярлыков любой ценой.
Мотивация включения ярлыков в язык
В первую очередь делает Дарт - лучшая цель для генерации кода.%
?


 Разрыв!


%
 состоит из
зарезервировано слово \BREAK{} и факультативная метка (\ref{labels}).

<break Statement> ::= < <идентификатор>? "


%
Пусть  будет утверждением .
Если  имеет форму \code{\BREAK{} ;}
Ошибка времени компиляции  не содержится в маркированном заявлении
с меткой  внутри самой внутренней функции, в которой  происходит.
Если  имеет форму ,
Ошибка времени компиляции  не закрепляется в одном
\code{{)  (),
< (), < (), \SWITCH{} ()
или  (перенаправлено с «P000787») заявление внутри
Внутренняя функция, в которой  происходит.

%
Исполнение a  заявление \code{{) ; Перерывы на этикетке 
(P000788).
Исполнение a  заявление  Разрыв без этикетки
(P000789).


 Продолжайте.


%
 Состоит из зарезервированного слова \CONTINUE{}
и факультативную этикетку ().


<продолжение заявления> ::= \CONTINUE{} <идентификатор>? "


%
Пусть  будет утверждением .
Если  имеет форму \code{\CONTINUE{} ;},
Ошибка времени компиляции, если  Ни в одном из них не содержится
\code{{)  (),
\DO{} (), \FOR{} (), или \WHILE ()
Заявление, обозначенное , или в \SWITCH{} Заявление с клаузулой дела
«Спасибо, если у вас есть все, что вы можете сделать, если у вас есть все, что у вас есть» (с) Удалось.
Если  Форма: 
Ошибка времени компиляции  не закрепляется в одном
\code{{)  ()
<< (), << (), или  (перенаправлено с «P000798») заявление
внутри самой внутренней функции, в которой  происходит.

%
Исполнение a  заявление \code{{) ; продолжать
на этикетке  (P000799).
Исполнение a \CONTINUE{} заявление  продолжать
без маркировки ().


 Доходность


%
 Добавить объект в
Результат функции генератора
().


<yieldStatement> ::= \YIELD{{}<выражение>> "


%
 быть заявлением о доходности формы .
 Для этого необходимо выполнить функцию .
Это ошибка времени компиляции, если такой функции нет.
Или это не генератор.
Это ошибка времени компиляции, если статический тип 
не может быть присвоен типу элемента 
().

%
Выполнение заявления  формы 
идет следующим образом:

%
Во-первых, выражение  Оценивается на объект .
Если функция ограждения  
()
и поток  Связанный с  был приостановлен,
ближайший асинхронный к петле
(),
если таковой имеется, приостанавливается и выполняется  приостановлено
до тех пор, пока  не будет возобновлен или отменен.

%
Далее  Добавляется в итерируемый или связанный поток
с непосредственной ограждающей функцией.

{%
Обратите внимание, что динамическая ошибка возникает, если динамический тип 
не является подтипом элемента типа упомянутого итерабельного или потокового.%
?

%
Если функция ограждения  и поток  Связано с  был отменен,
< Заявление возвращается без объекта
(),
В противном случае она завершается нормально.

{%
Поток, связанный с асинхронным генератором, может быть отменен
любым кодом со ссылкой на этот поток в любой точке
где генератор был пассивным.
Такая отмена представляет собой безвозвратную ошибку для генератора.
На данный момент единственным правдоподобным действием для генератора является
Чтобы очистить себя через \FINALLY{} Оговорки.%
?

%
В противном случае, если прилагаемая функция  
()
После этого функция может быть приостановлена,
в каком случае ближайшая замкнутая асинхронная для петли
(),
Если таковые имеются, то сначала останавливают.

{%
Если  Происходит внутри бесконечной петли
и функция прикрытия никогда не приостанавливается,
Возможно, у потребителей замкнутого потока не будет возможности
для запуска и доступа к данным в потоке.
Затем поток может накапливать неограниченное количество элементов.
Такая ситуация несостоятельна.
Таким образом, мы допускаем приостановку прилагающей функции.
когда к связанному потоку добавляется новый объект.
Однако это не обязательно (и на самом деле может быть довольно дорогостоящим).
Приостановить функцию на каждом .
Реализация свободна решать
Как часто приостанавливать функцию ограждения.
Единственным требованием является то, что потребители не блокируются на неопределенный срок.
?

%
Если функция ограждения   () затем:



Выполнение функции  немедленно закрывается  приостановлено
до нулевого метода  На него ссылаются
Итератор используется для инициирования текущего вызова .

Текущий призыв к  Возвращается .



 Доходность


%
 Добавляет ряд объектов к
Результат функции генератора
().


<yieldEachStatement> ::= \YIELD{} '*' <выражение>'; "


%
 (b) быть заявлением о доходности по форме .
 быть функцией непосредственного окружения .
Это ошибка времени компиляции, если такой функции нет.
Или это не генератор.

%
 Тип элемента 
(),
 Статический тип .
Если  является синхронным генератором,
Ошибка времени компиляции  не может быть назначен
.
В противном случае  является асинхронным генератором,
Ошибка времени компиляции, если  не может быть назначен
.

%
Исполнение заявления  формы  "
идет следующим образом:

%
Выражение  оценивается на предмет .

%
Если функция немедленного закрытия  
(),
затем:



% Эта ошибка может произойти из-за неявных слепков.
Ошибка динамического типа
Если класс  не является подтипом .
Иначе

Метод  На него ссылаются  Возвращение объекта .

  Метод  на него ссылаются
Без аргументов.
Если  Возвращение  >> Выполнение  завершено.
Иначе

Геттер  Называется на .
Если призыв бросает
(),
Выполнение  бросает тот же самый объект исключения и след стека
(P000813).
В противном случае результат  Призыв Геттера добавляется кИтерация, связанная с .
{%
Обратите внимание, что динамическая ошибка возникает, если динамический тип  это
не является подтипом элемента типа упомянутого итерабельного.%
?
Выполнение функции  немедленно закрывается  приостановлено
до нулевого метода  на него ссылаются
на итераторе, используемом для инициирования текущего вызова ,
в какой момент выполнения  Продолжается на .

Текущий призыв к  Возвращается .


%
<<<<<<<<<<<<<<<<<<<<<<<<<<<<<<<<<<<<<<<<<<<<<<<<<<<<<<<<<<<<<<<<<<<<<<<<<<<<<<<<<<<<<<<<<<<<<<<<<<<<<<<<<<<<<<<<<<<<<<<<<<<<<<<<<<<<<<<<<<<<<<<<<<<<<<<<<<<<<<<<<<<<<<<<<<<<<<<<<<<<<<<<<<<<<<<<<<<<<<<<<<<<<<<<<<<<<<<<<<<<<<<<<<<<<<<<<<<<<<<<<<<<<<<<<<<<<<<< Для справки: <<<<<<<<<<<<<<<<<<<<<<<<<<<<<<<<<<<<<<<<<<<<<<<<<<<<<<<<<<<<<<<<<<<<<<<<<<<<<<<<<<<<<<<<<<<<<<<<<<<<<<<<<<<<<<<<<<<<<<<<<<<<<<<<<<<<<<<<<<<<<<<<<<<<<<<<<<<<<<<<<<<<<<<<<<<<<<<<<<<<<<<<<<<<<<<<<<<<<<<<<<<<<<<<<<<<<<<<<<<<<<<<<<<<<<<<<<<<<<<<<<<<<<<<<< (на английском языке), а именно:



% Эта ошибка может произойти из-за неявных слепков.
Ошибка динамического типа, если класс 
не является подтипом .
Иначе

Ближайший замкнутый асинхронный цикл (),
Если есть, то приостанавливается.

 Прослушивается поток, создавая подписку ,
и для каждого события  или ошибки  со следом стека , из :



Если поток  Связанный с  был приостановлен,
Выполнение  Приостановлено до  возобновляется или отменяется.

Если поток  Связанный с  был отменен,
 Отменяется оценкой \code{\AWAIT{} v.cancel()}
где  Это новая переменная, ссылающаяся на подписку на поток .
Если отмена завершится нормально,
Потоковое исполнение  Возвращение без объекта
(P000817).

В противном случае, $x$ или $e$ с $t$ добавляются к
Поток, связанный с  в порядке, указанном в .
{%
Обратите внимание, что динамическая ошибка возникает, если  добавляется
и динамический тип  не является подтипом
Тип элемента указанного потока.%
?
Функция  может приостановиться.


Если поток  выполняется, выполняется  обычно завершается.



 Утверждение


%
 Используется для нарушения нормального исполнения
если данное булевое состояние не выдерживает


<assert Statement> ::= <утверждение>; "

<утверждение> ::= \ASSERT{} '(' <выражение> (',' <выражение>)? ','?' "


%
Грамматика позволяет запятую до закрытия скобки,
По аналогии со списком аргументов.
Эта запятая, если она присутствует, не имеет никакого эффекта.
Утверждение с задней запятой эквивалентно утверждению с удаленной запятой.

%
Утверждение формы  является эквивалентным
a утверждение формы \code{\ASSERT(, \NULL{})}.

%
Выполнение заявления об утверждении выполняет утверждение, как описано ниже
и завершается так же, как и утверждение.

%
Когда не допускается утверждение,
Выполнение утверждения немедленно завершается нормально
(P000818).
 То есть никакие подвыражения утверждения не оцениваются. ?
Когда выдвигаются требования,
выполнение утверждения \code{{}(,  идет следующим образом:

%
Выражение  оценивается на предмет .
% Эта ошибка может произойти из-за неявных отливок и нулевых.
Ошибка динамического типа, если  не является типом .
{%
Это ошибка времени компиляции, если такая ситуация возникает.
во время оценки утверждения в  Запрос конструктора.%
?
Если   Затем исполнение заявления об утверждении завершается обычно
(P000819).В противном случае  оценивается на предмет 
а затем исполнение утверждения бросает
()
  и с
след стека, соответствующий текущему состоянию исполнения при утверждении.

%
Ошибка времени компиляции, если тип 
Не может быть присвоено .

{%
Почему это утверждение, а не встроенный вызов функции?
Потому что с ним обращаются магически
Таким образом, он не имеет никакого эффекта и никаких накладных расходов, когда утверждения отключены.
При отсутствии окончательных методов,
Нельзя было предотвратить его переохлаждение.
(хотя в этом нет никакого реального вреда).
Его нельзя рассматривать как вызов функции, который оптимизируется.
Потому что аргументы могут иметь побочные эффекты.
?


 Библиотеки и сценарии


%
А. Программа Dart состоит из одной или нескольких библиотек.
может быть построено из одного или нескольких .
Блок компиляции может быть библиотекой или частью ().

%
Библиотека состоит из (возможно, пустого) набора импорта, набора экспорта.
и ряд деклараций высшего уровня.
Декларация верхнего уровня представляет собой класс ().
декларация псевдонима типа (),
Функция ()
или переменной декларацией ().
Члены библиотеки  Это декларации верхнего уровня, приведенные в .


<topLevelDeclaration> ::= <classDeclaration>
 <mixinDeclaration>
< <расширенная декларация>
 <enumType>
 <typeAlias>
 \EXTERNAL{} <функциональная подпись>>; '
 \EXTERNAL{{getterSignature}} "
\alt \EXTERNAL{{item.group_date}} "
 <функциональная подпись> <Функциональное тело>
 <getterSignature> <functionBody>
 <setterSignature> <functionBody>
 (\FINAL{} | \CONST{}) <type>? <staticFinalDeclarationList>; "
 < \FINAL{} <type>? <initializedIdentifierList>> "
 \LATE? <varOrType> <initializedIdentifierList>> "

<libraryDeclaration> ::= <<
<scriptTag>. <название библиотеки>. <importOrExport> <partDirective> *
\gnewline{} (<metadata> <topLevelDeclaration>)* <EOF>

<scriptTag> ::=#! () \>\>)* <LINE\_ Брейк>

<libraryName> ::= <metadata> \LIBRARY{{}<dottedIdentifierList>> "

<importOrExport> ::= <libraryImport>
 <libraryExport>

<dottedIdentifierList> ::= <identifier> ('.' <identifier>) *


%
Библиотека содержит строку, полученную из {libraryDeclaration}.

{%
Мы можем сказать, что {библиотечная декларация}

грамматики,
Потому что отсюда начинается синтаксический вывод любой библиотеки.
В отличие от традиционной контекстной свободной грамматики,
Грамматика Дарта не имеет точно одного начального символа.
В частности, \synt{partDeclaration}
()
Используется таким же образом для получения содержимого части.
Может быть больше, например,
Гипотетический цикл Dart REPL (read-eval-print loop)
\synt{выражение} или {утверждение} в качестве начального символа.
Таким образом, нет грамматики для программ Дарта как таковых.
только для некоторых строительных блоков, которые используются для создания программ Dart.
?

%
Библиотеки могут быть названы явно или неявно.
 начинается со слова 
(возможно, предваряемый любыми применимыми аннотациями метаданных),
Затем следует квалифицированный идентификатор, который дает название библиотеки.

{%
Технически каждая точка и идентификатор являются отдельным токеном.
Пространства между ними вполне приемлемы.
Тем не менее, фактическое название библиотеки является конкатенацией
Простые идентификаторы и точки, не содержащие пробелов.%?

%
Неявно названная библиотека имеет пустую строку в качестве своего названия.

{%
Название библиотеки используется для привязки к
отдельно составленные части библиотеки (называемые частями) и
Может использоваться для печати и, в более общем смысле, для отражения.
Название может иметь отношение к дальнейшей эволюции языка.
?

{%
Библиотеки, предназначенные для широкого использования, должны избегать столкновений имен.
Dart's  Система управления пакетами обеспечивает
механизм для этого.
Каждому пабу гарантировано уникальное название,
Эффективное обеспечение глобального пространства имен.%
?

%
Библиотека может необязательно начинаться с .
Теги сценариев предназначены для использования со скриптами ().
Тег сценария может использоваться для идентификации интерпретатора сценария для
В какую бы вычислительную среду не встроен сценарий.
Тег сценария должен появиться перед любым белым пространством или комментариями.
Сценарий начинается с \lit{\#!} и заканчивается в конце строки.
Любые символы, следующие за \lit{\#! Сценарий игнорируется
Реализация Дарта.

%
Библиотеки являются единицами приватности.
Частная декларация, объявленная в библиотеке 
Доступ возможен только с помощью кода внутри .

{%
Поскольку частные лица высшего уровня не импортируются,
Использование приватов высшего уровня другой библиотеки никогда не возможно.
?

%
 Библиотека  Пространство имен, которое отображает
Название каждой публичной декларации членов высшего уровня   .
 Библиотека  Пространство имен, которое отображает
имена, вводимые каждой декларацией верхнего уровня 
к соответствующей декларации.
 Библиотека  является самой внешней областью в ,
и его пространство имен является пространством имен библиотеки 
(P000828).

%
Настройщик не может быть сопряжен с функцией, классом и т.д.
Это ошибка времени компиляции, если локальное пространство имен библиотеки 
имеет две декларации с одинаковым именем,
Если только они не являются геттерами и сеттерами.

{%
Два разных имени  и  может иметь только одно и то же имя
Когда они имеют форму  и ,
Таким образом, этот вид конфликта  >> всегда включает сеттера.
Но другой декларацией может быть функция, класс и т.д.
?


 Импорт


%
 Библиотека, чье экспортируемое пространство имен
(или подмножество его отображений) доступно в текущей библиотеке.


<libraryImport> ::= <metadata> <importSpecification>

<importspecification> ::= <
< <configurableUri>
(). \AS{} <typeIdentifier>)?
<комбинатор>*; "


%
Интерпретация настраиваемых URI описана в другом месте.
().
Импорт определяет URI 
где находится декларация импортируемой библиотеки.
Ошибка времени компиляции, если указанный URI импорта
Это не относится к библиотечной декларации.

%
 В настоящее время ведется работа над библиотекой.
Импорт изменяет пространство имен текущей библиотеки.
в порядке, определяемом импортной библиотекой и
факультативные элементы импорта.

%
Импорт может быть отложен или немедленным.
А.
\IndexCustom{deferred import}{import!deferred}
Отличается возникновением
встроенный идентификатор {} после URI.
Ан
{немедленный импорт}{import!immediate}
Это импорт, который не откладывается.

%
Немедленная директива по импорту  может необязательно включать
а  формы , используемой для префикса
Имя привезено по адресу .
В данном случае мы говорим, что  ,
Или просто «P000618».

{%
Обратите внимание, что грамматика обеспечивает отсроченный импорт.
включает в себя пункт префикса,
Таким образом, мы можем ссылаться на пункт префикса {the} отложенного импорта.%
?

%
Это ошибка времени компиляции, если префикс используется в отложенном импорте.
Он также используется в качестве приставки к другому пункту импорта.
Ошибка времени компиляции  является префиксом импорта,
и текущая библиотека объявляет члена верхнего уровня с именем базы .

%
Директива по импорту  может необязательно включать положения комбинатора пространства имен
Используется для ограничения набора имен, импортируемых по .
Их синтаксис, использование и влияние на пространства имен описаны в другом месте.
(, ,
.

%
Библиотека ядра дротика 
Неявно вносится в каждую библиотеку дарта, кроме самой себя.
через импортную оговорку формы
{{} 'dart:core;}
Если библиотека не импортирует .
любой импорт ,
Если вы ограничены через ,  или ,
Предотвращает автоматический импорт.

<%
Было бы неплохо, если бы в P000650 не было ничего особенного.
Однако его использование является распространенным,
Это приводит к решению импортировать его автоматически.
С другой стороны, некоторые библиотеки  может пожелать определить
с именами, используемыми 
(что легко сделать, поскольку имена, объявленные библиотекой, имеют приоритет).
Другие библиотеки могут пожелать использовать ,
и может использовать членов  Это противоречит основной библиотеке.
без использования префикса и без возникновения ошибок.
Вышеприведенное правило делает это возможным.
по существу отменить  Специальный режим
С помощью еще одного специального правила.%
?

 Импортируемое пространство имен


%
Ниже мы указываем импортное пространство имен библиотеки ,
{{import}}
и использовать это для определения пространства имен, которое определяет объем библиотеки .

%
Мы должны ввести системные библиотеки, потому что у них есть особые правила.
A  - это библиотека, которая является частью реализации Dart.
Любая другая библиотека представляет собой .

{%
Системную библиотеку можно распознать, если
URI, начинающийся с .%
?

{%Особые правила для системных библиотек существуют по следующей причине.
Нормальные конфликты решаются во время развертывания.
Функциональность системной библиотеки
вводится в приложение во время выполнения,
Они могут меняться с течением времени по мере обновления платформы.
Возможны конфликты с системной библиотекой.
вне контроля застройщика.
Чтобы избежать взлома развернутых приложений,
Конфликты с системными библиотеками рассматриваются специально.
?

%
 Будьте импортной директивой и позвольте  Библиотека, импортированная по адресу .
The
{пространство имен, предоставленное}{%
Пространство имен! в соответствии с импортной директивой
Директива по импорту  Пространство имен, полученное при применении
Комбинаторы пространства имен 
в экспортируемое пространство имен 
(P000537).

<%
Обратите внимание, что пространство имен предусмотрено директивой импорта  не является
То же самое, что и пространство имен, импортируемое из .
Последнее включает в себя разрешение конфликтов.
и определено ниже в этом разделе.%
?

%
Пусть \NamespaceName{{local}} будет локальным пространством имен $L$
(P000538).
Пусть {I}{1}{m} будут директивами импорта ,
И пусть {L}{1}{m} быть библиотеками
Например,   для всех .
{%
Не является ошибкой иметь многократный импорт одной и той же библиотеки.
?

%
Пусть .

%
 Шаг первый.
На первом этапе мы вычисляем пространство имен, полученное из каждой импортируемой библиотеки.
соответственно каждой группе библиотек, импортируемых с одинаковым префиксом,
{{one},i}.

%
 Шаг первый для импорта без приставки
Когда  не имеет префикса, \NamespaceName{{one},i}
из пространства имен, предусмотренного директивой импорта 
Удаление всех привязок к имени, базовое имя которого
то же самое, что и базовое имя декларации верхнего уровня в ,
или имя основания которого является префиксом импортной директивы в .

{%
Этот шаг гарантирует, что \SHOW{} и \HIDE{} принимаются во внимание директивы,
Префикс импорта, а также местная декларация
Импортное имя.%
?


%
 Шаг первый для префиксированного импорта
На этом этапе мы решаем конфликты имен.
Среди названий, импортируемых с тем же приставкой.
Когда  имеет приставку ,
давать \List{I'}{1}{k} является подсписком {I}{1}{m} с префиксом ,
И пусть  L'{1}{k} - соответствующие библиотеки.
{(который является подсписком {L}{1}{m})}.
Пусть \NamespaceName{{экспорт},j}, ,
быть пространством имен, предусмотренным директивой импорта 
{, которое \SHOW{} и {} в расчете)}.

%
Когда  является неотложенным импортом:
Пусть {{prefix},i} будет пространством имен, полученным при применении
Конфликт сливается в
 NS}}{\metavar{exported},1}{{exported},k}
().

%
Когда  Отсроченный импорт:
В данном случае  1.
{(в противном случае произошла бы ошибка времени компиляции)},
и \NamespaceName{{prefix},i} - пространство имен, полученное
Добавление связывания имени 
неявно вызванное объявление функции с подписью

{{экспортировано},1}.

%
Тогда \NamespaceName{{one},i} является пространством имен, которое имеет
одно связывание, которое отображает  в \NamespaceName{\metavar{prefix},i}.

%
В этой ситуации мы говорим, что \NamespaceName{\metavar{one},i} является
,
Потому что он отображает префикс библиотеки в пространстве имен.
{%
Пространство имен префикса не является объектом Дарта.
Это просто устройство, которое используется для управления
Выражения формы {qualifiedName} и что они означают
во время статического анализа.
Следовательно, любая попытка использовать префикс пространства имен в качестве объекта
Ошибка времени компиляции, например,
Это не может быть результатом оценки выражения.

Обратите внимание, что если  и  являются такими, что
 и  Оба имеют одинаковую приставку 
затем
${{one},l} = {{one},q} =
{{one},l}  {{one},q}$%
?


%
 Шаг второй.
На втором этапе мы решаем конфликты на высшем уровне.
среди пространств имен, полученных на первом этапе.

%
 , \NamespaceName{{import}}
Это результат применения конфликта, сливающегося с пространствами имен.
 NS}}{\metavar{one},1}{<<< P000992>>>{one},m}.

%
 быть набором заявлений о продлении
()
со следующими членами:
Декларация о продлении  является членом 
Если есть  такой, что
 находится в пространстве имен, предусмотренном директивой импорта 
{%
(обратите внимание, что это {} и {} с учетом,
и включает в себя расширения, импортированные как без, так и с префиксом
?
 В действующей библиотеке не указан
{(что может иметь место, если он импортирует сам)}.
Пусть {{расширения}} будет пространством имен, которое
для каждого расширения   На карте новое имя .

{%
\NamespaceName{\metavar{расширения}} дает новое имя
предоставление имплицитного доступа к каждому расширению, экспортируемому импортируемой библиотекой;
и не удаляется \HIDE{{} или ,
Даже те, которые не могут быть доступны, используя их заявленное имя,
Из-за столкновения имени.%
?

%
  тогда
$\NamespaceName{\metavar{local}}\cup
\NamespaceName{{import}}\cup
\NamespaceName{{расширения}}$.

%
Пусть .
Если  импортируется без приставки,
{именное пространство импортировано из}{библиотека! Именная область импортируется из
 Конфликт сузил пространство имен
()
\NamespaceName{\metavar{one},i}.
%
Иначе  импортируется с приставкой ,
В этом случае {именное пространство импортировано из} 
Конфликт сузил пространство имен
{{экспортировано},i}.

{%
Имя пользователя:  содержит привязки, экспортируемые ,
За исключением тех, которые удалены комбинаторами пространства имен,
За исключением тех, которые были устранены путем слияния конфликтов.%
?

%
 Библиотека с импортным пространством имен .
Мы говорим, что имя
\IndexCustom{imported by $L$}{imported!name)
Если имя является ключом .
Мы говорим, что декларация является
\IndexCustom{imported by $L$}{imported!declaration}}
если декларация является значением .
Мы говорим, что имя
<\IndexCustom с префиксом $p$}{imported!name, с приставкой}
если имя является ключом {}{p}.
Мы говорим, что декларация является
 с префиксом {%}
Импорт! Декларация с префиксом
если декларацией является значение {}{p}.


{Семантика импорта}


%
 быть директивой импорта, которая относится к URI через строку .
Семантика  Указывается следующее:

%
 Семантика отложенного импортаЕсли  является отложенным импортом с префиксом , связыванием  до

\NamespaceName{\metavar{deferred}}
Присутствующее в библиотеке пространство имен текущей библиотеки .
Пусть \NamespaceName{\metavar{import},i}
пространство имен, импортируемое из библиотеки, указанной в ,
как определено ранее
().
\NamespaceName{\metavar{deferred}} имеет следующие связывания:


 Оригинальное название: P000656 привязанный к
Функция с подписью .
Эта функция возвращает будущее .
При вызове функция вызывает
Немедленный импорт  быть исполнены в будущем,
где  Происходит от  Удаление слова 
Добавить   Комбинаторное предложение.
Выполнение немедленного импорта может привести к неудаче.
по конкретным причинам реализации.
{%
Например,  Библиотека отличается от той, что
что указанный URI упоминается во время компиляции;
или ошибка чтения файла уровня ОС; и т.д.%
?
Мы говорим, что вызов 
\IndexCustom{succeeds}{loadLibrary!succeeds} если $f$ завершается ценностью,
и что призыв
\IndexCustom{fails}{loadLibrary!fails} если $f$ завершается ошибкой.

Для каждой функции верхнего уровня  Оригинальное название: P000755 в
\NamespaceName{{import},i}
Соответствующая функция называется  с такой же подписью, как .
% Эта ошибка может произойти, потому что загрузка является динамическим свойством.
Вызов функции приводит к динамической ошибке, которая возникает до
Оцениваются любые реальные аргументы.
Закрытие функции
()
Это также приводит к динамической ошибке.

Для каждого получателя высшего уровня  Оригинальное название: P000757 в
\NamespaceName{{import},i}
Соответствующий геттер по имени  с той же подписью, что и .
% Эта ошибка может произойти, потому что загрузка является динамическим свойством.
Вызов геттера приводит к динамической ошибке.

Для каждого наборщика верхнего уровня  Оригинальное название: P000660 в
\NamespaceName{{import},i}
Соответствующий сеттер по имени  с той же подписью, что и .
% Эта ошибка может произойти, потому что загрузка является динамическим свойством.
Вызов установщика приводит к динамической ошибке, которая возникает до
Фактический аргумент оценивается.

Для каждого класса, смесителя, перечня и типа алиаса декларация под названием  в
\NamespaceName{{import},i}
соответствующий геттер с именем  с возвратным типом .
% Эта ошибка может произойти, потому что загрузка является динамическим свойством.
Вызов геттера приводит к динамической ошибке.


{%
Для того, чтобы члены импортной библиотеки могли
\NamespaceName{{отложенный}
чтобы обеспечить использование членов, которые еще не были загружены
может быть решена нормально и имеет четко определенное поведение;
Что вызывает ошибки.%
?

%
При обращении к  Успешно,
имя <<p000096> отображается в неотложенный префикс run-time namespace
\NamespaceName{\metavar{loaded}}
с привязками, как описано ниже, для немедленного импорта.
Кроме того, \NamespaceName{\metavar{loaded}}
Карты  для функции с той же подписью, что и раньше.
Таким образом, можно ссылаться  Опять же,
Которая всегда будет успешной.
Если вызов не удался, библиотека не была загружена.
и у вас есть возможность вызвать  Опять.
Неоднократный звонок по адресу  Успехи будут разными,
как описано ниже.

{%
Обратите внимание, что это ошибка времени компиляции для отложенного префикса.
используется более чем в одном импорте;
что означает обновление связывания не вмешивается в другие импортные операции.%
?

%
Эффект повторного вызова  выглядит следующим образом:



При повторном вызове  Уже удалось,
Неоднократные призывы также увенчались успехом.
В противном случае

При повторном вызове  Не удалось:



Если ошибка вызвана ошибкой компиляции,
Повторное обращение не удается по той же причине.

Если неудача вызвана другими причинами,
Повторяющиеся призывы ведут себя так, как будто никакого предыдущего призыва не было.



{%
Другими словами, успешный  гарантируют, что
Доступ к импорту можно получить через префикс импорта.
Если ссылка на  Инициируется после
Успешно завершился очередной призыв,
Это также гарантирует успешное завершение (успех является идемпотентным).
Мы не указываем, какой объект возвращает будущее.%
?


%
{Семантика немедленного импорта}
Когда  является немедленным импортом с библиотекой URI, ,
в виде строки :
 библиотека, полученная из исходного кода, обозначенного .
Затем мы говорим, что URI  Обозначает библиотеку .
Весь импорт и экспорт одного и того же URI в программе Dart обозначает
той же библиотеки,
Импорт и экспорт различных URI обозначают различные библиотеки.

%
Пусть {{import}, i} будет пространством имен, импортированным из .
Пространство имен времени выполнения {i} будет иметь следующие привязки:



Для каждой функции верхнего уровня, getter или setter  Название: P000113 в
\NamespaceName{{import}, i},
привязка от  к указанной функции, геттеру или сеттеру.

Для каждого имени  в \NamespaceName{\metavar{import}, i}
который связан с классом, смешиванием, перечислением или объявлением псевдонима типа
Введение типа ,
a привязка от  к составленному представлению .


%
Если  имеет приставку ,
Пространство имен времени выполнения текущей библиотеки
Карты  в неотложенное пространство имен во время выполнения \NamespaceName{p}
Содержит отображения {i}.

{%
\NamespaceName{p} могут иметь дополнительные карты
Потому что может быть несколько импортов с одним и тем же префиксом.
До тех пор, пока все они являются немедленными.%
?

%
В противном случае, когда  не имеет префикса,
Пространство имен библиотеки времени выполнения текущей библиотеки
Содержит каждое отображение в {i}.



{Экспорт}


%
Библиотека  экспортирует пространство имен (), что означает, что
Заявления в пространстве имен доступны другим библиотекам.
Если вы решили импортировать  ().
Пространство имен, которое  Экспорт известен как
{экспортируемое пространство имен}{библиотека!экспортируемое пространство имен}.

%
Библиотека всегда экспортирует все имена и декларации в публичное пространство имен.
Кроме того, библиотека может реэкспортировать дополнительные библиотеки.
, часто называемый просто :


<libraryExport> ::= <metadata> \EXPORT{{} <configurableUri> <combinator>> "


%
Интерпретация настраиваемых URI описана в другом месте.
(P000546).
Экспорт определяет URI 
где находится декларация экспортируемой библиотеки.
Ошибка времени компиляции, если указанный URI
Это не относится к библиотечной декларации.

%
The

Библиотека определяется следующим образом, в два этапа.
 Будьте библиотекой,давать  E}{1}{m} являются экспортными директивами ,
И пусть {L}{1}{m) - библиотеки
  для всех .
Пусть \NamespaceName{\metavar{public}} будет публичным пространством имен $L$
(P000547).

{%
Обратите внимание, что приватные имена и префиксы импорта отсутствуют.
в \NamespaceName{{public}}%
?

%
На первом этапе вычисляется пространство имен, предоставляемое каждой экспортируемой библиотекой:
Для каждого 
\NamespaceName{\metavar{экспорт},i}
Пространство имен, полученное при применении
Комбинаторы пространства имен  то
Произошло это из-за того, что «Perfect»
(),
и удаление каждого связывания имени  такой, что
\NamespaceName{{public}} определяется по адресу ,
 и  Имеют одинаковое имя.
 Потому что местные декларации будут тенью для реэкспорта.

%
The
{namespace re-exported from}{library!namespace re-exported from}
 \NamespaceName{\metavar{экспортировано},i}.

%
На втором этапе вычисляется экспортируемое пространство имен .
Пусть \NamespaceName{\metavar{merged}} будет результатом применения
слияние конфликтов
()
 NS}}{\metavar{экспортировано},1}{<<< P001107>>>{экспорт},m}.
Ошибка времени компиляции возникает, если имя
\NamespaceName{{merged}}
является конфликтным
().

<%
Это правило является более строгим, чем соответствующее правило для импорта:
Когда две импортные декларации имеют столкновение названий, это всего лишь ошибка.
{использовать} противоречивое название,
Не является ошибкой то, что слово «столкновение» существует.
С экспортированными именами, это ошибка, что столкновение имени существует.
Причина этого различия заключается в том, что
Конфликт может бесшумно разорвать импортеров нынешней библиотеки , - говорится в сообщении.
Если мы будем использовать тот же подход к экспорту, что и к импорту:
Если библиотека  Импорт 
и использует имя  который реэкспортируется  из 
Затем добавьте декларацию под названием 
В реэкспортируемую библиотеку
будет использовать  в $L'$ Ошибка,
и хранителей  может быть не в состоянии изменить .%
?

%
The
{экспортируемое пространство имен}{библиотека!экспортируемое пространство имен}
 это
.

%
Мы говорим, что имя 
если имя находится в экспортируемом пространстве имен библиотеки.
Мы говорим, что заявление 
если декларация находится в экспортируемом пространстве имен библиотеки.

%
Для данного ,
Мы говорим, что 

, а также 

\NamespaceName{\metavar{экспортировано},i}.
Когда не может возникнуть путаницы, мы можем просто сказать:
  , или
 \NoIndex{re-exports} \NamespaceName{{exported},i}.


 Комбинаторы пространства имен


%
Импорт () и экспорт () полагаться

Для настройки пространства имен
()
Управляйте столкновениями имен.
Поддерживаемыми комбинаторами пространства имен являются \SHOW{} и \HIDE.


<комбинатор> ::= \SHOW{} <identifierList> | \HIDE{} <identifierList>

<identifierList> ::= <identifier> (',' <identifier>) *


%
Мы определяем несколько операций, которые вычисляют пространства имен.
The
{союз двух пространств имен}{namespace!union},
\IndexCustom{>>>{%)
$\cup$@$\NamespaceName{a} \cup \NamespaceName{b}$,
Определяется, когда каждый ключ, где оба определены, отображается на одно и то же значение.
Этот союз - пространство имен, которое отображает
Каждый ключ  \NamespaceName{a}
соответствующего значения \Namespace{a}{n}и каждый ключ  \NamespaceName{b}
{b}{n}.

<%
Обратите внимание, что это <{является} пространством имен, потому что союз только определен.
когда отображение любого заданного ключа является однозначным.%
?

%
Функция
\IndexCustom{{l}{{a}}}
hide(l,NSa)@\Hide{l}{{}}}
Выберите список идентификаторов  и пространство имен \NamespaceName{a},
и создает пространство имен \NamespaceName{b} Это то, что
идентично  Кроме того, что для каждого идентификатора  в ,
{b} не определено в  .

%
Функция
\IndexCustom{{l}{{a}}}
show(l,NSa)@\Show{l}{{}}}
Возьмите список идентификаторов  и пространство имен {a},
и создает пространство имен  это
Карты каждого идентификатора  
где \NamespaceName{a}
то {a}{n}.
Кроме того, для каждого идентификатора  
где \NamespaceName{a} определено в ,
\NamespaceName{b} карты \code{\id=} до \Namespace{a}{}.
Наконец, {b} не определено во всех других именах.

%
Пусть \List{C}{1}{n} будет последовательностью комбинаторных оговорок.
 Они будут поступать из импортной или экспортной директивы. ?
Результатом
{применение комбинаторных оговорок}{%
Комбинаторные предложения! Приложение для пространства имен
в заданное пространство имен \NamespaceName{\metavar{start}}
является пространством имен \NamespaceName{\metavar{end}},
который рассчитывается следующим образом.

%
Пусть \NamespaceName{0} будет \NamespaceName{\metavar{старт}}.
Для каждого комбинатора , :
Если  имеет форму \code{\SHOW{}  Тогда пусть
.
Если  имеет форму \code{\HIDE{}  Тогда пусть
.
Тогда \NamespaceName{\metavar{end}} является \NamespaceName{n}.

{%
\NamespaceName{\metavar{end}} всегда соглашается с
\NamespaceName{\metavar{старт}} везде, где оба определены,
и набор имен, где определен \NamespaceName{\metavar{end}}
всегда является подмножеством того из {{старт}}.
В этом смысле это процедура {сужение}.%
?

{%
Обратите внимание, что можно использовать \HIDE{} или {} на идентификаторе
которого нет в названном пространстве.%
?
{%
Это позволяет избежать ситуаций, когда, например,
удаление декларации из библиотеки  вызвать
Библиотека, которая импортирует %
?


 Слияние пространств имен.


%
В этом разделе мы определяем операцию на пространствах имен.
которые устраняют определенные привязки имен в случае столкновений имен,
и отслеживает такие столкновения имен, используя особое значение,
.

%
Когда имя  отображается на {} по пространству имен,
Мы говорим:  это
\IndexCustom{conflicted}{конфликтное название}.
Поиск противоречивого имени будет успешным
Нормальные правила для пространств имен и лексического определения,
Но это ошибка времени компиляции в месте, где используется имя.
Это название является противоречивым.

%
Пусть \List{{NS}}{1}{m} будет списком пространств имен
{(список может содержать или не содержать дубликатов)}.
The
{конфликтное слияние}{именное пространство!конфликтное слияние}
 NS}}{1}{m} является
$\NamespaceName{{конфликт}}\cup
\NamespaceName{\metavar{сокращенно},1} 
\ldots\ \NamespaceName{\metavar{narrowed},m}$,
где пространство имен \NamespaceName{{конфликт}} и
\NamespaceName{\metavar{narrowed},i}, ,
Указывается следующее:

%\NamespaceName{\metavar{конфликт}} пустой, за исключением
Обязательства, упомянутые ниже.
Для каждого ,
\NamespaceName{\metavar{суженный},i} идентичен \NamespaceName{i},
За исключением того, что первое не определено в каждом имени  где
Последняя определяется,
и одно из следующих условий:


<<<p001212> Существует  Название: P000186 такой, что
{j} определено в ,
 и  иметь одинаковые базовые имена,
{i}{n} — декларация в системной библиотеке,
и {j}{n'} является декларацией в несистемной библиотеке.
<{%
Заявление из несистемной библиотеки теней
Заявления из системных библиотек.%
?
 В противном случае существует  Название: P000191 такой, что
\NamespaceName{j} определено в ,
 и  иметь одинаковые базовые имена,
\Namespace{i}{n} и {j}{n'} не являются одним и тем же заявлением
И не одноимённое пространство
а не геттер и сеттер, объявленные в той же библиотеке,
Ни один из них, ни оба
\Namespace{i}{n} и \Namespace{j}{n''
Деклараций в системной библиотеке.
%
{%
С двумя разными объявлениями с одинаковым именем,
Оба они устранены,
За исключением случаев, когда это геттер и сеттер из одной и той же библиотеки.
Обратите внимание, что \Namespace{i}{n} и \Namespace{j}{n'')
оба должны быть объявлениями или оба должны быть пространствами имен;
потому что слияние конфликтов применяется к пространствам имен, где
импортные декларации, противоречащие импортному префиксу
Они уже ликвидированы.
Когда они оба пространства имен, нет никакого конфликта.
Потому что это одно и то же имя.%
?

В этой ситуации  Соответствует 
\NamespaceName{{конфликт}}
{%
Мы можем поменять все имена и снова использовать правило.
 Конфликт.%
?


{%
Короче говоря, слияние конфликтов берет список пространств имен и производит
Пространство имен, которое является объединением нескольких несвязанных пространств имен
(Поскольку все названия были уничтожены):
Пространство имен конфликтов, которое записывает все неразрешенные столкновения имен.
и список пространств имен, где были удалены сталкивающиеся имена.
Некоторые столкновения названий были решены путем предпочтения.
Заявления из несистемных библиотек
по заявлениям системных библиотек.%
?

%
Полезно уметь ссылаться на результат сужения.
Пусть \List{{NS}}{1}{m} будет списком пространств имен, $i \in 1 .. m$,
и рассматривать слияние конфликта, как указано выше.
The
{конфликт сузил пространство имен}{namespace!conflict narrowed}
Из {i} получается \NamespaceName{{narrowed},i}.


 Части.


%
Библиотеку можно разделить на ,
Каждый из них может храниться в отдельном месте.
Библиотека идентифицирует свои части, перечисляя их через {} директивы.

%
 Укажите URI, где
а Блок компиляции Dart, который должен быть включен в текущую библиотеку
Может быть найден.


<partDirective> ::= <metadata> {{i}>> "

<part Header> ::= <metadata> {} \OF{} (<dottedIdentifierList> | <uri>); "

<parteclaration> ::=
<partHeader> (<metadata> <topLevelDeclaration>)* <EOF>


%
Часть содержит строку, полученную из {partDeclaration}.

{%
Таким образом, можно сказать, что {partDeclaration} является начальным символом грамматики.
Как описано в статье.~\ref{librariesAndScripts}.%
?

%
 начинается с  < сопровождаемый
Название библиотеки, к которой принадлежит часть,
или {uri}, обозначающий указанную библиотеку.
Декларация части состоит из заголовка части, за которым следуетПоследовательность деклараций высшего уровня.

%
Составление части директивы формы \code{\PART{) ;
Это приводит к тому, что система Дарта пытается собрать содержимое
URI - значение .
Объявления верхнего уровня на этом URI затем компилируются компилятором Dart.
в рамках действующей библиотеки.
Ошибка времени компиляции, если содержимое URI не
декларация действительной части.
Это ошибка времени компиляции, если указанная часть декларации  имена
библиотека, отличная от текущей библиотеки, в которой  принадлежит.

%
Ошибка времени компиляции, если библиотека содержит
Две части директивы с одинаковым URI.

%
Мы говорим, что библиотека  является  Библиотека  если
все нижеследующее верно (, :


  и  — одна и та же библиотека.
  импортирует или экспортирует библиотеку  и  Доступно из .


%
 Будьте библиотекой, пусть  Будь Ури,
 и  Отдельные библиотеки, доступные из .
Ошибка времени компиляции, если  и  Оба содержат
директива с URI .

{%
В частности, ошибкой является использование одной и той же части дважды.
В той же программе
().
Обратите внимание, что относительный URI интерпретируется как относительный по отношению к местоположению URI.
библиотека (), что означает  и  могут оба
имеют часть, идентифицированную по , но они не одинаковы
URI, если  и  имеют одинаковое месторасположение.%
?


{Сценарии}


%
 Библиотека, экспортировавшая пространство имен () включает
Декларация функций верхнего уровня 
Для этого требуется либо ноль, либо два аргумента.

Сценарий  Выполняется следующим образом:

%
Во-первых,  составляется как библиотека, как указано выше.
Функция верхнего уровня, определяемая как 
в экспортируемом пространстве имен  (перенаправлено с «P000560»)
следующим образом:
Если  может быть вызвана двумя позиционными аргументами,
Он ссылается на следующие два фактических аргумента:


 Объект, тип времени выполнения которого реализует .
 Объект, указанный в текущем Изоляторе  был создан,
Например, посредством ссылки  что породило ,
или нулевого объекта (), если такой объект не был поставлен.


%
Если  не может быть вызван двумя аргументами,
Но его можно назвать одним позиционным аргументом.
он вызывается объектом, чей тип времени выполнения реализует 
Как единственный аргумент.
Если  не может быть вызван одним или двумя аргументами,
На него ссылаются без аргументов.

{%
Обратите внимание, что если  требует более двух аргументов,
Библиотека не считается скриптом.%
?

{%
Программа Dart обычно выполняется путем выполнения сценария.%
?

%
Ошибка времени компиляции, если объем экспорта библиотеки содержит декларацию
Называется она «P000687», а библиотека — не сценарий.
{%
Это ограничение гарантирует, что все топ-уровни  заявления
Введите главную функцию сценария, чтобы не было
Геттер верхнего уровня или поле
не может быть функцией, требующей более двух
аргументы. Ограничение позволяет инструментам рано выходить из строя на недействительном 
без необходимости знать, будет ли библиотека использоваться в качестве входа
Смысл программы Дарт. Возможно, это ограничение будет снято.
в будущем.%
? УРИ


%
URI определяются с помощью струнных букв:


<uri>:= <stringLiteral>

<configurableUri> ::= <uri> <конфигурация Ури*

<configurationUri> ::= \IF{} '(' <uriTest> ''

<uriTest> ::= <dottedIdentifierList> ('==' <stringLiteral>)?


%
Это ошибка времени компиляции, если строка буквально описывает URI.
или строка, которая используется в 
содержит интерполяцию строк.

%
  в форме
\code{{uri)   
{указывает URI}{указывает URI} следующим образом:



Пусть  будет \metavar{uri}.

Для каждой из следующих конфигураций URI формы
\code{{) () ,
В порядке источника сделайте следующее.



Если < \code{\metavar{ids}}
без  Оговорка, это
эквивалентно \code{\metavar{ids} == "истинное"}.

Если $\metavar{test}_i$ \code{\metavar{ids} == \metavar{string}}
Затем создайте строку \metavar{key} из \metavar{ids}
путем конкатенирования идентификаторов и точек,
Опустить любые пробелы между ними, которые могут возникнуть в источнике.

Смотрите {key} в доступном разделе
.
{%
Среда компиляции предоставляется платформой.
Он отображает некоторые струнные ключи для значений строк,
могут быть получены программным путем с использованием
 Конструктор.
Инструменты могут выбирать только некоторые части среды компиляции.
Доступен для выбора конфигурации URIs.%
?

Если окружение содержит запись для {ключ} и
ассоциированная стоимость равна, как постоянная струнная стоимость, значению
строка буквальная \metavar{string},
  Перестаньте повторять URI конфигурации.

В противном случае переходите к следующей конфигурации URI.


URI, указанный в , является .


%
Эта спецификация не обсуждает интерпретацию URI.
со следующими исключениями.

{%
Интерпретация URI в основном оставлена на
окружающей вычислительной среды.
Например, если Dart работает в веб-браузере,
Этот браузер, скорее всего, интерпретирует некоторые URI.
Хотя может показаться привлекательным указать, скажем, что
URI интерпретируются в отношении стандарта, такого как IETF RFC 3986.
На практике это обычно зависит от браузера.
На него нельзя положиться.%
?

%
URI формы  толкуется как
Ссылка на системную библиотеку () .

%
URI формы  интерпретируется в
конкретного способа реализации.

{%
Цель состоит в том, чтобы во время разработки программисты Dart могли полагаться на
менеджер пакетов для поиска элементов своей программы.%
?

%
В противном случае любой относительный URI интерпретируется как относительный.
Местонахождение нынешней библиотеки.
Вся дальнейшая интерпретация URI зависит от реализации.

<<<p001293>> Это означает, что он зависит от времени выполнения.


 Типы


%
Dart поддерживает статическую типизацию на основе типов интерфейсов.

{%
Типовая система является звуковой в том смысле, что
если переменная типа  относится к объекту типа  во время бега,
 Это подтип .
Другими словами, содержимое кучи удовлетворяет ожиданиям.
выражается статической типизацией.

Однако параметры типа являются ковариантными.
(например, \SubtypeNE{\code{List<int>}}{\code{List<num>}})Это означает, что некоторые операции подвергаются динамическим проверкам типа.
(например, , который будет бросаться во время бега)
Если  объявленный тип ,
На самом деле это P000699.
Таким образом, программа может быть свободна от ошибок времени компиляции.
И он все еще может повлечь за собой ошибку типа во время выполнения.

Этот выбор был сделан намеренно в первые дни Дарта.
(И это, несомненно, спорно).
Представляет собой компромисс, где
Потенциал для ошибок типа времени выполнения - это стоимость.
Преимуществом являются более простые программы.
В последние годы Dart развивался, чтобы иметь все больше и больше безопасности статического типа.
Например, нулевая безопасность, и эта тенденция, вероятно, продолжится.
?


 Статические типы


%
Аннотации типов могут встречаться в переменных декларациях (),
включая формальные параметры (),
в возвратных типах функций (),
и в границах переменных типа ().
Аннотации типов используются во время статической проверки и при запуске программ.
Типы определяются с использованием следующих правил грамматики.

{%
{typeIdentifier} - идентификатор, который может быть именем типа,
То есть,
обозначает  Идентификатор, который не является  BUILT В  Идентификатор
()%
?

{%
Нетерминалы с названиями формы \synt{\ldots{}NotFunction}
Выведите термины, которые являются типами, которые не являются типами функций.
Обратите внимание, что в нем {does} выводится тип ,
который сам по себе не является функцией,
Но это наименьшая верхняя граница всех типов функций.
?



<type> ::= <functionType> '?'?
 <typeNotFunction>

<typeNotVoid> ::= <functionType> '?'
 <typeNotVoidNotFunction> '?'

<typeNotFunction> ::= <
 <typeNotVoidNotFunction>>?

<typeNotVoidNotFunction> ::= <typeName> <typeArguments>
 (<typeIdentifier>>)? \FUNCTION{}

<typeName> ::= <typeIdentifier> ('.' <typeIdentifier>)?

<typeArguments> ::= <<typeList>> "

<typeList> ::= <type> (',' <type>)*

<typeNotVoidNotFunctionList>:= <
<typeNotVoidNotFunction> (',' <typeNotVoidNotFunction>) *

<functionType> ::= <functionTypeTails>
 <typeNotFunction> <functionTypeTails>

<функция TypeTails>::= <functionTypeTail> '?'? <functionTypeTails>
 <functionTypeTail>

<функция TypeTail ::= \FUNCTION{} <typeParameters> <parameterTypeList>

<parameterTypeList> ::= ''' "
<<normalParameterTypes>> <optionalParameterTypes>> "
 <<типы нормального параметра>> "
 <("Необязательные параметры">> "

<normalParameterTypes> ::= <
<normalParameterType> (>normalParameterType>)*

<normalParameterType> ::= <metadata> <typedIdentifier
 <metadata> <type>

<optionalParameterTypes> ::= <optionalPositionalParameterTypes>
 <namedParameterTypes>

<optionalPositionalParameterTypes> ::= '[' <normalParameterTypes> ','?' "

<namedParameterTypes> ::=
\{> <namedParameterType> ('', <namedParameterType>''), \} "

<namedParameterType> ::=
<metadata>  <typedIdentifier>

<typedIdentifier> ::= <type> <identifier>


%
А. Реализация Dart должна обеспечивать статическую проверку, которая обнаруживает и сообщает
именно те ситуации, которые эта спецификация идентифицирует как ошибки времени компиляции;
И только эти ситуации.
Аналогично, статическая шашка должна выдавать статические предупреждения.
По крайней мере, для ситуаций, указанных как таковые в данной спецификации.

{%
Ничто не исключает дополнительных инструментов, которые реализуют альтернативный статический анализ.
(например, толкование аннотаций существующего типа в звуковой форме)
такие как невариантные дженерики или вывод о дисперсии на основе декларации
из фактических деклараций.Однако использование этих инструментов не должно препятствовать
Успешная компиляция и выполнение кода Dart.%
?

%
Тип   Если [f]:



 Имеет форму  или форму \code{\metavar{префикс}.\id},
И это не означает объявление типа.

 обозначает переменную типа,
Это происходит в сигнатуре или теле статического элемента.

 является параметризованным типом формы ,
 является неполноценным,
 не является общим типом,
 является родовым типом,
Об этом сообщает  Типовые параметры и ,
 Для некоторых .

 является функциональным типом формы

 \FUNCTION{}<$X_1\ \EXTENDS\ B_1, \ldots,\ X_m\ \EXTENDS\ B_m$>




или в форме

 \FUNCTION{}$X_1\ \EXTENDS\ B_1, \ldots,\ X_m\ \EXTENDS\ B_m$>

 $T_1\ x_1, \ldots,\ T_k\ x_k,\ $\{$T_{k+1}\ x_{k+1}, \ldots,\ T_n\ x_n$\}]


где каждый  который не является поименованным параметром, может быть опущен;
 является деформированным для некоторых ,
 Для некоторых из них .

 Деклараций, которые были импортированы из нескольких пунктов импорта.


%
Любое появление неправильного типа в библиотеке является ошибкой времени компиляции.

%
Тип  \IndexCustom{deferred}{type!deferred}
Если он имеет форму   Это отложенный префикс.
Ошибка времени компиляции для использования отложенного типа
в аннотации типа, тесте типа, отливке типа или в качестве параметра типа.
Тем не менее, все другие ошибки времени компиляции должны быть выпущены.
Все отложенные библиотеки были успешно загружены.

% Теперь, когда передается дженерик, p.T также должен рассматриваться как динамический -
В противном случае мы должны немедленно потерпеть неудачу. Где мы это скажем? И как это делается
% это согласуется с идеей о том, что как объект типа он терпит неудачу? Стоит ли говорить, что
%-ный аксессуар на p возвращает динамику вместо отказа? Различают ли их использование
% в конструкторе против его использования в аннотации? Не то чтобы мы оценивали тип
% объектов в конструкторских аргах — они не могут представлять параметризованные типы.


 Тип поощрения


%
Система статического типа приписывает статический тип каждому выражению.
В некоторых случаях тип локальной переменной
{(который может быть формальным параметром)}
могут быть повышены из заявленного типа, исходя из управляющего потока.

%
Мы говорим, что переменная  Известно, что тип 
Каждый раз, когда мы разрешаем рекламировать тип .
Точные обстоятельства, когда разрешено продвижение по типу, приведены в
соответствующие разделы спецификации
(,  и .

%
Продвижение типа для переменной  Допускается только тогда, когда мы можем сделать вывод, что
Такое продвижение действительно на основе анализа определенных булевых выражений.
В таких случаях мы говорим, что
Булево выражение  Показывает, что  имеет тип .
Как правило, для всех переменных  и типы , булево выражение
Не показывает, что  имеет тип .
Те ситуации, когда выражение действительно показывает, что переменная имеет тип
прямо упомянутых в соответствующих разделах настоящей спецификации
( и ).


 Система динамического типа


%  Когда имя класса появляется как тип, %
% это имя обозначает интерфейс класса. Так что мы можем сказать
% что динамический тип экземпляра является негенерическим классом или
% генерическая инстанциация генерического класса.

%
Пусть  будет примером.  Класс, который определяется
Для ситуации, когда  был получен в качестве нового примера
(,
, , , .

{%
В частности, динамический тип экземпляра никогда не меняется.
Иногда указывается только, что данный класс реализует
определенный тип, например, для списка буквальный.
В этих случаях динамический тип зависит от реализации.
За исключением, конечно, того, что упомянутое ограничение суперинтерфейса должно быть удовлетворено.
?

%
Динамические типы запущенной программы Дарта эквивалентны
Статические типы в отношении субтипирования.

{%
Определенные динамические проверки типа выполняются во время выполнения.
(,
,
,
,
,
,
,
,
,
,
,
,
,
,
,
.
Как указано в этих местах,
Эти динамические проверки основаны на динамических типах экземпляров.
и фактические виды деклараций
()%
?

%
Когда типы уточнены как экземпляры встроенного класса ,
Эти объекты перекрывают  оператор
Унаследовано от  класс, так что
 Объекты равны по оператору \lit{=======================================================================================================================================================================================================================================================
f соответствующие типы являются подтипами друг друга.

{%
Например,  объекты для типов
\DYNAMIC{{} и  равны между собой.
 Вы должны оценить до .
На встроенную функцию не накладываются ограничения ,
 может быть {} или .

Точно так же  Примеры заявлений о различных типах псевдонимов
Объявление имени для одного и того же типа функций равно: %
?



%
{%
Случаи  могут быть получены различными способами,
Например, используя рефлексию,
Читать далее «P000712» Для объекта,
или путем оценки выражения .%
?

%
Выражение представляет собой  [тип буквальный], если оно является идентификатором,
квалифицированный идентификатор,
который обозначает класс, смесь, enum или тип alias Declaration, или это
идентификатор, обозначающий параметр типа общего класса или функции.
Это {постоянный тип буквальный} если он не обозначает параметр типа,
и не квалифицируется отложенным префиксом.
{%
Постоянная буквальная форма - это постоянное выражение ().%
?


 Типовые алиазы


%
 Назовите имя какого-нибудь типа,
или для отображения от аргументов типа к типам.

{%
Для такого заявления обычно используется фраза "a typedef",
Из-за заметного появления токена %
?


<typeAlias> ::= \gnewline{}
{} <typeIdentifier> <typeParameters>? '=' <type>'; "
 < <functionTypeAlias>

<functionTypeAlias> ::= <functionPrefix> <formalParameterPart>; "

<functionPrefix> ::= <type>? <identifier>


%
Рассмотрим заявление типа alias  в форме


\code{{) \id\TypeParametersStd> = $T$;


Объявлено в библиотеке .
Эффект  Предлагается ввести  в библиотечный объем .
Когда 
{(где тип псевдонима не является общим)}
 Связано с .
Когда 
\commentary{(где тип псевдонима является общим)}
 привязан к отображению из списка аргументов типа
{U}{1}{}
к типу
.

%
В предположении, что {X}{1}{s} являются такими типами, что, для всех ,
Ошибка времени компиляции  не является регулярным,
и это ошибка времени компиляции, если какой-либо тип встречается в  Он не очень хорошо ограничен.

{%
Это означает, что границы, объявленные для
Формальные параметры типа общего псевдонима типа
должны быть такими, чтобы, когда они удовлетворены,
границы, относящиеся к телу, также удовлетворены;
и тип, происходящий как субтермин тела, может нарушать его границы;
но только в том случае, если это правильный сверхограниченный тип.
?

%
Кроме того,
 быть типами
 быть параметризованным типом 
в месте, где  обозначает .
Ошибка времени компиляции .
Ошибка времени компиляции  не является хорошо ограниченным
().

%
По историческим причинам тип псевдонима может иметь еще две формы.
производное использование {функцияТипАлиас}.
Пусть {S?} будет термином, который является пустым или полученным из {тип},
и аналогично для {T_j?} для любого .
Затем две более старые формы определяются следующим образом:

%
Типовой псевдоним формы


\code{{)  \id<\TypeParametersStd>()





рассматривается как


\code{{) <

<<< 001393>


%
Типовой псевдоним формы


\code{{)  <\TypeParametersStd>()





рассматривается как


\code{{} <




%
В этих правилах для каждого $j$, если $T_j?$ пустой
затем {T_j} - ,
иначе  .

{%
Это означает, что старые формы позволяют опустить тип параметра.
В этом случае оно считается ,
Но имя параметра не может быть опущено.
Такое поведение подвержено ошибкам,
и, следовательно, рекомендуется более новая форма (с {=}).
Например, заявление \code{\TYPEDEF{} (int);
Уточняется, что  обозначает тип
\code{{) \FUNCTION(\DYNAMIC)
документируется, но технически игнорируется,
Этот параметр имеет имя .
Очень вероятно, что читатель неправильно это прочитает.
Предположим, что  Тип параметра.%
?

%
Это ошибка времени компиляции, если значение по умолчанию указано для
формальный параметр в этих старых формах,
или если формальный параметр имеет модификатор .

<{%
Обратите внимание, что старые формы могут определять только типы функций.
Они не могут определить тип общей функции.
Когда такой тип псевдонима имеет параметры типа,
Она всегда выражает семейство негенерических типов функций.
Эти ограничения существуют, потому что синтаксис был определен.
До добавления общих функций в Dart.%
?

%
Пусть  быть типом alias заявление формы


\code{{) $F$<\TypeParametersStd> = $T$;


Если  или  Для некоторых  является или содержит
а \synt{typeName}} $G$ обозначение заявления типа alias ,
Мы говорим, что
 Зависит от  [тип псевдонима]! зависимость.

%%TODO(Eernst): Расслабление, когда возникает зависимость типа псевдонима
% % через , но не через , является более разрешительным и позволяет
% % тип псевдонимов, которые должны быть объявлены, что будет ошибкой с вышеуказанным правилом.
% % Очень вероятно, что можно использовать более расслабленное правило. Это а
% возможного улучшения в будущем и непреодолимых изменений.

%
 быть типом псевдонима декларации,И пусть  Промежуточное закрытие
Заявления типа alias, которые  Зависит от.
Ошибка времени компиляции возникает, если .

{%
Другими словами, это ошибка для объявления псевдонима типа.
зависеть от себя, прямо или косвенно.%
?

{%
Эта ошибка может возникнуть, когда аргументы типа
Опущены в программе, но добавляются при статическом анализе
Через инстанциацию связываться
()
или с помощью типоразмера,
который будет указан позднее
()%
?

%
Когда  является типом alias декларации формы


\code{\TYPEDEF{} $F$<\TypeParametersStd> = $T$;


Параметризованный тип  в форме
< U}{1}{}>}
в области, где  Выберите 
{alias расширяется в один шаг до}{%
Типовой псевдоним! alias расширяется в один шаг
.

{%
Обратите внимание, что  может быть равно нулю, в этом случае  не является общим,
и мы просто заменяем имя псевдонима на его тело.
?

%
Если  Это тип, который мы можем многократно заменять каждый субтермин 
который является параметризованным типом, применяющим псевдоним типа к некоторым аргументам типа
Расширение в один шаг
{(в том числе необычный случай, когда нет аргументов типа)}.
Когда дальнейшие шаги невозможны, мы говорим, что полученный тип
{расширение транзитивного псевдонима}{%
Типовой псевдоним! Промежуточное расширение псевдонима
.

{%
Обратите внимание, что транзитивное расширение псевдонима существует,
потому что это ошибка для объявления псевдонима типа
Зависить от самого себя.
Однако он может быть большим.
Например, \code{\TYPEDEF{}  = Карта  доходность
увеличение экспоненциального размера для \code{$F$<$F$<\ldots$F$<<\ldots{}>{}>>.%
?

%
%
 быть типом alias заявление формы


\code{{) < = ;


И пусть  U - тип формы  или 
в области, где этот термин обозначает .
Предположим, что переходное псевдонимное расширение \code{$F$<\List{ X}{1}{s}>} это
%
Тип одной из форм  или ,
необязательно с последующим \synt{typeArguments},
где  является идентификатором, обозначающим префикс импорта,
  обозначает класс или смесь
{(в частности,  не может быть переменной типа.
Предположим, что  встречается в выражении  в форме
. \id\;\metavar{args}} "
где \metavar{args} происходит от \syntax{<argumentPart>?},
 Это имя
Статический член  соответственно .
Выражение  Затем рассматривается как
. \id\;\metavar{args}} "
соответственно
. . \id\;\metavar{args}}.

{%
Это означает, что можно использовать псевдоним типа
Для вызова статического члена класса или смесителя.
Например: %
?



{%
Обратите внимание, что аргументы типа перешли на  соответственно 
стираются,
Таким образом, полученное выражение может быть правильным статическим вызовом члена.
Если будущая версия Dart позволяет передавать аргументы типа через класс
при обращении к статичному члену,
Предпочтительнее сохранить эти аргументы.
В то время,
существующие призывы, когда аргументы типа передаются таким образомНе сломается,
потому что в настоящее время для статического элемента является ошибкой зависеть от значения
формальный параметр типа закрывающего класса.%
?

%
%
 быть типом alias заявление формы


\code{{) $F$<\TypeParametersStd> = $T$;


где .
%
1990-1474-й год. Y}{1}{s} - переменные свежего типа,
Предполагалось, что он удовлетворяет пределам .
Мы говорим, что 
{расширяется до переменной типа}
Если транзитивный псевдоним расширение
< Y}{1}{}>}
 для некоторых .

%
%
 быть параметризованным < T}{1}{}>}
(перенаправлено с «P000401») может быть нулевой
в области, где  обозначение типа alias Declaration .
Мы говорим, что 
< Типовой псевдоним! Используется как класс
когда  происходит одним из следующих способов:



 происходит в случае создания выражения формы
 или форму ,
где  или \NEW{} или \CONST{}
().
{%
Обратите внимание, что, например,  может быть творением,
Но затем он рассматривается как \code{\NEW{} $T$(42)} или {{} $T$(42)},
Что означает, что он включен здесь
() %
?

 или  является перенаправлением в перенаправлении заводского конструктора,
().

 Используется в качестве суперкласса, смеси или суперинтерфейса.
в классовой декларации
(, , ,
или как  Тип или суперинтерфейс в смешанной декларации
().
% Его можно использовать как  Тип заявления о продлении:
% В этом положении он используется как тип, а не как класс.

 используется для вызова статического члена класса или смесителя,
Как описано выше.


%
Ошибка времени компиляции возникает, если  расширяется до переменной типа,
и  использует  как класс.
Ошибка времени компиляции возникает, если используется  как класс,
 не является регулярным.
 Например: }



%
Когда мы используем такие фразы, как Let  быть классом или предположить, что  Это смешение,
Следует понимать, что это включает случай, когда
 {typeName} обозначает псевдоним типа, или
 Параметризованный тип формы
\syntax{<typeName> <typeArguments>>
где имя типа обозначает псевдоним типа,
и расширение транзитивного псевдонима  Класс обозначает соответственно смесь.


 Подтипы


%
Этот раздел определяет, когда тип является  другого типа.
Ядром этого раздела является набор правил, определенных в
Рисунок ~,
Но сначала нужно будет ввести несколько понятий.
Чтобы понять, что означают эти правила.

{%
Читатель, прочитавший множество научных работ об объектно-ориентированных системах
Смысл данного обозначения может быть очевиден,
Но нам все еще нужно прояснить некоторые детали о том, как справиться с этим.
синтаксически разные обозначения одного и того же типа,
Как выбрать правильную начальную среду (P000431).
%
Для читателя, который не знаком с обозначением, используемым в этом разделе,
Объяснений, приведенных здесь, должно быть достаточно, чтобы понять, что это значит.
Ссылаясь на объяснения естественного языка, приведенные в конце
раздел для получения интуиции о значении.%
?

%
Этот раздел посвящен отношениям подтипов между типами.во время статического анализа
а также отношения подтипа, запрашиваемые в динамических проверках,
Типовые испытания
(),
и типовые отливки
().

{%
Вариант правил, описанных здесь, показан в приложении.
(),
Демонстрация того, что субтипирование Dart может быть решено эффективно.
?

%
% % TODO(Eernst): Введение этих специализированных типов пересечений
% в подходящем месте, где указан тип продвижения.
Типы формы
\IndexCustom{$X \& S$>>>{тип!of} Форма %
\IndexExtraEntry{\&@$X \& S$>
возникают во время статического анализа из-за продвижения по типу
(P000611).
Они никогда не происходят во время казни.
Они никогда не являются аргументом другого типа.
ни тип возврата, ни тип формального параметра,
И это всегда так, что  является подтипом границы .
{%
Оригинальное название: P000437 является то, что он представляет
тип локальной переменной 
тип которого объявлен переменной типа ,
который, как известно, имеет тип  Из-за повышения.
Точно так же  может рассматриваться как тип пересечения,
что является подтипом  А также подтип .
Типы пересечений {не} поддерживаются в целом.
Только в этом случае.%
?
Любая другая форма может возникнуть во время статического анализа.
Как и во время исполнения,
Взаимоотношения между подтипами всегда определяются одинаково.

% Номер правил подтипа
\newcommand{}{1}
\newcommand{}{2}
\newcommand{}{3}}
\newcommand{}{4}
\newcommand{}{5}
\newcommand{}{6}
\newcommand{}{7}
\newcommand{}{8]
\newcommand{}{9]
\newcommand{}{10]
\newcommand{}{11}
\newcommand{}{12]
\newcommand{}{13}
\newcommand{}{14}
\newcommand{}{15}
\newcommand{}{16}
\newcommand{}{17}
\newcommand{}{18}

[p]
\def\VSP{\vspace{4 мм)
\def\ExtraVSP{\vspace{2 мм)
\def\Axiom#1#2#3#4{\centerline{\inference[#1]{}{{#3}{#4}}}<<<< P001551>>>
\def\Rule#1#2#3#4#5#6{\centerline{[#1]{{#3}{#4}}{{#5}{#6}}}<<<< P001558>>>
#1#2#3#4#5#6#7#8
\centerline{[#1]{{#3}{#4} & \SubtypeStd{#5}{#6}}{{#7}{#8}}}\VSP
\def\RuleRaw#1#2#3#4#5
\centerline{[#1]{#3}{{#4}{#5}}}\VSP]
\def\RuleRawRaw#1#2#3#4{\centerline{\inference[#1]{#3}{#4}}\VSP]
%
[c]{0.49}
{ Рефлексивность.
\Axiom{\SrnBottom>>> Левый дно { T

[c]{0.49}
{ Правый верх.  \code{Object} , \VOID\}}{S}{T}
\RuleRaw{\SrnNull>>> Левый нуль. <<<p001592>> \bot}{}{T}



{ Тип Alias Left {%}
\code{{) <\TypeParametersNoBounds{X}{}> = U} &
<< S_1/X_1,\ldots,S_s/X_s]U}{T}}{\code{<\List{S}{1}{}}}{T}
\RuleRaw{\SrnRightTypeAlias>>> Тип Alias Right {%}
\code{{) <\TypeParametersNoBounds{X}{}> = U} &
\SubtypeStd{S}{[T_1/X_1,,T_s/X_s]U}}{S}{<<<<<<<<<<<<<<<<<<<<<<<<<<<<<<<<<<<<<<<<<<<<<<<<<<<<<<<<<<<<<<<<<<<<<<<<<<<<<<<<<<<<<<<<<<<<<<<<<<<<<<<<<<<<<<<<<<<<<<<<<<<<<<<<<<<<<<<<<<<<<<<<<<<<<<<<<<<<<<<<<<<<<<<<<<<<<<<<<<<<<<<<<<<<<<<<<<<<<<<<<<<<<<<<<<<<< P001611>>>{$F$<\List T}{1}{}>}}

[c]{0.49}
<\RuleTwo{\SrnLeftFutureOr>>> Оставленное будущее {S}{T}{%}
\code{Future<>}}{T}{}{T}\RuleTwo{\SrnRightPromotedVariable>>> Право Продвигаемая Переменная {S}{X}{S}{T}{%
S {\displaystyle X}  T
<\Rule{\SrnRightFutureOrB>>> Правильное будущее Или B}{S}{T}{S}{\code{FutureOr<>}
<\Rule{\SrnLeftVariableBound>>> Левая переменная граница {\Delta(X)}{T}{X}{T}

[c]{0.49}
\Axiom{\SrnTypeVariableReflexivityA>>> Переменная A}{X  S}{X}
\Rule{\SrnRightFutureOrA>>> Правильное будущее Или A}{S}{}{%
S {
<\Rule{\SrnLeftPromotedVariable>>> Левая продвигаемая переменная B}{S}{T}{X  S {T }
\RuleRaw{\SrnRightFunction>>> Правильная функция {T Тип функции {%}
T}{

%

<{ Позиционные функции {%}
\Delta = \Delta\uplus\{X_ i\mapsto{}B_i1  i \leq s\} и
\Subtype{'}{S_0}{T_0}
n_1  n_2
n_1 + k_1 < n_2 + k_2
 j  1.. n_2 + k_2\!:\;\Subtype{\Delta'} T_j} {S_j}}


\RuleRawRaw{\SrnNamedFunctionType}{Названные типы функций}{%
\Delta = \Delta\uplus\{X_ i\mapsto{}B_i\, |\,1  i \leq s\} и
\Subtype{}{S_0}{T_0} &
 j  1.. n\!:\;\Subtype{\Delta'}{T_j}{S_j}
 \List{y}{n+1}{n+k_2}<<<<<<<<<<<<<<<<<<<<<<<<<<<<<<<<<<<<<<<<<<<<<<<<<<<<<<<<<<<<<<<<<<<<<<<<<<<<<<<<<<<<<<<<<<<<<<<<<<<<<<<<<<<<<<<<<<<<<<<<<<<<<<<<<<<<<<<<<<<<<<<<<<<<<<<<<<<<<<<<<<<<<<<<<<<<<<<<<<<<<<<<<<<<<<<<<<<<<<<<<<<<<<<<<<<<<<<<<<<<<<<<<<<<<<<<< P001761>>>  \{\,\List{x}{n+1}{n+k_1}\, 
 >> p  1.. k_2, q \in 1.. k_1:\quad
y_{n+p} = x_{n+q}\quad\Rightarrow\quad\Subtype{\Delta>{T_{n+p}}{S_{n+q}}}{%

%

\RuleRaw{\SrnCovariance>>> Ковариантность
\code{{} $C$\TypeParametersNoBounds{X}{}>\,\{\} и
 j  1..s\!:\;\SubtypeStd{S_j}{T_j}}{%
< S}{1}{}>}}{<<< P001691>>>{$C$< T}{1}{}>}}

<\RuleRaw{\SrnSuperinterface>>> Суперинтерфейс (%)
\code{{) $C$\TypeParametersNoBounds{X}{}>\,\ldots\,\{\}\\
\Superinterface{{<\List{T}{1}{m}> и
<< S_1/X_1,\ldots,S_s/X_s]\code{<\List{T}{1}{m}>}}
< S}{1}{}>}\;}{ T
%
 Правила подтипа.




<\subsubsection{{{{{{{{{{{{{{{{{{{{{{{{{{{{{{{{{{{{{{{{{{{{{{{{{{{{{{{{{{{{{{{{{{{{{{{{{{{{{{{{{{{{{{{{{{{{{{{{{{{{{{{{{{{{{{{{{{{{{{{{{}}{{{{{{{{{{{{{{{{{{{{{{{{{{{{{{{{{{{{{{{{{{{{{{{{{{{{{{{{{{{{{{{{{{{{{{{{{{{{{{{{{{{{{{{{{{{{{{{{{{{{{{{{{{{{{{{{{{{{{{{{{{{{{ Мета-варианты.


%
 Символ, обозначающий синтаксическую конструкцию
удовлетворяет некоторым статическим семантическим требованиям.

{%
Например,  является мета-вариабельным, обозначающим
идентификатор ,
Но только если  обозначает переменную типа, объявленную в замкнутом пространстве.
В приведенных ниже определениях мы определяем это, говоря, что
 Диапазоны по переменным типа.
Точно так же  является мета-вариабельным, обозначающим
a \synt{typeName}, например, ,
Но только если  обозначает класс в заданном объеме.
Мы определяем это как  Диапазон классов.%
?

%
В этом разделе мы используем следующие мета-варианты:


  Диапазоны по переменным типа.
  Диапазон классов,
  Диапазоны по типам псевдонимов.  и  диапазон по типам, возможно с индексом, таким как  или .
  диапазоны по типам, опять же, возможно, с индексом;
Используется только как переменная типа.



 Правила подтипа


%
В этом разделе мы определяем несколько правил о субтитрах.
Всякий раз, когда правило содержит одну или несколько метавариабель,
Это правило может быть использовано
\IndexCustom{instantiating}{instantiation!subtype rule}
то есть путем последовательной замены
Каждое появление данной метавариабельности
конкретный синтаксис, обозначающий один и тот же тип.

{%
Это означает, что два или более случая
заданная мета-вариабельность в правиле
обозначает идентичные части синтаксиса,
и введение правила продолжается как
Простая операция поиска и замены.
Например,
Правило ~\SrnReflexivity На рисунке ~
Можно использовать для заключения
\Subtype{}{}{}
где  обозначает пустую среду
(Любой среды будет достаточно, потому что переменные типа не встречаются).

Однако вышеприведенная формулировка "обозначение одного и того же типа" охватывает
Дополнительные ситуации также:
Например, мы можем использовать правило ~\SrnReflexivity
Чтобы показать, что  является подтипом
, когда  Класс, объявленный в
библиотека  которые импортируются библиотеками  и  и
используется в декларациях,
 и  Импортируются с префиксами
  по текущей библиотеке.
Важно отметить, что все случаи одной и той же мета-вариабельности
в данном правиле инстанциация обозначает один и тот же тип,
Даже в том случае, если этот тип не обозначен
Один и тот же синтаксис в обоих случаях.

И наоборот, мы можем использовать одно и то же правило, чтобы сделать вывод.
 является подтипом 
в случае, когда первый обозначает класс, объявленный в библиотеке 
и последний обозначает класс, объявленный в , с .
Эта ситуация может возникнуть без ошибок времени компиляции, например,
Если  и  Косвенный ввоз в нынешнюю библиотеку
и два «смысла»  используются
аннотации типа на переменные или формальные параметры функций
задекларировано в промежуточных библиотеках импортом  соответственно .
Неспособность доказать
\Subtype{}{}{] " "
Это произойдет, например, в ситуации, когда мы проверим,
такая переменная может быть передана в качестве фактического аргумента такой функции;
Потому что два случая  не обозначают один и тот же тип.%
?

%
Предполагается, что каждый {typeName}, используемый в типе, упомянутом в этом разделе,
Не имеет ошибки компиляции и обозначает тип.

{%
То есть, никакие отношения субтипирования не могут быть доказаны.
Тип, который содержит или содержит неопределенное имя
или имя, которое обозначает нечто иное, чем тип.
Обратите внимание, что для определения отношений субтипирования нет необходимости.
Каждый тип удовлетворяет заявленным границам.
Отношение субтипирования не зависит от границ.
Однако, если предпринимается попытка доказать связь подтипа
и один или более {typeName}s получает фактический список аргументов типа
длина которого не соответствует декларации
(включая случаи, когда приводятся аргументы определенного типа)
неродовой класс,
и случай, когда возникает общий класс,
Но никаких аргументов не приводится.
Тогда попытка доказать отношения просто проваливается.
?

%
Правила на рисунке ~ использовать
Символ  Чтобы обозначить данное знание о
границы типовых переменных.
 Это частичная функция, которая отображает переменные типа на типы.
В заданном месте, где переменные типа в области применения
<
{(как указано в прилагаемых классах и/или функциях)},
Мы определяем окружающую среду следующим образом:
.
{%
То есть,  и так далее,
 не определен при применении к переменной типа 
которого нет в .%
?
Когда правила используются, чтобы показать, что существует определенный подтип отношений,
Это начальное значение .

%
Если встречается общий тип функции, расширение  используется,
Как показано в правилах ~\SrnPositionalFunctionType
и ~
Изображение: P000514.
Расширение сред использует оператор ,
который является оператором, производящим объединение разрозненных множеств,
и отдает приоритет правой руке операнда в случае конфликтов.

<%
Так
$   ,     
    
   ,   ,    
и
$   ,   FutureOr<List<double>{}>  
   =
   ,    $.
Обратите внимание на оператора  Это связано с масштабами и тенью,
без связи, например, с подтипами или методом экземпляра overriding.%
?

%
В данной спецификации мы часто ссылаемся на
отношения подтипа и присваиваемость
Не говоря уже об окружающей среде,
 {T}}
Это делается только тогда, когда конкретное место в коде находится в фокусе.
Окружающая среда – это то, что получается.
путем отображения каждой переменной типа в объеме в этом месте
Объявлено его обязательным.

%
Каждое правило на рисунке  имеет горизонтальную линию,
Слева от которого находится  указывается;
Под горизонтальной линией находится суждение, которое
{заключение}{правило!заключение}
Правило,
над горизонтальной линией находится ноль или более
{предпосылки}{правило!предпосылка}
Правило,
которые, как правило, также являются подтипом суждения.
Если это не относится к определенной предпосылке
Мы четко указываем значение.

{%
Обоснование вышеупомянутого правила,
обозначает последовательную замену метавариабельных
посредством фактических синтаксических терминов, обозначающих типы повсюду в правиле,
То есть, как в помещениях, так и в заключении, одновременно.
?

 Быть подтипом


%
Тип  Показано, что  другого типа 
в окружающей среде  путем предоставления
Инстанциация правила  чей вывод является
 S}{T}}{@\SubtypeStd S}{T}}
наряду с инсталляциями правил, показывающими
каждого из помещений ,
Продолжается до тех пор, пока не будет достигнуто правило, не имеющее помещения.

{%
Для правила  обратите внимание, что  тип
является подтипом всех не- типы,
Хотя на самом деле он не расширяет и не реализует эти типы.
Другие типы эффективно рассматриваются так, как если бы они были недействительными.
Что делает нулевой объект () Присваивается им. %
?

%
Первая предпосылка в
Правила ~\SrnLeftTypeAlias и ~
Это своеобразная декларация.
Эта предпосылка удовлетворяется в каждой из следующих ситуаций:


 Объявляется псевдоним недженерического типа под названием .
В данном случае  равна нулю,
Нет никаких предположений о существовании
любых формальных типовых параметров,
Списки аргументов фактического типа опускаются везде в правиле.
 Мы можем выбрать  и {X}{1}{s} таким образом, чтобы следующее:
Общий псевдоним типа  объявляется,
Формальные параметры типа  X}{1}{s}
{%
Каждый формальный параметр типа  Может быть, у него есть связь,
Но границы никогда не используются в этом контексте.
Поэтому мы не вводим метавариабельности для них.%
?


%
Правило ~\SrnRightFunction В качестве предпосылки, что  является функциональным типом.
Это означает, что  является одной из форм, введенных в
Раздел .
{%
Это то же самое, что и формы, которые встречаются на высшем уровне.
в выводах из
Правило ~\SrnPositionalFunctionType{}
Правило ~\SrnNamedFunctionType%
?

%
В правилах ~\SrnCovariance и ~,
Первое условие – это классовая декларация.
Эта предпосылка удовлетворяется в каждой из следующих ситуаций:


 Объявляется необычный класс с именем .
В данном случае  равна нулю,
Нет никаких предположений о существовании
любых формальных типовых параметров,
Списки аргументов фактического типа опускаются везде в правиле.
 Мы можем выбрать  и {X}{1}{s} таким образом, чтобы следующее:
Общий класс, названный  объявляется,
Формальные параметры типа  X}{1}{s}
<<<<<<<<<<<<<<<<<<<<<<<<<<<<<<<<<<<<<<<<<<<<<<<<<<<<<<<<<<<<<<<<<<<<<<<<<<<<<<<<<<<<<<<<<<<<<<<<<<<<<<<<<<<<<<<<<<<<<<<<<<<<<<<<<<<<<<<<<<<<<<<<<<<<<<<<<<<<<<<<<<<<<<<<<<<<<<<<<<<<<<<<<<<<<<<<<<<<<<<<<<<<<<<<<<<<<<<<<<<<<<<<<<<<<<<<<<<<<<<<<<<<<<<<
Каждый формальный параметр типа  Может быть, у него есть связь,
Но границы никогда не используются в этом контексте.
Поэтому мы не вводим метавариабельности для них.%
?


%
Вторая предпосылка правила ~ указывает, что
параметризированный тип \code{<\ldots{}>> принадлежит
\IndexCustom{{C}}{суперинтерфейсы(C)@{ С.
Семантическая функция \Superinterfaces{\_} применяется к общему классу  доходность
Набор прямых суперинтерфейсов 
(P000518).

{%
Обратите внимание, что один из прямых суперинтерфейсов  это
интерфейс суперкласса ,
Это может быть приложение для смешивания
(),
в каком случае  В правиле есть
синтетический класс, определяющий
Семантика этого смешанного приложения.%
?

{%
Последняя предпосылка правила ~
Подменяет аргументы фактического типа \List{S}{1}{s}
Формальные параметры типа \List{X}{1}{}
 T}{1}{m} может содержать формальные параметры типа.%
?

{%
Правила  и ~\SrnSuperinterface
применимы к интерфейсам,
Также их можно использовать с классами.Потому что негенерический класс  который используется как тип
обозначает интерфейс ,
Для параметризированного типа
< T}{1}{k}>
где  обозначает общий класс.%
?


 Неофициальные описания правил подтипа


{%
В этом разделе дается неформальное и ненормативное описание естественного языка.
Каждое правило на рисунке .

В описаниях используются номера правил, чтобы сделать соединение явным,
и добавляет имена к правилам, которые могут быть полезны для понимания
Роль, которую играет каждое правило.

В следующем, многие правила содержат мета-вариабельность.
()
как  и ,
И всегда можно говорить о произвольных типах.
Например, правило ~\SrnRightFutureOrA{} Говорит, что
Тип   ,
Это означает, что для любых произвольных типов  и ,
Показать  является подтипом  Достаточно показать, что  это
Подтип .

Другой пример — формулировка в правиле ~\SrnReflexivity{}:
\ldots{} в любой среде ,
что указывает на то, что правило может применяться независимо от того, какие привязки
Переменные типа к границам существуют в окружающей среде.
Следует отметить, что окружающая среда имеет значение даже при соблюдении правил.
где оно просто заявлено как простое  в заключение
в одном или нескольких помещениях,
Поскольку доказательство этих предпосылок может прямо или косвенно
Включает применение правила, в котором используется окружающая среда.

#1#2{[#1]{{#2:}}}

\Item{ Рефлексивность
Каждый тип является подтипом самого себя, в любой среде .
В следующих правилах, кроме нескольких,
Это правило действует в любой среде.
Окружающая среда никогда не используется открыто.
Поэтому мы не будем повторять этого.
\Item{\SrnTop}} Вершина!
Каждый тип является подтипом ,
Каждый тип является подтипом ,
Каждый тип является подтипом .
Это означает, что эти типы эквивалентны
в соответствии с отношением подтипа.
Мы обозначаем эти типы,
и другие с таким же свойством (например, ),
как лучшие типы
().
\Item{}
Каждый тип является супертипом P000056.
\Item{}
Все типы, кроме  Это супертип P000592.
\Item{\SrnLeftTypeAlias> Тип Alias Left
Применение псевдонима типа к некоторым фактическим аргументам типа
Подтип другого типа 
если расширение псевдонима типа на тип, который он обозначает
Это подтип .
Обратите внимание, что псевдоним необщего типа обрабатывается с помощью разрешения .
\Item{\SrnRightTypeAlias}} Тип Alias Right
Тип  является подтипом приложения псевдонима типа
Если  является подтипом
расширение псевдонима типа на тип, который он обозначает.
Обратите внимание, что псевдоним необщего типа обрабатывается с помощью разрешения .
\Item{\SrnLeftFutureOr}} Левое будущее или
Тип  является подтипом данного типа 
Если  является подтипом  и  является подтипом ,
для каждого типа  и .
\Item{\SrnTypeVariableReflexivityA}} Продвигаемая переменная
Тип  Это подтип .
{ Правильно повышен Переменная А}
Тип  является подтипом  если
 является подтипом как , так и .
{ Правильное будущее Или А.Тип  является подтипом  если
 Это подтип .
\Item{}} Правильное будущее Или Б.
Тип  является подтипом  если
 Это подтип .
\Item{\SrnLeftPromotedVariable}} Продвигаемая левая переменная B}
Тип  является подтипом  если
 является подтипом .
\Item{\SrnLeftVariableBound}} Левая переменная граница
Переменная типа  является подтипом типа  если
Связь 
(как указано в текущей среде )
Это подтип .
\Item{\SrnRightFunction}} Правильная функция
Каждый тип функции является подтипом типа .
\Item{\SrnPositionalFunctionType}} Тип позиционной функции
Тип функции  с позиционными дополнительными параметрами
является подтипом
Другой тип функции  с опциональными параметрами
Если первый имеет самое
количество требуемых параметров, равное последнему,
Последний имеет по крайней мере
общее количество параметров, равное первому;
Тип возврата  является подтипом подтипа ;
и каждого типа параметров  {супертип}
соответствующий тип параметров , если таковой имеется.
Обратите внимание, что отношение к типам функций без дополнительных параметров
и отношения между типами функций без дополнительных параметров,
Покрывается разрешением  соответственно .
Для каждого подтипа отношения, рассматриваемого в этом правиле,
Формальные параметры типа  и  должны быть приняты во внимание
(как это отражено в использовании расширенной среды ).
Мы можем предполагать без потери общности.
что имена переменных типов попарно идентичны,
Мы считаем, что типы общих функций эквивалентны в
последовательное переименование
().
Короче говоря, «во время доказательства мы переименуем их по мере необходимости».
Наконец, обратите внимание, что взаимосвязь между негенерическими типами функций
Покрывается разрешением .
\Item{\SrnNamedFunctionType}} Названы функции типа
Тип функции  с названными факультативными параметрами является подтипом
Другой тип функции  с названными дополнительными параметрами
если они имеют одинаковое количество требуемых параметров,
и набор имен названных параметров для последнего является подмножеством
Для тех, кто раньше
Тип возврата  является подтипом подтипа ;
и каждый параметр типа  {супертип}
соответствующий тип параметров , если таковой имеется.
Обратите внимание, что отношение к типам функций без дополнительных параметров
и отношения между типами функций без дополнительных параметров,
Покрывается разрешением  соответственно ,
и что последний случай идентичен правилу, полученному из
Правило ~\SrnPositionalFunctionType
в отношении подтипирования между типами функций без дополнительных параметров.
Как правило, ,
Мы можем принять без потери общности.
Имена переменных типов попарно идентичны.
Аналогично, негенерические функции охватываются разрешением .
\Item{}} Класс Ковариантность
Параметризованный тип на основе общего класса  является подтипом
параметризованный тип на основе того же класса  если
Каждый действительный тип аргумента первого является подтипом
соответствующего фактического типа аргумента последнего.
Это правило может иметь значение  и охватывают также необычный класс,
Но это избыточно, потому что это уже покрыто
Правило .
\Item{\SrnSuperinterface}} СуперинтерфейсРассмотрим случай, когда  и  Во-первых,
параметризованный тип на основе негенерического класса  является подтипом
параметризованный тип на основе другого негенерического класса  если
 Это прямой суперинтерфейс .
Когда  или , это правило описывает отношение подтипа
который включает один или несколько общих классов,
В этом случае мы должны дать названия формальным параметрам типа ,
и указать, как они используются в спецификации суперинтерфейса
На основе .
С этими частями на месте, мы можем указать отношения подтипа.
Существует между двумя параметризованными типами на основе  и .
%
% % TODO(Eernst): Обратите внимание, что спецификация того, как передавать аргументы типа
% % в  Это неправильно, и это должно быть
% полностью переписаны для интеграции новой конструкции смесителя.
Случай, когда суперкласс представляет собой приложение для смешивания, рассматривается через
эквивалентность с объявлением
обычный (возможно, общий) суперкласс
(),
Это означает, что может быть несколько шагов подтипа от
a заданное объявление класса классу, указанному в {} пункт.
%
?


{дополнительные концепции субтипирования}


%
   в данной среде ,
написанный {S}{T}
если {f}  T.S.

%
Тип 

Тип  в среде ,
написанный {S}{T}
Если {f} {S}{T} или  T.S.
В этом случае мы говорим, что типы  и  есть
.

{%
Это правило может удивить читателей, привыкших к обычной проверке типа.
Намерение  отношение
Это не означает, что задание гарантирует динамический успех.
Вместо этого он направлен только на флаги.
Это почти наверняка будет ошибочным.
Не исключая заданий, которые могут работать.

Например, назначение объекта статического типа 
переменной со статическим типом ,
Несмотря на то, что это не гарантировано,
Может быть хорошо, если значение времени выполнения оказывается строкой.

Статический анализатор или компилятор
может поддерживать более строгие статические проверки в качестве опции.
?


{Функциональные типы}


%
{Функциональные типы}{тип!функция}
Выпускаются в двух вариантах:



Типы функций, которые имеют только позиционные параметры.
Они имеют общую форму
{T}.

Виды функций с названными параметрами.
Они имеют общую форму
{T}.


{%
Обратите внимание, что необычный случай охватывается наличием ,
в этом случае объявления о параметрах типа опускаются
(P000526).
Случай без факультативных параметров покрывается наличием ;
Обратите внимание, что все правила, касающиеся типов функций двух типов
совпадает в данном случае.%
?

%
Два типа функций считаются равными при последовательном переименовании типа.
Параметры могут сделать их идентичными.

{%
Обычный способ сказать, что мы не различаем
функции, которые являются альфа-эквивалентными.
Правило субтипирования ниже означает, что мы можем предположить, что
Подходящее переименование уже состоялось.
В тех случаях, когда это невозможно
Поскольку количество типовых параметров в двух типах различается
Или же границы разные,
Подтипа не существует.%
?

%
Объект функции всегда является экземпляром некоторого класса  который реализует
класс < (),
и который имеет метод, названный ,сигнатурой которого является тип функции  себя.
{%
Следовательно, все типы функций являются подтипами 
()%
?


\subsection{Тип \FUNCTION}}


%
Встроенный класс {} является супертипом всех типов функций.
().
Невозможно расширить, реализовать или смешать в классе .

%
Если заявление класса или приложение для смешивания имеет {} в качестве суперкласса,
Вместо этого он использует P000600 в качестве суперкласса.

%
Если класс или смешанная декларация реализуется , это не имеет никакого эффекта.
Это как если бы P001007 Он был удален из  пункт
(и если это единственный реализованный интерфейс, весь пункт удаляется).
Полученный класс или интерфейс микширования
не имеет {} в качестве суперинтерфейса.

%
Если приложение для смешивания смешивает {} на суперкласс, следует за
Нормальные правила применения миксина, но так как результат этого миксинга
Приложение эквивалентно классу с , и
Это положение не имеет эффекта, полученный класс также не
осуществлять .
{%
{} класс не объявляет конкретных членов,
Таким образом, приложение микширования создает подкласс
Суперкласс без новых членов и без новых интерфейсов.
?

{%
После использования \FUNCTION{} Таким образом, это не имеет никакого эффекта, это было бы
разумно полностью запретить его, как мы расширяем, реализуем или
смешивание в таких типах, как  или .
Для обратной совместимости с программами Dart 1
Синтаксис разрешается сохранять, даже если он не имеет эффекта.
Инструменты могут предупреждать пользователей о том, что их код не имеет эффекта.
?


\subsection{Тип 


%
Тип {} является статичным типом, который является супертипом всех других типов.
Как будто ,
отличается от других видов тем, что статический анализ
Предполагает, что каждый член имеет соответствующий член
с подписью, допускающей данный доступ.

{%
Например,
когда приемник в обычном способе вызова имеет тип ,
Можно использовать любое название метода,
с любым количеством аргументов или без них,
с любым количеством аргументов,
И всякие именованные доводы,
любого типа,
Без ошибок.
Обратите внимание, что вызов все равно вызовет ошибку времени компиляции.
Если есть ошибка в одном или нескольких аргументах или других терминах.
?

%
Предполагается, что процентный вывод имел место, поэтому тип не был выведен.
Если аннотация статического типа не была предоставлена,
Типовая система рассматривает декларации как имеющие тип .
% % TODO(Eernst): Изменение при добавлении спецификации instantiate-to-bound.
Если используется общий тип, но аргументы типа не приводятся,
Аргументы типа по умолчанию вводят .

{%
% % TODO(Eernst): Изменить при добавлении спецификации мгновенного подключения.
Это означает, что при наличии общей декларации
,
где  является формальной декларацией параметров типа ,
Тип  является эквивалентным


%
?

%
Встроенная декларация типа ,
который объявлен в библиотеке ,
обозначает \DYNAMIC{} Тип.
Когда имя \DYNAMIC{} Экспорт:  оценивается
как выражение,
Для сравнения:  Объект, представляющий собой \DYNAMIC{} тип,
Даже если {} не является классом.

{%
 объект должен быть равен соответствующему 
 и  По словам оператора \lit{==================================================================================================================
()%
?

%
Для повышения точности статических типов
Член получает доступ к приемнику типа \DYNAMIC{} который относится к
Декларации встроенного класса дается статический тип, соответствующий этим декларациям
Всякий раз, когда это происходит, это звучит хорошо.



 быть выражением формы ,
который не следует за аргументом,
где статический тип  ,
 имя геттера, объявленное в ;
если тип возврата   затем
Статический тип  является .
{%
Например,  имеет тип 
 имеет тип ,%
?


 быть выражением формы ,
который не следует за аргументом,
где статический тип  ,
 это название метода, объявленного в 
сигнатура метода которого имеет тип 
{(который является функциональным типом)}.
Статический тип  В этом случае .
{%
Например,  имеет тип .%
?


 быть выражением, которое имеет форму
\code{$d$.\id(\metavar{аргументы})}
или форму \code{$d$.\id<\metavar{typeArguments}>(\metavar{arguments})},
где статический тип  ,
 имя геттера, объявленное в  Тип возврата ,
\metavar{аргументы} происходит от \synt{аргументы} и
\metavar{typeArguments} являются производными от \synt{typeArguments}, если они присутствуют.
Статический анализ обрабатывается  как функция выражения призыв
где объект статического типа  Применяется к данной аргументной части.
<<<p001052>{%
Это всегда ошибка времени компиляции.
Например,  Ошибка времени компиляции,
Потому что он проверяется как
Выражение функции вызова, где
сущность статического типа  это
Призвано. Обратите внимание, что это действительно может быть успешным: Преобладающая реализация
 Можно ли вернуть пример
динамический тип которого является подтипом
 Имеет метод {}.
Мы решили сделать это ошибкой, потому что это может быть ошибкой.
Особенно в таких случаях, как 
где разработчик мог забыть, что  Геттер.%
?


 быть выражением формы {. \id(\metavar{аргументы})}
где статический тип  , <{аргументы}
фактический список аргументов, полученный из {аргументов},
 Это название метода, объявленного в 
чей метод подписи имеет тип .
Если число позиционных фактических аргументов в 
меньше числа требуемых позиционных аргументов 
или больше числа позиционных аргументов в ,
или если {аргументы} включает любые названные аргументы
с именем, которое не указано в ,
Тип  .
В противном случае тип  Тип возврата в .
{%
 имеет тип < тогда как
 имеет тип .
Обратите внимание, что призывы, которые «не подходят» статически
Заявления не являются ошибками, они просто получают тип возврата. %
?


 быть выражением формы
. \id<\metavar{typeArguments}>({arguments})} где
Статический тип  - ,
{typeArguments} - список актуальных
Аргументы типа, полученные из {typeArguments} и{аргументы} - это фактический список аргументов, полученный из \synt{аргументы}.
Ошибка времени компиляции  Это имя
необычный метод, указанный в .
{%
Общие методы не указаны в .
Поэтому мы не утверждаем, что должно быть
Статически необходимое количество аргументов фактического типа и
Они должны соответствовать границам.
В противном случае это был бы последовательный подход.
Потому что призыв гарантированно потерпит неудачу, когда любой из этих
нарушены требования,
Но обобщения этого механизма должны включать такие правила.
?


Например, вызов метода  (в том числе призывы геттеров,
сеттеры и операторы), где приемник имеет статический тип <{}
 не соответствует ни одному из вышеперечисленных случаев статического типа  это
.
Когда выполняется выражение, полученное из 
геттер или метод вызова, который соответствует одному из приведенных выше случаев,
Применяется соответствующий статический анализ и ошибки компиляции времени.
{%
Например,  Ошибка.%
?


{%
Обратите внимание, что только очень немногие формы вызова метода экземпляра с
приемник типа \DYNAMIC{} Это может быть ошибка времени компиляции.
Конечно, некоторые выражения, такие как  Синтаксические ошибки
Хотя они также могут считаться «призывами»,
Субэкспрессии проверяются отдельно.
Любой конкретный аргумент может быть ошибкой времени компиляции.
Но почти любая форма списка аргументов может быть обработана через
,
Любой аргумент может быть принят, потому что любой
Формальный параметр в основной декларации может иметь свой тип
Аннотация была изменена на .
Это естественное следствие принципа
это
a Получатель {} допускает почти все вызовы метода экземпляра.
Несколько случаев, когда метод инстанций вызывает
Получатель типа {} Это ошибка
Либо гарантированно проваливаются во время бега,
или они очень, очень вероятно, будут ошибки разработчика.
?


 Тип будущего или


%
Встроенная декларация типа ,
который вывозится библиотекой ,
определяет общий тип с одним параметром типа ().
Тип  Это неклассовый тип
Для всех, кто имеет регулярные ограничения .

{%
Отношения подтипа с участием  указывается в других
().
Обратите внимание, однако, что они влекут за собой определенные полезные свойства:


 
T.
<<<p001090>> 
 T T.
[]
Если  и  T$>} <: S$, затем $\code{FutureOr<$ T.


То есть  является в некотором смысле
Союз  и соответствующего будущего типа.
Последний пункт гарантирует, что
 <:,
А также  является ковариантным по своему типу параметру,
Как и типы классов:
Если  <:, то  <: 
?

%
Если аргументы типа перешли на  будет иметь ошибки компиляции времени
если применяется к обычному общему классу с одним параметром типа,
Ошибки времени компиляции выдаются для .
Оригинальное название: P000656 как выражение
обозначает  объект, представляющий тип .

{%
 Тип представляет собой случай, когда объект может быть
Или один из экземпляров типа 
или тип .
Такие случаи происходят естественным образом в асинхронном коде.
Доступной альтернативой будет использование верхнего типа (например, ). позволяет использовать некоторые инструменты для
Более точный анализ типов.%
?

%
Тип  имеет интерфейс, идентичный этому
.
{%
То есть только члены, которые  можно ли ссылаться
на объекте со статическим типом ,%
?

{%
Мы допускаем только призывы членов, которые унаследованы от
Обычный супертип как , так и .
В большинстве случаев единственным распространенным супертипом является .
Исключения, такие как 
Что такое P000669 как обычный супертип,
малочисленны и практически не полезны,
Поэтому сейчас мы выбираем только призывы к
члены, унаследованные от .%
?

%
Определяем вспомогательную функцию
 T}}{futureOrBase(t)@\emph{futureOrBase}$(T)$}
следующим образом:


 Если   Для некоторых 
.
 Иначе .



 Тип Пустота.


%
Специальный тип \VOID{} Используется для обозначения
Значение выражения бессмысленно и должно быть отброшено.

{%
Типичным случаем является тип \VOID{} Используется как тип возврата
Функция, которая «ничего не возвращает».
Технически всегда будет объект {some}
который возвращается
().
Но очень важно иметь функцию
чьей единственной целью является создание побочных эффектов,
так, что {любое} использование возвращаемого объекта
Это было бы неправильно.
%
Это не значит, что есть что-то неправильное
Возвращенный объект как таковой.
Это может быть любой объект.
Но разработчик, выбравший тип возврата 
Это указывает на то, что это недоразумение
придать этому объекту какой-либо смысл,
или использовать его в любых целях.%
?

{%
Тип  является высшим типом
(),
\VOID{{} и  являются подтипами друг друга
(),
Это означает, что любой объект может быть
значение выражения типа .
%
Следовательно, любой экземпляр типа  который соответствует типу \VOID{}
должны сравниваться равными (по оператору {==} )
Для любого случая  который соответствует типу 
(P000537).
Не гарантируется, что  Оценивается как правда.
На самом деле, не рекомендуется, чтобы реализации стремились достичь этого.
Потому что может быть более важным обеспечить, чтобы диагностические сообщения
(включая следы стека и динамические сообщения об ошибках)
Хранить достаточно информации, чтобы использовать слово «пустота» при обращении к типам
которые указаны как таковые в исходном коде.%
?

%
В поддержку этого понятия
значение выражения со статическим типом \VOID{} должны быть отброшены,
Такие объекты могут использоваться только в конкретных ситуациях:
Возникновение выражения типа \VOID{} является ошибкой времени компиляции
Если это не разрешено в соответствии с одним из следующих правил.



В {выражение ,  Может иметь тип .
 Значение  отбрасывается.

При инициализации и инкрементных выражениях петли,
\code{{) (; ;  {\ldots}}
 может иметь тип ,
и каждое из выражений в списке выражений  может иметь тип .
 Значения  и  отбрасываются.

В типе литья ,  может иметь тип .
{%
Таким образом, разработчики получают возможность на значениях с статическим типом .
Это означает, что такие ценности не должны быть отброшены,
Но они могут быть использованы только тогда, когда это необходимо.
Указывается явно.%
?

В скобчатом выражении ,  Может иметь тип .
{%
Обратите внимание, что  Сам по себе тип ,
Это означает, что оно должно происходить в определенном контексте.
где это не ошибка, чтобы иметь его.%
?

В условном выражении ,
 и  может иметь тип .
{%
Статический тип условного выражения тогда ,
Даже если одна из ветвей имеет другой тип,
Это означает, что условное выражение должно произойти снова.
В каком-то контексте, где это не ошибка, чтобы иметь его.
?

В нулевом слитном выражении ,
 Может иметь тип .
{%
Статический тип нулевого коалесцирующего выражения тогда ,
Что, в свою очередь, ограничивает, где это может произойти.
?

В выражении формы ,  может иметь тип .
{%
Это правило было принято, потому что это было существенное изменение.
Превратить эту ситуацию в ошибку
В то время, когда была изменена обработка .
Инструменты могут дать подсказку в таких случаях.%
?

{%
В ответном заявлении ,
 может иметь тип {} в ряде ситуаций
()%
?

{%
В теле функции стрелы ,
Возвращенное выражение  может иметь тип <
в ряде ситуаций
()%
?

Начальное выражение для переменной типа \VOID{}
может иметь тип .
 Использование этой переменной ограничено.

Фактическое выражение аргумента, соответствующее формальному параметру
Статически известный тип аннотации 
может иметь тип .
{%
Использование этого параметра в теле вызывающего
Статически ожидается, что они будут ограничены наличием типа .
Смотрите дискуссию о звуке ниже
()%
?
Это правило также применяется к операторам  и .
{%
Например,  и $e_2$ может иметь тип <
В выражении формы 
когда параметры статически известного оператора 
Оба имеют тип {.%}
?
Наконец, это правило относится и к сеттерам.
{%
Например, с выражением формы 
Который обозначает призывник,
Статически известный параметр типа ,
 может иметь тип \VOID{.%}
?

В выражении формы 
где  является {назначаемый выражение, обозначающее локальную переменную
{(который также может быть формальным параметром)}
тип ,
 может иметь тип .
{%
Использование этой переменной ограничено, поскольку она имеет тип .
Смотрите дискуссию о звуке ниже
()%
?

 быть выражением, заканчивающимся на 
в форме ,
где  имеет форму


\syntax{(<cascadeSelector> <argumentPart>*)
(<assignableSelector> <argumentPart>*)*)


 Имеет форму {выражение БезКаскада}.

Если  является {назначенный выборщик}
форма  или где статический тип идентификатора  ,
 может иметь тип .
В противном случае, если  {назначаемый выборщик} формы
 где статический тип
Первый формальный параметр в статически известной декларации
оператор  ,
 может иметь тип .
Кроме того, если статический тип второго формального параметра ,
 Может иметь тип .


%
Наконец, мы должны рассмотреть ситуации, связанные с неявным использованием
Объект, статический тип которого может быть .
%
Это ошибка времени компиляции для заявления for-in, чтобы иметь итератор.
выражение типа  {Итератор\VOID{}>}
Наиболее специфичная инстанциация 
Это суперинтерфейс , если только
Переменная итерации имеет тип .
%
Это ошибка времени компиляции для асинхронного заявления for-in.
Иметь потоковое выражение типа 
{Stream\VOID{}>} является наиболее специфическим
Инстанциация  Это суперинтерфейс ,
Если переменная итерации не имеет типа .

 Вот несколько примеров:



{%
Однако в примерах, которые не являются ошибками
Использование  В петле тело ограничено,
Потому что имеет тип %
?


 Звуковая пустота


%
Ограничения на использование объекта, полученные в результате оценки
выражение со статическим типом \VOID{}
не строго соблюдаются.

{%
Использование «пустого значения» не является вопросом здравости, то есть
нет инвариантов, необходимых для правильного выполнения программы Дарта
могут быть нарушены из-за такого использования.
?

{%
Можно сказать, что тип  используется
Помогать разработчикам поддерживать определенную дисциплину
О том, что некоторые объекты не предназначены для использования.
%
В силу того, что принуждение не является необходимым,
и из-за обработки {} в более ранних версиях Дарта,
Язык использует подход {лучшие усилия} для обеспечения
что значение выражения типа \VOID{}
не используется.%
?

<{%
На самом деле, есть много способов в дополнение к типу литья.
где разработчик может получить доступ к такому объекту: %
?



<<%
В (1) переменная  тип < >>{}
общая функция ,
что допускается, так как фактический тип аргумента \VOID{}
и разрешается передавать фактический аргумент типа <<
Формальный параметр того же типа.
%
Тем не менее, никакой специальной обработки не дается, когда выражение имеет тип.
которая является или содержит переменную типа, значение которой может быть ,
Поэтому нам разрешено вернуться  в теле ,
Это означает, что мы косвенно получаем доступ к стоимости.
выражение типа , под статичным типом .

В пункте 2 мы косвенно получаем доступ к значению
Переменная  с типом ,
потому что мы используем задание, чтобы получить доступ к экземпляру 
который был создан с аргументом типа  по типу
.
Обратите внимание, что \code{A<Object>} и \code A<\VOID{}> являются подтипами друг друга,
То есть они эквивалентны по правилам подтипа,
Таким образом, ни статические, ни динамические проверки не будут работать.

В пункте 3 мы косвенно получаем доступ к стоимости
Переменная  с типом <
Статический тип ,
потому что статически известный способ подписи 
имеет тип параметра ,Но фактическая реализация  на который ссылается
является оверрайдом, тип параметра которого ,
что допускается, потому что  и \VOID{} являются двумя верхними типами.
?

{%
Очевидно, что поддержка для предотвращения разработчиков от использования объектов
полученных из выражений типа \VOID{} Далеко не звук,
В том смысле, что есть много способов обойти правило
Такой объект следует выбросить.

Однако мы решили сосредоточиться на простом использовании первого порядка.
(где выражение имеет тип , и используется значение),
И мы оставили дела более высокого порядка в значительной степени неконтролируемыми,
использование дополнительных инструментов, таких как linters, для проведения анализа;
Покрывает косвенные потоки данных.

Конечно, можно было бы определить разумные правила.
Таким образом, значение выражения типа \VOID{}
будет гарантированно выброшен после некоторых переводов.
от одной переменной или параметра к следующей, все с типом ,
явно или в качестве значения параметра типа.
В частности, мы можем потребовать, чтобы переопределение метода
Никогда не отменяйте тип возврата  Тип возврата ,
Типы параметров в противоположном направлении;
параметризованные типы с аргументом типа \VOID{} не может быть назначен
переменных, где аргумент соответствующего типа является чем-либо иным, чем
 и т.д.\ etc.

Но это было бы совершенно непрактично.
В частности, необходимость предотвращения большого количества переменных типа.
из когда-либо имевших значение ,
или предотвращение использования определенных значений, тип которых является такой переменной типа,
или в котором содержится такая переменная типа,
Это сильно ограничило бы большую часть всего кода Дарта.

Поэтому мы решили помочь разработчикам поддерживать эту самонавязанную дисциплину.
в простых и прямых случаях,
и оставить его для специальных рассуждений или отдельных инструментов для обеспечения
Косвенные случаи рассматриваются так же тщательно, как это необходимо на практике.
?


 Параметризованные типы


% TODO (Eernst): Возможно, мы хотели бы добавить новую концепцию применения
% общая функция для аргументов фактического типа (возможно, это дополнительный вид аргументов).
% "параметризованный тип", но он отличается от общего класса, поскольку
% мы можем иметь динамические вызовы общей функции. Но это делает
% не возникает как самостоятельная организация до того, как мы введем общие отрывы
% (), или если мы позволим ему возникнуть на неявной основе
% на вывод. Эта новая концепция, вероятно, должна быть добавлена в этот раздел.

%
{параметризированный тип} — синтаксическая конструкция.
где применяется наименование декларации общего типа
Список фактических аргументов.
{общая инстанциация} - операция, при которой
Общий тип применяется к фактическим аргументам типа.

{%
Параметризованный тип — это синтаксическая концепция, которая соответствует
Семантическая концепция родовой инстанциации.
При использовании первого мы часто оставляем последний неявным.
?

%
 быть параметризованным .

%
Ошибка времени компиляции  не является общим типом,
 является универсальным типом,
но количество формальных параметров типа в декларации  Это не P000291.
В противном случае пусть

быть формальными параметрами типа , и

быть соответствующими верхними границами, используя  когда не объявлялось никаких обязательств.

%
  Если {f}
 Для одного или нескольких ,
 Это не очень хорошо (P000542).

%
Ошибка времени компиляции  Он ограничен.

%
 оценивается следующим образом.
 быть результатом оценки , для . Затем оценивает дженерикальную инстанциацию
где  Применяется к .

%
 быть параметризованным типом формы

Предположим, что  Он не деформирован и не деформирован.
Если  является статическим типом элемента  ,
Статический тип элемента  Выражение типа

это
,
где  Формальные параметры типа .

% % TODO(Eernst): Это место, где мы можем указать, что каждый из
% % аргументы типа получателя  должен быть
% % заменено нижним типом («Null», на данный момент) в местах расположения члена
Тип %%, где это происходит противоположно. Например, "c.f" должен иметь
% % статический тип «пустая функция» когда "c" имеет статический тип "C<T>" для любого
% % 'T', и мы имеем 'класс C<X> { void Function(X) f; }'.


{Актуальные типы}


% Примечание (Eernst): Аргументы фактического типа в этом разделе являются динамическими объектами.
% не синтаксис (используется понятие «фактический тип» и «фактическая связь»)
% определяют динамическую семантику, включая динамические ошибки. Мы используем 
% для обозначения аргументов такого типа, как и все те места, где
Используется понятие %, а не  который часто используется для обозначения
% синтаксис аргумента реального типа.
%
% Дело в том, что  никогда не будет содержать формальный параметр типа, поэтому
% не нужно беспокоиться о необходимости повторять замену, которая находит
% фактической границы. Например, нет проблем с определением фактической границы.
% для  в обращении к {void foo< X расширяет C<X>>() ...
%, даже если это рекурсивное обращение по форме  в
% тела.

%
Пусть  Термин, полученный из {тип}.
Пусть {X}{1}{s} будут формальными параметрами типа в области применения
в месте, где  происходит.
В контексте, где фактический тип аргументов соответствует
 X}{1}{s} являются {t}{1}{}
  тогда
.

%
 быть декларацией с именем  и введите аннотацию .
  и  В данном контексте является
Тогда фактическое значение  в этом контексте.

{%
В необычный случай, когда ,
фактический тип равен заявленному типу,
В том смысле, что мы используем простые термины, такие как , для обозначения обоих.
Обратите внимание, что {X}{1}{s} может быть объявлено несколькими объектами, например,
одна или несколько прилагаемых общих функций и прилагаемый общий класс.%
?

%
 <  Формальное объявление параметров типа.
  В данном контексте является
Фактическое значение  в этом контексте.

{%
Обратите внимание, что даже если  Может возникнуть в 
не происходит в фактическом значении ,
Поскольку ни один тип не имеет фактического значения, которое включает переменную типа.
?


 Наименее верхние границы


% TODO (Eernst): Этот раздел был обновлен для выполнения общих функций.
%, но никаких других изменений не произошло. Следовательно, мы все еще
% необходимо обновить этот раздел, чтобы использовать правила Dart 2 для LUB.

%
В некоторых случаях это расходится?
Дано два интерфейса  и ,
 быть набором суперинтерфейсов ,
 Набор суперинтерфейсов 
И пусть .
Кроме того,
 Для любого конечного числа 
где  Число шагов на самом длинном пути наследования
От  до .%TODO(lrn): Укажите: «путь наследования» — это путь в графе суперинтерфейса.
 Самое большое число, такое, что  Кардинальность одна,
Он должен существовать, потому что  является .
Наименьшая верхняя граница  и  является единственным элементом .

%
Наименьшая верхняя граница  и любой тип  является .
Наименьшая верхняя граница  и любой тип  является .
Наименьшая верхняя граница  Любой тип  является .
 быть переменной типа с верхней границей .
Наименьшая верхняя граница  Тип  это
наименьшая верхняя граница  и .

%
Операция с наименьшей верхней границей является коммутативной и идемпотентной.
Но она не является ассоциативной.

% Типы функций

% % TODO(Eernst): Этот раздел должен использовать новый синтаксис для типов функций.
%% ({} и т.д.); эти обновления сделаны CL 84908.

%
Наименее верхняя граница типа функции и типа интерфейса  это
наименьшая верхняя граница {} и .
 и  быть функциональными типами.
Если  и  отличаются по количеству требуемых параметров,
Тогда наименьшая верхняя граница  и  составляет .
В противном случае:


 Если


 и





где  наименьшая верхняя граница  и  это





где  является наименьшей верхней границей  и .
 Если


,


 $X_1\ B_1, \ldots,\ X_s\ B_s$> <  


наименьшая верхняя граница  и  это





где  является наименьшей верхней границей  и .
 Если


 $X_1\ B_1, \ldots,\ X_s\ B_s$>$T_1, \ldots,\ T_r,\ $\{$T_{r+1}\ p_{r+1}, \ldots,\ T_f\ p_f$\} ,


 $X_1\ B_1, \ldots,\ X_s\ B_s$>$S_1, \ldots,\ S_r,\ $\{$S_{r+1}\ q_{r+1}, \ldots,\ S_g\ q_g$\} 

Тогда пусть

и пусть  быть наименьшей верхней границей типов  в  и
.
Тогда наименьшая верхняя граница  и  это


< $L_1, \ldots,\ L_r,\ $\{$X_m\ x_m, \ldots,\ X_n\ x_n$\} 

где  является наименьшей верхней границей  и 


{%
Обратите внимание, что необычный случай охватывается использованием ,
в этом случае объявления параметров типа опущены ().%
?


 Ссылка.



{Лексические правила}


%
Исходный текст Dart представлен в виде последовательности точек кода Unicode.
Эта последовательность сначала преобразуется в последовательность токенов
в соответствии с лексическими правилами, приведенными в данной спецификации.
В любой момент процесса токенизации,
Признается самый длинный токен.


<<<p001280>> Зарезервированные слова


%
 Может использоваться только в синтаксических положениях.
определяется грамматикой.
В частности, ошибка времени компиляции возникает при использовании зарезервированного слова.
где ожидается идентификатор.

{%
Обратите внимание, что зарезервированные слова встречаются жирными и нецитируемыми в правилах грамматики.
(например, {})
Даже если последовательное обозначение будет использовать цитаты
(например, {assert}).
Это злоупотребление происходит потому, что мы считаем, что
Это делает правила грамматики более читаемыми.
?

<RESERVED Слово>:= \ASSERT{} | \BREAK{} | \CASE{} | \CATCH{} !
\CLASS{} | \CONST{}
\alt\hspace{-3mm} \CONTINUE{} | \DEFAULT{} | \DO{} | \ELSE{} | \ENUM{} !
\EXTENDS{} | \FALSE{} | \FINAL{} | \FINALLY{} | \FOR{}
\alt\hspace{-3mm} \IF{} | \IN{} | \IS{} | \NEW{} | \NULL{} | \RETHROW{} !
\RETURN{} | \SUPER{} | \SWITCH{} | \THIS{}
\alt\hspace{-3mm} \TRUE{} | {} | {} | \VOID{} | \WHILE{} | {}


%
В грамматике должно происходить правило для зарезервированных слов выше.
перед правилом для  BUILT В  Идентификатор
(P000544).

{%
Это гарантирует, что \synt{IDENTIFIER} и {IDENTIFIER No ДОЛЛАР. Не надо.
Они выводят любые зарезервированные слова и не выводят никаких встроенных идентификаторов.
?


 Комментарии


%
\IndexCustom{Комментарии}{комментарий}
Это разделы текста программы, которые используются для документации.


< LINE Комментарий> ::= <
 (<LINE\_BREAK>)* (перенаправлено с «P001507») Дырнуть?

<MULTI LINE Комментарий> ::= \gnewline{}
"/*" (<MULTI) LINE Комментарий> | \gtilde{} '*/''** "


%
Dart поддерживает как однострочные, так и многострочные комментарии.
 начинается с токена .
Все, что между  и конец строки
Компилятор Dart должен быть проигнорирован
Если комментарий не является комментарием к документации.

%
 начинается с токена 
И кончается токеном .
Все, что между * и /
Компилятор Dart должен быть проигнорирован
Если комментарий не является комментарием к документации.
Комментарии могут гнездиться.

%
\IndexCustom{Комментарии к документации}{комментарии к документации}
Комментарии, которые начинаются с токенов  или .
Комментарии к документации предназначены для обработки
Инструмент, который производит человеческую читаемую документацию.

%
Текущий объем комментария к документации, непосредственно предшествующий
Объявление класса  Это
{Объем тела}{Сфера действия! для органа по заявлениям
.

%
Текущий объем комментария к документации, непосредственно предшествующий
Объявление неперенаправляющего генеративного конструктора 
с параметрами инициализации - формальный параметр инициализации 
().

%
В противном случае текущая сфера применения комментария к документации, непосредственно предшествующего
Объявление функции  Формальный диапазон параметров 
(P000546).


 Прецедент оператора.


%
Преимущество оператора косвенно определяется грамматикой.

{%
Следующая ненормативная таблица может быть полезной


%
?


 Оригинальное название: Algorithmic Subtyping


% Номер правил подтипа
\newcommand{\AppSrnReflexivity>>>{ 1_{\scriptsize\mbox{algo}}}}
\newcommand{\AppSrnTypeVariableReflexivityB}{. 1)
\newcommand{\AppSrnTypeVariableReflexivityC}{. 2.
\newcommand{}{.3.
\newcommand{\AppSrnRightFutureOrC}{\SrnRightFutureOrB. 1)
\newcommand{\AppSrnRightFutureOrD}{\SrnRightFutureOrB. 2.

[h!]
\def\VSP{\vspace{3 мм)
\def\ExtraVSP{\vspace{1mm)
\def\Axiom#1#2#3#4{{[#1]{}{{#3}{#4}}}<<< P001370>>>\def\Rule#1#2#3#4#5#6{\centerline{[#1]{{#3}{#4}}{{#5}{#6}}}<<<< P001377>>>
\def#1#2#3#4#5#6#7#8
<<<<<<<<<<<<<<<<<<<<<<<<<<<<<<<<<<<<<<<<<<<<<<<<<<<<<<<<<<<<<<<<<<<<<<<<<<<<<<<<<<<<<<<<<<<<<<<<<<<<<<<<<<<<<<<<<<<<<<<<<<<<<<<<<<<<<<<<<<<<<<<<<<<<<<<<<<<<<<<<<<<<<<<<<<<<<<<<<<<<<<<<<<<<<<<<<<<<<<<<<<<<<<<<<<<<<<<<<<<<<<<<<<<<<<<<<<<<<<<<<<<<<<<<<<<<<<<< P001385>>>
#1#2#3#4#5
\centerline{[#1]{#3}{{#4}{#5}}}<<< P001391>>>
#1#2#3#4 [#1]{#3}{#4}}\VSP]
%
\begin{minipage}[c]{0.49}
\RuleRaw{\AppSrnReflexivity>>> Рефлексивность {S\mbox{ Не композитный. S}{S}
\Rule{\AppSrnTypeVariableReflexivityC>>> Тип переменной рефлексивности B}{X}{T}{X}{X  T
\Rule{\AppSrnRightFutureOrC>>> Правильное будущее Или C}{\Delta(X)}{}{X}{}

[c]{0.49}
\Axiom{\AppSrnTypeVariableReflexivityB>>> Переменная рефлексивность типа {X}{X}
<\Rule{\AppSrnTypeVariableReflexivityD>>> Переменная рефлексивность типа C {X}  S}{T}{X  S {\displaystyle X}  T
<\Rule{\AppSrnRightFutureOrD>>> Правильное будущее Или D}{S}{{X}  S}{\code{FutureOr<>}

%
 Правила алгоритмического подтипа.
Правила \SrnTop--\SrnSuperinterface{} неизменны и поэтому здесь опущены.



%
Текст этого приложения не является частью спецификации языка Дарта.
Тем не менее, мы по-прежнему используем запись, где точная информация
использует стиль, связанный с нормативным текстом в спецификации (этот стиль);
{в то время как примеры и объяснения используют стиль комментариев (как этот)}.

%
В настоящем добавлении представлен вариант правил подтипа, приведенных
На рисунке ~ на странице ~<\pageref{фиг:правила подтипа}.

{%
Правила докажут тот же набор отношений подтипа,
Но приведенные здесь правила показывают, что существует эффективная реализация.
Это будет определять, будет ли {S}{T} держать,
для любых типов  и .
Легко видеть, что алгоритмические правила докажут самое большее.
одинаковые отношения подтипа,
Все приведенные здесь правила могут быть доказаны.
С помощью правил на рисунке ~.
Также относительно легко нарисовать доказательства.
Алгоритмические правила могут доказать, по крайней мере, те же отношения подтипа.
также при соблюдении следующих ограничений по порядку и прекращению.%
?

%
Единственным измененным правилом является число ~.
Изменяется на .
Это только меняет применимость правила:
Это правило используется только для типов, которые не являются атомными.
{атомный} Тип {type!atomic}
Это тип, который не является переменной типа.
не является продвигаемой переменной типа,
не функциональный тип,
Не параметризованный тип.

{%
Иными словами, правило {} используется для
Особые типы, такие как ,  и ,
Используется для негенеральных классов.
Но он не используется для любого типа операций.
Для этого требуется более одного сравнения, чтобы определить,
Это то же самое, что и любой другой тип.
%
Дело в том, что оставшиеся правила заставят
в любом случае, при необходимости,
Таким образом, мы можем просто опустить начальный структурный переход.
которые могут предпринять много шагов только для того, чтобы сообщить об этих двух терминах большого типа.
Не совсем идентичны.%
?

%
Правила упорядочены по их номерам правил:
Приведенное здесь правило пронумеровано  вставляется сразу после правила ,
Далее следует правило , и так далее.
Далее следует правило, число которого .
{%
Так что приказ
, \SrnTop--\SrnTypeVariableReflexivityA,
, ,
,
 и так далее.%
?

%
Теперь мы определим процедуру, которая используется для определения того,
{S}{T} держит,для некоторых конкретных типов  и :
Выберите первое правило  чьи синтаксические ограничения удовлетворены
по данным типам  и ,
и далее показать, что его помещения держатся.
Если это так, мы заканчиваем и заключаем, что отношения подтипа имеют место.
В противном случае мы прекращаем и завершаем
что отношения подтипа не поддерживаются,
За исключением случаев, когда  это
, ,
, или .
{%
В частности, для исходного запроса \SubtypeStd Сюда!
Мы не отказываемся от использования правила, которое
более высокое число правил, чем у ,
Но мы можем попробовать все
Правила с  вправо.%
?

{%
Помимо того, что вся сложность субтипирования
потенциально возникает каждый раз, когда проверяется, имеет ли предпосылка
Проверки, применяемые для каждого правила, связаны с количеством работы.
которые являются постоянными для всех правил, кроме следующих:
Во-первых, группа правил
, ,
 и 
Это может привести к обратному ходу.
Далее, правила \SrnPositionalFunctionType-\SrnCovariance{}
требует работы, пропорциональной размерам  и ,
из-за количества помещений, которые необходимо проверить.
Наконец, правило ~\SrnSuperinterface требует работы
пропорционально размеру ,
Это также может повлечь за собой затраты на поиск всего набора данных.
прямые и косвенные суперинтерфейсы подтипа кандидата ,
До тех пор, пока не будет показана соответствующая предпосылка для одного из них,
Если есть.

Применяются дополнительные оптимизации.
Например,
Сразу можно сделать вывод, что отношения подтипа не
когда мы собираемся проверить правило ~
Если  Это переменная типа или тип функции.
Для некоторых других типов, например,
- переменная продвигаемого типа,
, , 
 для любого  или ,
Известно, что это никогда не произойдет как  для правила ~,
Что означает, что этот, казалось бы, дорогой шаг
может быть ограничено до некоторой степени.%
?


 Добавление: Целые варианты осуществления


{%
 Тип представляет целые числа.
Спецификация написана с целыми числами дополнения 64-битных двух в качестве
предполагаемой реализации. Но когда Дарт компилируется в JavaScript,
Использование  Вместо этого он будет использовать тип номера JavaScript.
и соответствующих операций JavaScript,
За исключением битовых операций, как описано ниже.

Это приводит к ряду различий:


 
Действительные значения JavaScript  IEEE-754 64-битный
Число с плавающей точкой без дробной части. Это включает
Положительный и отрицательный , который может быть достигнут
переполнение (целое деление на ноль по-прежнему является динамической ошибкой).
В противном случае действительные целые буквы (включая любой ведущий знак минус)
Представляет собой недействительный JavaScript  Значения составляют время компиляции
Ошибки при компиляции на JavaScript. Операции на целых числах могут
теряют точность, потому что операнды и результат представлены
64-битные числа с плавающей запятой, ограниченные 53-значными
кусочки.
 
JavaScript  Также используются экземпляры  и
целочисленное значение  Примеры также реализуют .
 и  Классы остаются отдельными подклассами
класса , но {инстанции} любого класса, который
представляют собой целое действие, как если бы они на самом деле были
Общий подкласс, использующий оба  и .
Фракционные числа реализуют только .
[]Битовые операции на целых числах (\lit{\&}, { |}} ,
 \lit{} и {} Все усеченные
операнды к целым числам комплемента 32-битных двух, выполняют 32-битные
операции на них, и полученные 32 бита интерпретируются как
32-битное неподписанное целое число. Например,  оценивать
4294967294 (перенаправлено с «P000760»). Правый
оператор смены, \lit{\gtgt}, выполняет подписанное правое смещение
32 бита, когда исходное число отрицательное, и неподписанное право
смещение, когда исходное число не является отрицательным. Оба вида сдвига
Читать далее  в положение , ; но подписано
сдвиг оставляет бит 31 неизменным, а неподписанный сдвиг записывает 0
Как раз 31. Например:

.
В этом примере мы отмечаем, что  и 
получить представление дополнения 32-битных двух ,
Но они имеют разные значения для бита знака IEEE 754.
[]
 Способ не может различать значения  и
, и он не может распознать любой  значение как идентичное
Для себя. Для эффективности:  Операция использует
JavaScript Оператор P000768.
%
?



\end{document}
